\chapter{Matrix Product Unitaries}\label{ch:mpu}

This chapter develops the algebraic prerequisites for matrix product unitaries.
The first result supplies the corrected nil-matrix bound used in the blocking
argument of \cite{Cirac2017MPU}.

\begin{theorem}[Dimension bound for a nilpotent Lie action]
    \label{thm:mpu_lcs_finrank}
    \lean{LieModule.lowerCentralSeries_finrank_eq_bot}
    \leanok
    Let a Lie algebra $L$ act nilpotently on a finite-dimensional vector space
    $V$ over a field $K$.  Write $C_r(L,V)$ for the $r$th term of the lower
    central series of this action.  Then
    \begin{align}
        C_{\dim_K(V)}(L,V) & = 0.
    \end{align}
\end{theorem}

\begin{proof}\leanok
    As long as $C_r(L,V)$ is nonzero, nilpotency prevents two consecutive
    terms from being equal: equality would make every later term equal, in
    contradiction with eventual vanishing.  Thus each nonzero step strictly
    lowers dimension, and at most $\dim_K(V)$ such steps are possible.
\end{proof}

\begin{theorem}[Finite-dimensional nil-matrix theorem]
    \label{thm:mpu_nil_matrix}
    \lean{NonUnitalSubalgebra.list_prod_eq_zero_of_forall_isNilpotent}
    \leanok
    Let $V$ be a finite-dimensional vector space over a field $K$, and let
    $A\subseteq\operatorname{End}_K(V)$ be a not necessarily unital
    associative subalgebra.  If every element of $A$ is nilpotent, then every
    ordered product of $\dim_K(V)$ elements of $A$ is zero.  The empty product
    in dimension zero is also zero because the identity endomorphism of the
    zero space is zero.
\end{theorem}

\begin{proof}\leanok
    \uses{thm:mpu_lcs_finrank}
    Regard $A$ as a Lie algebra under the commutator bracket.  Engel's theorem
    makes $V$ a nilpotent Lie module, so the dimension bound gives
    \begin{align}
        C_{\dim_K(V)}(A,V) & = 0.
    \end{align}
    For every nonnegative integer $r$, every ordered family
    $a_1,\ldots,a_r$ in $A$, and every vector $v$ in $V$, induction on $r$
    gives
    \begin{align}
        (a_1\cdots a_r)v & \in C_r(A,V).
    \end{align}
    Taking $r=\dim_K(V)$ proves that the product annihilates every vector and
    is therefore zero.
\end{proof}

\begin{theorem}[Nil matrix subalgebras]
    \label{thm:mpu_nil_matrices}
    \lean{NonUnitalSubalgebra.matrix_list_prod_eq_zero_of_forall_isNilpotent}
    \leanok
    Let $A$ be a not necessarily unital associative subalgebra of
    $n\times n$ matrices over a field.  If every matrix in $A$ is nilpotent,
    then every ordered product of $n$ matrices from $A$ vanishes.
\end{theorem}

\begin{proof}\leanok
    \uses{thm:mpu_nil_matrix}
    Transport the matrix algebra to the endomorphisms of the coordinate space
    through the standard algebra equivalence.  Nilpotency, membership, and the
    ordered product are preserved, so the finite-dimensional nil-matrix
    theorem applies; injectivity of the equivalence returns the zero matrix.
\end{proof}

The length $n$ is a weak bound, not a strict bound below $n$; the latter is
false for strictly upper-triangular matrices.  This corrects the intermediate
strict inequality in \cite{Cirac2017MPU} without changing its
final strict $D^4$ estimate.

\section{The matrix product unitary condition}

Let $\mathcal U=\{U^{ij}\}_{i,j=0}^{d-1}$ be a matrix product operator tensor of
bond dimension $D$, and write $U^{(N)}$ for the periodic operator obtained by
closing a chain of $N$ copies of $\mathcal U$.

\begin{theorem}[Unitarity under simultaneous reindexing]
    \label{thm:mpu_reindex_unitary}
    \lean{Matrix.reindex_mem_unitaryGroup}
    \leanok
    Let $I$ and $J$ be finite sets, let $e\colon I\to J$ be a bijection, and
    let $A$ be a unitary matrix indexed by $I$.  The matrix indexed by $J$
    whose $(j,k)$ entry is
    \begin{align}
        A_{e^{-1}(j),e^{-1}(k)}
    \end{align}
    is unitary.
\end{theorem}

\begin{proof}\leanok
    Simultaneous reindexing preserves the identity and matrix multiplication,
    and it commutes with the conjugate transpose.  Reindex the equation
    $AA^\dagger=\Id$.
\end{proof}

\begin{definition}[Matrix product unitary tensor]
    \label{def:mpu}
    \lean{MPOTensor.IsMPU}
    \lean{MPOTensor.IsMPU.mpo_mem_unitaryGroup}
    \leanok
    \uses{def:mpo_family}
    The tensor $\mathcal U$ generates matrix product unitaries if $U^{(N)}$ is
    unitary for every integer $N>1$.  This is the matrix product unitary
    condition in \cite[Section~III, equation~(10)]{Cirac2017MPU}.
\end{definition}

\begin{lemma}[First matrix product unitarity equation]
    \label{lem:mpu_mul_adjoint}
    \lean{MPOTensor.IsMPU.mpo_mul_conjTranspose_mpo}
    \leanok
    \uses{def:mpu}
    If $\mathcal U$ generates matrix product unitaries and $N>1$, then
    \begin{align}
        U^{(N)}\bigl(U^{(N)}\bigr)^\dagger & = \Id.
    \end{align}
\end{lemma}

\begin{proof}\leanok
    \uses{def:mpu}
    Expand the unitary-group condition for $U^{(N)}$.
\end{proof}

\begin{lemma}[Second matrix product unitarity equation]
    \label{lem:mpu_adjoint_mul}
    \lean{MPOTensor.IsMPU.conjTranspose_mpo_mul_mpo}
    \leanok
    \uses{def:mpu}
    If $\mathcal U$ generates matrix product unitaries and $N>1$, then
    \begin{align}
        \bigl(U^{(N)}\bigr)^\dagger U^{(N)} & = \Id.
    \end{align}
    This is the unitarity equation displayed in
    \cite[Section~III, equation~(12)]{Cirac2017MPU}.
\end{lemma}

\begin{proof}\leanok
    \uses{def:mpu}
    Use the equivalent left-multiplied characterization of a unitary matrix.
\end{proof}

\begin{theorem}[Composition preserves matrix product unitarity]
    \label{thm:mpu_composition}
    \lean{MPOTensor.IsMPU.mulTensor}
    \leanok
    \uses{def:mpu, def:mpdo_operator_product, thm:mpdo_operator_product_mpo}
    Let $\mathcal U$ and $\mathcal V$ be matrix product unitary tensors with the
    same physical dimension.  Contracting the output index of $\mathcal U$
    with the input index of $\mathcal V$ gives a matrix product unitary tensor
    whose periodic operator is $U^{(N)}V^{(N)}$, in that order.  This is the
    concatenation tensor $\mathcal W$ in the proof of
    \cite[Theorem~IV.6(ii)]{Cirac2017MPU}.
\end{theorem}

\begin{proof}\leanok
    \uses{def:mpu, thm:mpdo_operator_product_mpo}
    For every $N>1$, both $U^{(N)}$ and $V^{(N)}$ are unitary.  The defining
    contraction of the product tensor gives the periodic operator
    $U^{(N)}V^{(N)}$, which is unitary because the unitary matrices form a
    group under multiplication.
\end{proof}

\begin{lemma}[Physical blocking preserves matrix product unitarity]
    \label{lem:mpu_blocking}
    \lean{MPOTensor.IsMPU.blockTensor}
    \leanok
    \uses{def:mpu}
    Suppose that $\mathcal U$ generates matrix product unitaries.  For every
    positive integer $L$, the tensor $\mathcal U^{[L]}$ obtained by blocking
    $L$ adjacent sites also generates matrix product unitaries.  This is the
    blocking operation of \cite[Definition~II.5 and equation~(9)]{Cirac2017MPU};
    preservation of the MPU property under blocking is used in the proof of
    \cite[Proposition~III.3]{Cirac2017MPU}.
\end{lemma}

\begin{proof}\leanok
    \uses{def:mpu, thm:mpu_reindex_unitary, thm:mpdo_closed_chain_physical_blocking}
    At system size $N$, closing the blocked tensor gives $U^{(NL)}$, up to the
    canonical bijection between $N$ blocked configurations and $NL$ unblocked
    configurations.  If $N>1$ and $L>0$, then $NL>1$, so $U^{(NL)}$ is
    unitary.  Simultaneous reindexing preserves unitarity.
\end{proof}

\section{Simple matrix product unitary tensors}

\begin{definition}[Double layer]
    \label{def:mpu_double_layer}
    \lean{MPOTensor.doubleLayerTensor,
        MPOTensor.doubleLayerTensor_apply}
    \leanok
    \uses{def:mpdo_physical_adjoint, def:mpdo_operator_product}
    For an MPO tensor $\mathcal U$, let $\mathcal W$ be the double-layer
    tensor obtained by placing the physical adjoint above $\mathcal U$ and
    contracting the adjacent physical legs:
    \begin{align}
        \mathcal W^{ik}
         & =\sum_j (\mathcal U^\sharp)^{ij}\otimes\mathcal U^{jk}.
    \end{align}
    The bond space of $\mathcal W$ is the product of the two original bond
    spaces.  The physical adjoint exchanges and conjugates the physical
    indices without reversing the order of the virtual chain, exactly as in
    the barred upper layer of Figures~\texttt{II\_Simple1.png} and
    \texttt{II\_Simple2.png} of \cite{Cirac2017MPU}.
\end{definition}

\begin{theorem}[Closed double layer]
    \label{thm:mpu_closed_double_layer}
    \lean{MPOTensor.mpo_doubleLayerTensor}
    \leanok
    \uses{def:mpu_double_layer}
    For every chain length $N$, closing the double layer gives
    \begin{align}
        W^{(N)} & =\bigl(U^{(N)}\bigr)^\dagger U^{(N)}.
    \end{align}
\end{theorem}

\begin{proof}\leanok
    \uses{def:mpdo_physical_adjoint, thm:mpdo_operator_product_mpo}
    The operator family is multiplicative under the local product tensor, and
    the physical adjoint generates the conjugate transpose without reflecting
    the chain.
\end{proof}

\begin{theorem}[Physical-adjoint identities for matrix product unitaries]
    \label{thm:mpu_physical_adjoint_identities}
    \lean{MPOTensor.physicalAdjointTensor_physicalAdjointTensor,
        MPOTensor.IsMPU.physicalAdjointTensor,
        MPOTensor.mpo_physicalAdjointTensor_eq_conjTranspose}
    \leanok
    \uses{def:mpu, def:mpdo_physical_adjoint}
    The physical adjoint is involutive. If $\mathcal U$ generates matrix
    product unitaries, then so does $\mathcal U^\sharp$, and for every chain
    length $N$,
    \begin{align}
        (\mathcal U^\sharp)^\sharp & =\mathcal U,
                                   & (\mathcal U^\sharp)^{(N)} & =\bigl(\mathcal U^{(N)}\bigr)^\dagger.
    \end{align}
\end{theorem}

\begin{proof}\leanok
    \uses{def:mpu, def:mpdo_physical_adjoint, lem:mpu_mul_adjoint}
    Applying the physical adjoint twice conjugates every entry twice and
    exchanges the physical indices twice. Closing a physical-adjoint tensor
    exchanges the two physical configurations and conjugates the resulting
    scalar. Hence its periodic operator is the conjugate transpose. The first
    MPU unitarity equation for $U^{(N)}$ is therefore the second unitarity
    equation for $(U^\sharp)^{(N)}$.
\end{proof}

\begin{theorem}[Output-first double-layer entry]
    \label{thm:mpu_output_first_double_layer_entry}
    \lean{MPOTensor.doubleLayerTensor_physicalAdjointTensor_apply}
    \leanok
    \uses{def:mpu_double_layer, def:mpdo_physical_adjoint}
    Order the left doubled bond as $(\alpha,\gamma)$ and the right doubled
    bond as $(\beta,\delta)$. Then the double layer whose lower tensor is
    $\mathcal U^\sharp$ has entries
    \begin{align}
        \mathcal W(\mathcal U^\sharp)^{ik}_{(\alpha,\gamma),(\beta,\delta)}
         & =\sum_j
        \mathcal U^{ij}_{\alpha,\beta}\,
        \overline{\mathcal U^{kj}_{\gamma,\delta}}.
    \end{align}
\end{theorem}

\begin{proof}\leanok
    \uses{def:mpu_double_layer, def:mpdo_physical_adjoint}
    Expand the double layer and then apply the physical adjoint to its upper
    factor. The two exchanges of physical indices leave the contracted index
    $j$ in the second physical position of both factors. The product-index
    convention keeps $(\alpha,\gamma)$ and $(\beta,\delta)$ in the displayed
    order.
\end{proof}

\begin{definition}[Simple matrix product unitary tensor]
    \label{def:mpu_simple}
    \lean{MPOTensor.IsMPUSimple,
        MPOTensor.IsMPUSimple.simple1,
        MPOTensor.IsMPUSimple.simple2}
    \leanok
    \uses{def:mpu_double_layer}
    The tensor $\mathcal U$ is \emph{simple} if there are vectors $a,b$ in
    the double-layer bond space such that, for all physical indices
    $i,j,k,\ell$,
    \begin{align}
        (a|W^{ij}|b) & =\delta_{ij},
        \tag{simple1}\label{eq:mpu_simple1}    \\
        W^{ij}W^{k\ell}
                     & =W^{ij}|b)(a|W^{k\ell}.
        \tag{simple2}\label{eq:mpu_simple2}
    \end{align}
    Thus (\ref{eq:mpu_simple1}) closes one double-layer letter to the
    physical identity, while (\ref{eq:mpu_simple2}) inserts the rank-one
    operator $|b)(a|$ between adjacent letters.  These are precisely the two
    diagrams in \cite[Definition~III.2]{Cirac2017MPU}.
\end{definition}

\begin{theorem}[Sequential contraction of a simple double layer]
    \label{thm:mpu_simple_closed_identity}
    \lean{MPOTensor.IsMPUSimple.mpo_doubleLayerTensor_eq_one}
    \leanok
    \uses{def:mpu_simple}
    If $\mathcal U$ is simple and $N>1$, then
    \begin{align}
        W^{(N)} & =\Id.
    \end{align}
\end{theorem}

\begin{proof}\leanok
    \uses{def:mpu_simple}
    Apply (\ref{eq:mpu_simple2}) sequentially between neighboring letters.
    The inserted rank-one operators split the periodic contraction into the
    product of the one-letter contractions $(a|W^{\sigma_x\tau_x}|b)$.
    Identity~(\ref{eq:mpu_simple1}) turns this product into
    $\prod_x\delta_{\sigma_x,\tau_x}$, which is the matrix element of the
    physical identity.
\end{proof}

\begin{theorem}[A simple tensor generates matrix product unitaries]
    \label{thm:mpu_simple_is_mpu}
    \lean{MPOTensor.IsMPUSimple.isMPU}
    \leanok
    \uses{def:mpu, def:mpu_simple}
    Every simple tensor generates matrix product unitaries.
    This is \cite[Proposition~III.3(i)]{Cirac2017MPU}.
\end{theorem}

\begin{proof}\leanok
    \uses{thm:mpu_closed_double_layer, thm:mpu_simple_closed_identity}
    For every $N>1$, the closed-double-layer identity and the sequential
    contraction give
    \begin{align}
        \bigl(U^{(N)}\bigr)^\dagger U^{(N)} & =\Id.
    \end{align}
    Hence $U^{(N)}$ is unitary.
\end{proof}


\section{Double-layer blocking and mixed physical traces}

\begin{theorem}[Physical adjoints and double layers commute with blocking]
    \label{thm:mpu_double_layer_blocking}
    \lean{MPOTensor.physicalAdjointTensor_blockTensor,
        MPOTensor.doubleLayerTensor_blockTensor}
    \leanok
    \uses{def:mpu_double_layer, thm:mpdo_closed_chain_physical_blocking}
    For every blocking length $L$, physical adjunction commutes with blocking,
    and consequently
    \begin{align}
        \mathcal W(\mathcal U^{[L]}) & =\mathcal W(\mathcal U)^{[L]}.
    \end{align}
    The virtual order is unchanged on both sides.
\end{theorem}

\begin{proof}\leanok
    Physical adjunction conjugates each matrix entry without reversing the
    virtual word.  The first identity follows entrywise.  The second follows
    because physical blocking commutes with the local operator product.
\end{proof}

\begin{theorem}[Blocked double-layer diagonal is a transfer power]
    \label{thm:mpu_blocked_double_layer_diagonal}
    \lean{MPOTensor.normalizedDiagonal_doubleLayerTensor,
        MPOTensor.normalizedDiagonal_blockTensor,
        MPOTensor.normalizedDiagonal_doubleLayerTensor_blockTensor,
        MPOTensor.IsMPU.normalizedDiagonal_doubleLayerTensor_blockTensor_eq_vecMulVec}
    \leanok
    \uses{def:mpu_double_layer, def:mpu_normalized_flattening, thm:mpu_double_layer_blocking}
    Assume $d>0$, and let $E$ be the normalized transfer matrix of
    $\mathcal U$.  After blocking $L$ sites, the normalized physical diagonal of the double layer is
    exactly $E^L$.  More precisely, if
    \begin{align}
        q:\Fin D\times\Fin D & \longrightarrow\Fin{D^2}
    \end{align}
    is the standard product encoding, then its entries are
    \begin{align}
        \bigl(E^L\bigr)_{q^{-1}(a),q^{-1}(b)}.
    \end{align}
    Thus the row and column reindexings have the same explicit orientation.
    If $\mathcal U$ is an MPU and $D>1$, then at the stabilization exponent
    $J$ this diagonal is the same reindexing of the rank-one projector
    $|\rho)(\Phi|$.
\end{theorem}

\begin{proof}\leanok
    \uses{thm:mpu_double_layer_blocking, thm:mpu_normalized_transfer_stabilization}
    On one site, expanding the diagonal double-layer contraction gives the
    transfer matrix after the stated product reindexing.  The factor
    $d^{-1}$ is the square of the local factor $d^{-1/2}$.  Summing blocked
    diagonal letters over words and distributing the ordered product shows
    that normalized diagonals turn blocking into powers.  Simultaneous
    reindexing preserves multiplication, giving the displayed formula.
    Transfer stabilization gives the final rank-one specialization.
\end{proof}

\begin{theorem}[Normalized diagonal as a diagonal-word sum]
    \label{thm:mpu_normalized_diagonal_word_expansion}
    \lean{MPOTensor.normalizedDiagonal_pow_eq_sum_evalWord}
    \leanok
    \uses{def:mpo_eval_word}
    Let $\mathcal W$ be an MPO tensor of physical dimension $d$, and
    let
    \begin{align}
        R & =\frac1d\sum_{i=0}^{d-1}W^{ii}
    \end{align}
    be its normalized physical diagonal.  Here and in the two diagonal-tail
    identities below, scalar inverses are the totalized complex-field inverses,
    so $0^{-1}=0$.  For every $K\geq0$,
    \begin{align}
        R^K
         & =d^{-K}\sum_{\tau\in\{0,\ldots,d-1\}^K}
        W^{\tau_0\tau_0}\cdots W^{\tau_{K-1}\tau_{K-1}}.
    \end{align}
    At $K=0$, the word product is the identity matrix.  This is the algebraic
    diagonal-tail expansion used in the source-factor networks of
    \cite[Theorem~III.8, Section~III.B, equations~(31)--(32)]{Cirac2017MPU}.
\end{theorem}

\begin{proof}\leanok
    \uses{def:mpo_eval_word}
    Expand the $K$th power of the finite sum defining $R$.  Distributivity
    indexes the resulting terms by physical words $\tau$, and each term is
    the ordered diagonal word displayed above.
\end{proof}

\begin{theorem}[A matrix basis adapted to the trace]
    \label{thm:identity_traceless_matrix_basis}
    \lean{Matrix.exists_identity_traceless_basis}
    \leanok
    If $d>0$, the identity matrix extends to a basis of $\MN d$ whose other
    elements are traceless.
\end{theorem}

\begin{proof}\leanok
    The trace functional is nonzero on the identity. Extend the identity by a
    basis of the kernel of the trace.
\end{proof}

\begin{theorem}[Identity plus traceless physical decomposition]
    \label{thm:mpu_WIsom}
    \lean{MPOTensor.exists_identity_add_traceless_decomposition,
        MPOTensor.residualSlice_eq_contractPhysical_tracelessPart}
    \leanok
    \uses{thm:identity_traceless_matrix_basis}
    Assume $d>0$.  For every double-layer tensor $\mathcal W$, there is a basis
    $\{\sigma_0,\sigma_\alpha\}$ of the physical matrix algebra and virtual
    matrices $E,S_\alpha$ such that $\sigma_0=\Id$ and, for every nonzero
    index $\alpha$, $\tr(\sigma_\alpha)=0$, while
    \begin{align}
        \mathcal W & =E\otimes\Id+
        \sum_{\alpha\ne0}S_\alpha\otimes\sigma_\alpha.
    \end{align}
    Here $E$ is the normalized physical diagonal.  Equivalently, contraction
    against a physical matrix $X$ separates as
    \begin{align}
        \operatorname{contr}_{\rm phys}(\mathcal W,X)
         & =\tr(X)E+R_{\mathcal W}(X),      \\
        R_{\mathcal W}(X)
         & =\operatorname{contr}_{\rm phys}
        \left(\mathcal W,
        X-\frac{\tr(X)}d\Id\right),
    \end{align}
    where the second physical argument is traceless.  This is equation
    \texttt{WIsom} in \cite[Proposition~III.3(ii)]{Cirac2017MPU}.
\end{theorem}

\begin{proof}\leanok
    The trace functional is nonzero on the identity.  Extend the identity to
    a basis adapted to the kernel of the trace; all remaining basis vectors
    are therefore traceless.  The bilinear pairing
    $(X,Y)\mapsto\tr(XY)$ is nondegenerate.  Its dual basis gives the virtual
    coefficients by physical contraction.  The dual vector paired with the
    identity is $d^{-1}\Id$, so its virtual coefficient is precisely $E$.
    Expanding matrix units in the dual basis yields the decomposition.
    Subtracting the trace component gives the equivalent residual formula.
\end{proof}

\begin{theorem}[Multilinear physical contraction]
    \label{thm:mpu_multilinear_physical_contraction}
    \lean{MPOTensor.trace_prod_contractPhysical_of_mpo_eq_one,
        MPOTensor.IsMPU.trace_prod_contractPhysical_doubleLayerTensor}
    \leanok
    \uses{def:mpu, def:mpu_double_layer}
    Suppose a length-$N$ closed tensor equals the physical identity.  For
    arbitrary physical matrices $X_0,\ldots,X_{N-1}$,
    \begin{align}
        \tr\!\left(\prod_{k=0}^{N-1}
        \operatorname{contr}_{\rm phys}(\mathcal W,X_k)\right)
         & =\prod_{k=0}^{N-1}\tr(X_k).
    \end{align}
    In particular this applies to the double layer of an MPU whenever $N>1$.
\end{theorem}

\begin{proof}\leanok
    Expand each contraction and distribute the ordered product.  The closed
    tensor coefficient is a Kronecker delta between the two physical words,
    so only equal words survive.  The remaining independent diagonal sums
    factor into the product of the physical traces.
\end{proof}

% Resolved local correction: docs/paper-gaps/mpu_one_letter_mixed_residual_trace.tex
\begin{theorem}[Positive-length mixed residual traces]
    \label{thm:mpu_positive_mixed_residual_trace}
    \lean{MPOTensor.IsMPU.trace_residualSlice_doubleLayerTensor_eq_zero,
        MPOTensor.IsMPU.trace_prod_mixedResidualFactor_doubleLayerTensor_eq_zero_of_pos}
    \leanok
    \uses{def:mpu, def:mpu_double_layer}
    Let $\mathcal U$ be an MPU with physical dimension $d>0$.  Every closed
    product of positive length that contains at least one residual coefficient
    has zero trace.  In particular, this includes a single residual
    coefficient.  This is the mixed-trace assertion preceding equation
    \texttt{ESE=0} in
    \cite[Proposition~III.3(ii)]{Cirac2017MPU}.
\end{theorem}

\begin{proof}\leanok
    \uses{thm:mpu_WIsom, thm:mpu_multilinear_physical_contraction, thm:mpu_shifted_zero_trace_power_nilpotent}
    At lengths greater than one, represent identity factors by
    $d^{-1}\Id$ and residual factors by traceless physical arguments.  The
    multilinear contraction identity then vanishes.  For a single residual
    coefficient $A$, repeat the same coefficient at every length $q>1$ to
    obtain $\tr(A^q)=0$.  Thus the one-letter case is recovered after the
    longer words: the shifted zero-trace-power theorem makes $A$ nilpotent,
    and therefore $\tr(A)=0$.
\end{proof}

\begin{theorem}[Controlled mixed residual traces and rank-one sandwich]
    \label{thm:mpu_controlled_mixed_residual_trace}
    \lean{MPOTensor.IsMPU.trace_prod_mixedResidualFactor_doubleLayerTensor_eq_zero,
        MPOTensor.IsMPU.normalizedDiagonal_mul_prod_residualSlice_mul_normalizedDiagonal_eq_zero,
        MPOTensor.IsMPU.blocked_normalizedDiagonal_sandwich_residual_prod_eq_zero}
    \leanok
    \uses{def:mpu, def:mpu_double_layer, thm:mpu_blocked_double_layer_diagonal}
    Let $\mathcal U$ be an MPU with physical dimension $d>0$, and let
    $N>1$.  Every closed product of $N$ identity and residual
    coefficients that contains at least one residual coefficient has zero
    trace.  If the normalized diagonal is rank one, then every nonempty
    residual word $P$ satisfies
    \begin{align}
        EPE & =0.
    \end{align}
    The same conclusion holds after blocking $J>0$ sites whenever the
    $J$th normalized transfer power is $|\rho)(\Phi|$, with both virtual
    indices reindexed by $q^{-1}:\Fin{D^2}\longrightarrow\Fin D\times\Fin D$.
\end{theorem}

\begin{proof}\leanok
    \uses{thm:mpu_WIsom, thm:mpu_multilinear_physical_contraction, thm:mpu_blocked_double_layer_diagonal}
    Represent identity factors by $d^{-1}\Id$ and residual factors by their
    traceless physical arguments.  Since $N>1$, the multilinear contraction
    identity applies, and one trace factor vanishes.  For a nonempty residual
    word, adjoin $E$ as one extra factor; its closed trace has length at least
    two and therefore vanishes.  Cyclicity of trace and
    \begin{align}
        |\rho)(\Phi|P|\rho)(\Phi|
         & =\tr\!\left(P|\rho)(\Phi|\right)|\rho)(\Phi|
    \end{align}
    give $EPE=0$.  The blocked statement follows from the transfer-power
    formula with the same product-index reindexing in both matrix indices.
\end{proof}

\begin{definition}[Residual nonunital algebra]
    \label{def:mpu_residual_algebra}
    \lean{MPOTensor.residualGeneratorSet,MPOTensor.residualAlgebra,
        MPOTensor.residualSlice_mem_residualAlgebra}
    \leanok
    Let $\mathcal U$ have bond dimension $D$.  Contract the double layer
    against arbitrary physical matrices, subtract the normalized diagonal
    coefficient, and let $\mathcal S$ be the resulting set of residual
    matrices in $M_{D^2}(\mathbb C)$.  The residual algebra
    $\mathfrak S$ is the not necessarily unital algebra generated by
    $\mathcal S$.  This is the algebra generated by the matrices
    $S'_\alpha$ in \cite[Proposition~III.3(ii)]{Cirac2017MPU}.
\end{definition}

\begin{theorem}[Trace and nilpotence of the residual algebra]
    \label{thm:mpu_residual_algebra_nilpotent}
    \lean{MPOTensor.IsMPU.trace_list_prod_eq_zero_of_mem_residualGeneratorSet,
        MPOTensor.IsMPU.trace_eq_zero_of_mem_residualAlgebra,
        MPOTensor.IsMPU.trace_pow_eq_zero_of_mem_residualAlgebra,
        MPOTensor.IsMPU.isNilpotent_of_mem_residualAlgebra}
    \leanok
    \uses{def:mpu, def:mpu_residual_algebra}
    If $\mathcal U$ is an MPU with positive physical dimension, then every
    nonempty product of residual generators has zero trace.  Consequently,
    for every $A\in\mathfrak S$ and every integer $N>1$,
    \begin{align}
        \tr(A) & =0,
               & \tr(A^N) & =0,
    \end{align}
    and every element of $\mathfrak S$ is nilpotent.
\end{theorem}

\begin{proof}\leanok
    \uses{thm:mpu_positive_mixed_residual_trace, thm:mpu_shifted_zero_trace_power_nilpotent}
    A nonempty product in the multiplicative semigroup generated by
    $\mathcal S$ is a nonempty residual word, so its trace vanishes by the
    positive-length mixed residual identity.  The generated nonunital algebra
    is the linear span of this semigroup.  Linearity of the trace therefore
    gives $\tr(A)=0$ for every $A\in\mathfrak S$.  Since every positive power
    of $A$ remains in $\mathfrak S$, the same argument gives
    $\tr(A^N)=0$ for $N>1$.  The shifted zero-trace-power theorem now implies
    that $A$ is nilpotent.
\end{proof}

\begin{theorem}[Length-$D^2$ residual products vanish]
    \label{thm:mpu_residual_product_d_sq}
    \lean{MPOTensor.IsMPU.residualAlgebra_list_prod_eq_zero,
        MPOTensor.IsMPU.prod_residualSlice_doubleLayerTensor_eq_zero}
    \leanok
    \uses{def:mpu, def:mpu_residual_algebra}
    Let $\mathcal U$ be an MPU of bond dimension $D$ and positive physical
    dimension.  Every ordered product of exactly $D^2$ elements of
    $\mathfrak S$ vanishes.  In particular, for arbitrary residual generators,
    \begin{align}
        S_{\alpha_1}\cdots S_{\alpha_{D^2}} & =0.
    \end{align}
    This is equation \texttt{Sprime=0} in
    \cite[Proposition~III.3(ii)]{Cirac2017MPU}, with the corrected exact length
    $D^2$ in place of the unsupported strict intermediate bound.
\end{theorem}

\begin{proof}\leanok
    \uses{def:mpu_residual_algebra, thm:mpu_nil_matrices, thm:mpu_residual_algebra_nilpotent}
    The algebra $\mathfrak S$ acts on a space of dimension $D^2$, and every
    element is nilpotent.  The finite-dimensional nil-matrix theorem therefore
    annihilates every ordered product of $D^2$ elements.  Each residual
    generator belongs to $\mathfrak S$, giving the stated specialization.
\end{proof}

\begin{definition}[Three-form residual subspace]
    \label{def:mpu_three_form_subspace}
    \lean{MPOTensor.projectorMulResidualSubmodule,
        MPOTensor.residualMulProjectorSubmodule,
        MPOTensor.residualMulProjectorMulResidualSubmodule,
        MPOTensor.threeFormSubmodule}
    \leanok
    \uses{def:mpu_residual_algebra}
    Let $E$ be the normalized diagonal after the first blocking and let
    $\mathfrak S$ be its residual nonunital algebra.  Define
    \begin{align}
        \mathcal T={} & \operatorname{span}\{EA:A\in\mathfrak S\}
        +\operatorname{span}\{AE:A\in\mathfrak S\}                     \\
                      & +\operatorname{span}\{AEB:A,B\in\mathfrak S\}.
    \end{align}
    This is the basis-independent sum of the three forms in equation
    \texttt{sprimeforms} of
    \cite[Proposition~III.3(ii)]{Cirac2017MPU}.
\end{definition}

\begin{theorem}[Ordered second-block expansion]
    \label{thm:mpu_ordered_second_block_expansion}
    \lean{MPOTensor.ordered_prod_smul_projector_add_sub_mem_threeFormSubmodule}
    \leanok
    \uses{def:mpu_three_form_subspace}
    Suppose $E^2=E$, $EAE=0$ for every $A\in\mathfrak S$, and every ordered
    product of $N>0$ elements of $\mathfrak S$ vanishes.  For scalars $c_k$
    and elements $S_k\in\mathfrak S$,
    \begin{align}
        \prod_{k=0}^{N-1}(c_kE+S_k)
        -\left(\prod_{k=0}^{N-1}c_k\right)E & \in\mathcal T.
    \end{align}
    Products are taken in increasing site order and need not commute.
\end{theorem}

\begin{proof}\leanok
    Expand the product as a sum over binary choices of $E$ or $S_k$, retaining
    the original order.  The all-$E$ choice is the subtracted term because
    $E^q=E$ for $q>0$.  The all-residual choice vanishes by the length-$N$
    hypothesis.  In every remaining word, contiguous copies of $E$ collapse.
    If two resulting copies of $E$ are separated by a nonempty residual word,
    the factor $EAE$ makes the word zero.  Thus each surviving word has one
    contiguous projector block and belongs to one of the three summands of
    $\mathcal T$.
\end{proof}

\begin{theorem}[Second-block residuals lie in the three-form subspace]
    \label{thm:mpu_second_block_residual_three_forms}
    \lean{MPOTensor.IsMPU.normalizedDiagonal_mul_mem_residualAlgebra_mul_normalizedDiagonal_eq_zero,
        MPOTensor.IsMPU.residualSlice_doubleLayerTensor_blockTensor_single_mem_threeFormSubmodule,
        MPOTensor.IsMPU.residualSlice_doubleLayerTensor_blockTensor_mem_threeFormSubmodule}
    \leanok
    \uses{def:mpu, def:mpu_residual_algebra, def:mpu_three_form_subspace}
    Let the first-block tensor be an MPU of bond dimension $D>0$.  Assume its
    normalized double-layer diagonal $E$ is a rank-one idempotent.  After a
    second blocking of exactly $D^2$ sites, every residual coefficient belongs
    to $\mathcal T$.

    The conclusion is membership in the sum of the three subspaces.  It does
    not assign one exclusive type to a coefficient in an arbitrary physical
    basis.
\end{theorem}

\begin{proof}\leanok
    \uses{thm:mpu_controlled_mixed_residual_trace, thm:mpu_residual_product_d_sq, thm:mpu_ordered_second_block_expansion}
    First extend $EPE=0$ from nonempty residual words $P$ to every element of
    $\mathfrak S$ by writing the generated algebra as the linear span of the
    multiplicative semigroup of residual generators.  For a blocked matrix
    unit, each local double-layer entry is
    \begin{align}
        \delta_{ij}E+R_{ij}.
    \end{align}
    The ordered expansion theorem applies with $N=D^2$.  Its all-$E$
    coefficient is the blocked Kronecker delta and is exactly removed by the
    residual slice.  Hence every blocked matrix-unit residual lies in
    $\mathcal T$.  Expanding an arbitrary physical matrix in matrix units and
    using linearity proves the result.
\end{proof}

\begin{theorem}[Three-form contraction identities]
    \label{thm:mpu_three_form_contractions}
    \lean{MPOTensor.mul_eq_mul_projector_mul_of_mem_threeFormSubmodule,
        MPOTensor.trace_mul_projector_eq_zero_of_mem_threeFormSubmodule}
    \leanok
    \uses{def:mpu_three_form_subspace}
    Suppose $E^2=E$ and $EAE=0$ for every $A\in\mathfrak S$.  For all
    $X,Y\in\mathcal T$,
    \begin{align}
        XY & =XEY,
           & \tr(XE) & =0.
    \end{align}
\end{theorem}

\begin{proof}\leanok
    It is enough by bilinearity and linearity of trace to consider the
    generators $EA$, $AE$, and $AEB$ of the three summands.  Representative
    products reduce as
    \begin{align}
        (EA)(BE)    & =E(AB)E=0,      \\
        (AE)(EB)    & =AEB=(AE)E(EB), \\
        \tr((AEB)E) & =\tr(A(EBE))=0.
    \end{align}
    The remaining generator pairs follow by the same use of associativity,
    $E^2=E$, and $EAE=0$ for $A\in\mathfrak S$.
\end{proof}

\begin{theorem}[Exact simplicity after the second blocking]
    \label{thm:mpu_second_block_simple}
    \lean{MPOTensor.IsMPU.blockTensor_sq_simple_contractions_of_normalizedDiagonal_eq_vecMulVec,
        MPOTensor.IsMPU.blockTensor_sq_isMPUSimple_of_normalizedDiagonal_eq_vecMulVec}
    \leanok
    \uses{thm:mpu_second_block_residual_three_forms, thm:mpu_three_form_contractions}
    Let $\mathcal V$ be an MPU of positive physical and bond dimensions whose
    normalized double-layer diagonal is
    \begin{align}
        E=|\rho)(\Phi|,
         &  & (\Phi|\rho)=1.
    \end{align}
    Then blocking $\mathcal V$ by exactly $D^2$ sites produces a simple tensor
    with witnesses $a=\Phi$ and $b=\rho$.
\end{theorem}

\begin{proof}\leanok
    The pairing normalization makes $E$ idempotent.  Write every blocked
    double-layer coefficient as $\delta_{ij}E+R_{ij}$, where
    $R_{ij}\in\mathcal T$.  The trace identity for $R_{ij}E$ gives the
    diagonal contraction, and
    $R_{ij}R_{k\ell}=R_{ij}ER_{k\ell}$, together with $E^2=E$, gives the
    rank-one insertion identity after expanding both coefficients.
\end{proof}

\begin{definition}[Physical-basis reindexing]
    \label{def:mpu_physical_reindexing}
    \lean{MPOTensor.reindexPhysical,
        MPOTensor.reindexPhysicalConfigEquiv}
    \leanok
    For a bijection $e$ between two physical index sets, the reindexed tensor
    and its action on length-$N$ configurations are
    \begin{align}
        (e^*\mathcal U)^{ij} & =\mathcal U^{e(i)e(j)},
                             & (e_*\sigma)(n)          & =e(\sigma(n)).
    \end{align}
\end{definition}

\begin{theorem}[Words under physical reindexing]
    \label{thm:mpu_eval_word_physical_reindexing}
    \lean{MPOTensor.evalWord_reindexPhysical}
    \leanok
    \uses{def:mpu_physical_reindexing, def:mpo_eval_word}
    For ket and bra words $i,j$, physical reindexing gives
    \begin{align}
        (e^*\mathcal U)[i,j] & =\mathcal U[e_*i,e_*j].
    \end{align}
\end{theorem}

\begin{proof}\leanok
    \uses{def:mpu_physical_reindexing, def:mpo_eval_word}
    Induct simultaneously along the two words. The empty pair gives the
    identity matrix, mismatched lengths give zero, and the nonempty case
    applies the defining equality to the first letters and the induction
    hypothesis to the tails.
\end{proof}

\begin{theorem}[Finite-size operators under physical reindexing]
    \label{thm:mpu_finite_size_physical_reindexing}
    \lean{MPOTensor.mpo_reindexPhysical}
    \leanok
    \uses{def:mpu_physical_reindexing, def:mpo_family}
    At every length $N$, the sitewise configuration bijection identifies
    \begin{align}
        (e^*\mathcal U)^{(N)} & \cong\mathcal U^{(N)}.
    \end{align}
\end{theorem}

\begin{proof}\leanok
    \uses{def:mpu_physical_reindexing, def:mpo_family, thm:mpu_eval_word_physical_reindexing}
    Apply the word identity to every pair of configurations and take the
    virtual trace. This is simultaneous row and column reindexing by $e_*$.
\end{proof}

\begin{theorem}[Double layer under physical reindexing]
    \label{thm:mpu_double_layer_physical_reindexing}
    \lean{MPOTensor.doubleLayerTensor_reindexPhysical}
    \leanok
    \uses{def:mpu_physical_reindexing, def:mpu_double_layer}
    For every physical-index bijection $e$, taking the double layer commutes
    with physical reindexing:
    \begin{align}
        W(e^*\mathcal U)^{ij} & =W(\mathcal U)^{e(i)e(j)}.
    \end{align}
\end{theorem}

\begin{proof}\leanok
    \uses{def:mpu_physical_reindexing, def:mpu_double_layer}
    This follows entrywise from the definitions of physical reindexing and the
    double layer.
\end{proof}

\begin{lemma}[Iterated physical blocking]
    \label{lem:mpu_iterated_physical_blocking}
    \lean{MPOTensor.reindexPhysical_blockTensor_blockTensor}
    \leanok
    \uses{def:mpu_physical_reindexing, def:mpdo_physical_blocking}
    Let $e$ be the canonical bijection which groups a word of length $mn$ into
    $n$ consecutive words of length $m$. Iterating blocking at lengths $m$ and
    $n$, then applying this physical reindexing, agrees with direct blocking at
    length $mn$:
    \begin{align}
        e^*((\mathcal U_m)_n) & =\mathcal U_{mn}.
    \end{align}
\end{lemma}

\begin{proof}\leanok
    \uses{def:mpu_physical_reindexing, def:mpdo_physical_blocking}
    Expand the two blocked words. The canonical regrouping concatenates their
    $n$ words of length $m$ into the direct word of length $mn$.
\end{proof}

\begin{theorem}[Bounded blocking to simplicity]
    \label{thm:mpu_bounded_simple_blocking}
    \lean{MPOTensor.IsMPU.exists_blockTensor_isMPUSimple_of_one_lt,
        MPOTensor.IsMPU.blockTensor_one_isMPUSimple_fin_one,
        MPOTensor.IsMPU.exists_blockTensor_isMPUSimple,
        MPOTensor.IsMPUSimple.blockTensor,
        MPOTensor.IsMPUSimple.reindexPhysical}
    \leanok
    \uses{def:mpu_physical_reindexing, thm:mpu_normalized_transfer_stabilization, thm:mpu_normalized_transfer_fin_one, thm:mpu_second_block_simple}
    Every MPU of positive physical and bond dimensions has a positive blocking
    length $k\leq D^4$ for which the blocked tensor is simple.  If $D>1$, one
    may choose
    \begin{align}
        k=JD^2<D^4,
         &  & 0<J\leq D^2-1.
    \end{align}
    If $D=1$, the length-one block is simple.  Simplicity is preserved under
    every further iterated blocking by a positive factor and under a bijective
    relabeling of the physical basis.
\end{theorem}

\begin{proof}\leanok
    \uses{lem:mpu_iterated_physical_blocking, thm:mpu_normalized_transfer_stabilization, thm:mpu_normalized_transfer_fin_one, thm:mpu_second_block_simple, def:mpu_physical_reindexing}
    At the stabilization exponent $J$, the normalized double-layer diagonal
    is the rank-one fixed projector $|\rho)(\Phi|$.  The preceding theorem
    makes the subsequent length-$D^2$ block simple.  If $e$ is the canonical
    grouping bijection, the two blocking stages satisfy
    \begin{align}
        e^*((\mathcal U_J)_{D^2})=\mathcal U_{JD^2}.
    \end{align}
    For $D>1$,
    \begin{align}
        JD^2 & <(D^2)D^2=D^4.
    \end{align}
    The one-dimensional transfer matrix is already the identity.  For a
    further block with letters $W_1,\ldots,W_L$ and
    $R=|\rho)(\Phi|$, repeated insertion gives
    \begin{align}
        W_1\cdots W_L & =W_1R W_2R\cdots R W_L,          \\
        (\Phi|W_1\cdots W_L|\rho)
                      & =\prod_{q=1}^{L}(\Phi|W_q|\rho).
    \end{align}
    The second identity multiplies the sitewise Kronecker deltas.
\end{proof}

\begin{theorem}[Matching contractions at the direct stabilized block]
    \label{thm:mpu_matching_direct_block_contractions}
    \lean{MPOTensor.IsMPU.blockTensor_mul_sq_simple_contractions_of_transfer_power}
    \leanok
    Let $\mathcal U$ be an MPU of positive physical and bond dimensions.
    Suppose that $\rho$ has trace one and that, for some positive integer $J$,
    the normalized transfer matrix satisfies
    \begin{align}
        E^J & =|\rho)(\Id|.
    \end{align}
    Put $K=JD^2$. After the standard product reindexing of the physical
    coordinates, the direct length-$K$ block satisfies the contractions
    (\ref{eq:mpu_simple1}) and (\ref{eq:mpu_simple2}) with the specific witnesses
    $\rho$ and $\Id$ from this transfer-power identity.
\end{theorem}

\begin{proof}\leanok
    \uses{thm:mpu_blocked_double_layer_diagonal, thm:mpu_second_block_simple, thm:mpu_bounded_simple_blocking, lem:mpu_blocking, thm:mpu_double_layer_physical_reindexing}
    First block $J$ sites.  The normalized diagonal of this block is the
    reindexed rank-one matrix $|\rho)(\Id|$, and the trace-one condition is
    exactly $(\Id|\rho)=1$.  The exact second-block contraction theorem then
    supplies both identities after a further $D^2$ sites.  The canonical
    equivalence between iterated blocking and direct length-$JD^2$ blocking
    transports each double-layer coefficient without changing either virtual
    witness.
\end{proof}

\begin{theorem}[Rank-one transfer power after the second blocking]
    \label{thm:mpu_blocked_transfer_power_rank_one}
    \lean{MPOTensor.normalizedDiagonal_blockTensor_mul_sq_eq_vecMulVec_of_transfer_power}
    \leanok
    \uses{def:mpu_double_layer, def:mpu_normalized_flattening, def:mpdo_physical_blocking}
    Let $\mathcal U$ be an MPO tensor with $d,D>0$, let $E$ be the transfer
    matrix of its normalized flattening, and suppose that for some $J\geq0$
    and some trace-one $\rho$,
    \begin{align}
        E^J & =|\rho)(\Id|.
    \end{align}
    Then the normalized diagonal of the double layer of the direct
    length-$JD^2$ block is the same rank-one matrix, after the standard product
    reindexing of the doubled bond.
\end{theorem}

\begin{proof}\leanok
    \uses{thm:mpu_blocked_double_layer_diagonal, thm:mpu_double_layer_blocking}
    Trace normalization gives $(\Id|\rho)=1$, so $|\rho)(\Id|$ is idempotent
    and every positive power of it equals itself.  The blocked normalized
    diagonal is the corresponding power of $E$, and $E^{JD^2}$ is the $D^2$-th
    power of $E^J$.
\end{proof}

\begin{theorem}[Normalized input-tail isometry]
    \label{thm:mpu_normalized_input_tail_isometry}
    \lean{MPOTensor.IsMPU.normalized_mpo_tail_isometry}
    \leanok
    \uses{def:mpu, def:mpo_family}
    Let $d>0$, let $\mathcal U$ be an MPU, and let $K\geq0$.  For
    $p,q\in\Fin d\times\Fin d$, separate the first two input sites from a
    common input tail $\tau\in\Fin K\to\Fin d$.  Then
    \begin{align}
        d^{-K}\sum_{\tau}\sum_{\eta}
        \overline{U^{(K+2)}_{\eta,(p,\tau)}}
        U^{(K+2)}_{\eta,(q,\tau)}
         & =\delta_{p,q}.
    \end{align}
    The global contraction is the input-first equation
    $(U^{(K+2)})^\dagger U^{(K+2)}=\Id$.  This is the input-oriented tail
    contraction supporting the network calculation in
    \cite[Theorem~III.8, Section~III.B, equations~(31)--(32)]{Cirac2017MPU}.
\end{theorem}

\begin{proof}\leanok
    \uses{lem:mpu_adjoint_mul}
    Apply input-first unitarity for each fixed tail.  The inner output sum is
    $\delta_{p,q}$.  Summing over the $d^K$ tails and multiplying by $d^{-K}$
    leaves the same Kronecker delta.
\end{proof}

% Formalization scope: identifying this normalized scalar with
% $(\texttt{sourceU})^\dagger\texttt{sourceU}$ is treated separately.


\begin{theorem}[Simplicity at every later direct blocking]
    \label{thm:mpu_all_later_simple_blockings}
    \lean{MPOTensor.IsMPU.isMPUSimple_of_simple2,
        MPOTensor.IsMPU.blockTensor_succ_isMPUSimple,
        MPOTensor.IsMPU.blockTensor_isMPUSimple_of_le,
        MPOTensor.blockTensor_succ_simple2_of_supplied,
        MPOTensor.IsMPU.simple1_of_simple2_supplied}
    \leanok
    \uses{def:mpu, def:mpu_simple}
    Let $\mathcal U$ generate an MPU, and suppose its direct length-$k$ block
    is simple for some $k>0$.  Then its direct length-$k'$ block is simple for
    every $k'\geq k$, without any divisibility condition on $k'$ by $k$.
    This is the corollary following
    \cite[Proposition~III.3]{Cirac2017MPU}.
\end{theorem}

\begin{proof}\leanok
    \uses{lem:mpu_blocking, lem:mpu_adjoint_mul, thm:mpu_closed_double_layer, thm:mpu_double_layer_blocking}
    First suppose vectors $a,b$ satisfy the rank-one insertion identity
    \begin{align}
        W^{ij}W^{k\ell} & =W^{ij}|b)(a|W^{k\ell}.
    \end{align}
    Set $x_{ij}=(a|W^{ij}|b)$.  Applying this identity around the constant
    two-site and three-site words, and then using MPU unitarity, gives
    \begin{align}
        x_{ij}^2 & =\delta_{ij}, \\
        x_{ij}^3 & =\delta_{ij}.
    \end{align}
    Hence $x_{ii}=1$, while $x_{ij}=0$ for $i\ne j$, so the same vectors obey
    the one-letter contraction $(a|W^{ij}|b)=\delta_{ij}$.

    For the successor step, take four words $I,J,K,L$ of length $k+1$ and
    write $W^{I,J}=W^{i_0j_0}\cdots W^{i_kj_k}$, with the analogous notation
    for $W^{K,L}$.  Apply the length-$k$ rank-one insertion identity to the suffixes
    $I_+=(i_1,\ldots,i_k)$, $J_+=(j_1,\ldots,j_k)$ and the prefixes
    $K_-=(k_0,\ldots,k_{k-1})$, $L_-=(\ell_0,\ldots,\ell_{k-1})$.  Then
    \begin{align}
        W^{I,J}W^{K,L}
         & =W^{i_0j_0}\bigl(W^{I_+,J_+}W^{K_-,L_-}\bigr)W^{k_k\ell_k} \\
         & =W^{I,J}|b)(a|W^{K,L}.
    \end{align}
    Thus the length-$(k+1)$ block has the same insertion witnesses, and the
    two- and three-site argument gives their one-letter contraction.  Induction
    on $k'-k$ proves the claim.
\end{proof}

\begin{theorem}[Common simple blocking along a continuous MPU path]
    \label{thm:mpu_common_pow_four_simple_blocking}
    \lean{MPOTensor.IsMPU.blockTensor_pow_four_isMPUSimple,
        MPOTensor.continuous_evalWord,
        MPOTensor.continuous_blockTensor,
        Continuous.blockTensor}
    \leanok
    \uses{def:mpu, def:mpu_simple, def:mpdo_physical_blocking}
    Let $d,D>0$.  Every bond-$D$ MPU tensor becomes simple after blocking
    exactly $D^4$ sites.  For each fixed $L$, the map
    \begin{align}
        \mathcal U & \longmapsto \mathcal U_L
    \end{align}
    is continuous.  Consequently, if $x\mapsto\mathcal W(x)$ is a continuous
    family of bond-$D$ MPU tensors, then
    $x\mapsto\mathcal W(x)_{D^4}$ is a continuous family of simple tensors.
    This is only the common-blocking step in the proof of
    Proposition~\ref{prop:continuity-index}; it does not choose reduced
    representatives or assert constancy of their source ranks.
    % Source lines: paper_v2.tex 747--752.
\end{theorem}

\begin{proof}\leanok
    \uses{thm:mpu_bounded_simple_blocking, thm:mpu_all_later_simple_blockings}
    Choose for each tensor a positive simple blocking length $k\leq D^4$.
    Simplicity persists at every later direct blocking, hence at $D^4$.
    For fixed blocked ket and bra words, every matrix entry of the blocked
    tensor is a finite sum of products of entries of the original tensor.
    These entries are continuous, and the finite product and sum operations
    preserve continuity.  Applying this fixed blocking map to the original
    path proves the path statement.
\end{proof}

\begin{theorem}[Supplied-projector simple2 contraction]
    \label{thm:mpu_supplied_projector_simple2}
    \lean{MPOTensor.IsMPUSimple.simple2_of_normalizedDiagonal_pow_eq_vecMulVec}
    \leanok
    \uses{def:mpu_simple}
    Let $\mathcal U$ be simple with positive physical dimension, let $W$ be
    its double layer, and let $E$ be the normalized diagonal of $W$.  Suppose
    that for some $J>0$,
    \begin{align}
        E^J & =|\rho)(\Phi|.
    \end{align}
    Then the simple insertion identity holds with these same supplied
    witnesses:
    \begin{align}
        W^{ij}W^{k\ell}
         & =W^{ij}|\rho)(\Phi|W^{k\ell}.
    \end{align}
\end{theorem}

\begin{proof}\leanok
    \uses{def:mpu_simple}
    Choose the witnesses $a,b$ from simplicity and put $R=|b)(a|$.  Expanding
    the normalized diagonal as the average of the diagonal letters, the two
    defining simple contractions give
    \begin{align}
        W^{ij} E W^{k\ell} & =W^{ij}W^{k\ell}.
    \end{align}
    Indeed, insert $R$ on both sides of each diagonal letter; its middle
    sandwich is $R W^{qq}R=R$ by the first contraction, and the normalized sum
    contains $d$ equal terms.  The same identity makes
    $W^{ij}E^N=W^{ij}E$ for every $N>0$.  Hence
    \begin{align}
        W^{ij}E^J W^{k\ell} & =W^{ij}W^{k\ell}.
    \end{align}
    Substituting the supplied equality $E^J=|\rho)(\Phi|$ proves the claim.
    No uniqueness of rank-one factorizations is used.
\end{proof}

\begin{theorem}[Fixed boundaries of a normalized rank-one power]
    \label{thm:mpu_normalized_diagonal_power_fixed_pair}
    \lean{MPOTensor.normalizedDiagonal_pow_fixed_pair_of_eq_vecMulVec}
    \leanok
    Let $R$ be the normalized physical diagonal of the double layer of an MPO
    tensor.  Suppose
    \begin{align}
        (\Phi\rvert\rho) & =1,
                         & R^J & =\lvert\rho)(\Phi\rvert.
    \end{align}
    Then the two boundaries of the stabilized diagonal tail are fixed:
    \begin{align}
        R^J\lvert\rho) & =\lvert\rho),
                       & (\Phi\rvert R^J & =(\Phi\rvert.
    \end{align}
    This is the fixed-pair normalization from the canonical-form-II
    fixed-point equations in
    \cite[equations~(6a)--(6b)]{Cirac2017MPU} and underlies the tail
    contractions in
    \cite[Theorem~III.8, Section~III.B, equations~(31)--(32)]{Cirac2017MPU}.
\end{theorem}

\begin{proof}\leanok
    Substitute $R^J=\lvert\rho)(\Phi\rvert$ in each expression and use
    $(\Phi\rvert\rho)=1$.
\end{proof}

\begin{theorem}[Rank-one insertion as a normalized diagonal tail]
    \label{thm:mpu_rank_one_insertion_diagonal_tail}
    \lean{MPOTensor.doubleLayerTensor_rankOne_mul_eq_normalizedDiagonal_tail}
    \leanok
    \uses{def:mpu_double_layer, thm:mpu_normalized_diagonal_word_expansion}
    Let $W$ be the double layer of an MPO tensor, and let $R$ be its
    normalized physical diagonal.  If
    \begin{align}
        R^J & =\lvert\rho)(\Phi\rvert,
    \end{align}
    then, for all physical indices $i,j,k,\ell$,
    \begin{align}
        W^{ij}\lvert\rho)(\Phi\rvert W^{k\ell}
        =d^{-J}\sum_{\tau\in\{0,\ldots,d-1\}^J}
        W^{ij}W^{\tau_0\tau_0}\cdots
        W^{\tau_{J-1}\tau_{J-1}}W^{k\ell}.
    \end{align}
    This is the exact diagonal-word tail between the two retained letters in
    \cite[Theorem~III.8, Section~III.B, equations~(31)--(32)]{Cirac2017MPU}.
\end{theorem}

\begin{proof}\leanok
    \uses{thm:mpu_normalized_diagonal_word_expansion}
    Replace the rank-one insertion by $R^J$, apply the diagonal-word
    expansion, and distribute the two retained letters across the normalized
    finite sum.
\end{proof}

\begin{theorem}[Simple two-letter product as a normalized diagonal tail]
    \label{thm:mpu_simple_two_letter_diagonal_tail}
    \lean{MPOTensor.doubleLayerTensor_mul_eq_normalizedDiagonal_tail_of_simple2}
    \leanok
    \uses{def:mpu_double_layer, thm:mpu_rank_one_insertion_diagonal_tail}
    In the setting of the preceding theorem, suppose in addition that the
    same fixed pair satisfies the supplied two-letter contraction
    \begin{align}
        W^{ij}W^{k\ell}
         & =W^{ij}\lvert\rho)(\Phi\rvert W^{k\ell}
    \end{align}
    for all physical indices.  Then
    \begin{align}
        W^{ij}W^{k\ell}
        =d^{-J}\sum_{\tau\in\{0,\ldots,d-1\}^J}
        W^{ij}W^{\tau_0\tau_0}\cdots
        W^{\tau_{J-1}\tau_{J-1}}W^{k\ell}.
    \end{align}
    Thus the supplied insertion and its stabilized transfer power occur in
    one aligned network, as required in the calculation of
    \cite[Theorem~III.8, Section~III.B, equations~(31)--(32)]{Cirac2017MPU}.
\end{theorem}

\begin{proof}\leanok
    \uses{thm:mpu_rank_one_insertion_diagonal_tail}
    Apply the supplied two-letter contraction and then replace its rank-one
    insertion by the normalized diagonal-word tail.
\end{proof}

\begin{theorem}[The fixed-point simple1 contraction]
    \label{thm:mpu_simple1}
    \lean{MPOTensor.IsMPU.exists_normalized_fixed_points_simple1,
        MPOTensor.IsMPU.simple1}
    \leanok
    \uses{def:mpu, thm:mpu_positive_mixed_residual_trace, thm:mpu_normalized_transfer_stabilization, thm:mpu_normalized_transfer_fin_one}
    For every MPU of positive physical and bond dimensions there are normalized
    left and right fixed vectors $\Phi$ and $\rho$ of the transfer matrix such
    that the unblocked double layer satisfies
    \begin{align}
        (\Phi|W^{ij}|\rho) & =\delta_{ij}.
    \end{align}
    The unblocked tensor need not satisfy the rank-one insertion identity.
\end{theorem}

\begin{proof}\leanok
    Let $E^J=|\rho)(\Phi|$ be a stabilized transfer power.  Decompose
    $W^{ij}=\delta_{ij}E+R_{ij}$.  Fixedness and normalization give
    \begin{align}
        \tr(E|\rho)(\Phi|) & =1.
    \end{align}
    The mixed residual trace identity applied to one residual factor followed
    by $J$ diagonal factors gives
    \begin{align}
        \tr(R_{ij}E^J) & =0.
    \end{align}
    Since $E^J=|\rho)(\Phi|$, these two trace identities are exactly the
    displayed contraction.
\end{proof}

\section{Trace-power lemma}\label{sec:mpu_trace_power}

For the transfer-matrix spectrum argument of
\cite[Proposition~1]{Cirac2017MPU} we first record the
algebraic fact that a matrix whose positive powers all have trace~$1$ must have
characteristic polynomial $X^{n-1}(X-1)$.  The source assumes only the powers
with exponent strictly greater than one; this weaker hypothesis determines the
nonzero spectrum as a set.

\begin{definition}[Rank-one diagonal idempotent]
    \label{def:mpu_diagonal_one_at}
    \lean{Matrix.diagonalOneAt}
    \leanok
    For an index $i_0$ in a finite set $I$ of cardinality $n$, let $M_{i_0}$ be the $n\times n$
    diagonal matrix with a $1$ at position $i_0$ and zeros elsewhere.
\end{definition}

\begin{lemma}[Properties of the rank-one diagonal idempotent]
    \label{lem:mpu_diagonal_one_at}
    \lean{Matrix.trace_diagonalOneAt,
        Matrix.diagonalOneAt_pow, Matrix.trace_pow_diagonalOneAt,
        Matrix.charpoly_diagonalOneAt}
    \leanok
    \uses{def:mpu_diagonal_one_at}
    The matrix $M_{i_0}$ from definition~\ref{def:mpu_diagonal_one_at} satisfies
    \begin{itemize}
        \item $M_{i_0}^k = M_{i_0}$ for every $k\ge 1$,
        \item $\tr(M_{i_0}^k) = 1$ for every $k\ge 1$,
        \item $\chi_{M_{i_0}}(X) = X^{\,n-1}(X-1)$.
    \end{itemize}
\end{lemma}

\begin{proof}\leanok
    The first two items follow from $1^k=1$ and $0^k=0$ for $k\ge 1$.
    The characteristic polynomial is computed from
    $\chi_{\diag(d)}(X) = \prod_i (X - d_i)$.
\end{proof}

\begin{theorem}[Finite-range trace-power lemma]
    \label{thm:mpu_trace_power_le_card}
    \lean{Matrix.charpoly_eq_X_pow_pred_mul_X_sub_one_of_trace_pow_eq_one_of_le_card}
    \leanok
    Let $A$ be an $n\times n$ complex matrix with $n\ge 1$.  If
    $\tr(A^k)=1$ for every $k$ with $1\le k\le n$, then
    \begin{align}
        \chi_A(X) & = X^{\,n-1}(X-1).
    \end{align}
\end{theorem}

\begin{proof}\leanok
    \uses{lem:mpu_diagonal_one_at}
    Choose any index $i_0$ (the hypothesis $n\ge1$ makes this possible)
    and let $B$ be the rank-one diagonal idempotent from
    Lemma~\ref{lem:mpu_diagonal_one_at}.  For $1\le k\le n$ we have
    $\tr(A^k)=1=\tr(B^k)$.  The finite-range half of the Newton--Girard
    the corresponding theorem in the companion quantum-channel volume~\cite[``Newton--Girard trace recursion'']{QICLean2026Blueprint} then yields $\chi_A = \chi_B =
        X^{\,n-1}(X-1)$.
\end{proof}

\begin{theorem}[All-positive-powers trace-power lemma]
    \label{thm:mpu_trace_power_all}
    \lean{Matrix.charpoly_eq_X_pow_pred_mul_X_sub_one_of_forall_trace_pow_eq_one}
    \leanok
    Let $A$ be an $n\times n$ complex matrix.  If $\tr(A^k)=1$ for
    every $k\ge 1$, then
    \begin{align}
        \chi_A(X) & = X^{\,n-1}(X-1).
    \end{align}
    (The hypothesis already forces $n\ge 1$ because the trace of a
    $0\times0$ matrix is $0$, contradicting $\tr(A)=1$.)
\end{theorem}

\begin{proof}\leanok
    \uses{thm:mpu_trace_power_le_card}
    If $n=0$ the hypothesis $\tr(A)=1$ contradicts the fact that the
    trace of an empty matrix is $0$; the conclusion follows vacuously.
    When $n\ge1$, the hypothesis on all positive powers implies the
    finite-range hypothesis of Theorem~\ref{thm:mpu_trace_power_le_card}.
\end{proof}

\begin{theorem}[Shifted moments under matrix powers]
    \label{thm:mpu_shifted_trace_power_moments}
    \lean{Matrix.forall_trace_pow_pow_eq_one_of_forall_trace_pow_eq_one_of_one_lt}
    \leanok
    Let $A$ be an $n\times n$ complex matrix such that $\tr(A^k)=1$ for every
    $k>1$, and let $m>1$.  Then $\tr((A^m)^r)=1$ for every $r\ge1$.
\end{theorem}

\begin{proof}\leanok
    For $r\ge1$, one has $(A^m)^r=A^{mr}$ and $mr>1$, so the hypothesis at
    exponent $mr$ gives the result.
\end{proof}

\begin{theorem}[Trace powers above one determine the nonzero spectrum]
    \label{thm:mpu_shifted_trace_power_spectrum}
    \lean{Matrix.spectrum_diff_zero_eq_singleton_of_forall_trace_pow_eq_one_of_one_lt,
        Matrix.eq_one_of_mem_spectrum_of_forall_trace_pow_eq_one_of_one_lt,
        Matrix.one_mem_spectrum_of_forall_trace_pow_eq_one_of_one_lt}
    \leanok
    Let $A$ be an $n\times n$ complex matrix. If $\tr(A^k)=1$ for every
    integer $k>1$, then
    \begin{align}
        \sigma(A)\setminus\{0\} & =\{1\}.
    \end{align}
    Thus $1$ occurs and every nonzero spectral value equals $1$. This is the
    set-spectrum consequence of \cite[Proposition~III.1]{Cirac2017MPU}; set
    equality alone does not assert algebraic multiplicity.
\end{theorem}

\begin{proof}\leanok
    \uses{thm:mpu_shifted_trace_power_moments, thm:mpu_trace_power_all}
    For every positive integer $r$, the shifted moments give
    \begin{align}
        \tr((A^2)^r) & =\tr(A^{2r})=1,
                     & \tr((A^3)^r)    & =\tr(A^{3r})=1.
    \end{align}
    Hence
    \begin{align}
        \chi_{A^2}(X) & =\chi_{A^3}(X)=X^{n-1}(X-1).
    \end{align}
    If $\mu\in\sigma(A)$ is nonzero, spectral mapping places $\mu^2$ in
    $\sigma(A^2)$ and $\mu^3$ in $\sigma(A^3)$, so
    \begin{align}
        \mu^2 & =\mu^3=1,
              & \mu       & =1.
    \end{align}
    Conversely, spectral mapping from $1\in\sigma(A^2)$ supplies a nonzero
    spectral value of $A$, which must be $1$.
\end{proof}

\begin{theorem}[Generalized eigenspace inclusion under squaring]
    \label{thm:mpu_generalized_eigenspace_one_square}
    \lean{Matrix.maxGenEigenspace_one_le_sq}
    \leanok
    For a complex square matrix $A$, write $G_1(A)$ for its maximal generalized
    eigenspace at the eigenvalue $1$. Then
    \begin{align}
        G_1(A) & \subseteq G_1(A^2).
    \end{align}
\end{theorem}

\begin{proof}\leanok
    The commuting factorization
    \begin{align}
        A^2-\Id & =(A+\Id)(A-\Id)
    \end{align}
    gives, for every nonnegative integer $k$,
    \begin{align}
        (A^2-\Id)^k & =(A+\Id)^k(A-\Id)^k.
    \end{align}
    Therefore every vector annihilated by a power of $A-\Id$ is annihilated by
    the corresponding power of $A^2-\Id$.
\end{proof}

\begin{theorem}[Multiplicity at one does not decrease under squaring]
    \label{thm:mpu_root_multiplicity_one_square}
    \lean{Matrix.rootMultiplicity_one_charpoly_le_sq}
    \leanok
    For every complex square matrix $A$,
    \begin{align}
        \operatorname{mult}_{1}(\chi_A)
         & \le \operatorname{mult}_{1}(\chi_{A^2}).
    \end{align}
\end{theorem}

\begin{proof}\leanok
    \uses{thm:mpu_generalized_eigenspace_one_square}
    With $G_1$ as above, algebraic multiplicity equals the dimension of the
    maximal generalized eigenspace. Hence
    \begin{align}
        \operatorname{mult}_{1}(\chi_A)
         & =\dim G_1(A)     \\
         & \le\dim G_1(A^2)
        =\operatorname{mult}_{1}(\chi_{A^2}).
    \end{align}
\end{proof}

\begin{theorem}[Trace powers above one determine the characteristic polynomial]
    \label{thm:mpu_shifted_trace_power_charpoly}
    \lean{Matrix.charpoly_rootMultiplicity_one_eq_one_of_forall_trace_pow_eq_one_of_one_lt,
        Matrix.charpoly_eq_X_pow_pred_mul_X_sub_one_of_forall_trace_pow_eq_one_of_one_lt}
    \leanok
    Let $A$ be an $n\times n$ complex matrix. If $\tr(A^k)=1$ for every
    integer $k>1$, then
    \begin{align}
        \operatorname{mult}_{1}(\chi_A) & =1,
                                        & \chi_A(X) & =X^{n-1}(X-1).
    \end{align}
    Thus $1$ is the sole nonzero eigenvalue and has algebraic multiplicity one.
    This is the exact shifted-moment conclusion of
    \cite[Proposition~III.1]{Cirac2017MPU}.
\end{theorem}

\begin{proof}\leanok
    \uses{thm:mpu_shifted_trace_power_moments, thm:mpu_trace_power_all, thm:mpu_shifted_trace_power_spectrum, thm:mpu_root_multiplicity_one_square}
    Applying the all-positive trace-power theorem to $A^2$ gives
    \begin{align}
        \chi_{A^2}(X) & =X^{n-1}(X-1),
                      & \operatorname{mult}_{1}(\chi_{A^2}) & =1.
    \end{align}
    The multiplicity inequality therefore bounds
    $\operatorname{mult}_{1}(\chi_A)$ by one. The set-spectrum theorem shows
    that $1$ occurs, so this multiplicity is positive and hence equals one.
    The same set-spectrum theorem excludes every other nonzero root. Splitting
    the monic characteristic polynomial now yields
    \begin{align}
        \operatorname{mult}_{1}(\chi_A) & =1,
                                        & \chi_A(X) & =X^{n-1}(X-1).
    \end{align}
\end{proof}

\begin{theorem}[Semisimple rank-one Jordan--Chevalley reduction]
    \label{thm:mpu_semisimple_scaled_trace_rank_one}
    \lean{Module.End.IsNilpotent.trace_add_pow_eq_trace_pow_of_commute,
        Module.End.finrank_range_eq_one_of_isSemisimple_of_scaled_trace_pow_eq_one,
        Module.End.IsNilpotent.mul_eq_zero_and_mul_eq_zero_of_commute_of_finrank_range_eq_one,
        Module.End.IsNilpotent.exists_add_pow_eq_right_of_mul_eq_zero,
        Module.End.IsNilpotent.products_eq_zero_and_eventual_add_pow_eq_of_scaled_trace}
    \leanok
    Let $V$ be a finite-dimensional complex vector space and let $S$ be a
    semisimple endomorphism of $V$. Suppose that $c\in\mathbb C$ is nonzero and
    \begin{align}
        \tr\bigl((cS)^k\bigr) & =1
    \end{align}
    for every integer $k>1$. Then
    \begin{align}
        \dim\operatorname{range}(S) & =1.
    \end{align}
    If $N$ is nilpotent and commutes with $S$, then, for every positive integer
    $k$,
    \begin{align}
        \tr\bigl((N+S)^k\bigr) & =\tr(S^k), \\
        NS                     & =SN=0.
    \end{align}
    Moreover, there is an integer $J>0$ such that, for every $k\geq J$,
    \begin{align}
        (N+S)^k & =S^k.
    \end{align}
    This is the generic Jordan--Chevalley step in arXiv:1706.07329v2,
    Proposition~20, lines 3839--3866.
\end{theorem}

\begin{proof}\leanok
    \uses{thm:mpu_shifted_trace_power_charpoly}
    Choose an arbitrary basis of $V$ and represent $cS$ by a matrix $A$.
    The shifted trace-power theorem gives
    \begin{align}
        \chi_A(X) & =X^{\dim V-1}(X-1).
    \end{align}
    Since $cS$ is semisimple, its generalized zero-eigenspace is its ordinary
    zero-eigenspace, namely its kernel. The dimension of the generalized
    zero-eigenspace is the root multiplicity of zero in the characteristic
    polynomial, so
    \begin{align}
        \dim\ker(cS) & =\dim V-1.
    \end{align}
    Rank--nullity and $c\ne0$ therefore give
    $\dim\operatorname{range}(S)=1$.

    For $k>0$, commutation gives the binomial expansion
    \begin{align}
        (N+S)^k & =\sum_{m=0}^k \binom{k}{m}N^mS^{k-m}.
    \end{align}
    If $m>0$, write $N^m=N N^{m-1}$. The first factor commutes with the
    remaining factors, so nilpotence of $N$ makes the entire summand
    nilpotent and its trace vanishes. Only the term $m=0$ remains, giving
    $\tr((N+S)^k)=\tr(S^k)$. This is the trace reduction used in
    arXiv:1706.07329v2, Proposition~20, line 3865.

    Commutation makes $\operatorname{range}(S)$ invariant under $N$. The
    restriction of $N$ to this one-dimensional space is both nilpotent and a
    scalar endomorphism, hence zero. Therefore $NS=0$, and commutation also
    gives $SN=0$. For every positive $k$, induction using these two identities
    yields
    \begin{align}
        (N+S)^k & =N^k+S^k.
    \end{align}
    Taking $J$ at least the nilpotency index of $N$ removes the first summand.
\end{proof}

\begin{theorem}[Vanishing positive trace powers determine the characteristic polynomial]
    \label{thm:mpu_zero_trace_power_charpoly}
    \lean{Matrix.charpoly_eq_X_pow_card_of_forall_trace_pow_eq_zero}
    \leanok
    Let $A$ be an $n\times n$ complex matrix. If $\tr(A^k)=0$ for every integer
    $k\geq1$, then
    \begin{align}
        \chi_A(X) & =X^n.
    \end{align}
\end{theorem}

\begin{proof}\leanok
    The zero matrix has the same positive trace powers as $A$ and characteristic
    polynomial $X^n$. Newton--Girard therefore identifies the two characteristic
    polynomials. This is the characteristic-polynomial consequence of the
    all-positive trace condition obtained in
    \cite[proof of Proposition~III.3]{Cirac2017MPU}.
\end{proof}

\begin{theorem}[Shifted vanishing moments under matrix powers]
    \label{thm:mpu_shifted_zero_trace_power_moments}
    \lean{Matrix.forall_trace_pow_pow_eq_zero_of_forall_trace_pow_eq_zero_of_one_lt}
    \leanok
    Let $A$ be an $n\times n$ complex matrix such that $\tr(A^k)=0$ for every
    $k>1$, and let $m>1$. Then $\tr((A^m)^r)=0$ for every $r\geq1$.
    This auxiliary result supports the generalization below to a weaker
    hypothesis; the source itself supplies vanishing for every positive exponent.
\end{theorem}

\begin{proof}\leanok
    For $r\geq1$, one has $(A^m)^r=A^{mr}$ and $mr>1$, so the hypothesis at
    exponent $mr$ gives the result.
\end{proof}

\begin{theorem}[Shifted zero trace powers force nilpotence]
    \label{thm:mpu_shifted_zero_trace_power_nilpotent}
    \lean{Matrix.eq_zero_of_mem_spectrum_of_forall_trace_pow_eq_zero_of_one_lt,
        Matrix.charpoly_eq_X_pow_card_of_forall_trace_pow_eq_zero_of_one_lt,
        Matrix.isNilpotent_of_forall_trace_pow_eq_zero_of_one_lt}
    \leanok
    Let $A$ be an $n\times n$ complex matrix. If $\tr(A^N)=0$ for every integer
    $N>1$, then $A$ is nilpotent and
    \begin{align}
        \chi_A(X) & =X^n.
    \end{align}
    The source obtains the all-positive trace hypothesis in
    \cite[proof of Proposition~III.3]{Cirac2017MPU}; its specialization is
    the elementwise nilpotence assertion used before the nil-matrix theorem. The
    result above is a stronger generalization because it assumes vanishing only
    for $N>1$. It is useful independently and requires no positivity, normality,
    or diagonalizability.
\end{theorem}

\begin{proof}\leanok
    \uses{thm:mpu_shifted_zero_trace_power_moments, thm:mpu_zero_trace_power_charpoly}
    Every positive power of $A^2$ has vanishing trace. Hence
    \begin{align}
        \chi_{A^2}(X) & =X^n.
    \end{align}
    If $\mu$ is a spectral value of $A$, spectral mapping shows that $\mu^2$ is
    a root of this characteristic polynomial. Therefore $\mu^2=0$, and hence
    $\mu=0$. Thus every root of the split
    monic polynomial $\chi_A$ is zero, so $\chi_A(X)=X^n$.
    Cayley--Hamilton now proves that $A$ is nilpotent.
\end{proof}

The standard-form construction uses the compact singular-value decomposition in
which the intermediate dimension is the rank rather than the smaller ambient
dimension.

\begin{theorem}[Compact rectangular singular-value decomposition]
    \label{thm:mpu_compact_svd}
    \lean{Matrix.exists_compactSVD}
    \leanok
    Let $M\in\MN{m,n}$ be a complex matrix of rank $r$. There are matrices
    $V\in\MN{r,m}$ and $U\in\MN{r,n}$ and strictly positive real numbers
    $s_0,\ldots,s_{r-1}$ such that, for $D=\diag(s_0,\ldots,s_{r-1})$,
    \begin{align}
        M & =V^\dagger D U, & VV^\dagger & =\Id_r, & UU^\dagger & =\Id_r.
        \label{eq:mpu_compact_svd}
    \end{align}
    Thus no zero singular values occur in the intermediate space. This is the
    compact decomposition used in the standard-form construction of
    \cite[Section~III]{Cirac2017MPU}.
\end{theorem}

\begin{proof}\leanok
    Apply the spectral theorem to $MM^\dagger$. Its nonzero eigenvalues are
    positive, and their number is
    $\operatorname{rank}(MM^\dagger)=\operatorname{rank}(M)=r$. For an
    orthonormal eigenvector $e_j$ with eigenvalue $\lambda_j>0$, put
    $w_j=M^\dagger e_j$. Then
    $\langle w_j,w_k\rangle=\lambda_j\delta_{jk}$. Hence the vectors
    $\lambda_j^{-1/2}w_j$ are orthonormal. Expanding each column of $M$ in the
    eigenbasis and discarding the zero-eigenvalue terms gives
    \begin{align}
        M & =\sum_{\lambda_j>0}
        \sqrt{\lambda_j}\,e_j
        \bigl(\lambda_j^{-1/2}w_j\bigr)^\dagger,
    \end{align}
    which has the form in~(\ref{eq:mpu_compact_svd}).
\end{proof}

We now introduce the two matrices used for the singular-value decompositions of
a Matrix Product Unitary tensor in Section~III of \cite{Cirac2017MPU}.  Write
$\mathcal U^i_{j,\alpha,\beta}$ for the entries of the tensor, where $i,j$ are
physical indices and $\alpha,\beta$ are auxiliary indices.

\begin{definition}[First matrix associated to a Matrix Product Unitary tensor]
    \label{def:mpu_source_matrix_one}
    \lean{MPOTensor.sourceCutM₁, MPOTensor.sourceCutM₁_apply,
        MPOTensor.sourceCutM₁Fin, MPOTensor.sourceCutM₁Fin_eq_reindex,
        MPOTensor.sourceCutM₁Fin_apply,
        MPOTensor.sourceCutM₁_rank_eq_sourceCutM₁Fin_rank}
    \leanok
    \uses{def:mpo_tensor, def:mps_independent_tensor_product}
    % Source-orientation audit: docs/paper-gaps/mpu_source_cut_orientation.tex.
    The matrix $M_1$ is obtained by combining the upper physical index with the
    right auxiliary index, and the left auxiliary index with the lower physical
    index. Thus
    \begin{align}
        (M_1)_{(i,\beta),(\alpha,j)}
         & = \mathcal U^i_{j,\alpha,\beta}.
    \end{align}
    Transporting the row and column indices along the standard product
    enumerations gives a matrix $\widetilde M_1$ indexed by
    $\{0,\ldots,dD-1\}\times\{0,\ldots,Dd-1\}$. Its entries are obtained by
    applying $M_1$ to the inverse images of the enumerated indices, and
    \begin{align}
        \operatorname{rank}(\widetilde M_1)
         & = \operatorname{rank}(M_1).
    \end{align}
\end{definition}

\begin{definition}[Second matrix associated to a Matrix Product Unitary tensor]
    \label{def:mpu_source_matrix_two}
    \lean{MPOTensor.sourceCutM₂, MPOTensor.sourceCutM₂_apply,
        MPOTensor.sourceCutM₂Fin, MPOTensor.sourceCutM₂Fin_eq_reindex,
        MPOTensor.sourceCutM₂Fin_apply,
        MPOTensor.sourceCutM₂_rank_eq_sourceCutM₂Fin_rank}
    \leanok
    \uses{def:mpo_tensor}
    The matrix $M_2$ is obtained by combining the left auxiliary index with the
    upper physical index, and the lower physical index with the right auxiliary
    index. Thus
    \begin{align}
        (M_2)_{(\alpha,i),(j,\beta)}
         & = \mathcal U^i_{j,\alpha,\beta}.
    \end{align}
    Transporting the row and column indices along the standard product
    enumerations gives a matrix $\widetilde M_2$ indexed by
    $\{0,\ldots,Dd-1\}\times\{0,\ldots,dD-1\}$. Its entries are obtained by
    applying $M_2$ to the inverse images of the enumerated indices, and
    \begin{align}
        \operatorname{rank}(\widetilde M_2)
         & = \operatorname{rank}(M_2).
    \end{align}
\end{definition}

\begin{definition}[Right rank]
    \label{def:mpu_right_rank}
    \lean{MPOTensor.rightRank, MPOTensor.rightRank_eq}
    \leanok
    \uses{def:mpu_source_matrix_one}
    The right rank is
    \begin{align}
        r & = \operatorname{rank}(M_1).
    \end{align}
\end{definition}

\begin{definition}[Left rank]
    \label{def:mpu_left_rank}
    \lean{MPOTensor.leftRank, MPOTensor.leftRank_eq}
    \leanok
    \uses{def:mpu_source_matrix_two}
    The left rank is
    \begin{align}
        \ell & = \operatorname{rank}(M_2).
    \end{align}
\end{definition}

% Auxiliary transformation identities from source cuts (paper lines 450--506)
% and the existing physical-adjoint tensor, not the time-reversal classification.
\begin{theorem}[Source cuts under physical adjunction]
    \label{thm:mpu_physical_adjoint_source_cuts}
    \lean{MPOTensor.sourceCutM₁_physicalAdjointTensor,
        MPOTensor.sourceCutM₂_physicalAdjointTensor}
    \leanok
    \uses{def:mpdo_physical_adjoint, def:mpu_source_matrix_one, def:mpu_source_matrix_two}
    Physical adjunction exchanges the two source cuts and takes their adjoints:
    \begin{align}
        M_1(\mathcal U^\sharp) & =M_2(\mathcal U)^\dagger, \\
        M_2(\mathcal U^\sharp) & =M_1(\mathcal U)^\dagger.
    \end{align}
\end{theorem}

\begin{proof}\leanok
    \uses{def:mpdo_physical_adjoint, def:mpu_source_matrix_one, def:mpu_source_matrix_two}
    Physical adjunction exchanges the upper and lower physical indices and
    complex conjugates every entry. Under the source-cut regroupings, it also
    exchanges rows and columns. Hence the first cut becomes the conjugate
    transpose of the second, and conversely.
\end{proof}

\begin{definition}[Physical-pair swap]
    \label{def:mpu_physical_pair_swap}
    \lean{MPOTensor.physicalPairSwapEquiv,
        MPOTensor.physicalPairSwapEquiv_finProdFinEquiv}
    \leanok
    On the flattened physical alphabet, define the involution
    \begin{align}
        s(i,j) & =(j,i).
    \end{align}
    This acts on physical coordinates and is distinct from the doubled-bond
    exchange used in reflected transfer calculations.
\end{definition}

\begin{theorem}[Normalized flattening under physical adjunction]
    \label{thm:mpu_physical_adjoint_normalized_flattening}
    \lean{MPOTensor.normalizedFlattening_physicalAdjointTensor}
    \leanok
    \uses{def:mpdo_physical_adjoint, def:mpu_normalized_flattening, def:mpu_physical_pair_swap}
    The normalized flattening of the physical adjoint satisfies
    \begin{align}
        \widetilde{\mathcal U^\sharp}^{\,(i,j)}
         & =\overline{\widetilde{\mathcal U}^{\,(j,i)}}.
    \end{align}
\end{theorem}

\begin{proof}\leanok
    \uses{def:mpdo_physical_adjoint, def:mpu_normalized_flattening, def:mpu_physical_pair_swap}
    For every pair of physical indices $(i,j)$,
    \begin{align}
        \widetilde{\mathcal U^\sharp}^{\,(i,j)}
         & =\frac{1}{\sqrt d}\,(\mathcal U^\sharp)^{(i,j)}
        =\frac{1}{\sqrt d}\,\overline{\mathcal U^{(j,i)}}
        =\overline{\widetilde{\mathcal U}^{\,(j,i)}}.
    \end{align}
    The first equality is the normalized flattening, the second is physical
    adjunction, and the third uses the reality of $d^{-1/2}$.
\end{proof}

\begin{definition}[Canonical-form-II data under physical adjunction]
    \label{def:mpu_physical_adjoint_cfii_data}
    \lean{MPSTensor.CPSVCanonicalFormIIData.physicalAdjointNormalizedFlattening}
    \leanok
    \uses{def:cpsv_cfii, def:mpu_physical_pair_swap, thm:mpu_physical_adjoint_normalized_flattening, thm:cpsv_map_star_preservation, thm:cpsv_physical_reindexing, thm:cpsv_map_star_transfer_fixed_point}
    Given canonical-form-II data for $\widetilde{\mathcal U}$, define the
    transformed tuple by
    \begin{align}
        r'             & =r,                      \\
        D'_k           & =D_k,                    \\
        \mu'_k         & =\overline{\mu_k},       \\
        (A'_k)^{(i,j)} & =\overline{A_k^{(j,i)}}, \\
        V'             & =\overline V.
    \end{align}
    These data reconstruct $\widetilde{\mathcal U^\sharp}$.
\end{definition}

\begin{theorem}[Canonical form II under physical adjunction]
    \label{thm:mpu_physical_adjoint_cfii}
    \lean{MPSTensor.IsCPSVCanonicalFormII.physicalAdjointNormalizedFlattening}
    \leanok
    \uses{def:cpsv_cfii, def:mpdo_physical_adjoint, def:mpu_normalized_flattening}
    If $\widetilde{\mathcal U}$ is in canonical form II, then
    $\widetilde{\mathcal U^\sharp}$ is in canonical form II.
\end{theorem}

\begin{proof}\leanok
    \uses{def:mpu_physical_adjoint_cfii_data, thm:cpsv_map_star_preservation, thm:cpsv_physical_reindexing}
    Theorem~\ref{thm:cpsv_map_star_preservation} independently preserves
    normality and left-canonical normalization under entrywise conjugation, and
    Theorem~\ref{thm:cpsv_physical_reindexing} preserves these properties under
    the physical-pair swap.  Conjugating the coisometry equation gives
    \begin{align}
        V'V'^\dagger
         & =\overline{VV^\dagger}=\Id.
    \end{align}
    Conjugating the reconstruction at the swapped physical index gives
    \begin{align}
        \widetilde{\mathcal U^\sharp}^{\,(i,j)}
         & =(V')^\dagger
        \left(\bigoplus_k\mu'_k(A'_k)^{(i,j)}\right)V'.
    \end{align}
    Finally, each diagonal positive matrix $\Lambda_k$ is real, and therefore
    \begin{align}
        \E_{\overline{A_k}}(\Lambda_k)
         & =\overline{\E_{A_k}(\Lambda_k)}
        =\overline{\Lambda_k}
        =\Lambda_k.
    \end{align}
    Physical relabeling leaves this transfer map unchanged, which proves the
    transformed fixed-point equation.
\end{proof}

\begin{theorem}[Full support under physical adjunction]
    \label{thm:mpu_physical_adjoint_full_support}
    \lean{MPSTensor.CPSVCanonicalFormIIData.hasFullSupport_physicalAdjointNormalizedFlattening}
    \leanok
    \uses{def:mpu_full_support, def:mpu_physical_adjoint_cfii_data}
    If canonical-form-II data for $\widetilde{\mathcal U}$ have full
    support, then the data from
    Definition~\ref{def:mpu_physical_adjoint_cfii_data} also have full
    support.
\end{theorem}

\begin{proof}\leanok
    \uses{def:mpu_full_support, def:mpu_physical_adjoint_cfii_data}
    Complex conjugation of the weights leaves the block indices and their bond
    dimensions unchanged, so the full-support sum is unchanged.
\end{proof}

% Scope: the preceding three entries concern literal CFII data only. They do not
% identify a selected compact SVD, source factor, sourceU, or sourceV witness.
% Downstream orientation: the reflected transfer theorem separately fixes the
% right-hand density as rho.transpose; no such orientation is asserted here.

\begin{theorem}[Source ranks under physical adjunction]
    \label{thm:mpu_physical_adjoint_source_ranks}
    \lean{MPOTensor.rightRank_physicalAdjointTensor,
        MPOTensor.leftRank_physicalAdjointTensor}
    \leanok
    \uses{def:mpdo_physical_adjoint, def:mpu_right_rank, def:mpu_left_rank}
    Physical adjunction exchanges the two source ranks:
    \begin{align}
        r(\mathcal U^\sharp)    & =\ell(\mathcal U), \\
        \ell(\mathcal U^\sharp) & =r(\mathcal U).
    \end{align}
\end{theorem}

\begin{proof}\leanok
    \uses{thm:mpu_physical_adjoint_source_cuts, def:mpu_right_rank, def:mpu_left_rank}
    Apply Theorem~\ref{thm:mpu_physical_adjoint_source_cuts} and use invariance
    of matrix rank under conjugate transpose.
\end{proof}

\begin{theorem}[Source cuts under physical reindexing]
    \label{thm:mpu_source_cuts_physical_reindexing}
    \lean{MPOTensor.sourceCutM₁_reindexPhysical,
        MPOTensor.sourceCutM₂_reindexPhysical}
    \leanok
    \uses{def:mpu_physical_reindexing, def:mpu_source_matrix_one, def:mpu_source_matrix_two}
    A bijective relabeling of the physical basis simultaneously relabels the
    physical component of every source-cut row and column. Hence
    \begin{align}
        M_1(e^*\mathcal U) & \cong M_1(\mathcal U),
                           & M_2(e^*\mathcal U)     & \cong M_2(\mathcal U),
    \end{align}
    with the right and left orientations unchanged.
\end{theorem}

\begin{proof}\leanok
    \uses{def:mpu_physical_reindexing, def:mpu_source_matrix_one, def:mpu_source_matrix_two}
    Evaluate each cut after relabeling. In both cuts the upper physical index
    lies in the row and the lower physical index lies in the column, with the
    virtual index occupying the other product coordinate. Applying $e$ to both
    physical coordinates gives the two displayed row-and-column equivalences.
\end{proof}

\begin{theorem}[Source ranks under physical reindexing]
    \label{thm:mpu_source_ranks_physical_reindexing}
    \lean{MPOTensor.rightRank_reindexPhysical,
        MPOTensor.leftRank_reindexPhysical}
    \leanok
    \uses{def:mpu_physical_reindexing, def:mpu_right_rank, def:mpu_left_rank}
    Bijective physical relabeling preserves both source ranks:
    \begin{align}
        r(e^*\mathcal U) & =r(\mathcal U),
                         & \ell(e^*\mathcal U) & =\ell(\mathcal U).
    \end{align}
\end{theorem}

\begin{proof}\leanok
    \uses{thm:mpu_source_cuts_physical_reindexing, def:mpu_right_rank, def:mpu_left_rank}
    Matrix rank is unchanged by bijective row and column relabelings. Apply this
    fact to the two source-cut identities.
\end{proof}

\begin{definition}[Independent tensor product of matrix product tensors]
    \label{def:mps_independent_tensor_product}
    \lean{finTupleProdEquiv, finTupleProdEquiv_apply,
        MPSTensor.tensorProduct, MPSTensor.tensorProduct_apply}
    \leanok
    \uses{def:mps_tensor}
    Let $A=\{A^i\}_{i\in\Fin d}$ and $B=\{B^k\}_{k\in\Fin e}$ have
    auxiliary dimensions $D$ and $E$, respectively.  Write
    \begin{align}
        \pi_{m,n}:\Fin m\times\Fin n & \longrightarrow\Fin{mn}
    \end{align}
    for the standard finite-product coordinate bijection.  Their independent
    tensor product $A\boxtimes B$ has physical dimension $de$, auxiliary
    dimension $DE$, and local matrices determined by
    \begin{align}
        (A\boxtimes B)^{\pi_{d,e}(i,k)}_{
                \pi_{D,E}(\alpha,\gamma),\pi_{D,E}(\beta,\delta)}
         & =A^i_{\alpha,\beta}B^k_{\gamma,\delta}.
        \label{eq:mps_independent_tensor_product_entries}
    \end{align}
    For every $N$, the same coordinates give the pointwise splitting
    \begin{align}
        (\Fin N\to\Fin{de})
                  & \simeq(\Fin N\to\Fin d)\times(\Fin N\to\Fin e), \\
        \sigma    & \longmapsto(\sigma_d,\sigma_e),
                  &
        \sigma(n) & =\pi_{d,e}(\sigma_d(n),\sigma_e(n)).
        \label{eq:mps_tensor_product_tuple_split}
    \end{align}
    This is infrastructure for the tensoring clause in the proof of the MPU
    Index Theorem~IV.6(ii)
    \cite[lines~824--845]{Cirac2017MPU}, not a separately stated theorem of
    that paper.  The notation $A=\{A^i\}_i$ agrees with the canonical-form block
    notation in \cite[Section~2.3, lines~214--245]{Cirac2016MPDO_arXiv}.
\end{definition}

\begin{theorem}[Word evaluation of an independent tensor product]
    \label{thm:mps_eval_word_independent_tensor_product}
    \lean{MPSTensor.evalWord_tensorProduct}
    \leanok
    \uses{def:mps_independent_tensor_product, def:eval_word}
    If $w$ has length $n$ and
    \begin{align}
        w   & =\bigl(\pi_{d,e}(i_1,k_1),\ldots,
        \pi_{d,e}(i_n,k_n)\bigr),               \\
        w_d & =(i_1,\ldots,i_n),
            &
        w_e & =(k_1,\ldots,k_n),
    \end{align}
    then
    \begin{align}
        (A\boxtimes B)^w
         & =\operatorname{reind}_{\pi_{D,E}}
        \left(A^{w_d}\otimes B^{w_e}\right).
        \label{eq:mps_eval_word_independent_tensor_product}
    \end{align}
\end{theorem}

\begin{proof}\leanok
    \uses{def:mps_independent_tensor_product, def:eval_word}
    Induct on the length of $w$.  For the empty word, use
    $\Id_D\otimes\Id_E=\Id_{DE}$ and invariance of the
    identity under the coordinate bijection.  For a leading letter
    $\pi_{d,e}(i,k)$, equation
    (\ref{eq:mps_independent_tensor_product_entries}) and the induction
    hypothesis reduce the claim to
    \begin{align}
        (A^i\otimes B^k)(X\otimes Y)
         & =(A^iX)\otimes(B^kY),
    \end{align}
    after transporting multiplication through the same coordinate bijection.
\end{proof}

\begin{theorem}[Homogeneous word spans of an independent tensor product]
    \label{thm:mps_word_span_tensor_product}
    \lean{MPSTensor.wordSpan_tensorProduct,
        MPSTensor.isNBlkInjective_tensorProduct,
        MPSTensor.isNormal_tensorProduct}
    \leanok
    \uses{def:mps_independent_tensor_product, def:word_span, def:normal}
    Let $S_N(A)=\spn_{\C}\{A^w:|w|=N\}$.  The pointwise splitting of
    length-$N$ words gives
    \begin{align}
        S_N(A\boxtimes B)
         & =\spn_{\C}\{\operatorname{reind}_{\pi_{D,E}}(A^u\otimes B^v):
        |u|=|v|=N\}.
        \label{eq:mpu_tensor_word_span}
    \end{align}
    In particular, if $S_N(A)=\MN{D}$ and $S_N(B)=\MN{E}$, then
    $S_N(A\boxtimes B)=\MN{DE}$.  If $A$ and $B$ are algebraically normal,
    then $A\boxtimes B$ is algebraically normal.

    This is the homogeneous-word infrastructure for the tensoring sentence in
    the proof of the MPU Index Theorem~IV.6(ii)
    \cite[lines~824--845]{Cirac2017MPU}, not a separate theorem stated there.
    It applies to the retained blocks in the canonical-form decomposition
    \cite[Section~2.3, lines~214--245]{Cirac2016MPDO_arXiv}.
\end{theorem}

\begin{proof}\leanok
    \uses{thm:mps_eval_word_independent_tensor_product, def:word_span, def:normal}
    For every product word $w$, let $(u,v)$ be its pointwise component words.
    The bijection between $w$ and $(u,v)$ and
    Theorem~\ref{thm:mps_eval_word_independent_tensor_product} give
    \begin{align}
        (A\boxtimes B)^w
         & =\operatorname{reind}_{\pi_{D,E}}(A^u\otimes B^v),
    \end{align}
    which proves~(\ref{eq:mpu_tensor_word_span}) after taking linear spans.
    Pure tensors span $\MN{D}\otimes\MN{E}$.  The Kronecker linear
    equivalence, followed by simultaneous reindexing along $\pi_{D,E}$, is an
    isomorphism onto $\MN{DE}$; this proves the fixed-length consequence.  If
    $S_{N_A}(A)=\MN{D}$ and $S_{N_B}(B)=\MN{E}$ for positive $N_A,N_B$, then
    persistence at positive multiples gives
    \begin{align}
        S_{N_A N_B}(A) & =\MN{D}, &
        S_{N_A N_B}(B) & =\MN{E}.
    \end{align}
    Applying the fixed-length consequence at $N_A N_B$ proves algebraic
    normality of $A\boxtimes B$.
\end{proof}

\begin{theorem}[Left-canonical normal blocks under independent tensor products]
    \label{thm:mps_normal_block_tensor_product}
    \lean{MPSTensor.leftCanonical_tensorProduct,
        MPSTensor.IsNormalTensor.tensorProduct_of_leftCanonical}
    \leanok
    \uses{def:mps_independent_tensor_product, def:left_canonical_tensor, def:nt_cpsv}
    Let $A$ and $B$ be CPSV normal tensors, and suppose that
    \begin{align}
        \sum_i(A^i)^\dagger A^i & =\Id_D, &
        \sum_k(B^k)^\dagger B^k & =\Id_E.
        \label{eq:mpu_tensor_left_canonical_inputs}
    \end{align}
    Then $A\boxtimes B$ is a CPSV normal tensor and is left-canonical:
    \begin{align}
        \sum_{i,k}
        \left((A\boxtimes B)^{\pi_{d,e}(i,k)}\right)^\dagger
        (A\boxtimes B)^{\pi_{d,e}(i,k)}
         & =\Id_{DE}.
        \label{eq:mpu_tensor_left_canonical_output}
    \end{align}

    This is the retained-normal-block infrastructure for the tensoring
    sentence in the proof of the MPU Index Theorem~IV.6(ii)
    \cite[lines~824--845]{Cirac2017MPU}, not a separate theorem stated there.
    The retained-normal-block terminology is that of the canonical-form
    decomposition in
    \cite[Section~2.3, lines~214--245]{Cirac2016MPDO_arXiv}; the two
    left-canonical identities are explicit infrastructure hypotheses, not an
    additional assertion attributed to that passage.
\end{theorem}

\begin{proof}\leanok
    \uses{thm:mps_word_span_tensor_product, thm:cpsv_normal_eventually_injective, thm:is_normal_tensor_of_is_normal_left_canonical}
    The adjoint and multiplication laws for Kronecker products, together
    with the fact that simultaneous reindexing preserves adjoints, products,
    and finite sums, give
    \begin{align}
         & \sum_{i,k}
        \left((A\boxtimes B)^{\pi_{d,e}(i,k)}\right)^\dagger
        (A\boxtimes B)^{\pi_{d,e}(i,k)}                 \\
         & \quad=\operatorname{reind}_{\pi_{D,E}}
        \left(\sum_{i,k}
        (A^i\otimes B^k)^\dagger(A^i\otimes B^k)\right) \\
         & \quad=\operatorname{reind}_{\pi_{D,E}}
        \left(
        \left(\sum_i(A^i)^\dagger A^i\right)\otimes
        \left(\sum_k(B^k)^\dagger B^k\right)
        \right)                                         \\
         & \quad=\operatorname{reind}_{\pi_{D,E}}
        (\Id_D\otimes\Id_E)=\Id_{DE},
    \end{align}
    proving~(\ref{eq:mpu_tensor_left_canonical_output}).  Spectral normality
    gives positive lengths $N_A,N_B$ at which the two homogeneous word spans
    are full.  Theorem~\ref{thm:mps_word_span_tensor_product} therefore makes
    $A\boxtimes B$ algebraically normal at the common length $N_A N_B$.
    The algebraically normal, left-canonical tensor is CPSV normal by
    Theorem~\ref{thm:is_normal_tensor_of_is_normal_left_canonical}.
\end{proof}

\begin{definition}[Independent tensor product of matrix product operator tensors]
    \label{def:mpu_independent_tensor_product}
    \lean{MPOTensor.tensorProduct, MPOTensor.tensorProduct_apply, finTupleProdEquiv}
    \leanok
    \uses{def:mpo_tensor}
    Let $\mathcal U$ and $\mathcal V$ have physical dimensions $d,e$ and
    auxiliary dimensions $D,E$, respectively. Their independent tensor
    product $\mathcal U\boxtimes\mathcal V$ has dimensions $de$ and $DE$ and
    entries
    \begin{align}
        (\mathcal U\boxtimes\mathcal V)^{(i,k),(j,l)}_{(\alpha,\gamma),(\beta,\delta)}
         & =\mathcal U^{i,j}_{\alpha,\beta}\mathcal V^{k,l}_{\gamma,\delta}.
    \end{align}
    A product-valued physical configuration splits sitewise into one
    configuration for $\mathcal U$ and one for $\mathcal V$.
\end{definition}

\begin{definition}[Blocked product-alphabet equivalence]
    \label{def:mpu_tensor_product_block_equiv}
    \lean{MPOTensor.tensorProductBlockEquiv}
    \leanok
    \uses{def:mpdo_physical_blocking}
    A blocked word of $L$ product letters separates canonically into two
    blocked words:
    \begin{align}
        \Fin{(de)^L} & \simeq\Fin{d^L e^L}.
    \end{align}
    The map decodes the length-$L$ word, splits every letter into its $d$ and
    $e$ components, and re-encodes the two resulting words.
\end{definition}

\begin{theorem}[Decoding the blocked product alphabet]
    \label{thm:mpu_tensor_product_block_decoding}
    \lean{MPOTensor.decodeBlock_tensorProductBlockEquiv_divNat,
        MPOTensor.decodeBlock_tensorProductBlockEquiv_modNat}
    \leanok
    \uses{def:mpu_tensor_product_block_equiv}
    If $q$ denotes the blocked product-alphabet equivalence, then each site
    $s$ satisfies
    \begin{align}
        \operatorname{dec}_d(q(I)_{1},s)
         & =\operatorname{dec}_{de}(I,s)_{1}, \\
        \operatorname{dec}_e(q(I)_{2},s)
         & =\operatorname{dec}_{de}(I,s)_{2}.
    \end{align}
\end{theorem}

\begin{proof}\leanok
    \uses{def:mpu_tensor_product_block_equiv}
    Each equality follows by composing the inverse product encoding at every
    site with the corresponding projection.
\end{proof}

\begin{theorem}[Finite-size operator of an independent tensor product]
    \label{thm:mpu_finite_tensor_product}
    \lean{MPOTensor.mpo_tensorProduct}
    \leanok
    \uses{def:mpu_independent_tensor_product, def:mpo_family}
    For every system size $N$, the sitewise configuration bijection identifies
    the finite-size operator with
    \begin{align}
        (\mathcal U\boxtimes\mathcal V)^{(N)}
         & \cong \mathcal U^{(N)}\otimes\mathcal V^{(N)}.
    \end{align}
\end{theorem}

\begin{proof}\leanok
    \uses{def:mpu_independent_tensor_product, def:mpo_family}
    Grouping the two auxiliary indices makes every closed-word matrix a
    Kronecker product. The identity
    \begin{align}
        \Tr(A\otimes B) & =\Tr(A)\Tr(B)
    \end{align}
    then gives the stated finite-size operator after the sitewise
    configuration bijection.
\end{proof}

\begin{theorem}[Independent tensor products preserve matrix product unitarity]
    \label{thm:mpu_tensor_product}
    \lean{MPOTensor.IsMPU.tensorProduct}
    \leanok
    \uses{def:mpu, def:mpu_independent_tensor_product}
    The independent tensor product of two matrix product unitary tensors is a
    matrix product unitary tensor. This is the tensoring operation in the proof
    of \cite[Theorem~IV.6(ii)]{Cirac2017MPU}.
\end{theorem}

\begin{proof}\leanok
    \uses{def:mpu_independent_tensor_product, thm:mpu_finite_tensor_product, thm:mpu_reindex_unitary, lem:mpu_mul_adjoint}
    The multiplication rule $(A\otimes B)(C\otimes D)=AC\otimes BD$ gives
    \begin{align}
         & (\mathcal U^{(N)}\otimes\mathcal V^{(N)})
        (\mathcal U^{(N)}\otimes\mathcal V^{(N)})^\dagger   \\
         & \quad=\mathcal U^{(N)}(\mathcal U^{(N)})^\dagger
        \otimes\mathcal V^{(N)}(\mathcal V^{(N)})^\dagger
        =\Id\otimes\Id=\Id.
    \end{align}
    Simultaneous reindexing preserves this unitarity equation.
\end{proof}

\begin{theorem}[Blocking an independent tensor product]
    \label{thm:mpu_block_tensor_product}
    \lean{MPOTensor.blockTensor_tensorProduct}
    \leanok
    \uses{def:mpdo_physical_blocking, def:mpu_physical_reindexing, def:mpu_independent_tensor_product, def:mpu_tensor_product_block_equiv}
    For every blocking length $L$, separating each blocked product letter by
    the canonical equivalence $q$ gives the exact identity
    \begin{align}
        (\mathcal U\boxtimes\mathcal V)^{[L]}
         & =q^*\bigl(\mathcal U^{[L]}\boxtimes\mathcal V^{[L]}\bigr).
    \end{align}
    The ordering of the two component words agrees with the independent
    tensor-product convention.
\end{theorem}

\begin{proof}\leanok
    \uses{def:mpdo_physical_blocking, def:mpu_physical_reindexing, def:mpu_independent_tensor_product, thm:mpu_tensor_product_block_decoding}
    Evaluate both sides on blocked ket and bra letters. Word multiplication of
    local Kronecker products gives
    \begin{align}
        (\mathcal U\boxtimes\mathcal V)[I,J]
         & =\mathcal U[I_1,J_1]\otimes\mathcal V[I_2,J_2].
    \end{align}
    The two decoding identities identify $I_1,I_2$ and $J_1,J_2$ with the
    corresponding letters of $q(I)$ and $q(J)$.
\end{proof}

\begin{definition}[Source-cut shuffle for an independent tensor product]
    \label{def:mpu_tensor_product_cut_shuffle}
    \lean{MPOTensor.tensorProductCutShuffle}
    \leanok
    For four finite index sets, the source-cut shuffle sends
    \begin{align}
        ((x,y),(x',y'))
         & \longmapsto ((x,x'),(y,y')).
    \end{align}
    It is applied to the row and column products with the orders appropriate to
    $M_1$ and $M_2$.
\end{definition}

\begin{theorem}[Source cuts of an independent tensor product]
    \label{thm:mpu_tensor_product_source_cuts}
    \lean{MPOTensor.sourceCutM₁_tensorProduct,
        MPOTensor.sourceCutM₂_tensorProduct}
    \leanok
    \uses{def:mpu_independent_tensor_product, def:mpu_tensor_product_cut_shuffle, def:mpu_source_matrix_one, def:mpu_source_matrix_two}
    The source cuts satisfy
    \begin{align}
        M_1(\mathcal U\boxtimes\mathcal V)
         & \cong M_1(\mathcal U)\otimes M_1(\mathcal V), \\
        M_2(\mathcal U\boxtimes\mathcal V)
         & \cong M_2(\mathcal U)\otimes M_2(\mathcal V).
    \end{align}
    Thus the paper's right and left orientations are unchanged.
\end{theorem}

\begin{proof}\leanok
    \uses{def:mpu_independent_tensor_product, def:mpu_tensor_product_cut_shuffle, def:mpu_source_matrix_one, def:mpu_source_matrix_two}
    Evaluate both source cuts entry by entry after the displayed row and column
    shuffles. Each entry is the product of the corresponding entries of the
    two source cuts.
\end{proof}

\begin{theorem}[Source-cut ranks of an independent tensor product]
    \label{thm:mpu_tensor_product_source_ranks}
    \lean{MPOTensor.rightRank_tensorProduct,
        MPOTensor.leftRank_tensorProduct}
    \leanok
    \uses{def:mpu_independent_tensor_product, def:mpu_right_rank, def:mpu_left_rank}
    The right and left source-cut ranks are multiplicative:
    \begin{align}
        r(\mathcal U\boxtimes\mathcal V)
         & =r(\mathcal U)r(\mathcal V),       \\
        \ell(\mathcal U\boxtimes\mathcal V)
         & =\ell(\mathcal U)\ell(\mathcal V).
    \end{align}
    This is the rank-multiplication ingredient for the tensoring case in the
    proof of \cite[Theorem~IV.6(ii)]{Cirac2017MPU}.
\end{theorem}

\begin{proof}\leanok
    \uses{def:mpu_independent_tensor_product, thm:mpu_tensor_product_source_cuts, def:mpu_right_rank, def:mpu_left_rank, thm:mpu_kronecker_rank}
    Matrix rank is unchanged by row and column bijections and satisfies
    $\operatorname{rank}(A\otimes B)=\operatorname{rank}(A)
        \operatorname{rank}(B)$.
\end{proof}

\begin{theorem}[One-site blocking bounds for the source-cut ranks]
    \label{thm:mpu_source_cut_rank_blocking_succ}
    \lean{MPOTensor.rightRank_blockTensor_succ_le,
        MPOTensor.leftRank_blockTensor_succ_le}
    \leanok
    \uses{def:mpdo_physical_blocking, def:mpu_source_matrix_one, def:mpu_source_matrix_two, def:mpu_right_rank, def:mpu_left_rank}
    Let $\mathcal U^{[L]}$ denote the direct blocking of $L$ sites of a tensor
    $\mathcal U$ with physical dimension $d$, and let $r_L$ and $\ell_L$ be
    its right and left source-cut ranks. Then
    \begin{align}
        r_{L+1}    & \leq d\,r_L,    \\
        \ell_{L+1} & \leq d\,\ell_L.
    \end{align}
    These are the one-site forms of the two upper bounds in the proof of
    \cite[Proposition~IV.2]{Cirac2017MPU}.
\end{theorem}

\begin{proof}\leanok
    \uses{def:mpu_source_svd_data}
    For the first cut, choose a compact factorization of the source matrix of
    $\mathcal U^{[L]}$ through a space of dimension $r_L$. Split each physical
    word of length $L+1$ into its first letter and its remaining word of length
    $L$. For suitable rectangular matrices $A_{1,L}$ and $B_{1,L}$, the
    successor source matrix factors as
    \begin{align}
        M_{1,L+1} & =A_{1,L}B_{1,L}.
    \end{align}
    The intermediate space is $\mathbb C^d\otimes\mathbb C^{r_L}$, and hence
    \begin{align}
        \operatorname{rank}(M_{1,L+1})
         & \leq \dim\bigl(\mathbb C^d\otimes\mathbb C^{r_L}\bigr)=d r_L.
    \end{align}
    Interchanging the upper and lower physical indices gives matrices
    $A_{2,L}$ and $B_{2,L}$ with
    \begin{align}
        M_{2,L+1} & =A_{2,L}B_{2,L}.
    \end{align}
    This factorization passes through
    $\mathbb C^d\otimes\mathbb C^{\ell_L}$, so
    \begin{align}
        \operatorname{rank}(M_{2,L+1})
         & \leq \dim\bigl(\mathbb C^d\otimes\mathbb C^{\ell_L}\bigr)=d\ell_L.
    \end{align}
\end{proof}

\begin{theorem}[Source-cut rank bounds between two blocking lengths]
    \label{thm:mpu_source_cut_rank_blocking}
    \lean{MPOTensor.rightRank_blockTensor_le_pow_mul,
        MPOTensor.leftRank_blockTensor_le_pow_mul}
    \leanok
    If $L_0\leq L$, then
    \begin{align}
        r_L    & \leq d^{L-L_0}r_{L_0},    \\
        \ell_L & \leq d^{L-L_0}\ell_{L_0}.
    \end{align}
    These are the two blocking upper bounds used in the proof of
    \cite[Proposition~IV.2]{Cirac2017MPU}.
\end{theorem}

\begin{proof}\leanok
    \uses{thm:mpu_source_cut_rank_blocking_succ}
    Induct on $L-L_0$. The case $L=L_0$ is equality. For the right rank, the
    successor bound and the induction hypothesis give
    \begin{align}
        r_{L+1}
         & \leq d r_L
        \leq d\bigl(d^{L-L_0}r_{L_0}\bigr)
        =d^{L+1-L_0}r_{L_0}.
    \end{align}
    The same multiplication by $d$ gives
    \begin{align}
        \ell_{L+1}
         & \leq d\ell_L
        \leq d\bigl(d^{L-L_0}\ell_{L_0}\bigr)
        =d^{L+1-L_0}\ell_{L_0}.
    \end{align}
\end{proof}

\begin{theorem}[Exact blocking growth from endpoint rank products]
    \label{thm:mpu_source_cut_rank_blocking_saturation}
    \lean{MPOTensor.blockingRanks_eq_pow_mul_of_products,
        MPOTensor.rightRank_blockTensor_eq_pow_mul_of_products,
        MPOTensor.leftRank_blockTensor_eq_pow_mul_of_products}
    \leanok
    Let $d>0$ and $L_0\leq L$. Suppose the endpoint source-cut ranks satisfy
    the supplied product identities
    \begin{align}
        r_{L_0}\ell_{L_0} & =d^{2L_0}, &
        r_L\ell_L         & =d^{2L}.
    \end{align}
    Then the blocking bounds are exact:
    \begin{align}
        r_L    & =d^{L-L_0}r_{L_0},    &
        \ell_L & =d^{L-L_0}\ell_{L_0}.
    \end{align}
    This is the saturation step in the proof of
    \cite[Proposition~IV.2]{Cirac2017MPU}.
    % The endpoint product identities are hypotheses here; they are not consequences
    % of the blocking bounds.
    % **Local fix (rank-product exponent):** Source line 703 prints $d^k$; the
    % consistent blocked identity is $d^{2k}$. See
    % docs/paper-gaps/mpu_blocking_rank_product_exponent.tex.
\end{theorem}

\begin{proof}\leanok
    \uses{thm:mpu_source_cut_rank_blocking}
    The blocking bounds give
    \begin{align}
        r_L    & \leq d^{L-L_0}r_{L_0},    &
        \ell_L & \leq d^{L-L_0}\ell_{L_0}.
    \end{align}
    Their upper bounds have product
    \begin{align}
        \bigl(d^{L-L_0}r_{L_0}\bigr)
        \bigl(d^{L-L_0}\ell_{L_0}\bigr)
         & =d^{2(L-L_0)}d^{2L_0}=d^{2L}=r_L\ell_L.
    \end{align}
    Since $d>0$, the latter product is positive, so both $r_L$ and $\ell_L$
    are positive. Hence neither upper bound can be strict, and both equalities
    follow.
\end{proof}

\begin{definition}[Product-index source weight]
    \label{def:mpu_source_weight}
    \lean{MPOTensor.sourceWeight}
    \leanok
    \uses{def:mpu_source_matrix_one}
    The row space of $M_1$ has product order
    $(\text{upper physical},\text{right auxiliary})$. For
    $\rho\in\MN{D}$, define
    \begin{align}
        W_\rho & =\Id_d\otimes\rho.
    \end{align}
    This is the source weight in
    \cite[equation labelled~\texttt{Y1Y1X1X1}]{Cirac2017MPU}.
    % Source lines: paper_v2.tex 479--487.
\end{definition}

\begin{theorem}[Positivity of the product-index source weight]
    \label{thm:mpu_source_weight_posdef}
    \lean{MPOTensor.sourceWeight_posDef}
    \leanok
    \uses{def:mpu_source_weight}
    If $\rho$ is positive definite, then $W_\rho$ is positive definite.
\end{theorem}

\begin{proof}
    \leanok
    \uses{def:mpu_source_weight}
    The Kronecker product of the positive-definite matrices $\Id_d$ and
    $\rho$ is positive definite.
\end{proof}

\begin{definition}[Source singular-value data and compressed source metric]
    \label{def:mpu_source_svd_data}
    \lean{MPOTensor.SourceCutSVD, MPOTensor.sourceSVD₁,
        MPOTensor.sourceSVD₂, MPOTensor.sourceGram₁}
    \leanok
    \uses{thm:mpu_compact_svd, def:mpu_source_matrix_one, def:mpu_source_matrix_two, def:mpu_right_rank, def:mpu_left_rank, def:mpu_source_weight}
    Choose compact singular-value decompositions
    \begin{align}
        M_i & =V_i^\dagger D_iU_i,
            & V_iV_i^\dagger       & =\Id,
            & U_iU_i^\dagger       & =\Id,
    \end{align}
    with intermediate dimensions $r$ and $\ell$, respectively. For
    $\rho\in\MN{D}$, define the compressed positive metric
    \begin{align}
        A & =V_1W_\rho V_1^\dagger.
    \end{align}
    The inverse square root $A^{-1/2}$, rather than $A$ itself, is the
    normalization factor used in the first weighted source factor. This is the
    source-factor construction of \cite[Section~III]{Cirac2017MPU}.
    % Source lines: paper_v2.tex 479--500; labels Y1Y1X1X1 and Z1Z2.
\end{definition}

\begin{theorem}[Positivity of the first compressed source metric]
    \label{thm:mpu_source_gram_posdef}
    \lean{MPOTensor.sourceGram₁_posDef}
    \leanok
    \uses{def:mpu_source_svd_data}
    If $\rho$ is positive definite, then $A$ is positive definite.
\end{theorem}

\begin{proof}
    \leanok
    \uses{def:mpu_source_svd_data, thm:mpu_source_weight_posdef}
    The map $V_1^\dagger$ is injective because $V_1V_1^\dagger=\Id_r$.
    Positive definiteness of $W_\rho$ is therefore preserved by this
    congruence.
\end{proof}

\begin{definition}[Weighted source factors]
    \label{def:mpu_weighted_source_factors}
    \lean{MPOTensor.sourceX₁, MPOTensor.sourceY₁,
        MPOTensor.sourceZ₁, MPOTensor.sourceX₂, MPOTensor.sourceY₂,
        MPOTensor.sourceZ₂, MPOTensor.sourceX₁_apply,
        MPOTensor.sourceY₁_apply, MPOTensor.sourceZ₁_apply,
        MPOTensor.sourceX₂_apply, MPOTensor.sourceY₂_apply,
        MPOTensor.sourceZ₂_apply, MPOTensor.SourceFactors,
        MPOTensor.sourceFactors}
    \leanok
    \uses{def:mpu_source_svd_data, thm:mpu_source_gram_posdef}
    Let $\rho\in\MN{D}$ be positive definite. Define
    \begin{align}
        X_1 & =V_1^\dagger A^{-1/2},
            & Y_1                    & =A^{1/2}D_1U_1,
            & Z_1                    & =U_1^\dagger D_1^{-1}A^{-1/2}, \\
        X_2 & =V_2^\dagger,
            & Y_2                    & =D_2U_2,
            & Z_2                    & =U_2^\dagger D_2^{-1}.
    \end{align}
    These are the factors in \cite[Section~III]{Cirac2017MPU}.
\end{definition}

\begin{proof}
    \leanok
    \uses{def:mpu_source_weight, def:mpu_source_svd_data, thm:mpu_source_gram_posdef}
    Since $A$ is positive definite, its support projection is the identity and
    \begin{align}
        A^{-1/2}A^{1/2} & =\Id_r,
                        & A^{-1/2}AA^{-1/2} & =\Id_r.
    \end{align}
    Therefore
    \begin{align}
        X_1Y_1
         & =V_1^\dagger A^{-1/2}A^{1/2}D_1U_1
        =V_1^\dagger D_1U_1=M_1,                      \\
        X_1^\dagger W_\rho X_1
         & =A^{-1/2}V_1W_\rho V_1^\dagger A^{-1/2}
        =A^{-1/2}AA^{-1/2}=\Id_r,                     \\
        Y_1Z_1
         & =A^{1/2}D_1U_1U_1^\dagger D_1^{-1}A^{-1/2}
        =A^{1/2}A^{-1/2}=\Id_r.
    \end{align}
    For the second cut, $V_2V_2^\dagger=U_2U_2^\dagger=\Id_\ell$ and
    $D_2D_2^{-1}=\Id_\ell$ give
    \begin{align}
        X_2Y_2          & =V_2^\dagger D_2U_2=M_2,              \\
        X_2^\dagger X_2 & =V_2V_2^\dagger=\Id_\ell,             \\
        Y_2Z_2          & =D_2U_2U_2^\dagger D_2^{-1}=\Id_\ell.
    \end{align}
\end{proof}

\begin{theorem}[Source-cut factorizations]
    \label{thm:mpu_source_factorizations}
    \lean{MPOTensor.sourceCutM₁_eq_sourceX₁_mul_sourceY₁,
        MPOTensor.sourceCutM₂_eq_sourceX₂_mul_sourceY₂}
    \leanok
    \uses{def:mpu_weighted_source_factors}
    The two factorizations are
    \begin{align}
        M_1 & =X_1Y_1, & M_2 & =X_2Y_2.
    \end{align}
    These are the source decompositions in
    \cite[Section~III]{Cirac2017MPU}.
\end{theorem}

\begin{proof}
    \leanok
    \uses{def:mpu_weighted_source_factors, def:mpu_source_svd_data}
    The source-factor construction supplies $M_1=X_1Y_1$. The second
    compact singular-value decomposition gives $M_2=X_2Y_2$ directly.
\end{proof}

\begin{theorem}[Entry forms of the source-cut factorizations]
    \label{thm:mpu_source_factorization_entries}
    \lean{MPOTensor.SourceFactors.X₁_mul_Y₁_apply,
        MPOTensor.SourceFactors.X₂_mul_Y₂_apply,
        MPOTensor.sourceX₁_mul_sourceY₁_apply,
        MPOTensor.sourceX₂_mul_sourceY₂_apply}
    \leanok
    \uses{thm:mpu_source_factorizations, def:mpu_source_matrix_one, def:mpu_source_matrix_two}
    The entries of the two products satisfy
    \begin{align}
        (X_1Y_1)_{(i,\beta),(\alpha,j)}
         & =\mathcal U^i_{j,\alpha,\beta}, \\
        (X_2Y_2)_{(\alpha,i),(j,\beta)}
         & =\mathcal U^i_{j,\alpha,\beta}.
    \end{align}
    These are the two graphical entry identifications in
    \cite[Section~III]{Cirac2017MPU}.
\end{theorem}

\begin{proof}
    \leanok
    \uses{thm:mpu_source_factorizations, def:mpu_source_matrix_one, def:mpu_source_matrix_two}
    Evaluate the two matrix factorizations at
    $((i,\beta),(\alpha,j))$ and $((\alpha,i),(j,\beta))$, respectively.
\end{proof}

\begin{theorem}[Source-factor normalizations]
    \label{thm:mpu_source_factor_normalizations}
    \lean{MPOTensor.sourceX₁_weighted_isometry,
        MPOTensor.sourceX₁_weighted_isometry_apply,
        MPOTensor.sourceX₂_isometry,
        MPOTensor.sourceX₂_isometry_apply}
    \leanok
    \uses{def:mpu_weighted_source_factors, def:mpu_source_weight}
    The left factors satisfy
    \begin{align}
        X_1^\dagger W_\rho X_1 & =\Id_r,
                               & X_2^\dagger X_2 & =\Id_\ell, \\
        \sum_{x,y}\overline{(X_1)_{x,r}}(W_\rho)_{x,y}(X_1)_{y,r'}
                               & =\delta_{r,r'},              \\
        \sum_{\beta,i}\overline{(X_2)_{(\beta,i),l}}(X_2)_{(\beta,i),l'}
                               & =\delta_{l,l'}.
    \end{align}
    These are \cite[Section~III]{Cirac2017MPU}.
\end{theorem}

\begin{proof}
    \leanok
    \uses{def:mpu_weighted_source_factors, def:mpu_source_svd_data}
    The source-factor construction supplies
    $X_1^\dagger W_\rho X_1=\Id_r$; evaluating at $(r,r')$ gives the displayed
    weighted sum. Since $X_2=V_2^\dagger$, the
    coisometry identity $V_2V_2^\dagger=\Id_\ell$ gives
    $X_2^\dagger X_2=\Id_\ell$; evaluating at $(l,l')$ gives the second
    displayed entry sum.
\end{proof}

\begin{theorem}[Rotated second source-cut Gram]
    \label{thm:mpu_source_y_two_rotated_gram}
    \lean{MPOTensor.sourceY₂_gram_eq_rotated_sourceCutM₂_gram}
    \leanok
    \uses{def:mpu_weighted_source_factors}
    For every tensor $\mathcal U$, physical indices $p,q$, and auxiliary
    indices $a,b$,
    \begin{align}
        \sum_{\beta,z}
        \overline{\mathcal U^z_{p,\beta,a}}\,
        \mathcal U^z_{q,\beta,b}
         & =\sum_l \overline{(Y_2)_{l,(p,a)}}(Y_2)_{l,(q,b)}.
    \end{align}
    Thus the factor indexed by $(p,a)$ is starred and the factor indexed by
    $(q,b)$ is unstarred. This is the second-cut rotation used in
    \cite[Theorem~III.8, Section~III.B, equations~(31)--(32)]{Cirac2017MPU}.
    The identity requires no Matrix Product Unitary, dimension, positivity, or
    simple-tensor hypothesis.
\end{theorem}

\begin{proof}
    \leanok
    \uses{def:mpu_weighted_source_factors, thm:mpu_source_factorization_entries, thm:mpu_source_factor_normalizations}
    Substitute the entry factorization $M_2=X_2Y_2$ into the left-hand side
    and reorder the finite sums. The contraction of the two $X_2$ factors is
    $X_2^\dagger X_2=\Id_\ell$, so the two source indices coincide and the
    remaining sum is the displayed Gram matrix of $Y_2$.
\end{proof}

\begin{theorem}[Recovery of the right source factors]
    \label{thm:mpu_source_factor_recovery}
    \lean{MPOTensor.sourceY₁_eq_sourceX₁_conjTranspose_mul_weight_mul_sourceCutM₁,
        MPOTensor.sourceY₂_eq_sourceX₂_conjTranspose_mul_sourceCutM₂}
    \leanok
    \uses{def:mpu_source_weight, def:mpu_weighted_source_factors, def:mpu_source_matrix_one, def:mpu_source_matrix_two}
    The right source factors are recovered from the two source matrices by
    \begin{align}
        Y_1 & =X_1^\dagger W_\rho M_1,
            & Y_2                      & =X_2^\dagger M_2.
    \end{align}
    Here $W_\rho=I_d\otimes\rho$ in the product-index order
    $(\text{upper physical},\text{right auxiliary})$.
\end{theorem}

\begin{proof}\leanok
    \uses{thm:mpu_source_factorizations, thm:mpu_source_factor_normalizations}
    Multiply $M_1=X_1Y_1$ on the left by $X_1^\dagger W_\rho$ and use
    $X_1^\dagger W_\rho X_1=I_r$. Similarly, multiply
    $M_2=X_2Y_2$ on the left by $X_2^\dagger$ and use
    $X_2^\dagger X_2=I_\ell$.
\end{proof}

\begin{theorem}[Right inverses of the source factors]
    \label{thm:mpu_source_factor_right_inverses}
    \lean{MPOTensor.sourceY₁_mul_sourceZ₁,
        MPOTensor.sourceY₂_mul_sourceZ₂}
    \leanok
    \uses{def:mpu_weighted_source_factors}
    The right factors satisfy
    \begin{align}
        Y_1Z_1 & =\Id_r, & Y_2Z_2 & =\Id_\ell.
    \end{align}
    This is \cite[Section~III]{Cirac2017MPU}.
\end{theorem}

\begin{proof}
    \leanok
    \uses{def:mpu_weighted_source_factors, def:mpu_source_svd_data}
    The source-factor construction supplies $Y_1Z_1=\Id_r$. For the
    second cut,
    \begin{align}
        Y_2Z_2
         & =D_2U_2U_2^\dagger D_2^{-1}
        =D_2D_2^{-1}=\Id_\ell.
    \end{align}
\end{proof}

\begin{definition}[Source tensors $u$ and $v$]
    \label{def:mpu_source_uv}
    \lean{MPOTensor.SourceFactors.sourceU,
        MPOTensor.SourceFactors.sourceV,
        MPOTensor.SourceFactors.sourceU_apply,
        MPOTensor.SourceFactors.sourceV_apply,
        MPOTensor.sourceFactors_sourceU,
        MPOTensor.sourceFactors_sourceV,
        MPOTensor.sourceU,MPOTensor.sourceV,
        MPOTensor.sourceU_apply,MPOTensor.sourceV_apply,
        MPOTensor.sourceU_conjTranspose_apply,
        MPOTensor.sourceV_conjTranspose_apply}
    \leanok
    \uses{def:mpu_weighted_source_factors}
    % Leg-order audit: docs/paper-gaps/mpu_source_cut_orientation.tex.
    Given source-cut decompositions $M_1=X_1Y_1$ and $M_2=X_2Y_2$, define
    \begin{align}
        u_{(l,r),(i_1,i_2)}
         & =\sum_\beta (Y_2)_{l,(i_1,\beta)}(Y_1)_{r,(\beta,i_2)},    \\
        v_{(j_1,j_2),(r,l)}
         & =\sum_\alpha (X_1)_{(j_1,\alpha),r}(X_2)_{(\alpha,j_2),l}.
    \end{align}
    Thus $u:\C^{d^2}\to\C^{\ell r}$ has source order $\ell\times r$,
    while $v:\C^{r\ell}\to\C^{d^2}$ has source order $r\times\ell$.
    These are the literal triangle contractions in
    \cite[Section~III, equations~\texttt{uuvv}--\texttt{vdagger}]{Cirac2017MPU}.
\end{definition}

\begin{theorem}[Paper-gate two-site factorization]
    \label{thm:mpu_source_gates_two_site_factorization}
    \lean{MPOTensor.trace_mul_eq_sourceV_mul_sourceU_swap,
        MPOTensor.mpo_two_pair_entry_eq_sourceV_mul_sourceU_swap,
        MPOTensor.mpo_two_reindex_eq_sourceV_mul_sourceU_swap}
    \leanok
    \uses{def:mpu_source_uv, def:mpo_family}
    Identify a two-site physical configuration with an ordered pair in
    $\Fin d\times\Fin d$. Then
    \begin{align}
        \operatorname{reindex}(U^{(2)})
         & =v\,\operatorname{swap}_{r,\ell}(u),
    \end{align}
    where $\operatorname{swap}_{r,\ell}$ reorders the source index of $u$ from
    $\ell\times r$ to the $r\times\ell$ order carried by the columns of $v$ and
    reverses the two physical input coordinates. This is the exact two-site
    factorization through the paper gates $u$ and $v$ of CPSV17
    \texttt{SVDforms2}--\texttt{uuvv}.
\end{theorem}

\begin{proof}\leanok
    \uses{thm:mpu_source_factorization_entries}
    Expand the periodic trace through $M_1=X_1Y_1$ and $M_2=X_2Y_2$.
    Reordering the four scalar finite sums gives
    \begin{align}
        \sum_{l,r}v_{(i_1,i_2),(r,l)}\,u_{(l,r),(j_2,j_1)}.
    \end{align}
    Reindexing both physical configuration spaces and swapping $(l,r)$ to
    $(r,l)$ gives the matrix identity.
\end{proof}

\begin{theorem}[Open two-site factorization through $v$]
    \label{thm:mpu_source_v_open_two_site}
    \lean{MPOTensor.SourceFactors.sum_sourceV_mul_Y₁,
        MPOTensor.SourceFactors.mul_apply_eq_sum_sourceV_mul_Y₁_mul_Y₂,
        MPOTensor.SourceFactors.blockTwo_apply_eq_sum_sourceV_mul_Y₁_mul_Y₂,
        MPOTensor.sum_sourceV_mul_sourceY₁,
        MPOTensor.mul_apply_eq_sum_sourceV_mul_sourceY₁_mul_sourceY₂,
        MPOTensor.mul_apply_eq_sum_sourceV_mul_sourceY₁_mul_sourceY₂_pair,
        MPOTensor.blockTwo_apply_eq_sum_sourceV_mul_sourceY₁_mul_sourceY₂}
    \leanok
    \uses{def:mpu_source_uv, def:mpdo_two_site_blocking,
        thm:mpu_source_factorization_entries}
    For $q=(q_1,q_2)$ and $j=(j_1,j_2)$, the untraced two-site product is
    \begin{align}
        (\mathcal U^{q_1j_1}\mathcal U^{q_2j_2})_{\alpha,\gamma}
         & =\sum_{r,l}v_{q,(r,l)}(Y_1)_{r,(\alpha,j_1)}
        (Y_2)_{l,(j_2,\gamma)}.
    \end{align}
    The intermediate contraction is
    \begin{align}
        \sum_r v_{(q_1,q_2),(r,l)}(Y_1)_{r,(\alpha,j_1)}
         & =\sum_\beta \mathcal U^{q_1}_{j_1,\alpha,\beta}
        (X_2)_{(\beta,q_2),l}.
    \end{align}
\end{theorem}

\begin{proof}\leanok
    \uses{def:mpu_source_uv, thm:mpu_source_factorization_entries}
    Expand $v$ as the contraction of $X_1$ with $X_2$. Contracting $X_1$
    with $Y_1$ gives the first tensor letter, while contracting $X_2$ with
    $Y_2$ gives the second. Reordering the finite sums proves the open product
    identity. The concrete blocked formula follows by encoding both physical
    pairs through $\Fin d\times\Fin d\simeq\Fin(d^2)$.
\end{proof}

\begin{theorem}[Open two-site factorization through $u$]
    \label{thm:mpu_source_u_open_two_site}
    \lean{MPOTensor.SourceFactors.mul_apply_eq_sum_X₂_mul_sourceU_mul_X₁,
        MPOTensor.SourceFactors.blockTwo_apply_eq_sum_X₂_mul_sourceU_mul_X₁,
        MPOTensor.blockTwo_apply_eq_sum_sourceX₂_mul_sourceU_mul_sourceX₁}
    \leanok
    \uses{def:mpu_source_uv, def:mpdo_two_site_blocking,
        thm:mpu_source_factorization_entries}
    The same untraced two-site product also satisfies
    \begin{align}
        (\mathcal U^{i_1j_1}\mathcal U^{i_2j_2})_{\alpha,\gamma}
         & =\sum_{l,r}(X_2)_{(\alpha,i_1),l}
        u_{(l,r),(j_1,j_2)}(X_1)_{(i_2,\gamma),r}.
    \end{align}
    The product indices of $u$ are $(l,r)$ on the source side and
    $(j_1,j_2)$ on the physical side.
\end{theorem}

\begin{proof}\leanok
    \uses{def:mpu_source_uv, thm:mpu_source_factorization_entries}
    Expand the first tensor letter through $M_2=X_2Y_2$ and the second through
    $M_1=X_1Y_1$. The common auxiliary index contracts $Y_2$ with $Y_1$ to
    give $u$. Summing the two remaining source indices gives the formula. The
    concrete blocked identity follows by decoding its physical indices into
    ordered pairs.
\end{proof}

\begin{theorem}[Closed double-layer trace for the normalized input tail]
    \label{thm:mpu_normalized_input_tail_closed_trace}
    \lean{MPOTensor.normalized_mpo_input_tail_eq_closed_doubleLayer_trace}
    \leanok
    \uses{def:mpo_family, def:mpu_double_layer}
    Assume $d>0$. For any MPO tensor $\mathcal U$, tail length $K$, and
    $p,q\in\Fin d\times\Fin d$, let $W$ be its double layer and let
    $E=d^{-1}\sum_i W^{ii}$ be the normalized physical diagonal. Then
    \begin{align}
        d^{-K}\sum_{\tau,\eta}
        \overline{U^{(K+2)}_{\eta,(p,\tau)}}
        U^{(K+2)}_{\eta,(q,\tau)}
        =\tr\!\left[W^{p_1q_1}W^{p_2q_2}E^K\right].
    \end{align}
    Thus $p$ labels the starred input and $q$ the unstarred input. This is
    the closed input-tail expression occurring on the right-hand side of
    \cite[equation~\texttt{uUnitary} and lines~550--556]{Cirac2017MPU}.
\end{theorem}

\begin{proof}\leanok
    \uses{thm:mpu_closed_double_layer, thm:mpu_normalized_diagonal_word_expansion}
    Closing the common output word gives the double-layer identity
    \begin{align}
        \sum_\eta
        \overline{U^{(K+2)}_{\eta,\sigma}}
        U^{(K+2)}_{\eta,\sigma'}
         & =W^{(K+2)}_{\sigma,\sigma'}.
    \end{align}
    Set $\sigma=(p,\tau)$ and $\sigma'=(q,\tau)$. Splitting off the two
    retained sites rewrites the right-hand side as
    \begin{align}
        \tr\!\left[
            W^{p_1q_1}W^{p_2q_2}
            W^{\tau_1\tau_1}\cdots W^{\tau_K\tau_K}\right].
    \end{align}
    Finally,
    \begin{align}
        E^K
         & =d^{-K}\sum_\tau
        W^{\tau_1\tau_1}\cdots W^{\tau_K\tau_K}.
    \end{align}
    Substitute this expansion and use linearity of the trace.
\end{proof}

\begin{definition}[Tensorized source dressing]
    \label{def:mpu_source_v_dressing}
    \lean{MPOTensor.sourceYTensor,MPOTensor.sourceZTensor,
        MPOTensor.sourceYTensor_apply,MPOTensor.sourceZTensor_apply,
        MPOTensor.sourceVRegroupEquiv,
        MPOTensor.sourceVRegroupEquiv_apply,
        MPOTensor.sourceVRegroupEquiv_symm_apply,
        MPOTensor.SourceVDressedGram}
    \leanok
    \uses{def:mpu_weighted_source_factors, def:mpu_source_uv}
    Put
    \begin{align}
        Y & =Y_1\otimes Y_2,
          & Z                & =Z_1\otimes Z_2,
    \end{align}
    with source-bond order $r\times\ell$. Thus $Y$ maps
    $(\C^D\otimes\C^d)\otimes(\C^d\otimes\C^D)$ to
    $\C^r\otimes\C^\ell$, and $Z$ maps in the opposite direction. The
    regrouping of the open indices is
    \begin{align}
        R\bigl((a,i),(j,b)\bigr)
         & =\bigl((i,j),\operatorname{flat}(a,b)\bigr),
    \end{align}
    where $\operatorname{flat}:\Fin D\times\Fin D\simeq\Fin(D^2)$ is the
    canonical product enumeration. The dressed source-$v$ Gram equation is
    \begin{align}
        Y^\dagger(v^\dagger v)Y & =Y^\dagger Y.
    \end{align}
    These are the tensorized factors surrounding $v^\dagger v$ in
    \cite[Theorem~III.8, Section~III.B, equations~(31)--(32)]{Cirac2017MPU}.
\end{definition}

\begin{theorem}[Four-letter source contractions]
    \label{thm:mpu_source_v_four_u_expansions}
    \lean{MPOTensor.sourceY₁_gram_eq_weighted_sourceCutM₁_gram,
        MPOTensor.doubleLayerTensor_mul_apply_four_u,
        MPOTensor.doubleLayerTensor_rankOne_mul_apply_four_u,
        MPOTensor.sourceYTensor_gram_eq_four_u_weighted}
    \leanok
    \uses{def:mpu_source_v_dressing, def:mpu_double_layer, thm:mpu_source_factorizations, thm:mpu_source_factor_normalizations, thm:mpu_source_y_two_rotated_gram}
    Let $p=(p_1,p_2)$ be the starred physical pair,
    $q=(q_1,q_2)$ the unstarred pair, and let $a=(a_1,a_2)$ and
    $b=(b_1,b_2)$ be virtual paths. The weighted first-cut contraction is
    \begin{align}
        \sum_r \overline{(Y_1)_{r,(a_1,p_1)}}(Y_1)_{r,(b_1,q_1)}
         & =\sum_{i_1,\beta,\beta'}
        \overline{U^{i_1p_1}_{a_1\beta}}
        \rho_{\beta\beta'}U^{i_1q_1}_{b_1\beta'}.
    \end{align}
    The entry of two ordinary double-layer letters from row $(a_1,b_1)$ to
    column $(a_2,b_2)$ is
    \begin{align}
        \sum_{\alpha,\beta,j_1,j_2}
        \overline{U^{j_1p_1}_{a_1\alpha}}
        U^{j_1q_1}_{b_1\beta}
        \overline{U^{j_2p_2}_{\alpha a_2}}
        U^{j_2q_2}_{\beta b_2}.
    \end{align}
    Set
    \begin{align}
        \rho'_x & =\operatorname{vec}(\rho)_{\operatorname{flat}^{-1}(x)},  &
        \Phi'_x & =\operatorname{vec}(\Id_D)_{\operatorname{flat}^{-1}(x)}.
    \end{align}
    Inserting $\ket{\rho'}\!\bra{\Phi'}$ between the two letters gives
    \begin{align}
         & \bigl(W^{p_1q_1}\ket{\rho'}\!\bra{\Phi'}W^{p_2q_2}\bigr)
        _{(a_1,b_1),(a_2,b_2)}                                      \\
         & \quad=\sum_{\alpha,\alpha',\beta,j_1,j_2}
        \overline{U^{j_1p_1}_{a_1\alpha}}
        U^{j_1q_1}_{b_1\alpha'}\rho_{\alpha'\alpha}
        \overline{U^{j_2p_2}_{\beta a_2}}U^{j_2q_2}_{\beta b_2}.
    \end{align}
    Finally, the Gram entry of $Y_1\otimes Y_2$ expands as
    \begin{align}
        \sum_{i_1,\beta,\beta',\delta,i_2}
        \overline{U^{i_1p_1}_{a_1\beta}}
        \rho_{\beta\beta'}U^{i_1q_1}_{b_1\beta'}
        \overline{U^{i_2p_2}_{\delta a_2}}
        U^{i_2q_2}_{\delta b_2}.
    \end{align}
\end{theorem}

\begin{proof}
    \leanok
    \uses{def:mpu_source_v_dressing, def:mpu_double_layer, thm:mpu_source_factorizations, thm:mpu_source_factor_normalizations, thm:mpu_source_y_two_rotated_gram}
    Expand the ordinary and rank-one-inserted matrix products entrywise.
    For the tensorized Gram entry, use
    $X_1^\dagger(\Id_d\otimes\rho)X_1=\Id_r$ on the first cut and
    $X_2^\dagger X_2=\Id_\ell$ on the second cut, then reorder the finite
    sums. No ambient identity $X_2X_2^\dagger=\Id$ is used.
\end{proof}

% This theorem supplies the closed factor of the complete source-$v$ network
% in Theorem~\ref{thm:mpu_source_v_complete_contraction}.
\begin{theorem}[Paper-$v$ open-leg contraction]
    \label{thm:mpu_source_v_mul_source_y_tensor_entry}
    \lean{MPOTensor.sourceV_mul_sourceYTensor_apply}
    \leanok
    \uses{def:mpu_source_v_dressing, def:mpu_source_uv,
        thm:mpu_source_factorization_entries}
    Let $\rho>0$. For $z=(z_1,z_2)$, $p=(p_1,p_2)$, and
    $a=(a_1,a_2)$,
    \begin{align}
        \bigl(v(Y_1\otimes Y_2)\bigr)_{z,((a_1,p_1),(p_2,a_2))}
        =\sum_\gamma U^{z_1p_1}_{a_1\gamma}U^{z_2p_2}_{\gamma a_2}.
    \end{align}
    This is the local open-leg identity for the literal gate
    $v=X_1\mathbin{-}X_2$.
\end{theorem}

\begin{proof}\leanok
    \uses{def:mpu_source_v_dressing, def:mpu_source_uv,
        thm:mpu_source_factorization_entries}
    Expand the matrix product and tensor product, reorder the finite sums, and
    contract $X_1Y_1$ and $X_2Y_2$ by the two source-cut factorizations.
\end{proof}

% Formalization scope: the input-tail, diagonal-tail, and one-sided source-$v$
% identities above do not prove a dressed Gram identity by themselves.  The
% complete same-tensor network is Theorem~\ref{thm:mpu_source_v_complete_contraction}.

\begin{theorem}[Tensorized source right inverse]
    \label{thm:mpu_source_yz_tensor}
    \lean{MPOTensor.sourceYTensor_mul_sourceZTensor}
    \leanok
    \uses{def:mpu_source_v_dressing}
    The tensorized source factors satisfy
    \begin{align}
        YZ & =\Id_{r\ell}.
    \end{align}
\end{theorem}

\begin{proof}
    \leanok
    \uses{thm:mpu_source_factor_right_inverses}
    Multiplicativity of the Kronecker product gives
    \begin{align}
        YZ & =(Y_1Z_1)\otimes(Y_2Z_2)
        =\Id_r\otimes\Id_\ell=\Id_{r\ell}.
    \end{align}
\end{proof}

\begin{theorem}[Regrouped entries of the dressed Gram equation]
    \label{thm:mpu_source_v_dressed_entries}
    \lean{MPOTensor.sourceVDressedGram_iff_regrouped_entries}
    \leanok
    \uses{def:mpu_source_v_dressing, def:mpu_source_uv}
    For every positive definite source weight $\rho$, and for
    $x=R^{-1}(p,a)$ and $y=R^{-1}(q,b)$, the dressed Gram equation is
    equivalent to
    \begin{align}
        \sum_t\left(\sum_s\overline{Y_{s,x}}
        \sum_z\overline{v_{z,s}}v_{z,t}\right)Y_{t,y}
         & =\sum_t\overline{Y_{t,x}}Y_{t,y}.
    \end{align}
    Here $p,q\in\Fin d\times\Fin d$, $a,b\in\Fin(D^2)$, the indices $s,t$
    range over $\Fin r\times\Fin\ell$, and $z$ ranges over
    $\Fin d\times\Fin d$.
\end{theorem}

\begin{proof}
    \leanok
    \uses{def:mpu_source_v_dressing}
    Evaluate both sides of $Y^\dagger(v^\dagger v)Y=Y^\dagger Y$ at
    $(x,y)$ and expand the three matrix products. Conversely, apply the
    displayed equality after regrouping each pair of open source-cut indices
    by $R$.
\end{proof}

\begin{theorem}[Right-inverse cancellation of a dressed Gram equation]
    \label{thm:mpu_dressed_gram_cancel}
    \lean{Matrix.conjTranspose_mul_mul_eq_conjTranspose_mul_iff_of_mul_eq_one}
    \leanok
    If $Y$ has a right inverse $Z$, so that $YZ=\Id$, then for every square
    matrix $G$ on the row space of $Y$,
    \begin{align}
        Y^\dagger GY=Y^\dagger Y
         & \quad\Longleftrightarrow\quad G=\Id.
    \end{align}
\end{theorem}

\begin{proof}
    \leanok
    Conjugate the first equality by $Z^\dagger$ on the left and $Z$ on the
    right. Since $Z^\dagger Y^\dagger=(YZ)^\dagger=\Id$, both outer factors
    cancel. The reverse implication follows by substituting $G=\Id$.
\end{proof}

\begin{theorem}[Cancellation of the source-$v$ dressing]
    \label{thm:mpu_source_v_dressed_gram_cancel}
    \lean{MPOTensor.sourceVDressedGram_iff_isIsometry}
    \leanok
    \uses{def:mpu_source_v_dressing, def:mpu_source_uv}
    For every positive definite source weight $\rho$,
    \begin{align}
        Y^\dagger(v^\dagger v)Y=Y^\dagger Y
         & \quad\Longleftrightarrow\quad v^\dagger v=\Id_{r\ell}.
    \end{align}
\end{theorem}

\begin{proof}
    \leanok
    \uses{thm:mpu_source_yz_tensor, thm:mpu_dressed_gram_cancel}
    Apply right-inverse cancellation with $G=v^\dagger v$ and
    $Z=Z_1\otimes Z_2$.
\end{proof}

\begin{theorem}[Right-rank bound]
    \label{thm:mpu_right_rank_bound}
    \lean{MPOTensor.rightRank_bound}
    \leanok
    \uses{def:mpu_right_rank}
    If the physical and auxiliary dimensions are $d$ and $D$, respectively,
    then
    \begin{align}
        r & \leq Dd.
    \end{align}
\end{theorem}

\begin{proof}
    \leanok
    \uses{def:mpu_source_matrix_one}
    The matrix $M_1$ has $dD$ columns. Its rank is at most its number of
    columns, hence at most $dD=Dd$.
\end{proof}

\begin{theorem}[Left-rank bound]
    \label{thm:mpu_left_rank_bound}
    \lean{MPOTensor.leftRank_bound}
    \leanok
    \uses{def:mpu_left_rank}
    If the physical and auxiliary dimensions are $d$ and $D$, respectively,
    then
    \begin{align}
        \ell & \leq Dd.
    \end{align}
\end{theorem}

\begin{proof}
    \leanok
    \uses{def:mpu_source_matrix_two}
    The matrix $M_2$ has $dD$ columns. Its rank is at most its number of
    columns, hence at most $dD=Dd$.
\end{proof}

\section{Full-support canonical representatives}

\begin{definition}[Full support]
    \label{def:mpu_full_support}
    \lean{MPSTensor.CPSVCanonicalFormData.HasFullSupport}
    \leanok
    \uses{def:cpsv_canonical_form}
    Let
    \begin{align}
        A^i = V^*\left(\bigoplus_{k=1}^r \mu_k A_k^i\right)V
    \end{align}
    be CPSV canonical-form data in ambient bond dimension $D$, where every
    $A_k$ has positive bond dimension $D_k$.  The data have \emph{full
        support} when
    \begin{align}
        \sum_{k=1}^rD_k=D.
    \end{align}
\end{definition}

\begin{theorem}[Reduced canonical representative of the normalized tensor]
    \label{thm:mpu_reduced_normalized_flattening_cfii}
    \lean{MPOTensor.IsMPU.exists_reduced_normalizedFlattening_cfii}
    \leanok
    \uses{def:mpu, def:mpu_normalized_flattening, def:cpsv_cfii, def:mpu_full_support, def:normal, def:same_mpv2_pos}
    Let $d,D>0$, let $U$ be an MPU tensor of physical dimension $d$ and
    bond dimension $D$, and let $\widetilde U=d^{-1/2}U$ be its normalized
    doubled-index tensor. There are an integer $D_{\mathrm{red}}>0$, a normal
    tensor $A_{\mathrm{red}}$ of bond dimension $D_{\mathrm{red}}$, and
    canonical-form-II data for $A_{\mathrm{red}}$ such that, for every
    integer $N>0$ and every word $\sigma$ of length $N$,
    \begin{align}
        D_{\mathrm{red}}  & \leq D,                        \\
        V^{(N)}(A_{\mathrm{red}})_\sigma
                          & =V^{(N)}(\widetilde U)_\sigma, \\
        \sum_{k=1}^{r}D_k & =D_{\mathrm{red}}.
    \end{align}
    This organizes the dimension-bounded canonical replacement used in
    \cite[lines~257--294 and 319--326]{Cirac2017MPU}. Its irreducible
    reduction and canonical-form-II gauge are those of
    \cite[lines~195--255 and 1058--1077]{Cirac2016MPDO_arXiv}. The shifted
    transfer argument is the one used in
    \cite[lines~344--356]{Cirac2017MPU}.
    % Source role: replacement infrastructure, not a separately named source theorem.
\end{theorem}

\begin{proof}\leanok
    \uses{def:cpsv_canonical_form, thm:nonzero_irreducible_block_decomposition, lem:overlap_eq_of_pos_mpv, thm:mpu_transfer_matrix_overlap, thm:mpu_normalized_transfer_trace, thm:mpu_transfer_matrix_eigenvalues, thm:mpu_dependent_block_inclusion_identities, lem:transfer_map_corner_intertwine, lem:irreducible_transfer_perron_pair, lem:irreducible_block_spectral_normalization, thm:mpu_shifted_trace_power_spectrum, thm:mpu_shifted_transfer_trace_single_block, thm:mpu_shifted_transfer_trace_primitive, thm:mpu_unique_full_support_normal, thm:cpsv_cf_to_cfii_gauge, thm:gauge_invariance}
    Remove the zero irreducible blocks of $\widetilde U$ and write
    \begin{align}
        A_0^i & =\bigoplus_{k=1}^{r}B_k^i,
              & \sum_{k=1}^{r}D_k              & \leq D, \\
        V^{(N)}(A_0)_\sigma
              & =V^{(N)}(\widetilde U)_\sigma,
    \end{align}
    for every $N>0$ and every word $\sigma$ of length $N$.
    At $N=2$ the common self-overlap is one, so the retained sum is nonempty
    and $D_0=\sum_kD_k>0$. For every integer $N>1$, overlap invariance gives
    \begin{align}
        \tr(E_{A_0}^{\,N})
         & =\sum_\sigma\left|V^{(N)}(A_0)_\sigma\right|^2
        =\sum_\sigma\left|V^{(N)}(\widetilde U)_\sigma\right|^2
        =1.
        \label{eq:mpu_reduced_trace}
    \end{align}
    Let $J_k:\mathbb C^{D_k}\to\mathbb C^{D_0}$ be the inclusion of the
    $k$th block. Then
    \begin{align}
        J_k^\dagger J_k & =\Id,
                        & A_0^iJ_k                    & =J_kB_k^i, \\
        \E_{A_0}(J_kXJ_k^\dagger)
                        & =J_k\E_{B_k}(X)J_k^\dagger.
    \end{align}
    Thus every nonzero transfer eigenvalue of $B_k$ is a nonzero eigenvalue of
    $E_{A_0}$ and hence equals one by
    (\ref{eq:mpu_reduced_trace}). Perron--Frobenius theory supplies
    $\rho_k>0$ and $s_k>0$ with
    $\E_{B_k}(\rho_k)=s_k\rho_k$. Therefore $s_k=1$, and every peripheral
    eigenvalue of $B_k$ is also one. Each $B_k$ is consequently normal.

    The direct sum now has literal canonical-form data with full support.
    Apply the nonsingular blockwise canonical-form-II gauge. It preserves the
    block dimensions and all periodic vectors, so the resulting tensor
    $A_{\mathrm{red}}$ still satisfies
    \begin{align}
        \sum_{k=1}^{r}D_k & =D_{\mathrm{red}},
                          & \tr(E_{A_{\mathrm{red}}}^{\,N}) & =1,
    \end{align}
    for every integer $N>1$.
    The shifted traces force $r=1$, spectral radius one, and primitivity.
    Full support identifies the unique block with the ambient bond space, so
    $A_{\mathrm{red}}$ is normal.
\end{proof}

\begin{theorem}[Reduced canonical representative of an MPU]
    \label{thm:mpu_reduced_cfii_representative}
    \lean{MPOTensor.IsMPU.exists_reduced_cfii_representative}
    \leanok
    \uses{def:mpu, def:mpu_normalized_flattening, def:cpsv_cfii, def:mpu_full_support}
    Let $d,D>0$ and let $U$ be an MPU tensor. There are
    $D_{\mathrm{red}}>0$ and an MPU tensor $U_{\mathrm{red}}$ of bond
    dimension $D_{\mathrm{red}}$ such that, for every integer $N>0$,
    \begin{align}
        D_{\mathrm{red}}       & \leq D,   \\
        U_{\mathrm{red}}^{(N)} & =U^{(N)}.
    \end{align}
    The normalized doubled-index tensor of $U_{\mathrm{red}}$ has
    canonical-form-II data whose retained block dimensions sum to
    $D_{\mathrm{red}}$. No equality is asserted at length zero, and no
    virtual gauge is asserted between the two potentially unequal-dimensional
    bond spaces.
    % Source role: MPU corollary of the canonical-replacement infrastructure.
\end{theorem}

\begin{proof}\leanok
    \uses{thm:mpu_reduced_normalized_flattening_cfii, lem:mpv_smul, lem:mpdo_three_first_site_contraction_coefficients}
    Apply Theorem~\ref{thm:mpu_reduced_normalized_flattening_cfii} and write
    its representative as $A_{\mathrm{red}}$. Define
    \begin{align}
        U_{\mathrm{red}}^{ij}
         & =\sqrt d\,A_{\mathrm{red}}^{(i,j)}.
    \end{align}
    Since $d>0$, normalization recovers $A_{\mathrm{red}}$ exactly. For every
    $N>0$ and every ket--bra pair $(\sigma,\tau)$, the doubled-index
    coefficient identity gives
    \begin{align}
        (U_{\mathrm{red}}^{(N)})_{\sigma,\tau}
         & =d^{N/2}V^{(N)}(A_{\mathrm{red}})_{(\sigma,\tau)} \\
         & =d^{N/2}V^{(N)}(\widetilde U)_{(\sigma,\tau)}     \\
         & =(U^{(N)})_{\sigma,\tau}.
    \end{align}
    Hence the periodic operators agree at every positive length. For $N>1$
    the right side is unitary, so $U_{\mathrm{red}}$ is an MPU.
\end{proof}

\begin{theorem}[Full support under physical blocking]
    \label{thm:mpu_full_support_blocking}
    \lean{MPSTensor.CPSVCanonicalFormData.hasFullSupport_blockTensor}
    \leanok
    \uses{def:mpu_full_support, thm:cpsv_literal_canonical_form_positive_blocking}
    Suppose that canonical-form data have full support:
    \begin{align}
        \sum_{k=1}^rD_k & =D.
    \end{align}
    Then every positive physical blocking of these data also has full
    support.
\end{theorem}

\begin{proof}\leanok
    \uses{def:mpu_full_support, thm:cpsv_literal_canonical_form_positive_blocking}
    Positive blocking supplies canonical-form data with the same block
    dimensions, so the full-support sum is unchanged.
\end{proof}

\begin{theorem}[Full support under physical relabeling]
    \label{thm:mpu_full_support_reindexing}
    \lean{MPSTensor.CPSVCanonicalFormData.hasFullSupport_reindexPhysical}
    \leanok
    \uses{def:mpu_full_support, thm:cpsv_physical_reindexing}
    A bijective relabeling of the physical alphabet preserves full
    support of canonical-form data.
\end{theorem}

\begin{proof}\leanok
    \uses{def:mpu_full_support, thm:cpsv_physical_reindexing}
    Physical relabeling leaves the blocks and their bond
    dimensions unchanged, so the full-support sum is unchanged.
\end{proof}

\begin{theorem}[An ambient-dimensional block is irreducible]
    \label{thm:mpu_block_ambient_irreducible}
    \lean{MPSTensor.CPSVCanonicalFormData.isIrreducibleTensor_of_dim_eq}
    \leanok
    \uses{def:cpsv_canonical_form}
    If a block has bond dimension $D$, then the ambient tensor $A$ is
    irreducible.
\end{theorem}

\begin{proof}\leanok
    \uses{def:cpsv_canonical_form, thm:irreducible_tensor_smul_conj}
    If $k$ is the block with $D_k=D$, the total-dimension bound gives
    \begin{align}
        D = D_k \leq \sum_{j=1}^r D_j \leq D.
    \end{align}
    Since every retained block dimension is positive, these equalities force
    $k$ to be the only retained block.  After reindexing its bond coordinates,
    the ambient coisometry $U$ is square, hence unitary, and reconstruction reads
    \begin{align}
        A^i = U^*\bigl(\mu_k B^i\bigr)U,
    \end{align}
    where $B$ is the reindexed normal block and $\mu_k\ne0$.  Thus
    Theorem~\ref{thm:irreducible_tensor_smul_conj} transports irreducibility of
    $B$ through the nonzero rescaling and unitary conjugation.
\end{proof}

\begin{theorem}[Unique block with full support]
    \label{thm:mpu_unique_full_support_irreducible}
    \lean{MPSTensor.CPSVCanonicalFormData.isIrreducibleTensor_of_r_eq_one_of_fullSupport}
    \leanok
    \uses{def:mpu_full_support}
    Canonical-form data with one block and full support reconstruct
    an irreducible ambient tensor.
\end{theorem}

\begin{proof}\leanok
    \uses{thm:mpu_block_ambient_irreducible}
    Let $k$ be the unique block.  The singleton sum and full-support
    identity give
    \begin{align}
        \sum_{j=1}^{r} D_j = D_k = D.
    \end{align}
    Thus $k$ has ambient bond dimension, so
    Theorem~\ref{thm:mpu_block_ambient_irreducible} applies.
\end{proof}

\begin{theorem}[Normality of a unique full-support block]
    \label{thm:mpu_unique_full_support_normal}
    \lean{MPSTensor.CPSVCanonicalFormData.isNormalTensor_of_r_eq_one_of_fullSupport}
    \leanok
    \uses{def:mpu_full_support}
    Suppose that canonical-form data have one block and full
    support.  If the transfer map has spectral radius one and is primitive,
    then the ambient tensor is a normal tensor.
\end{theorem}

\begin{proof}\leanok
    \uses{thm:mpu_unique_full_support_irreducible}
    Combine the preceding irreducibility conclusion with the supplied
    spectral-radius and primitivity properties.
\end{proof}

\section{Transfer-matrix identities}

Let $D$ be a positive integer and let
$T:M_D(\mathbb C)\to M_D(\mathbb C)$ be a linear map. The transfer matrix
$\widehat T$ of $T$ under the column-stacking identification
$M_D(\mathbb C)\cong\mathbb C^{D^2}$ is defined in the companion
quantum-channel volume~\cite[``Transfer matrix of a matrix endomorphism'']{QICLean2026Blueprint}, the matrix-units
special case of the one-family transfer matrix of the corresponding
definition in the companion quantum-channel
volume~\cite[``Transfer matrix in one family'']{QICLean2026Blueprint}.


The corresponding generic result is proved in the companion quantum-channel volume~\cite[``Functoriality of transfer matrices'']{QICLean2026Blueprint}.




The corresponding generic result is proved in the companion quantum-channel volume~\cite[``Trace preservation fixes the vectorized identity on the left'']{QICLean2026Blueprint}.



The corresponding generic result is proved in the companion quantum-channel volume~\cite[``Powers and traces of transfer matrices'']{QICLean2026Blueprint}.


\begin{theorem}[Eigenvalues of a transfer matrix]
    \label{thm:mpu_transfer_matrix_eigenvalues}
    \lean{transferMatrix_hasEigenvalue_iff}
    \leanok
    A complex number $\mu$ is an eigenvalue of $T$ if and only if it is an
    eigenvalue of $\widehat T$.
\end{theorem}

\begin{proof}\leanok

    Column stacking is a linear bijection and intertwines $T$ with
    multiplication by $\widehat T$.  It therefore carries nonzero
    eigenmatrices to nonzero eigenvectors and conversely.
\end{proof}

For a matrix product tensor $A$, write $V^{(N)}(A)_\sigma$ for the
coefficient of its length-$N$ periodic vector at the configuration $\sigma$.

\begin{theorem}[Transfer-matrix powers as periodic overlaps]
    \label{thm:mpu_transfer_matrix_overlap}
    \lean{MPSTensor.trace_transferMatrix_mixedMapLM_pow_eq_mpvOverlap,
        MPSTensor.trace_transferMatrix_transferMap_pow_eq_mpvOverlap}
    \leanok
    \uses{def:mixed_transfer, def:transfer_map, def:mpv_overlap, def:mps_tensor}
    Let $A$ and $B$ be matrix product tensors with the same positive bond
    dimension, and let $F_{A,B}(X)=\sum_i A^iX(B^i)^\dagger$ be their mixed
    transfer map.  Then
    \begin{align}
        \tr\!\left(\widehat{F_{A,B}}^{\,N}\right)
         & =\sum_\sigma V^{(N)}(A)_\sigma
        \overline{V^{(N)}(B)_\sigma}.
    \end{align}
    In particular, taking $B=A$ identifies the trace of the $N$th power of
    the ordinary transfer matrix with the length-$N$ self-overlap.
\end{theorem}

\begin{proof}\leanok
    \uses{thm:overlap_trace, thm:mixed_self}
    Replace the matrix power by the transfer matrix of the corresponding
    operator power, identify matrix trace with operator trace, and expand the
    latter over periodic words.  For the ordinary self-transfer statement,
    specialize to $B=A$ and use the identity $F_{A,A}=\mathcal E_A$ from
    Theorem~\ref{thm:mixed_self}.
\end{proof}

\begin{definition}[Normalized flattening of an MPO tensor]
    \label{def:mpu_normalized_flattening}
    \lean{MPOTensor.normalizedFlattening}
    \leanok
    \uses{def:mpo_tensor, def:mps_tensor, def:mpo_to_mps}
    For an MPO tensor $U$ of physical dimension $d>0$, its normalized
    flattening is the doubled-physical-index MPS tensor
    \begin{align}
        \widetilde U^{ij} & =\frac{1}{\sqrt d}\,U^{ij}.
    \end{align}
    This is the normalization used in the normal-tensor argument of
    \cite[Section~III, equation~(13)]{Cirac2017MPU}.
\end{definition}

\begin{definition}[Doubled product-coordinate equivalence]
    \label{def:mpu_doubled_product_coordinate_equiv}
    \lean{finDoubledProdEquiv, finDoubledProdEquiv_apply,
        finDoubledProdEquiv_symm_apply}
    \leanok
    Write
    $\pi_{m,n}:\Fin m\times\Fin n\to\Fin{mn}$ for the standard product
    coordinate bijection.  For natural numbers $d,e$, the canonical doubled
    product-coordinate equivalence
    $s_{d,e}:\Fin{(de)^2}\to\Fin{d^2e^2}$ and its inverse satisfy
    \begin{align}
        s_{d,e}\!\left(\pi_{de,de}
        \bigl(\pi_{d,e}(i,k),\pi_{d,e}(j,\ell)\bigr)\right)
         & =\pi_{d^2,e^2}
        \bigl(\pi_{d,d}(i,j),\pi_{e,e}(k,\ell)\bigr), \notag \\
        s_{d,e}^{-1}\!\left(\pi_{d^2,e^2}
        \bigl(\pi_{d,d}(i,j),\pi_{e,e}(k,\ell)\bigr)\right)
         & =\pi_{de,de}
        \bigl(\pi_{d,e}(i,k),\pi_{d,e}(j,\ell)\bigr).
        \label{eq:mpu_tensor_flattening_shuffle}
    \end{align}
\end{definition}

\begin{theorem}[Normalized flattening of an independent tensor product]
    \label{thm:mpu_normalized_flattening_tensor_product}
    \lean{MPOTensor.normalizedFlattening_tensorProduct}
    \leanok
    \uses{def:mpu_doubled_product_coordinate_equiv,
        def:mpu_normalized_flattening,
        def:mpu_independent_tensor_product,
        def:mps_independent_tensor_product}
    Let $\mathcal U$ and $\mathcal V$ have physical dimensions $d,e$ and
    auxiliary dimensions $D,E$, respectively.  With the coordinate
    equivalence $s_{d,e}$ of
    Definition~\ref{def:mpu_doubled_product_coordinate_equiv}, normalized
    flattening obeys
    \begin{align}
        \widetilde{\mathcal U\boxtimes\mathcal V}
         & =s_{d,e}^{*}
        \bigl(\widetilde{\mathcal U}\boxtimes
        \widetilde{\mathcal V}\bigr), \notag            \\
        \left(\widetilde{\mathcal U\boxtimes\mathcal V}\right)^{
                \pi_{d,e}(i,k),\pi_{d,e}(j,\ell)}_{
                \pi_{D,E}(\alpha,\gamma),\pi_{D,E}(\beta,\delta)}
         & =\widetilde{\mathcal U}^{i,j}_{\alpha,\beta}
        \widetilde{\mathcal V}^{k,\ell}_{\gamma,\delta}.
        \label{eq:mpu_tensor_normalized_flattening}
    \end{align}
    For positive physical dimensions this uses the normalization of
    \cite[Section~III, equation~(13)]{Cirac2017MPU}.  With the totalized
    inverse convention $0^{-1}=0$, the algebraic identity also holds when
    $d=0$ or $e=0$.  This identity makes explicit the normalized-flattening
    compatibility corresponding to the tensoring clause in
    \cite[Theorem~IV.6(ii), lines~824--845]{Cirac2017MPU}; the paper does not
    state it separately.
\end{theorem}

\begin{proof}\leanok
    \uses{def:mpu_doubled_product_coordinate_equiv,
        def:mpu_normalized_flattening,
        def:mpu_independent_tensor_product,
        def:mps_independent_tensor_product}
    Apply the two product-coordinate entry formulas after the regrouping in
    (\ref{eq:mpu_tensor_flattening_shuffle}).  The scalar normalizations agree
    because
    \begin{align}
        \bigl(\sqrt{de}\bigr)^{-1}
         & =\bigl(\sqrt d\sqrt e\bigr)^{-1}
        =\bigl(\sqrt d\bigr)^{-1}\bigl(\sqrt e\bigr)^{-1}.
    \end{align}
    Multiplying this identity by
    $\mathcal U^{i,j}_{\alpha,\beta}
        \mathcal V^{k,\ell}_{\gamma,\delta}$ gives
    (\ref{eq:mpu_tensor_normalized_flattening}).
\end{proof}

\begin{theorem}[Canonical-form-II data under independent tensor products]
    \label{thm:mpu_tensor_product_cfii_data}
    \lean{MPSTensor.transferMap_tensorProduct_kronecker,
        MPSTensor.CPSVCanonicalFormIIData.tensorProduct,
        MPSTensor.CPSVCanonicalFormIIData.hasFullSupport_tensorProduct,
        MPSTensor.CPSVCanonicalFormIIData.hasFullSupport_cast,
        MPOTensor.tensorProductCFIIData,
        MPOTensor.hasFullSupport_tensorProductCFIIData}
    \leanok
    \uses{def:cpsv_cfii, def:mpu_full_support,
        thm:mps_normal_block_tensor_product,
        thm:mpu_normalized_flattening_tensor_product,
        thm:cpsv_physical_reindexing}
    Let $A$ and $B$ have chosen canonical-form-II data with retained blocks
    $A_a$ and $B_b$, dimensions $D_a$ and $E_b$, weights $\mu_a$ and $\nu_b$,
    diagonal positive-definite fixed points $\Lambda_a$ and $\Gamma_b$, and
    ambient coisometries $V_A$ and $V_B$. Write $\pi_{m,n}$ for the standard
    product-coordinate equivalence and let $\chi$ be the retained-coordinate
    equivalence
    \begin{align}
        \chi\bigl((a,\alpha),(b,\beta)\bigr)
         & =\bigl(\pi_{r_A,r_B}(a,b),
        \pi_{D_a,E_b}(\alpha,\beta)\bigr).
    \end{align}
    The product ambient coisometry is
    \begin{align}
        W
         & =\operatorname{reind}_{\chi,\pi_{D,E}}(V_A\otimes V_B).
    \end{align}
    For each retained pair $(a,b)$, the chosen product data are
    \begin{align}
        D_{(a,b)} & =D_aE_b,
    \end{align}
    \begin{align}
        \omega_{(a,b)} & =\mu_a\nu_b,
    \end{align}
    \begin{align}
        C_{(a,b)}^{\pi_{d,e}(i,k)}
         & =\operatorname{reind}_{\pi_{D_a,E_b},\pi_{D_a,E_b}}
        (A_a^i\otimes B_b^k),
    \end{align}
    and
    \begin{align}
        \Lambda_{(a,b)}
         & =\operatorname{reind}_{\pi_{D_a,E_b},\pi_{D_a,E_b}}
        (\Lambda_a\otimes\Gamma_b).
    \end{align}
    The exact reconstruction is
    \begin{align}
        A\boxtimes B
         & =W^\dagger
        \left(\bigoplus_{a,b}\mu_a\nu_b
            (A_a\boxtimes B_b)\right)W.
        \label{eq:mpu_tensor_product_cfii_reconstruction}
    \end{align}
    Thus the retained direct sum is conjugated by the reindexed Kronecker
    ambient coisometry, rather than identified directly with the ambient
    tensor. For MPO tensors $\mathcal U$ and $\mathcal V$, the doubled physical
    equivalence of
    Theorem~\ref{thm:mpu_normalized_flattening_tensor_product} identifies these
    data with canonical-form-II data for
    $\widetilde{\mathcal U\boxtimes\mathcal V}$.

    Full support is a supplied reduced-representative restriction. If the two
    chosen input data satisfy
    \begin{align}
        \sum_a D_a & =D,
    \end{align}
    and
    \begin{align}
        \sum_b E_b & =E,
    \end{align}
    then the chosen product data satisfy
    \begin{align}
        \sum_{a,b}D_aE_b
         & =\left(\sum_aD_a\right)\left(\sum_bE_b\right)
        =DE.
    \end{align}
    Identifying canonical-form-II data along an equality of ambient tensors
    leaves this retained-dimension sum unchanged. This full-support assertion
    is not stated by the tensoring sentence in
    \cite[Theorem~IV.6(ii)]{Cirac2017MPU}; it is the supplied restriction
    documented in \cite{gap:mpu_canonical_form_full_support}.

    The data considered here do not include the separate weight-normalization
    condition $|\mu_a|\leq1$ together with the existence of a unit-modulus
    weight. For the same independent tensor product, the raw source ranks
    satisfy
    $r(\mathcal U\boxtimes\mathcal V)=r(\mathcal U)r(\mathcal V)$ and
    $\ell(\mathcal U\boxtimes\mathcal V)=\ell(\mathcal U)\ell(\mathcal V)$ by
    Theorem~\ref{thm:mpu_tensor_product_source_ranks}. These rank identities do
    not require the chosen canonical-form-II data or full-support hypotheses
    above, while the canonical-form construction does not use source ranks.
    Tensor-product additivity of the index is a separate conclusion.
\end{theorem}

\begin{proof}\leanok
    \uses{thm:mps_normal_block_tensor_product,
        thm:mpu_normalized_flattening_tensor_product,
        thm:cpsv_physical_reindexing}
    For physical labels $i$ and $k$, set
    \begin{align}
        \mathsf B_A^i & =\bigoplus_a\mu_aA_a^i,
    \end{align}
    and
    \begin{align}
        \mathsf B_B^k & =\bigoplus_b\nu_bB_b^k.
    \end{align}
    Regrouping retained coordinates by $\chi$ gives the weighted-block-diagonal
    identity
    \begin{align}
        \bigoplus_{a,b}\mu_a\nu_b
        C_{(a,b)}^{\pi_{d,e}(i,k)}
         & =\operatorname{reind}_{\chi,\chi}
        (\mathsf B_A^i\otimes\mathsf B_B^k).
        \label{eq:mpu_tensor_product_weighted_block_diagonal}
    \end{align}
    Substitute $A^i=V_A^\dagger\mathsf B_A^iV_A$ and
    $B^k=V_B^\dagger\mathsf B_B^kV_B$. The mixed-product law for Kronecker
    products and the definition of $W$ then give
    equation~\eqref{eq:mpu_tensor_product_cfii_reconstruction}. Moreover,
    $W W^\dagger=1$ follows from
    $V_AV_A^\dagger=1$ and $V_BV_B^\dagger=1$. Left-canonicality and normality
    follow from Theorem~\ref{thm:mps_normal_block_tensor_product}.

    For arbitrary matrices $X$ and $Y$, direct expansion of the finite Kraus
    sums gives
    \begin{align}
        \mathcal E_{A\boxtimes B}
        \left(\operatorname{reind}_{\pi_{D,E},\pi_{D,E}}(X\otimes Y)\right)
         & =\operatorname{reind}_{\pi_{D,E},\pi_{D,E}}
        \bigl(\mathcal E_A(X)\otimes\mathcal E_B(Y)\bigr).
    \end{align}
    Taking $X=\Lambda_a$ and $Y=\Gamma_b$ proves the fixed-point equation;
    Kronecker products and bijective reindexing preserve positive definiteness
    and diagonality. The finite double-sum identity proves the supplied
    full-support conclusion. Finally apply the normalized-flattening identity
    and physical relabeling to obtain the MPO statement.
\end{proof}

\begin{lemma}[Normalized flattening of a composition tensor]
    \label{lem:mpu_normalized_flattening_composition}
    \lean{MPOTensor.normalizedFlattening_mulTensor_apply}
    \leanok
    \uses{def:mpu_normalized_flattening, def:mpdo_operator_product}
    Let $\mathcal U$ and $\mathcal V$ be MPO tensors with the same positive
    physical dimension $d$, and let
    \begin{align}
        W^{ik} & =\sum_{j=0}^{d-1}U^{ij}\otimes V^{jk}
    \end{align}
    be their composition tensor. Under the canonical identification of its
    auxiliary space with $\C^{D_U}\otimes\C^{D_V}$, the normalized
    flattenings satisfy
    \begin{align}
        \widetilde W^{ik}
         & =\sqrt d\sum_{j=0}^{d-1}
        \widetilde U^{ij}\otimes\widetilde V^{jk}.
    \end{align}
    This identity combines the raw composition tensor used in the proof of
    \cite[Theorem~IV.6(ii)]{Cirac2017MPU} with the normalization in
    \cite[Section~III, equation~(13)]{Cirac2017MPU}; it is not a separately
    stated result of that proof.
\end{lemma}

\begin{proof}\leanok
    \uses{def:mpu_normalized_flattening, def:mpdo_operator_product}
    Substitute $U^{ij}=\sqrt d\,\widetilde U^{ij}$ and
    $V^{jk}=\sqrt d\,\widetilde V^{jk}$ into the defining contraction for
    $W^{ik}$, and then multiply the result by $d^{-1/2}$.
\end{proof}

\begin{theorem}[Normalized flattening under blocking]
    \label{thm:mpu_normalized_flattening_blocking}
    \lean{MPOTensor.normalizedFlattening_blockTensor}
    \leanok
    \uses{def:mpu_normalized_flattening, thm:mpdo_blocking_doubled_index}
    Let $\mathcal U$ be an MPO tensor of physical dimension $d$, and let
    $L\geq0$. For the normalized-tensor interpretation assume $d>0$, as in
    Definition~\ref{def:mpu_normalized_flattening}; with the convention
    $0^{-1}=0$, the algebraic identity below also holds for $d=0$. Pairing
    the ket and bra letters at each site gives a canonical bijection
    \begin{align}
        \vartheta_L:\operatorname{Fin}(d^L)^2
         & \longrightarrow \operatorname{Fin}(d^2)^L.
    \end{align}
    If $\widetilde{\mathcal U}=d^{-1/2}\mathcal U$ denotes normalized
    flattening, then
    \begin{align}
        \widetilde{\mathcal U^{[L]}}
         & =\vartheta_L^*\bigl(\widetilde{\mathcal U}^{[L]}\bigr).
        \label{eq:mpu_normalized_flattening_blocking}
    \end{align}
    Thus normalized flattening commutes with blocking, up to the canonical
    doubled-index relabeling. This is the normalized form of the blocking
    operation in \cite[Definition~II.5 and equation~(9)]{Cirac2017MPU}.
\end{theorem}

\begin{proof}\leanok
    \uses{def:mpu_normalized_flattening, thm:mpdo_blocking_doubled_index}
    A blocked letter is a product of $L$ original letters. The identity
    \begin{align}
        \bigl(\sqrt{d^L}\bigr)^{-1}
         & =\bigl((\sqrt d)^{-1}\bigr)^L
    \end{align}
    identifies its normalization with the product of the $L$ one-site
    normalizations. Theorem~\ref{thm:mpdo_blocking_doubled_index}
    identifies the two physical indexings for every $L\geq0$; at $L=0$ both
    blocked words are empty and both local matrix products are the identity.
\end{proof}

\begin{theorem}[Left-canonical normalization under blocking and relabeling]
    \label{thm:mpu_left_canonical_blocking_reindexing}
    \lean{MPSTensor.leftCanonical_reindexPhysical_equiv}
    \leanok
    \uses{def:blocked_tensor, def:left_canonical_tensor, thm:left_canonical_blocking}
    Suppose that an MPS tensor $A$ is left-canonical:
    \begin{align}
        \sum_i (A^i)^\dagger A^i & =\Id.
    \end{align}
    Then for every $L\geq0$ its blocked tensor is left-canonical:
    \begin{align}
        \sum_I \bigl((A^{[L]})^I\bigr)^\dagger(A^{[L]})^I & =\Id.
        \label{eq:mpu_left_canonical_blocking}
    \end{align}
    If $e$ is a bijection of physical alphabets, then $e^*A$ is
    left-canonical if and only if $A$ is:
    \begin{align}
        \sum_j \bigl((e^*A)^j\bigr)^\dagger(e^*A)^j=\Id
        \quad\Longleftrightarrow\quad
        \sum_i(A^i)^\dagger A^i=\Id.
        \label{eq:mpu_left_canonical_reindexing}
    \end{align}
\end{theorem}

\begin{proof}\leanok
    \uses{def:blocked_tensor, def:left_canonical_tensor, thm:left_canonical_blocking}
    Expand~(\ref{eq:mpu_left_canonical_blocking}) as a sum over
    length-$L$ words and sum one letter at a time, repeatedly applying the
    one-site left-canonical identity. For $L=0$, the unique empty word
    evaluates to the identity. The equivalence
    (\ref{eq:mpu_left_canonical_reindexing}) follows by reindexing the finite
    sum along $e$.
\end{proof}

\begin{definition}[Canonical-form-II data under blocking and relabeling]
    \label{def:mpu_cfii_blocking_reindexing}
    \lean{MPSTensor.CPSVCanonicalFormIIData.blockTensor,
        MPSTensor.CPSVCanonicalFormIIData.reindexPhysical,
        MPOTensor.blockTensorCFIIData}
    \leanok
    \uses{def:cpsv_cfii, def:mpu_normalized_flattening, thm:cpsv_literal_canonical_form_positive_blocking, thm:cpsv_physical_reindexing, lem:blocked_irreducibility_direct_support, thm:mpu_left_canonical_blocking_reindexing, thm:mpu_normalized_flattening_blocking}
    Let $A$ have canonical-form-II data with retained weights $\mu_k$, blocks
    $A_k$, and positive definite diagonal fixed matrices $\Lambda_k$. For
    every positive integer $L$, the blocked tensor $A^{[L]}$ has
    canonical-form-II data with weights $\mu_k^L$, blocks $A_k^{[L]}$, and
    the same matrices $\Lambda_k$. The retained blocks remain
    left-canonical:
    \begin{align}
        \sum_I \bigl((A_k^{[L]})^I\bigr)^\dagger(A_k^{[L]})^I & =\Id.
    \end{align}
    A bijective relabeling of the physical alphabet likewise preserves
    canonical-form-II data, including both the matrices $\Lambda_k$ and the
    left-canonical identities.

    Consequently, if the normalized flattening $\widetilde{\mathcal U}$ of an
    MPO tensor has canonical-form-II data, then for every $L>0$ the normalized
    flattening of $\mathcal U^{[L]}$ has the data obtained by first blocking
    and then relabeling by $\vartheta_L$ from
    Theorem~\ref{thm:mpu_normalized_flattening_blocking}.
    Theorem~\ref{thm:mpu_left_canonical_blocking_reindexing} preserves the
    left-canonical identity under both operations. For each retained block,
    Lemma~\ref{lem:blocked_irreducibility_direct_support} gives
    \begin{align}
        \E_{A_k^{[L]}} & =\E_{A_k}^{L},
                       & \E_{A_k}^{L}(\Lambda_k) & =\Lambda_k.
    \end{align}
    The second identity follows by iterating
    $\E_{A_k}(\Lambda_k)=\Lambda_k$. Thus the same positive definite diagonal
    matrix is fixed after blocking. A physical bijection leaves the transfer
    map unchanged. The
    final construction applies these two operations to
    (\ref{eq:mpu_normalized_flattening_blocking}).
\end{definition}

\begin{theorem}[Self-overlap scaling]
    \label{thm:mpu_self_overlap_scaling}
    \lean{MPSTensor.mpvOverlap_smul_self,
        MPOTensor.mpvOverlap_normalizedFlattening_self}
    \leanok
    \uses{def:mpu_normalized_flattening, def:mpv_overlap, def:mpo_to_mps, def:mpo_tensor, def:mps_tensor}
    For a matrix product tensor $A$, a complex number $c$, and a nonnegative
    integer $N$, scaling every local matrix by $c$ gives
    \begin{align}
        \sum_\sigma V^{(N)}(cA)_\sigma
        \overline{V^{(N)}(cA)_\sigma}
         & =(c\overline c)^N
        \sum_\sigma V^{(N)}(A)_\sigma
        \overline{V^{(N)}(A)_\sigma}.
    \end{align}
    Consequently, if $d>0$ and $U$ is an MPO tensor of physical dimension
    $d$, then
    \begin{align}
        \sum_\sigma V^{(N)}(\widetilde U)_\sigma
        \overline{V^{(N)}(\widetilde U)_\sigma}
         & =d^{-N}
        \sum_\sigma V^{(N)}(U^{\mathrm{MPS}})_\sigma
        \overline{V^{(N)}(U^{\mathrm{MPS}})_\sigma}.
    \end{align}
\end{theorem}

\begin{proof}\leanok
    \uses{lem:eval_word_smul}
    Each length-$N$ coefficient is multiplied by $c^N$.  Multiplication by
    its complex conjugate gives the first identity.  Taking
    $c=1/\sqrt d$ gives the second.
\end{proof}

\begin{theorem}[Normalized MPU transfer trace]
    \label{thm:mpu_normalized_transfer_trace}
    \lean{MPOTensor.IsMPU.trace_transferMatrix_normalizedFlattening_pow_eq_one}
    \leanok
    \uses{def:mpu, def:mpo_tensor, def:mps_tensor, def:mpu_normalized_flattening, def:transfer_map}
    Let $U$ be an MPU tensor with positive physical and bond dimensions, and
    let $E$ be the transfer matrix of its normalized flattening.  For every
    integer $N>1$,
    \begin{align}
        \tr(E^N) & =1.
    \end{align}
    This is the trace identity in the proof of the normal-tensor result
    \cite[Proposition~III.1]{Cirac2017MPU}.
\end{theorem}

\begin{proof}\leanok
    \uses{thm:mpu_transfer_matrix_overlap, thm:mpu_self_overlap_scaling, thm:mpdo_doubled_overlap_hilbert_schmidt, lem:mpu_mul_adjoint}
    Scaling both tensor copies by $1/\sqrt d$ multiplies the length-$N$
    self-overlap by $d^{-N}$.  The doubled-index overlap identity and MPU
    unitarity give the complete chain
    \begin{align}
        \tr(E^N)
         & =\sum_\sigma V^{(N)}(\widetilde U)_\sigma
        \overline{V^{(N)}(\widetilde U)_\sigma}                \\
         & =d^{-N}\sum_\sigma V^{(N)}(U^{\mathrm{MPS}})_\sigma
        \overline{V^{(N)}(U^{\mathrm{MPS}})_\sigma}            \\
         & =d^{-N}\tr\!\left(U^{(N)}(U^{(N)})^\dagger\right)   \\
         & =d^{-N}\tr(\Id)                                     \\
         & =d^{-N}d^N                                          \\
         & =1.
    \end{align}
    Here the fourth equality uses $N>1$, so that the MPU equation applies.
\end{proof}

\section{Block transfer multiplicity}

\begin{definition}[Flattened dependent-block inclusion]
    \label{def:mpu_dependent_block_inclusion}
    \lean{MPSTensor.blockInclusion}
    \leanok
    Let the blocks be indexed by $k\in\Fin{r}$, let $n_k$ be their
    nonnegative dimensions, and put $n=\sum_{k\in\Fin{r}}n_k$.  The standard
    block-order bijection is
    \begin{align}
        \phi\colon \bigsqcup_{k\in\Fin{r}}\Fin{n_k}
                  & \longrightarrow \Fin{n}, \\
        \phi(k,a) & =\sum_{j=0}^{k-1}n_j+a.
    \end{align}
    For each $k\in\Fin{r}$, define $J_k:\mathbb C^{n_k}\to\mathbb C^n$ by
    \begin{align}
        (J_k)_{x,a} & =\begin{cases}
                           1, & x=\phi(k,a),   \\
                           0, & x\ne\phi(k,a).
                       \end{cases}
    \end{align}
\end{definition}

\begin{theorem}[Dependent-block inclusion identities]
    \label{thm:mpu_dependent_block_inclusion_identities}
    \lean{MPSTensor.blockInclusion_apply,
        MPSTensor.blockInclusion_conjTranspose_mul_self,
        MPSTensor.blockInclusion_conjTranspose_mul_eq_zero,
        MPSTensor.toTensorFromBlocks_mul_blockInclusion}
    \leanok
    \uses{def:mpu_dependent_block_inclusion}
    The inclusions satisfy the following identities, with the second holding
    for distinct $k$ and $\ell$:
    \begin{align}
        J_k^\dagger J_k    & =\Id, \\
        J_k^\dagger J_\ell & =0.
    \end{align}
    If $B^i=\bigoplus_k\mu_k A_k^i$ is a weighted block-diagonal tensor,
    then
    \begin{align}
        B^iJ_k & =J_k(\mu_kA_k^i).
    \end{align}
\end{theorem}

\begin{proof}\leanok
    \uses{lem:matrix_sigma_block_inclusion}
    Reindex the canonical dependent-summand inclusion along the equivalence
    between dependent pairs and flattened coordinates.  Transporting its
    isometry, cross-block orthogonality, and block-diagonal intertwining
    identities gives the three displayed formulas.
\end{proof}

\begin{definition}[Ambient inclusion of a canonical block]
    \label{def:mpu_ambient_block_inclusion}
    \lean{MPSTensor.CPSVCanonicalFormData.ambientBlockInclusion}
    \leanok
    \uses{def:cpsv_canonical_form, def:mpu_dependent_block_inclusion}
    If $U$ is the reconstruction coisometry and $J_k$ is the inclusion of the
    $k$th retained block, define
    \begin{align}
        V_k & =U^\dagger J_k.
    \end{align}
\end{definition}

\begin{theorem}[Ambient block isometry and intertwining]
    \label{thm:mpu_ambient_block_isometry_intertwining}
    \lean{MPSTensor.CPSVCanonicalFormData.ambientBlockInclusion_conjTranspose_mul_self,
        MPSTensor.CPSVCanonicalFormData.mul_ambientBlockInclusion}
    \leanok
    \uses{def:cpsv_canonical_form, def:mpu_ambient_block_inclusion}
    For every retained block and physical index $i$,
    \begin{align}
        V_k^\dagger V_k & =\Id,              \\
        A^iV_k          & =V_k(\mu_k A_k^i).
    \end{align}
\end{theorem}

\begin{proof}\leanok
    \uses{thm:mpu_dependent_block_inclusion_identities}
    The coisometry $UU^\dagger=\Id$ and $J_k^\dagger J_k=\Id$ give
    \begin{align}
        V_k^\dagger V_k
         & =J_k^\dagger UU^\dagger J_k
        =J_k^\dagger J_k
        =\Id.
    \end{align}
    The canonical reconstruction and block-diagonal intertwining give
    \begin{align}
        A^iV_k
         & =U^\dagger\Bigl(\bigoplus_j\mu_jA_j^i\Bigr)J_k
        =U^\dagger J_k(\mu_kA_k^i)
        =V_k(\mu_kA_k^i).
    \end{align}
\end{proof}

\begin{definition}[Weighted-block transfer eigenvalue]
    \label{def:mpu_transfer_eigenvalue}
    \lean{MPSTensor.CPSVCanonicalFormData.transferEigenvalue}
    \leanok
    \uses{def:cpsv_canonical_form}
    For a block $k$, define its weighted transfer eigenvalue by
    \begin{align}
        q_k & =\mu_k\overline{\mu_k}=\lvert\mu_k\rvert^2.
    \end{align}
\end{definition}

\begin{theorem}[Shifted traces normalize every weight]
    \label{thm:mpu_transfer_eigenvalue_one}
    \lean{MPSTensor.CPSVCanonicalFormData.transferEigenvalue_eq_one}
    \leanok
    \uses{def:mpu_transfer_eigenvalue}
    If the ambient transfer matrix satisfies $\tr(E^N)=1$ for every integer
    $N>1$, then every block satisfies
    \begin{align}
        q_k & =1.
    \end{align}
\end{theorem}

\begin{proof}\leanok
    \uses{thm:mpu_ambient_block_isometry_intertwining, thm:mpu_shifted_trace_power_spectrum, thm:mpu_transfer_matrix_eigenvalues}
    Choose a nonzero fixed matrix $X_k$ for the unweighted normal block and set
    $Y_k=V_kX_kV_k^\dagger$. The isometry and intertwining identities give
    \begin{align}
        Y_k & \ne0,
            & E(Y_k) & =q_kY_k.
    \end{align}
    Thus $q_k$ lies in the ambient transfer spectrum; it is nonzero because
    $\mu_k\ne0$. The shifted trace-power spectrum is $\{0,1\}$, hence
    $q_k=1$.
\end{proof}

\begin{theorem}[Unique full-support block geometry]
    \label{thm:mpu_unique_full_support_geometry}
    \lean{MPSTensor.CPSVCanonicalFormData.dim_eq_of_r_eq_one_of_fullSupport,
        MPSTensor.CPSVCanonicalFormData.ambientBlockInclusion_mul_conjTranspose_eq_one}
    \leanok
    \uses{def:mpu_full_support, def:mpu_ambient_block_inclusion}
    If there is exactly one block $k$ and the blocks fill the
    ambient bond space, then
    \begin{align}
        D_k & =D,
            & V_kV_k^\dagger & =\Id.
    \end{align}
\end{theorem}

\begin{proof}\leanok
    \uses{thm:mpu_ambient_block_isometry_intertwining}
    The full-support sum reduces to the unique summand:
    \begin{align}
        D & =\sum_{j=1}^{r}D_j=D_k.
    \end{align}
    Therefore $V_k$ is a square matrix. Since $V_k^\dagger V_k=\Id$, finite
    equal-dimensional inverse symmetry gives $V_kV_k^\dagger=\Id$.
\end{proof}

\begin{theorem}[Shifted transfer traces determine one block]
    \label{thm:mpu_shifted_transfer_trace_single_block}
    \lean{MPSTensor.CPSVCanonicalFormData.r_eq_one_of_shifted_transfer_trace}
    \leanok
    \uses{def:cpsv_canonical_form}
    Let $A$ have literal CPSV canonical-form data with $r$ blocks.  If its
    transfer matrix $E$ satisfies
    $\tr(E^N)=1$ for every integer $N>1$, then
    \begin{align}
        r & =1.
    \end{align}
    This is the block multiplicity consequence used in
    \cite[Proposition~III.1]{Cirac2017MPU}.
\end{theorem}

\begin{proof}\leanok
    \uses{lem:transfer_map_corner_intertwine, thm:mpu_dependent_block_inclusion_identities, thm:mpu_shifted_trace_power_moments, thm:mpu_shifted_trace_power_spectrum, thm:mpu_trace_power_all, thm:mpu_transfer_matrix_eigenvalues}
    Each normal block $A_k$ has a nonzero transfer fixed vector $X_k$.
    Let $U$ be the reconstruction coisometry and let $J_k$ be the flattened
    inclusion of the $k$th block.  Define the ambient inclusion and transported
    matrix by
    \begin{align}
        V_k & =U^\dagger J_k,
            & Y_k             & =V_kX_kV_k^\dagger.
    \end{align}
    The matrix $Y_k$ is nonzero and is an ambient transfer eigenmatrix with
    eigenvalue $q_k=\mu_k\overline{\mu_k}\ne0$.
    Theorem~\ref{thm:mpu_shifted_trace_power_spectrum} gives $q_k=1$.

    For distinct $k$ and $\ell$, the inclusion identities give
    \begin{align}
        V_k^\dagger Y_kV_k     & =X_k, \\
        V_k^\dagger Y_\ell V_k & =0.
    \end{align}
    Hence the matrices $Y_k$, for $1\le k\le r$, are linearly independent.
    Vectorizing gives
    \begin{align}
        \operatorname{vec}(Y_k) & \in\ker(E^2-\Id).
    \end{align}
    For every $r\ge1$, the shifted traces and
    Theorem~\ref{thm:mpu_trace_power_all} give
    \begin{align}
        \tr((E^2)^r)  & =1,              \\
        \chi_{E^2}(X) & =X^{D^2-1}(X-1).
    \end{align}
    The eigenspace dimension is bounded by the algebraic multiplicity, so
    \begin{align}
        r
         & \le \dim\ker(E^2-\Id)
        \le \mult_{1}(\chi_{E^2})
        =1.
    \end{align}
    Finally, if $r$ were zero, the retained direct sum and therefore $A$ would
    vanish, contradicting $\tr(E^2)=1$.  Consequently $r=1$.
\end{proof}

\section{Normalized transfer spectrum and conditional normality}

\begin{theorem}[Shifted transfer traces imply radius one and primitivity]
    \label{thm:mpu_shifted_transfer_trace_primitive}
    \lean{spectralRadius_eq_one_and_isPrimitive_of_transferMatrix_shifted_trace}
    \leanok
    Let $T:M_D(\mathbb C)\to M_D(\mathbb C)$ be a linear map with $D>0$.
    If $\tr(\widehat T^{\,N})=1$ for every integer $N>1$, then
    \begin{align}
        r(T) & =1,
             & \operatorname{peripheral}(T) & =\{1\}.
    \end{align}
    This is the radius-and-primitivity part of
    \cite[Proposition~III.1]{Cirac2017MPU}.
\end{theorem}

\begin{proof}\leanok
    \uses{thm:mpu_shifted_trace_power_spectrum, thm:mpu_transfer_matrix_eigenvalues}
    The shifted traces give
    \begin{align}
        \sigma(\widehat T)\setminus\{0\} & =\{1\}.
    \end{align}
    Column stacking identifies the eigenvalues of $T$ and $\widehat T$.
    Consequently every spectral value of $T$ is zero or one, and one occurs,
    so
    \begin{align}
        r(T) & =\max_{\lambda\in\sigma(T)}|\lambda|=1.
    \end{align}
    Restricting to eigenvalues of modulus one gives
    \begin{align}
        \operatorname{peripheral}(T) & =\{1\},
    \end{align}
    which is primitivity.
\end{proof}

\begin{theorem}[Exact spectrum of the normalized MPU transfer matrix]
    \label{thm:mpu_normalized_transfer_nonzero_spectrum}
    \lean{MPOTensor.IsMPU.normalizedFlattening_nonzero_spectrum,
        MPOTensor.IsMPU.normalizedFlattening_charpoly}
    \leanok
    \uses{def:mpu, def:mpu_normalized_flattening, def:transfer_map}
    Let $U$ be an MPU tensor with positive physical and bond dimensions, and
    let $E$ be the transfer matrix of its normalized flattening. Then
    \begin{align}
        \sigma(E)\setminus\{0\} & =\{1\},
                                & \operatorname{mult}_{1}(\chi_E) & =1,
                                & \chi_E(X)                       & =X^{D^2-1}(X-1).
    \end{align}
    This is the transfer-matrix conclusion of
    \cite[Proposition~III.1]{Cirac2017MPU}. The first set equality alone does
    not assert algebraic multiplicity.
\end{theorem}

\begin{proof}\leanok
    \uses{thm:mpu_normalized_transfer_trace, thm:mpu_shifted_trace_power_charpoly}
    For every integer $N>1$, the normalized transfer identity gives
    \begin{align}
        \tr(E^N) & =1.
    \end{align}
    Applying the exact shifted trace-power theorem in dimension $D^2$ yields
    \begin{align}
        \chi_E(X) & =X^{D^2-1}(X-1),
                  & \operatorname{mult}_{1}(\chi_E) & =1,
                  & \sigma(E)\setminus\{0\}         & =\{1\}.
    \end{align}
\end{proof}

\begin{theorem}[Finite blocking eliminates the zero-primary transfer component]
    \label{thm:mpu_normalized_transfer_stabilization}
    \lean{MPOTensor.IsMPU.exists_normalized_transfer_stabilizes_to_rank_one}
    \leanok
    \uses{def:mpu, def:mpu_normalized_flattening, def:transfer_map, thm:mpu_normalized_transfer_nonzero_spectrum}
    Let $U$ be an MPU tensor with physical dimension $d>0$ and bond
    dimension $D>1$, and let $E$ be the transfer matrix of its normalized
    flattening. There is an integer
    $J$ satisfying
    \begin{align}
        0<J\leq D^2-1
    \end{align}
    and vectors $\rho,\Phi$ satisfying
    \begin{align}
        (\Phi\mid\rho) & =1,
                       & E\rho       & =\rho,
                       & (\Phi\mid E & =(\Phi\mid,
    \end{align}
    such that
    \begin{align}
        E^J & =\lvert\rho)(\Phi\rvert.
    \end{align}
    Moreover, for every integer $k\geq J$,
    \begin{align}
        E^k & =E^J.
    \end{align}
    Thus blocking $J$ sites removes the entire generalized zero-eigenspace.
    The one-dimensional case is stated separately below.

    This is the Jordan-elimination step in
    \cite[proof of Proposition~III.2(ii)]{Cirac2017MPU}.
\end{theorem}

\begin{proof}\leanok
    \uses{thm:mpu_normalized_transfer_nonzero_spectrum, thm:mpu_normalized_transfer_trace}
    The exact characteristic polynomial is
    \begin{align}
        \chi_E(X) & =X^{D^2-1}(X-1).
    \end{align}
    Cayley--Hamilton gives $E^{D^2}=E^{D^2-1}$. Hence every power from
    $J=D^2-1$ onward is equal, and $P=E^J$ is idempotent. Since $J>1$, the
    normalized transfer trace identity gives $\tr(P)=1$. The range of an
    idempotent has dimension equal to its trace, so $P$ has rank one and
    factors as $P=\lvert\rho)(\Phi\rvert$. Taking the factors with
    $(\Phi\mid\rho)=1$ gives the displayed fixed-point identities from
    $EP=PE=P$.
\end{proof}

\begin{theorem}[Uniqueness of the normalized fixed-pair factorization]
    \label{thm:mpu_stabilized_fixed_pair_uniqueness}
    \lean{Matrix.StabilizedRankOneData.power_eq_vec_mul_vec_of_fixed}
    \leanok
    Let $P=E^J$ be a stabilized rank-one transfer power with normalized left
    and right factors. If $\rho'$ and $\Phi'$ are any other fixed vectors with
    $(\Phi'\mid\rho')=1$, then
    \begin{align}
        P=\lvert\rho')(\Phi'\rvert.
    \end{align}
\end{theorem}

\begin{proof}\leanok
    Applying $P$ to $\rho'$ and applying $\Phi'$ to $P$ shows that the two
    pairs of rank-one factors differ by reciprocal scalars. Their normalized
    pairings force the product of those scalars to be one.
\end{proof}

\begin{theorem}[Identification with chosen reduced-CFII fixed points]
    \label{thm:mpu_reduced_cfii_transfer_stabilization}
    \lean{MPOTensor.IsMPU.normalized_transfer_power_eq_vecMulVec_of_reduced_cfii}
    \leanok
    \uses{def:mpu, def:mpu_normalized_flattening, def:cpsv_cfii, def:mpu_full_support, def:mpu_ambient_block_inclusion}
    Suppose $d>0$, $D>1$, and $\mathcal U$ is an MPU whose normalized
    flattening is presented by chosen canonical-form-II data with blocks
    filling the bond space. There is a
    unique block $k$. Normalize its diagonal positive definite fixed
    matrix $\Lambda_k$ to have trace one, and let $V_k$ be its inclusion into
    the ambient bond space. Then
    \begin{align}
        \rho         & =V_k\Lambda_kV_k^\dagger,
                     & \rho                      & >0,
                     & \tr(\rho)                 & =1,           \\
        E\lvert\rho) & =\lvert\rho),
                     & (\Phi\rvert E             & =(\Phi\rvert,
                     & (\Phi\mid\rho)            & =1,
    \end{align}
    where $(\Phi\rvert$ is the vectorized identity. Consequently,
    \begin{align}
        E^{D^2-1}=\lvert\rho)(\Phi\rvert.
    \end{align}
    % Source labels: Erightleft and prop:normal-tensor.
    % Source lines: paper_v2.tex 269--280 and 397--405.
    The diagonal matrix is the block matrix $\Lambda_k$ in the canonical-form-II
    fixed-point equations in
    \cite[equations~(6a)--(6b)]{Cirac2017MPU}; its ambient transport
    $\rho$ need not be diagonal. This is a statement about the chosen reduced
    representative and the blocking step in
    \cite[proof of Proposition~III.2(ii)]{Cirac2017MPU}; it does not assert
    that every ambient MPU tensor is itself in canonical form II.
\end{theorem}

\begin{proof}\leanok
    \uses{thm:mpu_normalized_transfer_trace, thm:mpu_shifted_transfer_trace_single_block, thm:mpu_transfer_eigenvalue_one, thm:mpu_ambient_block_isometry_intertwining, thm:mpu_unique_full_support_geometry, thm:matrix_posdef_coisometric_transport, lem:transfer_map_corner_intertwine, thm:mpu_normalized_transfer_stabilization, thm:mpu_stabilized_fixed_pair_uniqueness}
    The shifted trace identities satisfy $\tr(E^N)=1$ for every integer
    $N>1$ and give
    \begin{align}
        \tr(E^N)=1
        \quad\Longrightarrow\quad
        r=1.
    \end{align}
    Let $k$ be this block. Canonical form II supplies a diagonal
    $\Lambda_0>0$ fixed by the unweighted block transfer map. Normalize it by
    \begin{align}
        \Lambda & =\tr(\Lambda_0)^{-1}\Lambda_0,
                & \tr(\Lambda)                   & =1.
    \end{align}
    If $\omega_k$ is its weight, then
    \begin{align}
        \lvert\omega_k\rvert^2 & =1,
                               & E_k(\Lambda) & =\Lambda,
    \end{align}
    where $E_k$ is the transfer map of the weighted block
    $\widetilde B_k=\omega_k B_k$. Full support makes the block inclusion
    $V=V_k$ both isometric and onto:
    \begin{align}
        V^\dagger V & =I,
                    & VV^\dagger & =I.
    \end{align}
    Hence $\rho=V\Lambda V^\dagger$ is positive definite with trace one, and
    the intertwining equations $A_iV=V\widetilde B_k(i)$ give the fixedness
    transport
    \begin{align}
        E(V\Lambda V^\dagger)
         & =V E_k(\Lambda)V^\dagger
        =V\Lambda V^\dagger
        =\rho.
    \end{align}
    The blockwise left-canonical identity transports similarly:
    \begin{align}
        \sum_i \widetilde B_k(i)^\dagger\widetilde B_k(i)=I
        \quad\Longrightarrow\quad
        \sum_i A_i^\dagger A_i=I
        \quad\Longrightarrow\quad
        (\Phi\rvert E=(\Phi\rvert.
    \end{align}
    Finally,
    \begin{align}
        (\Phi\mid\rho) & =\tr(\rho)=1,
                       & E^{D^2-1}     & =\lvert\rho)(\Phi\rvert,
    \end{align}
    where the last equality is the uniqueness of the normalized fixed-pair
    factorization for the stabilized power.
\end{proof}

\begin{theorem}[Reduced-CFII data for the forced simple block]
    \label{thm:mpu_reduced_cfii_forced_block_simple_contractions}
    \lean{MPOTensor.IsMPU.exists_reduced_cfii_forced_block_simple_contractions}
    \leanok
    \uses{def:mpu, def:mpu_normalized_flattening, def:cpsv_cfii, def:mpu_full_support, def:mpdo_physical_blocking, def:mpu_double_layer}
    Let $d>0$, $D>1$, and let $\mathcal U$ be an MPU whose normalized
    flattening has chosen canonical-form-II data with full support. Put
    $J=D^2-1$ and $W=\mathcal U_{JD^2}$. There is a matrix $\rho>0$ with
    $\tr(\rho)=1$ such that
    \begin{align}
        E\lvert\rho) & =\lvert\rho),
                     & (\Phi\rvert E & =(\Phi\rvert,
                     & E^J           & =\lvert\rho)(\Phi\rvert.
    \end{align}
    With the corresponding vectorized pair $\rho',\Phi'$, the same blocked
    tensor $W$ satisfies the two aligned simple contractions
    \begin{align}
        \Phi'\mathbin{\boldsymbol\cdot}W^{ij}\rho'
         & =\delta_{ij},                              \\
        W^{ij}W^{k\ell}
         & =W^{ij}\lvert\rho')(\Phi'\rvert W^{k\ell}.
    \end{align}
    This theorem supplies only the fixed pair and the two simple contractions.
    It does not assert a source-cut factorization or source-$v$ isometry.
    % Source lines: paper_v2.tex 397--427, 550--556, 589, and 592.
\end{theorem}

\begin{proof}\leanok
    \uses{thm:mpu_matching_direct_block_contractions, thm:mpu_reduced_cfii_transfer_stabilization}
    Theorem~\ref{thm:mpu_reduced_cfii_transfer_stabilization} gives the
    positive trace-one fixed matrix and the identity
    $E^J=\lvert\rho)(\Phi\rvert$. Applying
    Theorem~\ref{thm:mpu_matching_direct_block_contractions} with this pair
    gives both contractions on the direct block $W=\mathcal U_{JD^2}$.
\end{proof}

\begin{theorem}[One-dimensional normalized transfer matrix]
    \label{thm:mpu_normalized_transfer_fin_one}
    \lean{MPOTensor.IsMPU.normalized_transfer_matrix_eq_one_fin_one}
    \leanok
    \uses{def:mpu, def:mpu_normalized_flattening}
    If the physical dimension satisfies $d>0$ and the bond dimension is $D=1$,
    then the normalized transfer matrix satisfies
    \begin{align}
        E & =1.
    \end{align}
    Consequently its powers stabilize at the positive exponent one.
\end{theorem}

\begin{proof}\leanok
    \uses{thm:mpu_normalized_transfer_nonzero_spectrum}
    In dimension one the characteristic polynomial identity reduces to
    $\chi_E(X)=X-1$, and Cayley--Hamilton gives $E=1$.
\end{proof}

\begin{theorem}[Normality of a full-support normalized MPU representative]
    \label{thm:mpu_normalized_full_support_normal}
    \lean{MPOTensor.IsMPU.isNormalTensor_normalizedFlattening_of_cpsv}
    \leanok
    \uses{def:mpu, def:mpu_normalized_flattening, def:cpsv_canonical_form, def:mpu_full_support, def:normal}
    Let $U$ be an MPU tensor with positive physical and bond dimensions.
    Suppose that the normalized flattening is equipped with CPSV
    canonical-form data whose blocks have full support in the ambient
    bond space. Then the normalized flattening is a normal tensor.

    This is the full-support representative form of
    \cite[Proposition~III.1]{Cirac2017MPU}. The full-support hypothesis is
    essential for this ambient representative: adjoining a zero virtual
    direct summand preserves every periodic operator but introduces a
    nontrivial invariant projection. Thus the theorem does not assert that an
    arbitrary, possibly nonminimal MPU representative is normal.
\end{theorem}

\begin{proof}\leanok
    \uses{thm:mpu_normalized_transfer_trace, thm:mpu_shifted_transfer_trace_single_block, thm:mpu_shifted_transfer_trace_primitive, thm:mpu_unique_full_support_normal}
    Write $\widetilde U$ for the normalized flattening, set
    $T=\mathcal E_{\widetilde U}$ for its transfer map, and let
    $E=\widehat T$ be the corresponding transfer matrix. Also write $r$ for
    the number of blocks and $D_k$ for the bond dimension of the block $k$.
    For every integer $N>1$, the normalized transfer identity gives
    \begin{align}
        \tr(E^N) & =1.
    \end{align}
    The block-multiplicity theorem and the transfer-spectrum theorem then give
    \begin{align}
        r & =1,
          & r(T)                         & =1,
          & \operatorname{peripheral}(T) & =\{1\}.
    \end{align}
    Full support supplies
    \begin{align}
        \sum_{k=1}^{r}D_k & =D.
    \end{align}
    Since $r=1$, the unique block has ambient dimension
    $D$ and therefore makes $\widetilde U$ irreducible. The displayed radius
    and peripheral-spectrum identities provide its remaining normality
    fields. Consequently $\widetilde U$ is a normal tensor.
\end{proof}

\section{Standard form, index, and symmetry classification}

\begin{definition}[Reduced full-support source datum]
    \label{def:mpu_reduced_full_support_source_datum}
    \notready
    % **Scope restriction (source data for the same tensor):** this datum is an
    % explicit hypothesis, not a construction of a reduced representative.
    % Documented in docs/paper-gaps/mpu_canonical_form_full_support.tex.
    \uses{def:mpu, def:mpu_normalized_flattening, def:cpsv_cfii, def:mpu_full_support, thm:mpu_reduced_cfii_transfer_stabilization, thm:mpu_normalized_transfer_fin_one}
    Let $\widetilde{\mathcal U}=d^{-1/2}\mathcal U$ be the normalized
    flattening, and let $E$ be its transfer matrix. A reduced full-support source
    datum for $\mathcal U$, with physical dimension $d>0$ and positive bond
    dimension, consists of an integer $J>0$, a diagonal positive definite matrix
    $\rho$ with $\tr(\rho)=1$, and
    $(\Phi\rvert=\operatorname{vec}(\Id_D)^{\mathsf T}$ such that
    \begin{align}
        E^J & =\lvert\rho)(\Phi\rvert.
    \end{align}
    The one-dimensional case is supplied directly: if $D=1$, take $J=1$
    and $\rho=\Id_1$ using
    Theorem~\ref{thm:mpu_normalized_transfer_fin_one}. If $D>1$, chosen
    canonical-form-II data with full support give a positive trace-one ambient
    fixed point with $J=D^2-1$ by
    Theorem~\ref{thm:mpu_reduced_cfii_transfer_stabilization}. Diagonality of
    this ambient fixed point is an additional hypothesis: the theorem gives a
    diagonal matrix in the unique block coordinates, but conjugation into the
    ambient coordinates need not preserve diagonality.
    For $D>1$, this states the fixed-pair convention in
    \cite[equations labelled~\texttt{Erightleft} and~\texttt{II\_UTransfer}, and
        proof of Proposition~III.2(ii)]{Cirac2017MPU}. The $D=1$ branch is the direct
    one-dimensional specialization. This definition does not assert that an
    arbitrary representative with $D>1$ admits the required canonical-form-II
    data and does not replace $\mathcal U$ by another tensor.
    % Source lines: paper_v2.tex 269--280 and 397--405.
\end{definition}

\begin{theorem}[Staggered Gram expansion of the source tensor $u$]
    \label{thm:mpu_source_u_staggered_gram}
    \lean{MPOTensor.SourceFactors.sourceU_gram_eq_staggered_four_u}
    \leanok
    \uses{def:mpu_source_uv, def:mpu_weighted_source_factors}
    Let $M_i=X_iY_i$ be supplied source factorizations for $\mathcal U$ and
    the weight $\rho$. For retained physical pairs $p=(p_1,p_2)$ and
    $q=(q_1,q_2)$, one has
    \begin{align}
        \sum_{l,r}u_{(l,r),q}\overline{u_{(l,r),p}}
        ={} & \sum_{\alpha,\alpha',\beta,\beta',\gamma,j_1,j_2}
        \overline{\mathcal U^{j_1}_{p_2,\alpha,\beta}}
        \rho_{\beta,\beta'}
        \mathcal U^{j_1}_{q_2,\alpha',\beta'} \notag            \\
            & \qquad\qquad\cdot
        \overline{\mathcal U^{j_2}_{p_1,\gamma,\alpha}}
        \mathcal U^{j_2}_{q_1,\gamma,\alpha'}.
    \end{align}
    The $p$ entries are starred and the $q$ entries are unstarred. This is
    precisely the staggered reflected--direct four-tensor expression obtained
    from $u=Y_2\mathbin{-}Y_1$ and the two identities in
    \cite[equations~\texttt{X1X2b} and~\texttt{uUnitary}]{Cirac2017MPU}.
\end{theorem}

\begin{proof}\leanok
    \uses{def:mpu_source_uv, thm:mpu_source_factorizations, thm:mpu_source_factor_normalizations}
    Since the first cut has rows $(i,\beta)$ and columns $(\alpha,j)$,
    its weighted normalization gives
    \begin{align}
        \sum_r\overline{(Y_1)_{r,(\alpha,p_2)}}
        (Y_1)_{r,(\alpha',q_2)}
         & =\sum_{\beta,\beta',j_1}
        \overline{\mathcal U^{j_1}_{p_2,\alpha,\beta}}
        \rho_{\beta,\beta'}
        \mathcal U^{j_1}_{q_2,\alpha',\beta'}.
    \end{align}
    The unweighted normalization of the second cut gives
    \begin{align}
        \sum_l\overline{(Y_2)_{l,(p_1,\alpha)}}
        (Y_2)_{l,(q_1,\alpha')}
         & =\sum_{\gamma,j_2}
        \overline{\mathcal U^{j_2}_{p_1,\gamma,\alpha}}
        \mathcal U^{j_2}_{q_1,\gamma,\alpha'}.
    \end{align}
    Expand both copies of
    $u_{(l,r),(i_1,i_2)}=\sum_\alpha
        (Y_2)_{l,(i_1,\alpha)}(Y_1)_{r,(\alpha,i_2)}$ and substitute these two
    Gram identities. Reordering the finite sums yields the stated formula.
\end{proof}

\begin{theorem}[Transpose fixed-pair trace closure]
    \label{thm:mpu_source_u_transpose_fixed_pair_trace}
    \lean{MPOTensor.SourceFactors.sourceU_gram_eq_transpose_fixed_pair_trace}
    \leanok
    \discussion{7632}
    \uses{def:mpu_source_uv, def:mpu_weighted_source_factors, def:mpu_double_layer}
    Let $\mathcal U$ be an MPO tensor, let $\rho$ be a bond matrix, and let
    $M_i=X_iY_i$ be supplied source factorizations for $\mathcal U$ and $\rho$.
    Construct the source tensor $u=Y_2\mathbin{-}Y_1$. For retained physical
    pairs $p=(p_1,p_2)$ and $q=(q_1,q_2)$, let $W^{ij}$ denote the $(i,j)$
    double-layer letter, let $\lvert\rho^{\mathsf T})$ be the column-stacking
    vector of $\rho^{\mathsf T}$, and set
    $(\Phi\rvert=\operatorname{vec}(\Id_D)^{\mathsf T}$. Then
    \begin{align}
        \sum_{l,r}u_{(l,r),q}\overline{u_{(l,r),p}}
         & =\operatorname{tr}\!\left(
        W^{p_2q_2}\lvert\rho^{\mathsf T})(\Phi\rvert W^{p_1q_1}
        \right).
    \end{align}
    Thus the fixed-pair weight is genuinely transposed at this stage. In the
    diagonal canonical-form-II coordinates fixed above,
    $\rho^{\mathsf T}=\rho$; this is the identification used in the complete
    contraction below.
\end{theorem}

\begin{proof}\leanok
    \uses{thm:mpu_source_u_staggered_gram, thm:mpu_source_v_four_u_expansions}
    Enumerate each doubled-bond coordinate as
    $\operatorname{flat}(\alpha,\alpha')$. Expanding a diagonal entry of the
    matrix under the trace gives
    \begin{align}
        \sum_{\beta,\beta',\gamma,j_1,j_2}
        \overline{\mathcal U^{j_1}_{p_2,\alpha,\beta}}
        \mathcal U^{j_1}_{q_2,\alpha',\beta'}
        (\rho^{\mathsf T})_{\beta',\beta}
        \overline{\mathcal U^{j_2}_{p_1,\gamma,\alpha}}
        \mathcal U^{j_2}_{q_1,\gamma,\alpha'}.
    \end{align}
    Since $(\rho^{\mathsf T})_{\beta',\beta}=\rho_{\beta,\beta'}$,
    summing over $\alpha,\alpha'$ gives the identity in
    Theorem~\ref{thm:mpu_source_u_staggered_gram}.
\end{proof}

\begin{theorem}[Retained-interior source-$u$ closure]
    \label{thm:mpu_source_u_retained_interior}
    \lean{MPOTensor.SourceFactors.sourceU_gram_eq_closed_doubleLayer_trace_of_isSymm,
        MPOTensor.SourceFactors.sourceU_gram_eq_normalized_input_tail}
    \leanok
    \discussion{7633}
    \uses{thm:mpu_source_u_transpose_fixed_pair_trace,
        thm:mpu_normalized_input_tail_closed_trace}
    Let $\mathcal U$ be an MPO tensor with physical dimension $d>0$, let $W$
    be its double layer, and put $E=d^{-1}\sum_iW^{ii}$. Let $\rho$ be
    symmetric, set $(\Phi\rvert=\operatorname{vec}(\Id_D)^{\mathsf T}$, and
    suppose that for some $K\geq0$,
    \begin{align}
        E^K & =\lvert\rho)(\Phi\rvert.
    \end{align}
    For supplied source factors built from the same weight $\rho$, the paper
    gate $u=Y_2\mathbin{-}Y_1$ satisfies
    \begin{align}
        \sum_{l,r}u_{(l,r),q}\overline{u_{(l,r),p}}
         & =\operatorname{tr}\!\left(W^{p_1q_1}W^{p_2q_2}E^K\right).
    \end{align}
    If $\rho$ is diagonal, the same quantity is
    \begin{align}
        d^{-K}\sum_{\tau,\eta}
        \overline{U^{(K+2)}_{\eta,(p,\tau)}}
        U^{(K+2)}_{\eta,(q,\tau)}.
    \end{align}
    Thus only the $K$-site interior is blocked; the two endpoint letters remain
    one-site tensors. The $p$ entries are conjugated and the $q$ entries are
    not. This is the supplied-fixed-pair contraction in
    \cite[Lemma~III.7]{Cirac2017MPU}.
    % Source locator: paper_v2.tex, lines 269--280 and 545--557.
    % Orientation: docs/paper-gaps/mpu_source_cut_orientation.tex.
\end{theorem}

\begin{proof}\leanok
    \uses{thm:mpu_source_u_transpose_fixed_pair_trace,
        thm:mpu_normalized_input_tail_closed_trace}
    Symmetry changes the inserted vector
    $\lvert\rho^{\mathsf T})$ to $\lvert\rho)$ in
    Theorem~\ref{thm:mpu_source_u_transpose_fixed_pair_trace}. Substitute the
    fixed pair and cycle the first endpoint letter to obtain the displayed
    trace. The normalized input-tail identity expands that trace into the
    second displayed sum.
\end{proof}

\begin{theorem}[Complete source-$u$ network contraction]
    \label{thm:mpu_source_u_complete_network}
    \lean{MPOTensor.SourceFactors.sourceU_gram_eq_normalized_mpo_input_tail_of_isDiag}
    \leanok
    \discussion{5982}
    \uses{def:mpu, def:mpu_normalized_flattening, def:mpu_source_uv, def:mpo_family}
    % Revalidation boundary: docs/paper-gaps/mpu_source_cut_orientation.tex.
    Let $\mathcal U$ be an MPU tensor with physical dimension $d>0$ and bond
    dimension $D>0$. Let $\rho>0$ be diagonal with $\tr(\rho)=1$, and put
    $(\Phi\rvert=\operatorname{vec}(\Id_D)^{\mathsf T}$. Suppose that for some
    $J>0$ the normalized transfer matrix satisfies
    \begin{align}
        E^J & =\lvert\rho)(\Phi\rvert.
    \end{align}
    Construct the source factors and $u$ from $\mathcal U$ using $\rho$, and set
    $K=JD^2$. For retained physical pairs
    $p,q\in\{0,\ldots,d-1\}^2$, the complete source-$u$ contraction satisfies
    \begin{align}
        \sum_{l,r}u_{(l,r),q}\overline{u_{(l,r),p}}
         & =d^{-K}\sum_{\tau,\eta}
        \overline{U^{(K+2)}_{\eta,(p,\tau)}}
        U^{(K+2)}_{\eta,(q,\tau)}.
    \end{align}
    Here $\tau$ ranges over input words of length $K$, and $\eta$ ranges over
    output words of length $K+2$. This is the supplied-fixed-pair form of the
    complete-network equality in \cite[Lemma~III.7]{Cirac2017MPU}.
    % Source locator: paper_v2.tex, lines 545--557.
\end{theorem}

\begin{proof}\leanok
    \uses{thm:mpu_source_u_retained_interior,
        thm:mpu_blocked_transfer_power_rank_one}
    Set $K=JD^2$. Trace normalization makes the supplied rank-one transfer
    matrix idempotent, and the normalized blocked-transfer identity therefore
    gives
    \begin{align}
        E^K & =\lvert\rho)(\Phi\rvert.
    \end{align}
    Diagonality gives $\rho^{\mathsf T}=\rho$, so
    Theorem~\ref{thm:mpu_source_u_retained_interior} applies at this $K$ and
    yields the displayed normalized input-tail sum.
\end{proof}

\begin{theorem}[Complete source-$v$ network]
    \label{thm:mpu_source_v_complete_contraction}
    \lean{MPOTensor.sourceYTensor_conjTranspose_sourceV_gram_apply,
        MPOTensor.sourceYTensor_gram_apply_eq_doubleLayerTensor_rankOne_mul_apply,
        MPOTensor.sourceVDressedGram_iff_simple2_transpose_fixed_pair}
    \leanok
    \uses{def:mpu_source_uv, def:mpu_source_v_dressing, def:mpu_double_layer}
    Let $\rho>0$, and compute both source cuts, the source operator $v$, and
    the right source factors $Y_1,Y_2$ from the same tensor $\mathcal U$ using
    $\rho$. Let $W$ be the double layer of $\mathcal U$, let
    $\lvert\rho^{\mathsf T})$ be the column-stacking vector of
    $\rho^{\mathsf T}$, and let
    $(\Phi\rvert=\operatorname{vec}(\Id_D)^{\mathsf T}$. For the starred
    physical pair $p=(p_1,p_2)$, the unstarred pair $q=(q_1,q_2)$, and
    virtual paths $a=(a_1,a_2)$ and $b=(b_1,b_2)$, put
    $x=((a_1,p_1),(p_2,a_2))$ and $y=((b_1,q_1),(q_2,b_2))$. Then
    \begin{align}
        \bigl((Y_1\otimes Y_2)^\dagger v^\dagger v(Y_1\otimes Y_2)\bigr)_{x,y}
         & =\bigl(W^{p_1q_1}W^{p_2q_2}\bigr)_{(a_1,b_1),(a_2,b_2)},
        \label{eq:mpu_source_v_dressed_two_letter}                              \\
        \bigl((Y_1\otimes Y_2)^\dagger(Y_1\otimes Y_2)\bigr)_{x,y}
         & =\bigl(W^{p_1q_1}\lvert\rho^{\mathsf T})(\Phi\rvert W^{p_2q_2}\bigr)
        _{(a_1,b_1),(a_2,b_2)}.
        \label{eq:mpu_source_v_undressed_inserted}
    \end{align}
    Consequently the dressed Gram equation
    \begin{align}
        (Y_1\otimes Y_2)^\dagger v^\dagger v(Y_1\otimes Y_2)
         & =(Y_1\otimes Y_2)^\dagger(Y_1\otimes Y_2)
    \end{align}
    holds if and only if
    \begin{align}
        W^{ij}W^{k\ell}
         & =W^{ij}\lvert\rho^{\mathsf T})(\Phi\rvert W^{k\ell}
    \end{align}
    for all $i,j,k,\ell$. The source gate and source factors use the same
    weight $\rho$ throughout. The transpose occurs only in the inserted
    column-stacking vector. In the source, $\rho$ is diagonal, so this vector
    equals $\lvert\rho)$ and the displayed equivalence is exactly
    \cite[equation~\texttt{vUnitary} and proof of Theorem~III.8]{Cirac2017MPU}.
    % Source lines: paper_v2.tex 577--600. The weight orientation is recorded
    % in docs/paper-gaps/mpu_source_cut_orientation.tex.
\end{theorem}

\begin{proof}\leanok
    \uses{thm:mpu_source_v_mul_source_y_tensor_entry, thm:mpu_source_v_four_u_expansions}
    Regrouping gives
    $(Y_1\otimes Y_2)^\dagger v^\dagger v(Y_1\otimes Y_2)
        =\bigl(v(Y_1\otimes Y_2)\bigr)^\dagger v(Y_1\otimes Y_2)$, and
    Theorem~\ref{thm:mpu_source_v_mul_source_y_tensor_entry} closes each
    factor by the two source-cut factorizations:
    \begin{align}
        \bigl(v(Y_1\otimes Y_2)\bigr)_{z,x}
         & =\sum_\gamma U^{z_1p_1}_{a_1\gamma}U^{z_2p_2}_{\gamma a_2}.
    \end{align}
    Summing the conjugate of this expression against the corresponding
    expression for $y$ over $z$ and reordering the finite sums gives the
    ordinary two-letter entry of
    Theorem~\ref{thm:mpu_source_v_four_u_expansions}, which
    proves~(\ref{eq:mpu_source_v_dressed_two_letter}). For
    (\ref{eq:mpu_source_v_undressed_inserted}), the same theorem expands the
    Gram entry of $Y_1\otimes Y_2$ with the weight $\rho_{\beta\beta'}$
    attached to the starred index $\beta$, and expands the inserted entry with
    $(\rho^{\mathsf T})_{\alpha'\alpha}=\rho_{\alpha\alpha'}$ attached to the
    starred index $\alpha$; the two finite sums agree after reordering. Both
    directions of the equivalence compare the two entry identities at every
    $x$ and $y$. The network is kept combined, and neither source cut is
    contracted separately.
\end{proof}

\begin{theorem}[Source-$v$ isometry of a simple tensor]
    \label{thm:mpu_simple_source_v_isometry}
    \lean{MPOTensor.IsMPUSimple.sourceV_isIsometry,
        MPOTensor.IsMPUSimple.rightRank_mul_leftRank_le}
    \leanok
    \uses{def:mpu_simple, def:mpu_source_uv, def:mpu_double_layer}
    Let $\mathcal U$ be simple with physical dimension $d>0$, let $W$ be its
    double layer with normalized diagonal $E$, and let the source's diagonal
    matrix $\rho>0$ satisfy $E^J=\lvert\rho)(\Phi\rvert$ for some $J>0$, where
    $(\Phi\rvert=\operatorname{vec}(\Id_D)^{\mathsf T}$. Let $v$ be the source
    operator of $\mathcal U$ built from that same weight $\rho$. Then
    $v^\dagger v=\Id_{r\ell}$, and consequently $r\ell\le d^2$. This is the
    implication $1\to2$ in \cite[proof of Theorem~III.8]{Cirac2017MPU},
    before the standing isometry of $u$ is used.
    % Source lines: paper_v2.tex 577--588.
\end{theorem}

\begin{proof}\leanok
    \uses{thm:mpu_supplied_projector_simple2, thm:mpu_source_v_complete_contraction, thm:mpu_source_v_dressed_gram_cancel}
    Theorem~\ref{thm:mpu_supplied_projector_simple2} gives
    $W^{ij}W^{k\ell}=W^{ij}\lvert\rho)(\Phi\rvert W^{k\ell}$. Since $\rho$
    is diagonal, $\rho^{\mathsf T}=\rho$, so
    Theorem~\ref{thm:mpu_source_v_complete_contraction} gives the dressed Gram
    equation for the factors built from the same weight $\rho$. Then
    Theorem~\ref{thm:mpu_source_v_dressed_gram_cancel} cancels the dressing. Since $v^\dagger v=\Id$ on
    $\C^r\otimes\C^\ell$, the map $v$ is injective, so $r\ell\le d^2$.
\end{proof}

\begin{theorem}[Simplicity from the source-$v$ isometry]
    \label{thm:mpu_simple_of_source_v_isometry}
    \lean{MPOTensor.IsMPU.isMPUSimple_of_sourceV_isIsometry}
    \leanok
    \uses{def:mpu, def:mpu_simple, def:mpu_source_uv}
    Let $\mathcal U$ be an MPU tensor, let $\rho>0$, and let $v$ be its
    source operator built from $\rho$. If $v^\dagger v=\Id_{r\ell}$, then
    $\mathcal U$ is simple, with witnesses $(\Phi\rvert$ and
    $\lvert\rho^{\mathsf T})$. This is the implication $4\to1$ in
    \cite[proof of Theorem~III.8]{Cirac2017MPU}.
    % Source lines: paper_v2.tex 596--600.
\end{theorem}

\begin{proof}\leanok
    \uses{thm:mpu_source_v_dressed_gram_cancel, thm:mpu_source_v_complete_contraction, thm:mpu_all_later_simple_blockings}
    By Theorem~\ref{thm:mpu_source_v_dressed_gram_cancel}, $v^\dagger v=\Id$
    gives the dressed Gram equation, and
    Theorem~\ref{thm:mpu_source_v_complete_contraction} turns it into
    $W^{ij}W^{k\ell}=W^{ij}\lvert\rho^{\mathsf T})(\Phi\rvert W^{k\ell}$.
    The same-witness implication in
    Theorem~\ref{thm:mpu_all_later_simple_blockings} supplies
    $(\Phi\rvert W^{ij}\lvert\rho^{\mathsf T})=\delta_{ij}$ from unitarity,
    so both simple contractions hold.
\end{proof}

\begin{lemma}[The source operator $u$ is an isometry]
    \label{lemuisometry}
    \notready
    \discussion{5708}
    \uses{def:mpu_source_uv}
    For any MPU tensor $\mathcal U$ under the source's standing
    canonical-form-II convention, let $r$ and $\ell$ be the two source-cut
    ranks and let $u$ be the source operator. Then
    \begin{align}
        u^\dagger u & =\Id.
    \end{align}
    Consequently $u$ is an isometry and $r\ell\geq d^2$. This is the
    unrestricted statement of \cite[Lemma~III.7]{Cirac2017MPU}.
    % Source lines: paper_v2.tex 545--557.
\end{lemma}

\begin{lemma}[Admissible source-$u$ isometry]
    \label{lem:mpu_admissible_source_u_isometry}
    \lean{MPOTensor.SourceFactors.sourceU_isIsometry_of_isMPU,
        MPOTensor.SourceFactors.mul_self_le_rightRank_mul_leftRank_of_isMPU,
        MPOTensor.IsMPU.sourceU_isIsometry}
    \leanok
    \discussion{5708}
    % **Scope restriction (supplied stabilized diagonal fixed pair):** this
    % lemma assumes the exact rank-one identity explicitly, whereas source
    % Lemma III.7 uses the preceding canonical-form-II convention. Documented
    % in docs/paper-gaps/mpu_canonical_form_full_support.tex and
    % docs/paper-gaps/mpu_source_cut_orientation.tex.
    \uses{def:mpu, def:mpu_source_uv, def:mpu_weighted_source_factors,
        def:mpu_double_layer}
    Let $\mathcal U$ be an MPU tensor with physical dimension $d>0$. Let $W$
    be its double layer, put $E=d^{-1}\sum_iW^{ii}$, and let $\rho>0$ be
    diagonal. Put $(\Phi\rvert=\operatorname{vec}(\Id_D)^{\mathsf T}$ and
    suppose that for some $K\geq0$,
    \begin{align}
        E^K & =\lvert\rho)(\Phi\rvert.
    \end{align}
    Let $r$ and $\ell$ be the source-cut ranks, and construct the compact-SVD
    source factors and $u$ from $\mathcal U$ using $\rho$. Then
    \begin{align}
        u^\dagger u & =\Id.
    \end{align}
    Moreover,
    \begin{align}
        d^2 & \leq r\ell.
    \end{align}
    Thus the compact-SVD paper gate is an isometry. This is the
    supplied-fixed-pair version of \cite[Lemma~III.7]{Cirac2017MPU}.
    % Source lines: paper_v2.tex 269--280 and 545--557.
\end{lemma}

\begin{proof}\leanok
    \uses{thm:mpu_source_u_retained_interior,
        thm:mpu_normalized_input_tail_isometry}
    The retained-interior identity and input-first MPU unitarity give, for all
    physical pairs $p,q$,
    \begin{align}
        (u^\dagger u)_{p,q}
         & =d^{-K}\sum_{\tau,\eta}
        \overline{U^{(K+2)}_{\eta,(p,\tau)}}
        U^{(K+2)}_{\eta,(q,\tau)}
        =\delta_{p,q}.
    \end{align}
    Hence $u^\dagger u=\Id$. Since $u$ maps a space of dimension $d^2$
    isometrically into one of dimension $r\ell$, one also has $d^2\leq r\ell$.
\end{proof}

\begin{theorem}[Simple-tensor equivalence theorem]
    \label{ThmFund1}
    \notready
    \discussion{5708}
    \uses{lemuisometry, def:mpu_simple, def:mpu_source_uv}
    Let $\mathcal U$ be an MPU tensor under the source's standing
    canonical-form-II convention. Let $r$ and $\ell$ be its source-cut ranks,
    and let $u$ and $v$ be its source operators. Then the following conditions
    are equivalent:
    \begin{enumerate}
        \item $\mathcal U$ is simple;
        \item $r\ell=d^2$;
        \item $u$ is unitary;
        \item $v$ is unitary.
    \end{enumerate}
    This is the unrestricted statement of
    \cite[Theorem~III.8]{Cirac2017MPU}.
    % Source lines: paper_v2.tex 563--601.
    % The supplied diagonal-fixed-pair specialization below is checked. The
    % unrestricted node remains not ready for the representative scope recorded
    % in docs/paper-gaps/mpu_canonical_form_full_support.tex.
\end{theorem}

\begin{theorem}[Simple-tensor equivalence with a supplied diagonal fixed pair]
    \label{thm:mpu_diagonal_fixed_pair_simple_tensor_equivalence}
    \lean{MPOTensor.IsMPU.isMPUSimple_tfae}
    \leanok
    \uses{def:mpu, def:mpu_simple, def:mpu_source_uv, def:mpu_double_layer}
    Let $\mathcal U$ be an MPU tensor with physical dimension $d>0$, let $W$
    be its double layer with normalized diagonal $E$, and let $\rho>0$ be
    diagonal. Suppose that $E^J=\lvert\rho)(\Phi\rvert$ for some $J>0$, where
    $(\Phi\rvert=\operatorname{vec}(\Id_D)^{\mathsf T}$. Compute the source-cut
    ranks $r$ and $\ell$ and both source operators $u$ and $v$ from this same
    tensor $\mathcal U$ and weight $\rho$. Then the following are equivalent:
    \begin{enumerate}
        \item $\mathcal U$ is simple;
        \item $r\ell=d^2$;
        \item $u$ is unitary;
        \item $v$ is unitary.
    \end{enumerate}
    This is the supplied diagonal-fixed-pair specialization of
    \cite[Theorem~III.8]{Cirac2017MPU}.
    % Source lines: paper_v2.tex 495 and 563--601.
\end{theorem}

\begin{proof}\leanok
    \uses{lem:mpu_admissible_source_u_isometry,
        thm:mpu_simple_source_v_isometry,
        thm:mpu_source_gates_two_site_factorization,
        thm:mpu_simple_of_source_v_isometry}
    The supplied fixed pair gives $u^\dagger u=\Id$ and
    $d^2\leq r\ell$ before any of the four conditions is assumed. Follow the
    cycle $1\to2\to3\to4\to1$.

    If $\mathcal U$ is simple, the source-$v$ contraction for the same weight
    $\rho$ gives $v^\dagger v=\Id$, hence $r\ell\leq d^2$ and therefore
    $r\ell=d^2$. Conversely, this equality makes the standing isometry $u$
    square, hence unitary. The two-site factorization gives
    \begin{align}
        \operatorname{reindex}(U^{(2)})
         & =v\,\operatorname{swap}_{r,\ell}(u).
    \end{align}
    Reindexing preserves unitarity. Since $U^{(2)}$ and $u$ are unitary,
    cancellation makes $v$ unitary. Finally, $v^\dagger v=\Id$ gives the
    second simple contraction, while the MPU identity gives the first, so
    $\mathcal U$ is simple.
\end{proof}

\begin{theorem}[Admissible simple-tensor equivalence theorem]
    \label{thm:mpu_admissible_simple_tensor_equivalence}
    \notready
    \discussion{5708}
    % **Scope restriction (all positive bond dimensions):** the $D=1$ branch
    % uses one-dimensional transfer stabilization; the $D>1$ branch uses chosen
    % reduced full-support canonical-form-II data.
    % Source Theorem III.8 is stated under the preceding representative
    % convention without these displayed hypotheses. Documented in
    % docs/paper-gaps/mpu_canonical_form_full_support.tex.
    \uses{def:mpu_simple, def:mpu_source_uv, def:mpu_reduced_full_support_source_datum}
    Let $\mathcal U$ be an MPU tensor equipped with a reduced full-support
    source datum, and let $\rho$ and $(\Phi\rvert$ be its normalized fixed pair.
    Compute the source-cut ranks $r$ and $\ell$ and the source operators $u$ and
    $v$ from this same tensor $\mathcal U$ using $\rho$. Then the following
    conditions are equivalent:
    \begin{enumerate}
        \item $\mathcal U$ is simple;
        \item $r\ell=d^2$;
        \item $u$ is unitary;
        \item $v$ is unitary.
    \end{enumerate}
    This is the reduced-source version of
    \cite[Theorem~III.8]{Cirac2017MPU}.
    % Source lines: paper_v2.tex 563--601.
\end{theorem}

\begin{proof}
    \uses{lem:mpu_admissible_source_u_isometry,
        thm:mpu_simple_source_v_isometry,
        thm:mpu_simple_of_source_v_isometry,
        thm:mpu_source_gates_two_site_factorization}
    Lemma~\ref{lem:mpu_admissible_source_u_isometry} gives
    $u^\dagger u=\Id$ before any of the four
    conditions is assumed. Thus $u$ is a standing isometry and
    $r\ell\geq d^2$ throughout the proof.

    Suppose first that $\mathcal U$ is simple. The reduced full-support datum
    requires $\rho$ to be diagonal, so
    Theorem~\ref{thm:mpu_simple_source_v_isometry} gives
    $v^\dagger v=\Id$ for the source operator built from the same weight
    $\rho$, hence $r\ell\leq d^2$.

    If $r\ell=d^2$, the standing isometry $u$ is square and therefore unitary.
    Theorem~\ref{thm:mpu_source_gates_two_site_factorization} gives the exact
    CPSV17 two-site factorization
    \begin{align}
        \operatorname{reindex}(U^{(2)})
         & =v\,\operatorname{swap}_{r,\ell}(u).
    \end{align}
    Since $\operatorname{reindex}(U^{(2)})$ and $u$, hence also
    $\operatorname{swap}_{r,\ell}(u)$, are unitary,
    \begin{align}
        v
         & =\operatorname{reindex}(U^{(2)})
        \,\operatorname{swap}_{r,\ell}(u)^\dagger.
    \end{align}
    Thus $v$ is a product of unitaries and is unitary. This step uses the
    corrected source-cut orientation recorded in
    \texttt{docs/paper-gaps/mpu\_source\_cut\_orientation.tex}.

    Finally, suppose that $v$ is unitary. Then $v^\dagger v=\Id$, and
    Theorem~\ref{thm:mpu_simple_of_source_v_isometry} shows that
    $\mathcal U$ is simple.
\end{proof}

\begin{definition}[Standard form]
    \label{SF}
    \notready
    \discussion{6019}
    \uses{ThmFund1, def:mpu_source_uv}
    Let $\mathcal U$ be a simple tensor under the source's standing
    canonical-form-II convention. Its standard form is the depth-two form of
    the two-site block $\mathcal U_2$ determined by the source unitaries $u$ and
    $v$. This is the unrestricted definition of
    \cite[Definition~III.9]{Cirac2017MPU}.
    % Source lines: paper_v2.tex 603--622.
\end{definition}

\begin{definition}[Admissible standard form]
    \label{def:mpu_admissible_standard_form}
    \notready
    \discussion{6019}
    % **Scope restriction (reduced full-support source datum):** this standard
    % form uses Theorem~\ref{thm:mpu_admissible_simple_tensor_equivalence} for a
    % tensor carrying the datum in
    % Definition~\ref{def:mpu_reduced_full_support_source_datum}.
    % Documented in docs/paper-gaps/mpu_canonical_form_full_support.tex.
    \uses{thm:mpu_admissible_simple_tensor_equivalence, def:mpu_source_uv, def:mpu_reduced_full_support_source_datum}
    Let $\mathcal U$ be a simple tensor equipped with a reduced full-support
    source datum. The admissible standard form of the two-site blocking
    $\mathcal U_2$ is the depth-two form determined by the source unitaries
    \begin{align}
        u & \colon\mathbb C^d\otimes\mathbb C^d
        \longrightarrow\mathbb C^\ell\otimes\mathbb C^r, \\
        v & \colon\mathbb C^r\otimes\mathbb C^\ell
        \longrightarrow\mathbb C^d\otimes\mathbb C^d.
    \end{align}
    On four consecutive blocked legs its unitary is
    \begin{align}
        U & =v_{23}v_{41}u_{12}u_{34}.
    \end{align}
    % Source lines: paper_v2.tex 603--622.
\end{definition}

\begin{theorem}[Fundamental Theorem of MPU]
    \label{FundamentalMPU}
    \notready
    \discussion{6020}
    \uses{SF}
    Given two simple tensors $\mathcal U$ and $\mathcal V$ with standard forms,
    they generate the same MPU for every $N$ if and only if there are unitaries
    $x$, $y$, and $z$ relating the two standard forms by the local gauge equations
    of \cite[Eqns.~(SFuu) and (SFvv)]{Cirac2017MPU}. This is the unrestricted
    statement of the source's Fundamental Theorem of MPU.
    % Source label: FundamentalMPU. Source lines: paper_v2.tex 624--648.
\end{theorem}

\begin{theorem}[Admissible Fundamental Theorem of MPU]
    \label{thm:mpu_admissible_fundamental}
    \notready
    \discussion{6020}
    % **Scope restriction (reduced full-support source data):** both tensors
    % carry the data required by
    % Definition~\ref{def:mpu_admissible_standard_form}. Documented in
    % docs/paper-gaps/mpu_canonical_form_full_support.tex.
    \uses{def:mpu_admissible_standard_form, thm:mpu_admissible_simple_tensor_equivalence, def:mpu_reduced_full_support_source_datum}
    Let $\mathcal U$ and $\mathcal V$ be simple tensors equipped with reduced
    full-support source data and written in admissible standard form, with
    source unitaries $(u_{\mathcal U},v_{\mathcal U})$ and
    $(u_{\mathcal V},v_{\mathcal V})$. They generate the same MPU for every
    $N$,
    \begin{align}
        U^{(N)} & =V^{(N)},
    \end{align}
    if and only if their admissible standard forms are related by the local
    unitary gauges
    of \cite[Eqns.~(SFuu) and (SFvv)]{Cirac2017MPU}: there are unitaries $x$,
    $y$, and $z$ on the three corresponding source and canonical-form-II legs.
    In particular, the source unitaries obey the local relation
    \begin{align}
        u_{\mathcal V} & =(x\otimes y)u_{\mathcal U},
    \end{align}
    while the two canonical-form-II tensors are related on their virtual leg by
    $z$ and on their source legs by $x$ and $y$.
\end{theorem}

\begin{theorem}[Rank of a Kronecker product]
    \label{thm:mpu_kronecker_rank}
    \lean{Matrix.rank_kronecker}
    \leanok
    Let $A$ and $B$ be finite matrices over a field. Then
    \begin{align}
        \operatorname{rank}(A\otimes B)
         & =\operatorname{rank}(A)\operatorname{rank}(B).
    \end{align}
\end{theorem}

\begin{proof}\leanok
    Identify the image of the Kronecker product with the tensor product of the
    two images and multiply their dimensions.
\end{proof}

\begin{definition}[Specified-tensor source-index value]
    \label{def:mpu_specified_source_index_value}
    \lean{MPOTensor.sourceIndexValue}
    \leanok
    \uses{def:mpu_right_rank, def:mpu_left_rank}
    For a specified tensor $\mathcal U$ whose right and left source-cut ranks
    $r$ and $\ell$ are positive, define
    \begin{align}
        \operatorname{ind}_{\mathrm{src}}(\mathcal U)
         & =\frac12(\log_2 r-\log_2\ell).
    \end{align}
    This definition makes no choice of a simple blocking.
    % Auxiliary to paper_v2.tex lines 681--688; issue 6255.
\end{definition}

\begin{theorem}[Auxiliary source-index identities]
    \label{thm:mpu_specified_source_index_identities}
    \lean{MPOTensor.sourceIndexValue_eq_of_common_rank_scale,
        MPOTensor.sourceIndexValue_eq_add_of_common_rank_product,
        MPOTensor.sourceIndexValue_eq_logb_rightRank_div,
        MPOTensor.sourceIndexValue_eq_neg_logb_leftRank_div}
    \leanok
    \uses{def:mpu_specified_source_index_value, def:mpu_right_rank, def:mpu_left_rank}
    Let $\mathcal U$ and $\mathcal V$ be specified tensors. Suppose that their
    source ranks obey
    \begin{align}
        r_{\mathcal V}    & =c r_{\mathcal U},   &
        \ell_{\mathcal V} & =c\ell_{\mathcal U},
    \end{align}
    where $c$, $r_{\mathcal U}$, and $\ell_{\mathcal U}$ are positive. Then
    \begin{align}
        \operatorname{ind}_{\mathrm{src}}(\mathcal V)
         & =\operatorname{ind}_{\mathrm{src}}(\mathcal U).
    \end{align}
    Let $\mathcal U$, $\mathcal V$, and $\mathcal W$ instead have positive
    source ranks and suppose that
    \begin{align}
        r_{\mathcal W}    & =c r_{\mathcal U}r_{\mathcal V},      \\
        \ell_{\mathcal W} & =c\ell_{\mathcal U}\ell_{\mathcal V},
    \end{align}
    where $c$ is a positive integer. Then
    \begin{align}
        \operatorname{ind}_{\mathrm{src}}(\mathcal W)
         & =\operatorname{ind}_{\mathrm{src}}(\mathcal U)
        +\operatorname{ind}_{\mathrm{src}}(\mathcal V).
    \end{align}
    If $d$, $r_{\mathcal U}$, and $\ell_{\mathcal U}$ are positive and a
    specified tensor has $r_{\mathcal U}\ell_{\mathcal U}=d^2$, then
    \begin{align}
        \operatorname{ind}_{\mathrm{src}}(\mathcal U)
         & =\log_2\frac{r_{\mathcal U}}d
        =-\log_2\frac{\ell_{\mathcal U}}d.
    \end{align}
    % Auxiliary to paper_v2.tex lines 681--704 and IndexTh lines 837--845;
    % issues 6255 and 6923.
\end{theorem}

\begin{proof}\leanok
    \uses{def:mpu_specified_source_index_value}
    Expand the definition. For common scaling, apply
    $\log_2(cx)=\log_2c+\log_2x$ to both ranks and cancel the two copies of
    $\log_2c$. For the three-tensor identity, apply the product formula twice
    to each rank and cancel the common term. For the equivalent formulas,
    apply the same product identity to $r_{\mathcal U}\ell_{\mathcal U}=d^2$
    and use $\log_2(x/y)=\log_2x-\log_2y$.
\end{proof}

\begin{theorem}[Conditional blocking invariance of the specified source-index value]
    \label{thm:mpu_specified_source_index_blocking_invariance}
    \lean{MPOTensor.sourceIndexValue_blockTensor_eq_of_products}
    \leanok
    \uses{def:mpu_specified_source_index_value}
    Let $d>0$ and $k_0\leq k$. Suppose that the endpoint source ranks satisfy
    \begin{align}
        r_{k_0}\ell_{k_0} & =d^{2k_0}, \\
        r_k\ell_k         & =d^{2k}.
    \end{align}
    Then the specified source-index values agree:
    \begin{align}
        \operatorname{ind}_{\mathrm{src}}(\mathcal U_k)
         & =\operatorname{ind}_{\mathrm{src}}(\mathcal U_{k_0}).
    \end{align}
    This statement assumes the endpoint products; it neither selects a simple
    blocking nor defines the public index.
    % Auxiliary to paper_v2.tex lines 690--704; issue 6958.
    % **Local fix (rank-product exponent):** Source line 703 prints $d^k$; the
    % consistent blocked identity is $d^{2k}$. See
    % docs/paper-gaps/mpu_blocking_rank_product_exponent.tex.
\end{theorem}

\begin{proof}\leanok
    \uses{thm:mpu_specified_source_index_identities, thm:mpu_source_cut_rank_blocking_saturation}
    Exact blocking growth gives the common positive factor $d^{k-k_0}$:
    \begin{align}
        r_k    & =d^{k-k_0}r_{k_0},    \\
        \ell_k & =d^{k-k_0}\ell_{k_0}.
    \end{align}
    Its logarithm occurs once in each rank term and therefore cancels from
    their difference.
\end{proof}

\begin{definition}[Index]
    \label{def:index}
    \notready
    \discussion{6014}
    \uses{ThmFund1, thm:mpu_bounded_simple_blocking}
    Let $k$ be any blocking length for which $\mathcal U_k$ is simple, and let
    $r$ and $\ell$ be its right and left source-cut ranks. The source index is
    \begin{align}
        \operatorname{ind}(\mathcal U)
         & =\frac12(\log_2 r-\log_2\ell).
    \end{align}
    This is the unrestricted definition of
    \cite[Definition~IV.1]{Cirac2017MPU}.
    % Source lines: paper_v2.tex 681--688.
\end{definition}

\begin{definition}[Admissible index]
    \label{def:mpu_admissible_index}
    \notready
    \discussion{6014}
    % **Scope restriction (admissible simple block):** the index is defined only
    % from a simple block carrying its own reduced full-support source datum.
    % Documented in docs/paper-gaps/mpu_canonical_form_full_support.tex.
    \uses{thm:mpu_admissible_simple_tensor_equivalence, thm:mpu_bounded_simple_blocking, thm:mpu_right_rank_bound, thm:mpu_left_rank_bound, def:mpu_reduced_full_support_source_datum}
    Let $k$ be a blocking length for which $\mathcal U_k$ is simple and is
    equipped with a reduced full-support source datum, and let $r$ and $\ell$
    be its right and left source-cut ranks. Define
    \begin{align}
        \operatorname{ind}(\mathcal U)
         & =\frac12(\log_2 r-\log_2\ell).
    \end{align}
    Since $r\ell=d^{2k}$, this is equivalently
    \begin{align}
        \operatorname{ind}(\mathcal U)
         & =\log_2\frac{r}{d^k}
        =-\log_2\frac{\ell}{d^k}.
    \end{align}
    % Source lines: paper_v2.tex 681--688.
\end{definition}

\begin{proposition}[Well-definedness of the index]
    \label{index-well-defined}
    \notready
    \discussion{6014}
    \uses{def:index, ThmFund1, thm:mpu_source_cut_rank_blocking, thm:mpu_source_cut_rank_blocking_saturation, thm:mpu_specified_source_index_blocking_invariance, thm:mpu_all_later_simple_blockings}
    The source index is independent of the chosen simple blocking. More
    precisely, if $k\geq k_0$ and both $\mathcal U_{k_0}$ and $\mathcal U_k$
    are simple, then
    \begin{align}
        r_k    & =d^{k-k_0}r_{k_0},    &
        \ell_k & =d^{k-k_0}\ell_{k_0}.
    \end{align}
    This is the unrestricted statement of
    \cite[Proposition~IV.2]{Cirac2017MPU}.
    % Source lines: paper_v2.tex 690--704.
\end{proposition}

\begin{proposition}[Well-definedness of the admissible index]
    \label{prop:mpu_admissible_index_well_defined}
    \notready
    \discussion{6014}
    % **Scope restriction (admissible simple blocks):** independence is proved
    % only between simple blocks carrying their own reduced full-support source
    % data. Documented in docs/paper-gaps/mpu_canonical_form_full_support.tex.
    \uses{def:mpu_admissible_index, thm:mpu_admissible_simple_tensor_equivalence, thm:mpu_source_cut_rank_blocking, thm:mpu_all_later_simple_blockings, def:mpu_reduced_full_support_source_datum}
    The admissible index is independent of the chosen admissible simple
    blocking. More precisely, if $k\geq k_0$, both $\mathcal U_{k_0}$ and
    $\mathcal U_k$ are simple, and each is equipped with a reduced full-support
    source datum, then
    \begin{align}
        r_k    & =d^{k-k_0}r_{k_0},    &
        \ell_k & =d^{k-k_0}\ell_{k_0}.
    \end{align}
    % Source lines: paper_v2.tex 690--704.
\end{proposition}

\begin{lemma}[Tensor-product additivity]
    \label{lem:mpu_index_tensor}
    \notready
    \discussion{6015}
    \uses{def:index, index-well-defined, thm:mpu_tensor_product_source_ranks}
    For MPU tensors under the source's standing canonical-form convention,
    \begin{align}
        \operatorname{ind}(\mathcal U\otimes\mathcal V)
         & =\operatorname{ind}(\mathcal U)+\operatorname{ind}(\mathcal V).
    \end{align}
    This is the unrestricted tensor-product statement in
    \cite[Theorem~IV.6(ii)]{Cirac2017MPU}.
    % Source lines: paper_v2.tex 824--833.
\end{lemma}

\begin{lemma}[Admissible tensor-product additivity]
    \label{lem:mpu_admissible_index_tensor}
    \notready
    \discussion{6015}
    % **Scope restriction (common admissible block):** all three indices are
    % computed from actual simple blocks carrying their own source data.
    % Documented in docs/paper-gaps/mpu_canonical_form_full_support.tex.
    \uses{def:mpu_admissible_index, prop:mpu_admissible_index_well_defined, thm:mpu_tensor_product_source_ranks, def:mpu_reduced_full_support_source_datum}
    Suppose there is one positive blocking length $k$ such that
    $\mathcal U_k$, $\mathcal V_k$, and
    $(\mathcal U\otimes\mathcal V)_k$ are simple and each carries its own
    reduced full-support source datum. Define all three indices from these
    blocks. Then
    \begin{align}
        \operatorname{ind}(\mathcal U\otimes\mathcal V)
         & =\operatorname{ind}(\mathcal U)+\operatorname{ind}(\mathcal V).
    \end{align}
    This is the common-block reduced-source version of
    \cite[Theorem~IV.6(ii)]{Cirac2017MPU}.
    % Source lines: paper_v2.tex 824--833.
\end{lemma}

\begin{lemma}[Four-site composition source-rank bounds]
    \label{lem:mpu_four_site_composition_source_rank_bounds}
    \lean{MPOTensor.rightRank_blockTensor_mulTensor_four_le,
        MPOTensor.leftRank_blockTensor_mulTensor_four_le}
    \leanok
    \uses{def:mpdo_operator_product, def:mpdo_physical_blocking, def:mpu_right_rank, def:mpu_left_rank}
    Let $\mathcal U$ and $\mathcal V$ be arbitrary matrix product operator
    tensors of common physical dimension $d$. For their product tensor,
    four-site blocking satisfies
    \begin{align}
        r[(\mathcal U\mathbin{\cdot}\mathcal V)_4]
         & \leq d^2 r[\mathcal U]r[\mathcal V],       \\
        \ell[(\mathcal U\mathbin{\cdot}\mathcal V)_4]
         & \leq d^2 \ell[\mathcal U]\ell[\mathcal V].
    \end{align}
    No unitarity or simplicity assumption is required.
    % Source lines: paper_v2.tex 837--845; issue 6999.
\end{lemma}

\begin{proof}\leanok
    \uses{thm:mpu_source_v_open_two_site,
        thm:mpu_source_u_open_two_site,
        thm:mpdo_blocking_product,
        thm:mpdo_block_two_as_positive_blocking,
        thm:mpu_source_ranks_physical_reindexing,
        lem:mpu_iterated_physical_blocking}
    Use the canonical physical reindexing to identify iterated two-site
    blocking with direct four-site blocking. For the right cut, write each
    two-site block in the $v$--$Y_1$--$Y_2$ form. Group $v$ with $Y_2$ on one
    side of the cut and retain $Y_1$ on the other. The cut therefore carries
    one physical leg and one right source leg from each tensor. Thus, for
    suitable rectangular matrices $A_{1,4}$ and $B_{1,4}$,
    \begin{align}
        M_1((\mathcal U\mathbin{\cdot}\mathcal V)_4) & =A_{1,4}B_{1,4},
    \end{align}
    where the columns of $A_{1,4}$ and the rows of $B_{1,4}$ are indexed by
    \begin{align}
        \mathcal K_1
         & =(\mathbb C^d\otimes\mathbb C^{r[\mathcal U]})
        \otimes
        (\mathbb C^d\otimes\mathbb C^{r[\mathcal V]}).
    \end{align}
    Consequently,
    \begin{align}
        r[(\mathcal U\mathbin{\cdot}\mathcal V)_4]
         & =\operatorname{rank}(A_{1,4}B_{1,4})
        \leq\dim(\mathcal K_1)
        =d^2r[\mathcal U]r[\mathcal V].
    \end{align}
    For the left cut, write each two-site block in the
    $X_2$--$u$--$X_1$ form. Retain $X_2$ on one side and group $u$ with $X_1$
    on the other. The cut carries one physical leg and one left source leg from
    each tensor. For suitable rectangular matrices $A_{2,4}$ and $B_{2,4}$,
    \begin{align}
        M_2((\mathcal U\mathbin{\cdot}\mathcal V)_4) & =A_{2,4}B_{2,4},
    \end{align}
    where the intermediate space is
    \begin{align}
        \mathcal K_2
         & =(\mathbb C^d\otimes\mathbb C^{\ell[\mathcal U]})
        \otimes
        (\mathbb C^d\otimes\mathbb C^{\ell[\mathcal V]}).
    \end{align}
    Hence
    \begin{align}
        \ell[(\mathcal U\mathbin{\cdot}\mathcal V)_4]
         & =\operatorname{rank}(A_{2,4}B_{2,4})
        \leq\dim(\mathcal K_2)
        =d^2\ell[\mathcal U]\ell[\mathcal V].
    \end{align}
    Both source-cut ranks are unchanged by the canonical physical relabeling.
\end{proof}

\begin{lemma}[Composition additivity]
    \label{lem:mpu_index_composition}
    \notready
    \discussion{6016}
    \uses{SF, def:index, index-well-defined, lem:mpu_four_site_composition_source_rank_bounds}
    Let $\mathcal W$ generate the composition of the MPUs generated by
    $\mathcal U$ and $\mathcal U'$. Under the source's standing canonical-form
    convention,
    \begin{align}
        \operatorname{ind}(\mathcal W)
         & =\operatorname{ind}(\mathcal U)+\operatorname{ind}(\mathcal U').
    \end{align}
    This is the unrestricted composition statement in
    \cite[Theorem~IV.6(ii)]{Cirac2017MPU}.
    % Source lines: paper_v2.tex 833--845.
\end{lemma}

\begin{lemma}[Admissible composition additivity]
    \label{lem:mpu_admissible_index_composition}
    \notready
    \discussion{6016}
    % **Scope restriction (six admissible blocks):** every actual block used in
    % the source rank computation carries its own source datum. Documented in
    % docs/paper-gaps/mpu_canonical_form_full_support.tex.
    \uses{def:mpu_admissible_standard_form, def:mpu_admissible_index, prop:mpu_admissible_index_well_defined, thm:mpu_source_cut_rank_blocking, lem:mpu_four_site_composition_source_rank_bounds, def:mpu_reduced_full_support_source_datum}
    Suppose there is a positive integer $k$ such that
    $\mathcal U_k$, $\mathcal U'_k$, and the composition tensor
    $\mathcal W_k$ are simple and each carries its own reduced full-support
    source datum. Suppose in addition that $\mathcal U_{2k}$,
    $\mathcal U'_{2k}$, and $\mathcal W_{4k}$ are simple and each carries its
    own reduced full-support source datum. Define the three indices initially
    from $\mathcal U_k$, $\mathcal U'_k$, and $\mathcal W_k$; the supplied data
    at $2k$ and $4k$ make the later comparison admissible. If $\mathcal W$
    generates the composition of the MPUs generated by $\mathcal U$ and
    $\mathcal U'$, then
    \begin{align}
        \operatorname{ind}(\mathcal W)
         & =\operatorname{ind}(\mathcal U)+\operatorname{ind}(\mathcal U').
    \end{align}
    This is the six-block reduced-source version of
    \cite[Theorem~IV.6(ii)]{Cirac2017MPU}.
    % Source lines: paper_v2.tex 833--845.
\end{lemma}

\begin{definition}[MPU canonical form]
    \label{def:mpu_canonical_form}
    \lean{MPSTensor.IsMPUCanonicalBlock,
        MPSTensor.MPUCanonicalFormData,
        MPSTensor.IsMPUCanonicalForm}
    \leanok
    \uses{def:transfer_map, def:block_diagonal_tensor}
    A tensor $A$ of physical dimension $d$ and bond dimension $D$ is in MPU
    canonical form if there are positive block dimensions $D_k$, weights
    $\mu_k\in\mathbb C$, tensors $A_k$ of bond dimension $D_k$, and a
    coisometry
    \begin{align}
        C\colon\mathbb C^D                & \longrightarrow
        \bigoplus_{k=1}^r\mathbb C^{D_k}, &
        CC^\dagger                        & =\Id,
    \end{align}
    such that
    \begin{align}
        A^i=C^\dagger\left(\bigoplus_{k=1}^r\mu_k A_k^i\right)C.
    \end{align}
    Every $A_k$ is irreducible and the spectral radius of its transfer map is
    one. The blocks may be periodic, so no uniqueness condition is imposed on
    their peripheral eigenvalues. This is the canonical form of
    \cite[lines~259--267]{Cirac2017MPU}; normal tensors form a strictly stronger
    class.
\end{definition}

\begin{definition}[Strict equivalence]
    \label{def:strictly-equivalent-tensors}
    \lean{MPOTensor.StrictlyEquivalent}
    \leanok
    \uses{def:mpu, def:mpu_physical_reindexing, def:mpu_canonical_form}
    Fix a virtual bond dimension $D$. Let $\mathcal U$ and $\mathcal V$ be
    canonical-form MPU tensors with physical dimensions $d_a$ and $d_b$. Given
    an equality $h\colon d_a=d_b$, reindex $\mathcal U$ from $d_a$ to $d_b$
    along $h^{-1}$. The tensors are strictly equivalent with respect to $h$ if
    there is a continuous path $\mathcal W(p)$, $p\in[0,1]$, in the fixed-bond
    MPU locus such that
    \begin{align}
        \mathcal W(0) & =\operatorname{reind}_{h^{-1}}(\mathcal U), \\
        \mathcal W(1) & =\mathcal V.
    \end{align}
    Only the canonical-form requirement is relaxed along the path. This is the
    fixed-bond interpretation of
    \cite[lines~708--714]{Cirac2017MPU}; it does not compare unequal raw bond
    dimensions.
\end{definition}

\begin{theorem}[Normalized transfer data under identity ancillas]
    \label{thm:mpu_identity_ancilla_reduced_cfii}
    \lean{MPOTensor.normalizedFlattening_tensorPhysicalId,
        MPOTensor.transferMap_normalizedFlattening_tensorPhysicalId,
        MPOTensor.tensorPhysicalIdCFIIData,
        MPOTensor.hasFullSupport_tensorPhysicalIdCFIIData}
    \leanok
    \uses{def:mpu, def:mpu_normalized_flattening, def:cpsv_cfii, def:mpu_full_support}
    Let $x>0$. After adjoining an identity ancilla of dimension $x$, the
    normalized flattening is the physical reindexing of the original flattening
    with a normalized diagonal ancilla alphabet. Its transfer map is exactly the
    original transfer map. Consequently, chosen canonical-form-II data with full
    support determine canonical-form-II data with full support for
    $\mathcal U^{(x)}=\mathcal U\otimes\Id_x$.

    The product-index equivalence has the orientation
    \begin{align}
        ((i,a),(j,b)) & \longmapsto((i,j),(a,b)).
    \end{align}
    This result preserves the normalized transfer data. It does not identify the
    old and new source cuts, ranks, or chosen compact-SVD factors.
    % Source lines: paper_v2.tex 706--724; normalized transfer lines 336--340.
\end{theorem}

\begin{proof}\leanok
    \uses{def:mpu_physical_reindexing}
    Only the $x$ diagonal ancilla letters are nonzero, each with coefficient
    $x^{-1/2}$. Hence their total transfer contribution is
    \begin{align}
        x\left(x^{-1/2}\right)^2 & =1.
    \end{align}
    Thus the transfer map and each diagonal positive fixed matrix are unchanged.
    The block dimensions, weights, and ambient coisometry are unchanged as well,
    so full support follows from the same block-dimension sum. The displayed
    physical reindexing gives the normalized flattening of the enlarged MPO.
\end{proof}

\begin{definition}[Equivalence after blocking and ancillas]
    \label{def:equivalent-tensors}
    \lean{MPOTensor.Equivalent}
    \leanok
    \uses{def:mpu, def:strictly-equivalent-tensors, lem:mpu_blocking}
    Fix a virtual bond dimension $D$. Let $\mathcal U$ and $\mathcal V$ be MPU
    tensors with physical dimensions $d_a$ and $d_b$. They are equivalent if
    there are a
    positive blocking length $k$ and coprime positive integers $p_a,p_b$ with
    \begin{align}
        p_a d_a & =p_b d_b
    \end{align}
    such that $(\mathcal U^{(p_a)})_k$ and $(\mathcal V^{(p_b)})_k$ are strictly
    equivalent. Here the positive-dimensional identity ancillas are attached
    before the common blocking, and
    \begin{align}
        \mathcal U^{(p)} & =\mathcal U\otimes\Id_p.
    \end{align}
    This is the fixed-bond interpretation of
    \cite[lines~716--724]{Cirac2017MPU}; it does not introduce a stabilization
    relation for unequal raw bond dimensions.
\end{definition}

\begin{definition}[Admissible equivalence and path datum]
    \label{def:mpu_admissible_equivalence_datum}
    \notready
    % **Scope restriction (no existence or preservation theorem):** this datum
    % records every reduced source hypothesis needed along a chosen path. It does
    % not assert preservation under blocking, ancillas, products, composition, or
    % continuous variation. Documented in
    % docs/paper-gaps/mpu_canonical_form_full_support.tex.
    \uses{def:strictly-equivalent-tensors, def:equivalent-tensors, def:mpu_reduced_full_support_source_datum}
    An admissible strict-equivalence datum for a path
    $[0,1]\ni p\mapsto\mathcal W(p)$ consists of a common positive blocking
    length $m$ such that the actual blocked path tensor $\mathcal W(p)_m$ is
    simple and carries a reduced full-support source datum for every $p$. The
    blocked tensors and their source data vary continuously with $p$, after
    fixed identifications of the finite-dimensional spaces involved. No
    equivalent or reduced representative may replace $\mathcal W(p)_m$. In
    particular, the actual blocked endpoint tensors carry their own source data.

    When the path is used under a specified symmetry $\mathcal S$, the datum also
    includes continuously varying reduced full-support source data for the fixed
    family of comparison paths
    \begin{align}
        \overline{\mathcal W(p)_m},\qquad
        \mathcal W(p)_m^\sharp,\qquad
        \mathcal W(p)_m^{\mathsf T},\qquad
        \mathcal S[\mathcal W(p)_m].
    \end{align}
    Every tensor in this family is simple. This requirement is part of the datum
    independently of which comparisons a later proof uses; no preservation under
    conjugation, physical adjoint, transposition, or $\mathcal S$ is assumed.

    Two tensors are admissibly strictly equivalent if their strict equivalence
    is witnessed by such a path. They are admissibly equivalent if the blocked,
    ancilla-enlarged endpoints in Definition~\ref{def:equivalent-tensors} are
    joined by an admissible strict-equivalence path. The source arguments
    motivating these additional path hypotheses are Proposition~IV.5,
    Theorem~IV.6, Corollary~IV.7, and the lemma following the
    symmetry-equivalence definitions in \cite{Cirac2017MPU}; the datum itself is
    introduced here and has no source-labelled counterpart.
    % Source lines: paper_v2.tex 733--910 and 1366--1380.
\end{definition}

\begin{theorem}[Generic two-projection angle block]
    \label{thm:mpu_two_projection_angle_block}
    \lean{LinearMap.IsSymmetricProjection.rangeCompression_isSymmetric,
        LinearMap.IsSymmetricProjection.angleDefect_block,
        LinearMap.IsSymmetricProjection.exists_orthonormal_angle_block}
    \leanok

    Let $P$ and $Q$ be orthogonal projections on a complex inner-product space.
    The compression of $Q$ to $\operatorname{ran}P$ is self-adjoint. If a unit
    vector $x\in\operatorname{ran}P$ satisfies
    \begin{align}
        PQx & =\mu x, & 0 & <\mu<1,
    \end{align}
    then there is a unit vector $w\perp x$ such that
    \begin{align}
        Px & =x,                             & Pw & =0, \\
        Qx & =\mu x+\sqrt{\mu(1-\mu)}\,w,    &
        Qw & =\sqrt{\mu(1-\mu)}\,x+(1-\mu)w.
    \end{align}
    Thus $\operatorname{span}\{x,w\}$ is the generic two-dimensional angle
    block. This lemma concerns one eigenvalue in $(0,1)$; it does not include the
    endpoint eigenspaces or an orthogonal direct-sum decomposition of the ambient
    space.
    % Formalization status: endpoint spaces and global Jordan assembly remain not ready.
    % Source: paper_v2.tex, Proposition IV.5 (prop:continuity-index), lines 756--764.
    % Local correction: line 762 repeats P where Q is intended; see
    % docs/paper-gaps/mpu_two_projection_projector_typo.tex.
\end{theorem}

\begin{proof}\leanok
    Set $d=Qx-\mu x$. Symmetry and idempotence give
    \begin{align}
        Pd               & =0,                    & \langle x,d\rangle & =0, &
        \lVert d\rVert^2 & =\mu(1-\mu),                                        \\
        Qd               & =\mu(1-\mu)x+(1-\mu)d.
    \end{align}
    Since $0<\mu<1$, the vector
    $w=d/\sqrt{\mu(1-\mu)}$ is well defined and has the stated properties.
\end{proof}

\begin{definition}[Compression spectral data]
    \label{def:mpu_two_projection_compression_spectral_data}
    \lean{LinearMap.IsSymmetricProjection.rangeCompressionEigenvalues,
        LinearMap.IsSymmetricProjection.rangeCompressionEigenbasis}
    \leanok
    \uses{thm:mpu_two_projection_angle_block}
    For orthogonal projections $P$ and $Q$ on a finite-dimensional complex
    inner-product space, let $(x_i)_i$ be an orthonormal eigenbasis of the
    compression of $Q$ to $\operatorname{ran}P$, with ordered real eigenvalues
    $(\mu_i)_i$.
    % Source lines: paper_v2.tex 759--762.
\end{definition}

\begin{theorem}[Compression-spectrum endpoint and angle-block trichotomy]
    \label{thm:mpu_two_projection_compression_spectrum}
    \lean{LinearMap.IsSymmetricProjection.rangeCompression_apply_eigenbasis,
        LinearMap.IsSymmetricProjection.rangeCompressionEigenbasis_apply_self,
        LinearMap.IsSymmetricProjection.rangeCompressionEigenbasis_apply_compression,
        LinearMap.IsSymmetricProjection.rangeCompressionEigenvalues_mem_unitInterval,
        LinearMap.IsSymmetricProjection.rangeCompressionEigenvalues_eq_zero_or_eq_one_or_exists_angle_block}
    \leanok
    \uses{def:mpu_two_projection_compression_spectral_data, thm:mpu_two_projection_angle_block}
    Let $P$ and $Q$ be orthogonal projections on a finite-dimensional complex
    inner-product space, with compression spectral data $(x_i,\mu_i)$. Then, for
    every $i$, $\mu_i\in[0,1]$ and exactly one of the following alternatives
    holds:
    \begin{enumerate}
        \item $\mu_i=0$ and $Qx_i=0$;
        \item $\mu_i=1$ and $Qx_i=x_i$;
        \item $0<\mu_i<1$, and $x_i$ belongs to the two-dimensional angle block
              described in Theorem~\ref{thm:mpu_two_projection_angle_block}.
    \end{enumerate}
    This theorem classifies the compression eigenvectors.
    % Source lines: paper_v2.tex 759--762.
    % Scope: no pairwise orthogonality of angle blocks, no ambient orthogonal
    % direct sum, and no commutation on the remaining summand.
\end{theorem}

\begin{proof}\leanok
    \uses{def:mpu_two_projection_compression_spectral_data, thm:mpu_two_projection_angle_block}
    Apply the finite-dimensional spectral theorem to the self-adjoint compression
    of $Q$ to $\operatorname{ran}P$. For a unit compression eigenvector $x$ with
    eigenvalue $\mu$, the angle-defect identity gives
    \begin{align}
        \lVert Qx-\mu x\rVert^2 & =\mu(1-\mu),
    \end{align}
    hence $0\leq\mu\leq1$. At $\mu=0$ or $1$ the defect vanishes, giving the two
    common one-dimensional blocks. For $0<\mu<1$, apply
    Theorem~\ref{thm:mpu_two_projection_angle_block}.
\end{proof}

\begin{definition}[Compression defect block]
    \label{def:mpu_two_projection_angle_defect_block}
    \lean{LinearMap.IsSymmetricProjection.angleDefect,
        LinearMap.IsSymmetricProjection.angleDefectBlock}
    \leanok
    \uses{def:mpu_two_projection_compression_spectral_data}
    For compression spectral data $(x_i,\mu_i)$, define
    \begin{align}
        d_i & =Qx_i-\mu_i x_i,                 \\
        B_i & =\operatorname{span}\{x_i,d_i\}.
    \end{align}
\end{definition}

\begin{theorem}[Orthogonality of compression angle blocks]
    \label{thm:mpu_two_projection_angle_block_orthogonality}
    \lean{LinearMap.IsSymmetricProjection.rangeCompressionEigenbasis_inner_angleDefect_eq_zero,
        LinearMap.IsSymmetricProjection.angleDefect_inner_rangeCompressionEigenbasis_eq_zero,
        LinearMap.IsSymmetricProjection.angleDefect_inner_angleDefect,
        LinearMap.IsSymmetricProjection.angleDefect_inner_angleDefect_eq_zero,
        LinearMap.IsSymmetricProjection.finrank_angleDefectBlock_le_two,
        LinearMap.IsSymmetricProjection.map_angleDefectBlock_le_first_projection,
        LinearMap.IsSymmetricProjection.map_angleDefectBlock_le_second_projection,
        LinearMap.IsSymmetricProjection.angleDefectBlock_isOrtho}
    \leanok
    \uses{def:mpu_two_projection_angle_defect_block}
    For all indices $i,j$,
    \begin{align}
        \langle x_i,d_j\rangle & =0,                                    \\
        \langle d_i,x_j\rangle & =0,                                    \\
        \langle d_i,d_j\rangle
                               & =\mu_j(1-\mu_j)\langle x_i,x_j\rangle.
    \end{align}
    In particular, $\langle d_i,d_j\rangle=0$ when $i\ne j$, including when
    $\mu_i=\mu_j$. The block $B_i$ has dimension at most two, is invariant under
    both $P$ and $Q$, and satisfies $B_i\perp B_j$ for distinct indices.
    % This does not assert an ambient direct sum or any property of the remaining
    % orthogonal complement.
    % Source lines: paper_v2.tex 759--762.
\end{theorem}

\begin{proof}\leanok
    \uses{def:mpu_two_projection_angle_defect_block, thm:mpu_two_projection_compression_spectrum, thm:mpu_two_projection_angle_block}
    Since $Px_i=x_i$, symmetry of $P$ and the compression eigenvalue equation give
    \begin{align}
        \langle x_i,d_j\rangle
         & =\langle x_i,Qx_j\rangle-\mu_j\langle x_i,x_j\rangle     \\
         & =\langle Px_i,Qx_j\rangle-\mu_j\langle x_i,x_j\rangle    \\
         & =\langle x_i,PQx_j\rangle-\mu_j\langle x_i,x_j\rangle=0.
    \end{align}
    Conjugate symmetry gives $\langle d_i,x_j\rangle=0$. For the exact defect
    pairing, symmetry of $Q$ and the local action formula give
    \begin{align}
        Qd_j & =\mu_j(1-\mu_j)x_j+(1-\mu_j)d_j,                      \\
        \langle d_i,d_j\rangle
             & =\langle Qx_i,d_j\rangle-\mu_i\langle x_i,d_j\rangle, \\
        \langle Qx_i,d_j\rangle-\mu_i\langle x_i,d_j\rangle
             & =\langle x_i,Qd_j\rangle,                             \\
        \langle x_i,Qd_j\rangle
             & =\mu_j(1-\mu_j)\langle x_i,x_j\rangle
        +(1-\mu_j)\langle x_i,d_j\rangle,                            \\
        \langle d_i,d_j\rangle
             & =\mu_j(1-\mu_j)\langle x_i,x_j\rangle.
    \end{align}
    Orthogonality of the generated spans follows by linearity. The action formulas
    \begin{align}
        Px_i & =x_i,                           \\
        Pd_i & =0,                             \\
        Qx_i & =\mu_i x_i+d_i,                 \\
        Qd_i & =\mu_i(1-\mu_i)x_i+(1-\mu_i)d_i
    \end{align}
    show that both projections preserve
    $B_i=\operatorname{span}\{x_i,d_i\}$. Its dimension bound follows from
    its two generators.
\end{proof}
\begin{definition}[Defect-block sum and complementary summand]
    \label{def:mpu_two_projection_angle_defect_sum}
    \lean{LinearMap.IsSymmetricProjection.angleDefectSum,
        LinearMap.IsSymmetricProjection.angleDefectComplement}
    \leanok
    \uses{def:mpu_two_projection_angle_defect_block}
    For the compression defect blocks $(B_i)_i$, define
    \begin{align}
        B & =\bigvee_i B_i, & K & =B^\perp.
    \end{align}
    The index $i$ ranges over every vector in the chosen compression eigenbasis,
    including repeated eigenvalues and the endpoint eigenvalues $0$ and $1$.
    % Source lines: paper_v2.tex 759--762.
\end{definition}

\begin{theorem}[Orthogonal defect-block sum and commuting remainder]
    \label{thm:mpu_two_projection_angle_defect_sum}
    \lean{LinearMap.IsSymmetricProjection.angleDefectBlock_orthogonalFamily,
        LinearMap.IsSymmetricProjection.angleDefectBlock_isInternal,
        LinearMap.IsSymmetricProjection.angleDefectSum_isCompl_angleDefectComplement,
        LinearMap.IsSymmetricProjection.range_le_angleDefectSum,
        LinearMap.IsSymmetricProjection.map_angleDefectSum_le_first_projection,
        LinearMap.IsSymmetricProjection.map_angleDefectSum_le_second_projection,
        LinearMap.IsSymmetricProjection.map_angleDefectComplement_le_first_projection,
        LinearMap.IsSymmetricProjection.map_angleDefectComplement_le_second_projection,
        LinearMap.IsSymmetricProjection.projections_comp_apply_eq_zero_of_mem_angleDefectComplement,
        LinearMap.IsSymmetricProjection.firstProjectionOnAngleDefectComplement,
        LinearMap.IsSymmetricProjection.secondProjectionOnAngleDefectComplement,
        LinearMap.IsSymmetricProjection.firstProjectionOnAngleDefectComplement_isSymmetricProjection,
        LinearMap.IsSymmetricProjection.secondProjectionOnAngleDefectComplement_isSymmetricProjection,
        LinearMap.IsSymmetricProjection.firstProjectionOnAngleDefectComplement_commute_secondProjection}
    \leanok
    \uses{def:mpu_two_projection_angle_defect_sum}
    The blocks $(B_i)_i$ form an orthogonal internal direct sum in $B$, and
    \begin{align}
        E & =B\mathbin{\perp\!\oplus}K.
    \end{align}
    Both $B$ and $K$ are invariant under $P$ and $Q$. Moreover,
    $\operatorname{ran}P\subseteq B$, so for every $z\in K$,
    \begin{align}
        Pz & =0, & P(Qz) & =0, & Q(Pz) & =0.
    \end{align}
    Consequently the restrictions $P|_K$ and $Q|_K$ are orthogonal projections
    and commute. No commutation of $P$ and $Q$ on the ambient space is assumed.
    % Source lines: paper_v2.tex 759--762.
    % Local correction: line 762 repeats P where Q is intended; see
    % docs/paper-gaps/mpu_two_projection_projector_typo.tex.
\end{theorem}

\begin{proof}\leanok
    \uses{thm:mpu_two_projection_angle_block_orthogonality}
    Pairwise orthogonality of the blocks implies independence, while their
    supremum is $B$ by definition. Hence they form an internal direct sum in
    $B$. The standard orthogonal-complement decomposition gives
    $E=B\mathbin{\perp\!\oplus}B^\perp$.

    The local action formulas show that each $B_i$ is invariant under both
    projections, hence so is $B$. Symmetry of each projection then makes
    $K=B^\perp$ invariant. The compression eigenvectors form a basis of
    $\operatorname{ran}P$, and each lies in its block $B_i$; therefore
    $\operatorname{ran}P\subseteq B$. Thus $P$ vanishes on $K$. Since $Qz\in K$
    whenever $z\in K$, both $P(Qz)$ and $Q(Pz)$ vanish. The two restricted
    projections therefore commute.
\end{proof}
\begin{definition}[Reduced projection]
    \label{def:mpu_two_projection_reduced_projection}
    \lean{LinearMap.IsSymmetricProjection.reducedProjection}
    \leanok

    Let $P$ and $Q$ be orthogonal projections on a finite-dimensional complex
    inner-product space. Define $\widetilde P$ as the orthogonal projection onto
    $\operatorname{ran}(PQ)$.
    % The source defines Ptilde blockwise; issue #6294 owns the identification
    % of that projector with this canonical construction.
    % Source lines: paper_v2.tex 764--770.
\end{definition}

\begin{theorem}[Reduced-projection identities]
    \label{thm:mpu_two_projection_reduced_projection}
    \lean{LinearMap.IsSymmetricProjection.reducedProjection_isSymmetric,
        LinearMap.IsSymmetricProjection.range_reducedProjection,
        LinearMap.IsSymmetricProjection.reducedProjection_range_triple,
        LinearMap.IsSymmetricProjection.comp_reducedProjection,
        LinearMap.IsSymmetricProjection.reducedProjection_comp,
        LinearMap.IsSymmetricProjection.comp_reducedProjection_left,
        LinearMap.IsSymmetricProjection.range_reducedProjection_comp}
    \leanok
    \uses{def:mpu_two_projection_reduced_projection}
    The reduced projection satisfies
    \begin{align}
        \operatorname{ran}(\widetilde P)   & =\operatorname{ran}(PQ),           \\
        \operatorname{ran}(PQP)            & =\operatorname{ran}(PQ),           \\
        P\widetilde P                      & =\widetilde P,                     \\
        \widetilde P Q                     & =PQ,                               \\
        Q\widetilde P                      & =QP,                               \\
        \operatorname{ran}(\widetilde P Q) & =\operatorname{ran}(\widetilde P).
    \end{align}
    The first two identities describe the canonical construction, while the next
    three establish the properties in source lines 767--769. The final range
    equality is the derived input to the right-factor theorem, not item~(iv) of
    the source.
    % Identifying the canonical construction with the source's blockwise
    % projector remains in issue #6294.
    % Source lines: paper_v2.tex 764--769.
\end{theorem}

\begin{proof}\leanok
    Set $A=QP$. The standard range identity for $A^\dagger A$ gives
    \begin{align}
        \operatorname{ran}(PQP)
         & =\operatorname{ran}(A^\dagger A)
        =\operatorname{ran}(A^\dagger)
        =\operatorname{ran}(PQ).
    \end{align}
    By definition, $\operatorname{ran}(\widetilde P)=\operatorname{ran}(PQ)$,
    which also gives $P\widetilde P=\widetilde P$. For all $x$ and $z$,
    \begin{align}
        \left\langle Qx-PQx,\,PQz\right\rangle & =0.
    \end{align}
    Hence $Qx-PQx\in\operatorname{ran}(PQ)^\perp$, so orthogonal projection
    onto this range yields $\widetilde P Q=PQ$. Taking adjoints gives
    $Q\widetilde P=QP$. Finally,
    \begin{align}
        \operatorname{ran}(\widetilde P Q)
         & =\operatorname{ran}(PQ)
        =\operatorname{ran}(\widetilde P).
    \end{align}
\end{proof}

\begin{theorem}[Transversality of the reduced range]
    \label{thm:mpu_two_projection_reduced_projection_transversality}
    \lean{LinearMap.IsSymmetricProjection.disjoint_ker_range_reducedProjection,
        Matrix.disjoint_ker_range_reducedProjection}
    \leanok
    \uses{def:mpu_two_projection_reduced_projection, thm:mpu_two_projection_reduced_projection}
    For orthogonal projections $P$ and $Q$, the reduced range contains no
    nonzero vector annihilated by $Q$:
    \begin{align}
        \ker Q\cap\operatorname{ran}(\widetilde P) & =\{0\}.
        \label{eq:mpu_reduced_projection_transversality}
    \end{align}
    % Source lines: paper_v2.tex 764--781.
\end{theorem}

\begin{proof}\leanok
    If $x=PQy$ and $Qx=0$, then $Px=x$ and symmetry of the two projections gives
    \begin{align}
        \langle x,x\rangle
         & =\langle PQy,PQy\rangle
        =\langle Qy,PQy\rangle
        =\langle y,QPQy\rangle
        =0.
    \end{align}
    Hence $x=0$.
\end{proof}

\begin{theorem}[Invertible right factor for the reduced projection]
    \label{thm:mpu_two_projection_reduced_projection_right_factor}
    \lean{LinearMap.IsSymmetricProjection.exists_reducedProjection_rightFactor}
    \leanok
    \uses{thm:mpu_two_projection_reduced_projection}
    There is an invertible linear map $Y$ such that
    \begin{align}
        \widetilde P QY=\widetilde P.
    \end{align}
    % Source line: paper_v2.tex 770.
\end{theorem}

\begin{proof}\leanok
    The equality
    \begin{align}
        \operatorname{ran}(\widetilde P Q) & =\operatorname{ran}(\widetilde P)
    \end{align}
    and the equal-ranges right-factor theorem give an invertible $Y$ satisfying
    $\widetilde P QY=\widetilde P$.
\end{proof}

\begin{definition}[Matrix reduced projection]
    \label{def:mpu_matrix_reduced_projection}
    \lean{Matrix.reducedProjection}
    \leanok
    \uses{def:mpu_two_projection_reduced_projection}
    In standard Euclidean coordinates, let $\widetilde P$ be the matrix of the
    reduced projection associated with two square matrices $P$ and $Q$.
\end{definition}

\begin{theorem}[Coordinate identities for the reduced projection]
    \label{thm:mpu_matrix_reduced_projection}
    \lean{Matrix.toEuclideanLin_reducedProjection,
        Matrix.reducedProjection_isOrthogonalProjection,
        Matrix.mul_reducedProjection,
        Matrix.reducedProjection_mul,
        Matrix.reducedProjection_mul_second,
        Matrix.second_mul_reducedProjection,
        Matrix.exists_reducedProjection_rightFactor}
    \leanok
    \uses{def:mpu_matrix_reduced_projection, thm:mpu_two_projection_reduced_projection, thm:mpu_two_projection_reduced_projection_right_factor}
    If $P$ and $Q$ are orthogonal projections, then $\widetilde P$ is an
    orthogonal projection and
    \begin{align}
        P\widetilde P  & =\widetilde P, &
        \widetilde P P & =\widetilde P, &
        \widetilde P Q & =PQ,           &
        Q\widetilde P  & =QP.
        \label{eq:mpu_reduced_projection_coordinate_identities}
    \end{align}
    There is also an invertible matrix $Y$ satisfying
    $\widetilde P QY=\widetilde P$.
    % The tagged Lean declarations prove that \widetilde P is an orthogonal
    % projection for arbitrary P and Q; the hypotheses above are those of the source.
    % Source lines: paper_v2.tex 764--770.
\end{theorem}

\begin{lemma}[Transfer map under virtual sandwiching]
    \label{lem:mpu_transfer_map_virtual_sandwich}
    \lean{MPOTensor.transferMap_virtualSandwich}
    \leanok
    \uses{def:mpo_transfer_map, def:mpu_virtual_sandwich}
    Let $U=(U^{ij})$ be an MPO tensor, and let
    $\widehat U^{ij}=AU^{ij}B$. Then
    \begin{align}
        E_{\widehat U}(X) & =A E_U(BXB^\dagger)A^\dagger.
    \end{align}
\end{lemma}

\begin{proof}\leanok
    Expand the transfer map and collect the factors $A$, $B$, $B^\dagger$, and
    $A^\dagger$ outside the finite sum.
\end{proof}

\begin{lemma}[Positive definite support compression]
    \label{lem:mpu_support_compression_posdef}
    \lean{Matrix.PosSemidef.compression_posDef_of_support_action_ne_zero}
    \leanok

    Let $\rho\succeq0$. If
    \begin{align}
        \operatorname{supp}(\rho)(Vx) & \neq0
        \qquad\text{for every }x\neq0,
    \end{align}
    then
    \begin{align}
        V^\dagger\rho V & \succ0.
    \end{align}
\end{lemma}

\begin{proof}\leanok
    If the quadratic form of $V^\dagger\rho V$ vanished at a nonzero vector $x$,
    positive semidefiniteness would give $\rho Vx=0$. The support projection would
    then annihilate $Vx$, contrary to the hypothesis.
\end{proof}

\begin{theorem}[Reduced representative under a supplied positive fixed pair]
    \label{thm:mpu_reduced_representative}
    \lean{MPOTensor.mpo_virtualSandwich_reducedProjection_eq,
        MPOTensor.transferMap_virtualSandwich_reducedProjection,
        MPOTensor.trace_compressed_fixed_pair_reducedProjection,
        MPOTensor.compressed_fixed_pair_posDef_reducedProjection}
    \leanok
    \uses{def:mpo_family, def:mpo_transfer_map, def:mpu_virtual_sandwich, def:mpu_matrix_reduced_projection, lem:mpu_transfer_map_virtual_sandwich, lem:mpu_support_compression_posdef, thm:mpu_matrix_reduced_projection, thm:mpu_two_projection_reduced_projection_transversality}
    Let $W=(W^{ij})$ be an MPO tensor. Let $L,R\succeq0$ satisfy
    \begin{align}
        E_W(X)                & =\operatorname{tr}(LX)R, &
        \operatorname{tr}(LR) & =1,
    \end{align}
    and let $P$ and $Q$ be their support projections. Assume, letter by letter,
    \begin{align}
        W^{ij}P & =W^{ij}, & QW^{ij} & =W^{ij}.
    \end{align}
    Set $T=\widetilde P$ and $\widetilde W^{ij}=TW^{ij}T$. Then, for every
    $N\geq0$,
    \begin{align}
        \operatorname{MPO}_N(\widetilde W)        & =\operatorname{MPO}_N(W),                    \\
        E_{\widetilde W}(X)
                                                  & =\operatorname{tr}\!\bigl((TLT)X\bigr)\,TRT, \\
        \operatorname{tr}\!\bigl((TLT)(TRT)\bigr) & =1.
    \end{align}
    Moreover, if $V^\dagger V=I$ and $VV^\dagger=T$, then
    \begin{align}
        V^\dagger LV & \succ0, & V^\dagger RV & \succ0.
    \end{align}
    % Scope restriction: the positivity, support, absorption, transfer, and
    % normalization assumptions are supplied explicitly. The source derives them
    % by blocking D^4 sites and using the isometry in the proof of Proposition IV.5.
    % See docs/paper-gaps/mpu_reduced_representative_supplied_fixed_pair.tex.
    % Source lines: paper_v2.tex 747--783.
\end{theorem}

\begin{proof}\leanok
    For every nonempty word $w$, idempotence of $T$ and
    (\ref{eq:mpu_reduced_projection_coordinate_identities}) give
    \begin{align}
        \widetilde W_w & =TW_wT.
    \end{align}
    Cyclicity of the trace, together with $QW_w=W_w=W_wP$, therefore gives
    $\operatorname{tr}(TW_wT)=\operatorname{tr}(W_w)$. Length zero is immediate.
    Lemma~\ref{lem:mpu_transfer_map_virtual_sandwich}, with $A=B=T$ and
    $T^\dagger=T$, gives
    \begin{align}
        E_{\widetilde W}(X)
         & =\operatorname{tr}(LTXT)\,TRT
        =\operatorname{tr}((TLT)X)\,TRT.
    \end{align}
    Since $P$ and $Q$ are the support projections of $L$ and $R$, respectively,
    \begin{align}
        PL & =LP=L, & QR & =RQ=R.
    \end{align}
    Combining these absorptions with
    (\ref{eq:mpu_reduced_projection_coordinate_identities}), $T^2=T$, and
    cyclicity of the trace gives
    \begin{align}
        \operatorname{tr}\!\bigl((TLT)(TRT)\bigr)
         & =\operatorname{tr}(LTRT)
        =\operatorname{tr}(LRT)
        =\operatorname{tr}(RTL)
        =\operatorname{tr}(RL)
        =\operatorname{tr}(LR)
        =1.
    \end{align}
    Let $x\neq0$ and set $u=Vx$. Since $V^\dagger V=I$, one has $u\neq0$;
    since $VV^\dagger=T$, one has $Tu=u$. Equation
    (\ref{eq:mpu_reduced_projection_coordinate_identities}) gives $Pu=u$, and
    (\ref{eq:mpu_reduced_projection_transversality}) gives $Qu\neq0$. Thus
    \begin{align}
        \operatorname{supp}(L)u & =Pu=u\neq0, &
        \operatorname{supp}(R)u & =Qu\neq0.
    \end{align}
    Applying Lemma~\ref{lem:mpu_support_compression_posdef} to $L$ and $R$ yields
    \begin{align}
        x^\dagger(V^\dagger LV)x & =u^\dagger Lu>0, &
        x^\dagger(V^\dagger RV)x & =u^\dagger Ru>0.
    \end{align}
\end{proof}

\begin{theorem}[Invertible factor for a reduced projection]
    \label{thm:reduced_projection_invertible_factor}
    \lean{Matrix.exists_units_supportInvExtension_reducedProjection_rightFactor}
    \leanok
    \uses{def:mpu_matrix_reduced_projection, def:mpu_ambient_support_inverse_extension}
    Let $L,R\succeq0$, let $P$ and $Q$ be their support projections, and let
    $\widetilde P$ be the reduced projection of $P$ and $Q$. There are
    invertible support-inverse extensions of $L$ and $R$, together with an
    invertible matrix $Y$, such that $PQY=\widetilde P$.
\end{theorem}

\begin{proof}\leanok
    Extend the inverses on the supports by the identity on their orthogonal
    complements, and choose an invertible right factor for the reduced
    projection.
\end{proof}

\begin{theorem}[Source-rank comparison of the hatted and reduced tensors]
    \label{thm:mpu_reduced_to_hat_source_ranks}
    \lean{MPOTensor.virtualSandwich_hat_eq_reducedProjection,
        MPOTensor.rightRank_hat_eq_reducedProjection,
        MPOTensor.leftRank_hat_eq_reducedProjection}
    \leanok
    \uses{def:mpo_family, def:mpu_virtual_sandwich, def:mpu_matrix_reduced_projection, def:mpu_right_rank, def:mpu_left_rank, def:mpu_ambient_support_inverse_extension, thm:reduced_projection_invertible_factor}
    Let $W=(W^{ij})$ be an MPO tensor, let $L$, $R$, $X$, $Z$, $Y$ be square
    matrices of the same size, let $P$ be an orthogonal projection, let $Q$
    be a square matrix of the same size, and set $\widetilde P$ equal to
    their reduced projection.
    Assume, letter by letter, $W^{ij}P=W^{ij}=QW^{ij}$. Set
    $\widehat W^{ij}=LW^{ij}R$ and
    $\widetilde W^{ij}=\widetilde P W^{ij}\widetilde P$. If $X,Z,Y$ satisfy
    $XL=P$, $RZ=Q$, and $PQY=\widetilde P$, then
    \begin{align}
        X\widehat W ZY             & =PWQY
        \label{eq:mpu_reduced_hat_chain_pwqy}                    \\
        PWQY                       & =PQWPQY
        \label{eq:mpu_reduced_hat_chain_pqwpqy}                  \\
        PQWPQY                     & =\widetilde P W\widetilde P
        \label{eq:mpu_reduced_hat_chain_reduced}                 \\
        \widetilde P W\widetilde P & =\widetilde W.
        \label{eq:mpu_reduced_hat_chain_tilde}
    \end{align}
    If in addition $X$, $Z$, and $Y$ are invertible, then
    \begin{align}
        r[\widehat W]    & =r[\widetilde W],\label{eq:mpu_reduced_hat_right_rank} \\
        \ell[\widehat W] & =\ell[\widetilde W].
        \label{eq:mpu_reduced_hat_left_rank}
    \end{align} If $L,R\succeq0$ and
    $P=\operatorname{supp}(L)$, $Q=\operatorname{supp}(R)$, such invertible
    matrices may be chosen as $X=L^++(\Id-P)$, $Z=R^++(\Id-Q)$, together
    with an invertible right factor $Y$ of the reduced projection.
    % Scope restriction: the fixed pair, the equations XL=P and RZ=Q, and the
    % letterwise absorptions are supplied hypotheses, not consequences of a bare
    % MPU assumption. See
    % docs/paper-gaps/mpu_reduced_representative_supplied_fixed_pair.tex.
    % The identity chain is formalized without the invertibility hypotheses
    % (the paper takes X, Y, Z invertible at line 792); invertibility of X, Z, Y
    % is used only for the rank conclusion.
    % Source lines: paper_v2.tex 786--804.
\end{theorem}

\begin{proof}\leanok
    \uses{def:mpu_virtual_sandwich, thm:mpu_matrix_reduced_projection, thm:mpu_virtual_sandwich_source_ranks, thm:mpu_ambient_support_inverse_extension}
    Expanding each virtual sandwich and using $XL=P$ and $RZ=Q$ gives
    (\ref{eq:mpu_reduced_hat_chain_pwqy}). The absorptions
    $QW^{ij}=W^{ij}=W^{ij}P$ give, letter by letter,
    \begin{align}
        PW^{ij}QY
         & =P(QW^{ij})QY=PQ(W^{ij}P)QY=PQW^{ij}PQY,
        \notag
    \end{align}
    which is (\ref{eq:mpu_reduced_hat_chain_pqwpqy}). Since
    $\widetilde P Q=PQ$ and $PQY=\widetilde P$,
    \begin{align}
        PQW^{ij}PQY
         & =(\widetilde P Q)W^{ij}\widetilde P
        =\widetilde P W^{ij}\widetilde P.
    \end{align}
    This proves (\ref{eq:mpu_reduced_hat_chain_reduced}) and
    (\ref{eq:mpu_reduced_hat_chain_tilde}). Thus $\widetilde W$ is obtained
    from $\widehat W$ by left multiplication by $X$ and right multiplication
    by $ZY$. Invertibility of these factors and the source ranks of a
    virtual sandwich (Theorem~\ref{thm:mpu_virtual_sandwich_source_ranks})
    give
    \begin{align}
        r[\widehat W]    & =r[X\widehat W(ZY)]=r[\widetilde W],
        \notag                                                        \\
        \ell[\widehat W] & =\ell[X\widehat W(ZY)]=\ell[\widetilde W],
        \notag
    \end{align}
    the first equality in each line being the invariance identity with
    invertible $X$ and $ZY$, and the second rewriting $X\widehat W(ZY)$
    through
    (\ref{eq:mpu_reduced_hat_chain_pwqy})--(\ref{eq:mpu_reduced_hat_chain_reduced})
    and (\ref{eq:mpu_reduced_hat_chain_tilde}) to $\widetilde W$, which gives
    (\ref{eq:mpu_reduced_hat_right_rank}) and
    (\ref{eq:mpu_reduced_hat_left_rank}). For positive semidefinite $L$ and
    $R$, their ambient support-inverse extensions supply $X$ and $Z$; the
    invertible right-factor identity for $\widetilde P Q$ supplies $Y$.
\end{proof}

\begin{theorem}[Lower semicontinuity of finite matrix rank]
    \label{thm:mpu_rank_lower_semicontinuous}
    \lean{Matrix.lowerSemicontinuous_rank,
        Continuous.lowerSemicontinuous_matrix_rank}
    \leanok
    For finite index sets $M,N$, matrix rank defines a lower-semicontinuous map
    \begin{align}
        \operatorname{rank}\colon
        \operatorname{Mat}_{M\times N}(\mathbb C) & \longrightarrow\mathbb N.
    \end{align}
    Consequently, if $A\colon X\to\operatorname{Mat}_{M\times N}(\mathbb C)$
    is continuous, then $x\mapsto\operatorname{rank}A(x)$ is lower
    semicontinuous.
    % Source lines: paper_v2.tex 813--814.
\end{theorem}

\begin{proof}\leanok
    For every $k\in\mathbb N$, the sublevel set is the closed determinantal set
    \begin{align}
        \{A:\operatorname{rank}A\leq k\}.
    \end{align}
    These closed sublevel sets characterize lower semicontinuity. Composition
    with a continuous map preserves lower semicontinuity.
\end{proof}

\begin{theorem}[Fixed positive product gives local constancy]
    \label{thm:mpu_fixed_product_locally_constant}
    \lean{isLocallyConstant_pair_of_lowerSemicontinuous_nat_mul_eq}
    \leanok
    Let $r,\ell\colon X\to\mathbb N$ be lower semicontinuous. Suppose that
    $c>0$ and
    \begin{align}
        r(x)\ell(x) & =c
    \end{align}
    for every $x\in X$. Then both $r$ and $\ell$ are locally constant.
\end{theorem}

\begin{proof}\leanok
    Fix $x\in X$. Since $c>0$, both $r(x)$ and $\ell(x)$ are positive. Lower
    semicontinuity gives a neighborhood $U$ of $x$ on which
    \begin{align}
        r(y) & \geq r(x), & \ell(y) & \geq\ell(x).
    \end{align}
    For $y\in U$, monotonicity of multiplication and the fixed-product identity
    give
    \begin{align}
        r(x)\ell(x)
         & \leq r(y)\ell(x)
        \leq r(y)\ell(y)
        =r(x)\ell(x).
    \end{align}
    Positivity permits cancellation, first of $\ell(x)$ and then of $r(x)$, so
    $r(y)=r(x)$ and $\ell(y)=\ell(x)$.
\end{proof}

\begin{theorem}[Fixed positive product on a preconnected set]
    \label{thm:mpu_fixed_product_preconnected}
    \lean{eq_of_lowerSemicontinuous_nat_mul_eq_of_isPreconnected,
        eq_of_lowerSemicontinuous_nat_mul_eq,
        unitInterval_eq_of_lowerSemicontinuous_nat_mul_eq}
    \leanok
    Under the hypotheses of
    Theorem~\ref{thm:mpu_fixed_product_locally_constant}, if $S\subseteq X$ is
    preconnected, then for all $x,y\in S$,
    \begin{align}
        r(x) & =r(y), & \ell(x) & =\ell(y).
    \end{align}
    In particular, both functions are constant when $X$ is preconnected and
    when $X=[0,1]$.
\end{theorem}

\begin{proof}\leanok
    \uses{thm:mpu_fixed_product_locally_constant}
    A locally constant function has open fibers. A preconnected set cannot meet
    two disjoint nonempty open fibers, so each of $r$ and $\ell$ has one value
    on $S$. The unit interval is connected, hence preconnected.
\end{proof}

\begin{definition}[Virtual sandwich]
    \label{def:mpu_virtual_sandwich}
    \lean{MPOTensor.virtualSandwich,MPOTensor.virtualSandwich_apply}
    \leanok
    For an MPO tensor $\mathcal U=(U^{ij})_{0\leq i,j<d}$ of bond dimension $D$
    and matrices $A,B\in\mathbb C^{D\times D}$, define
    \begin{align}
        \widehat{\mathcal U}^{ij} & =A U^{ij}B.
    \end{align}
    % Source lines: paper_v2.tex 786--804.
\end{definition}

\begin{theorem}[Continuity of the virtual sandwich]
    \label{thm:mpu_virtual_sandwich_continuous}
    \lean{MPOTensor.continuous_virtualSandwich}
    \leanok
    \uses{def:mpu_virtual_sandwich}
    If $x\mapsto A(x)$, $x\mapsto\mathcal U(x)$, and $x\mapsto B(x)$ are
    continuous families, then
    \begin{align}
        x\longmapsto A(x)\mathcal U(x)B(x)
    \end{align}
    is continuous.
    % Source lines: paper_v2.tex 807--812.
\end{theorem}

\begin{proof}\leanok
    \uses{def:mpu_virtual_sandwich}
    Every matrix entry is a finite sum of products of entries of $A(x)$,
    $\mathcal U(x)$, and $B(x)$, hence is continuous.
\end{proof}

\begin{definition}[Doubled-virtual contraction and factor-free sandwich]
    \label{def:mpu_factor_free_sandwich}
    \lean{MPOTensor.doubledVirtualContraction,
        MPOTensor.doubledVirtualContraction_apply,
        MPOTensor.factorFreeSandwich,
        MPOTensor.factorFreeSandwich_apply}
    \leanok
    \uses{def:mpu_double_layer}
    Let $q:\Fin D\times\Fin D\to\Fin{D^2}$ be the standard enumeration of a
    pair. For an MPO tensor $\mathcal W$ and a doubled-virtual matrix
    $E\in\mathbb C^{D^2\times D^2}$, define
    \begin{align}
        C(\mathcal W,E)^{ij}_{ab}
         & =\sum_{c,e=0}^{D-1}E_{q(b,e),q(c,a)}W^{ij}_{ce}.
        \label{eq:mpu_factor_free_contraction}
    \end{align}
    If $E(\mathcal W)$ denotes the normalized physical diagonal of the
    physical-adjoint double layer of $\mathcal W$, define
    \begin{align}
        \widehat{\mathcal W}_{\mathrm{ff}}
         & =C\bigl(\mathcal W,E(\mathcal W)\bigr).
        \label{eq:mpu_factor_free_sandwich}
    \end{align}
    This is the contraction shown in Figure~\texttt{IV\_index4.png} of
    \cite[Proposition~IV.5]{Cirac2017MPU}.
    % Source lines: paper_v2.tex 807--812.
\end{definition}

\begin{theorem}[Rank-one identity and continuity of the factor-free sandwich]
    \label{thm:mpu_factor_free_sandwich_continuous}
    \lean{MPOTensor.doubledVirtualContraction_vecMulVec,
        MPOTensor.doubledVirtualContraction_vecMulVec_vec,
        MPOTensor.continuous_doubledVirtualContraction,
        MPOTensor.continuous_doubleLayerTensor,
        MPOTensor.continuous_normalizedDiagonal,
        MPOTensor.continuous_normalizedDiagonal_doubleLayerTensor,
        MPOTensor.continuous_factorFreeSandwich,
        Continuous.factorFreeSandwich}
    \leanok
    \uses{def:mpu_factor_free_sandwich, def:mpu_virtual_sandwich}
    Let $L,R\in\mathbb C^{D\times D}$ and let $\operatorname{vec}$ stack
    matrix columns, so
    \begin{align}
        \operatorname{vec}(L)_{q(c,a)} & =L_{ac}, \\
        \operatorname{vec}(R)_{q(b,e)} & =R_{eb}.
    \end{align}
    Then the rank-one doubled-virtual matrix satisfies
    \begin{align}
        C\!\left(\mathcal W,
        \operatorname{vec}(R)\operatorname{vec}(L)^T\right)^{ij}
         & =L W^{ij}R.
        \label{eq:mpu_factor_free_rank_one}
    \end{align}
    No adjoint or normalization factor occurs in the contraction result. The
    binary map
    $(\mathcal W,E)\mapsto C(\mathcal W,E)$, the map
    $\mathcal W\mapsto E(\mathcal W)$, and consequently the map
    $\mathcal W\mapsto\widehat{\mathcal W}_{\mathrm{ff}}$ are continuous.
    These assertions neither choose $L$ and $R$ continuously nor infer their
    positivity. They make no claim about reduced representatives or constancy
    of ranks.
    % Source lines: paper_v2.tex 807--812.
\end{theorem}

\begin{proof}\leanok
    \uses{def:mpu_factor_free_sandwich}
    Expanding the left-hand side of
    (\ref{eq:mpu_factor_free_rank_one}) and using column stacking gives
    \begin{align}
        C\!\left(\mathcal W,
        \operatorname{vec}(R)\operatorname{vec}(L)^T\right)^{ij}_{ab}
         & =\sum_{c,e}R_{eb}L_{ac}W^{ij}_{ce} \\
         & =\sum_{c,e}L_{ac}W^{ij}_{ce}R_{eb}
        =(L W^{ij}R)_{ab}.
    \end{align}
    This proves (\ref{eq:mpu_factor_free_rank_one}). Each entry in
    (\ref{eq:mpu_factor_free_contraction}) is a finite sum of products, so the
    binary contraction is continuous. The double-layer entries are finite sums
    of products of entries of $\mathcal W$ and their complex conjugates, and the
    normalized physical diagonal is another finite sum. Hence
    $\mathcal W\mapsto E(\mathcal W)$ is continuous. Substitution into
    (\ref{eq:mpu_factor_free_sandwich}) proves continuity of the factor-free
    sandwich.
\end{proof}

\begin{theorem}[Source cuts of a virtual sandwich]
    \label{thm:mpu_virtual_sandwich_source_cuts}
    \lean{MPOTensor.sourceCutM₁_virtualSandwich,
        MPOTensor.sourceCutM₂_virtualSandwich}
    \leanok
    \uses{def:mpu_virtual_sandwich, def:mpu_source_matrix_one, def:mpu_source_matrix_two}
    The two raw source cuts satisfy the literal formulas
    \begin{align}
        M_1(\widehat{\mathcal U})
         & =(I_d\otimes B^{\mathsf T})M_1(\mathcal U)
        (A^{\mathsf T}\otimes I_d),                      \\
        M_2(\widehat{\mathcal U})
         & =(A\otimes I_d)M_2(\mathcal U)(I_d\otimes B).
    \end{align}
    The transposes in the first identity result from the corrected orientation
    of $M_1$; no complex conjugation occurs.
    % Source lines: paper_v2.tex 786--804.
\end{theorem}

\begin{proof}\leanok
    \uses{def:mpu_virtual_sandwich, def:mpu_source_matrix_one, def:mpu_source_matrix_two}
    Expand an entry of each right-hand side. The identity factors collapse the
    two physical sums. For the first cut, the entries of $B^{\mathsf T}$ and
    $A^{\mathsf T}$ put $B_{\delta\beta}$ on the left and
    $A_{\alpha\gamma}$ on the right; commutativity of scalar multiplication
    gives
    \begin{align}
        \sum_{\gamma,\delta}
        A_{\alpha\gamma}U^{ij}_{\gamma\delta}B_{\delta\beta}
         & =(A U^{ij}B)_{\alpha\beta}.
    \end{align}
    The second cut gives the same entry directly.
\end{proof}

\begin{theorem}[Source ranks of a virtual sandwich]
    \label{thm:mpu_virtual_sandwich_source_ranks}
    \lean{MPOTensor.rightRank_virtualSandwich,
        MPOTensor.leftRank_virtualSandwich}
    \leanok
    \uses{def:mpu_virtual_sandwich, def:mpu_right_rank, def:mpu_left_rank}
    If $A$ and $B$ are invertible, then
    \begin{align}
        r[\widehat{\mathcal U}]    & =r[\mathcal U],    &
        \ell[\widehat{\mathcal U}] & =\ell[\mathcal U].
    \end{align}
    This assertion concerns the raw source ranks only.
    % Source lines: paper_v2.tex 786--804.
\end{theorem}

\begin{proof}\leanok
    \uses{thm:mpu_virtual_sandwich_source_cuts}
    The determinant identities for Kronecker products show that
    $A\otimes I_d$ and $I_d\otimes B$ are invertible. Left and right
    multiplication by these matrices preserves rank. Apply this observation to
    each source-cut formula in
    Theorem~\ref{thm:mpu_virtual_sandwich_source_cuts}.
\end{proof}

\begin{definition}[Ambient support-inverse extension]
    \label{def:mpu_ambient_support_inverse_extension}
    \lean{Matrix.PosSemidef.supportInvExtension,
        Matrix.PosSemidef.supportInvExtensionInverse}
    \leanok

    Let $A\succeq0$ have support projection $P$, and let $A^+$ denote its
    inverse on the support. Define
    \begin{align}
        X_A & =A^++(\Id-P).
    \end{align}
    Define also
    \begin{align}
        Z_A & =A+(\Id-P).
    \end{align}
    This construction gives explicit witnesses for the invertible matrices
    required in Proposition~IV.5 of \cite{Cirac2017MPU}; the paper requires
    those matrices but does not state these formulas.
    % Source requirement: paper_v2.tex, Proposition IV.5, lines 786--797.
    % The displayed construction is a formal witness, not a formula from the paper.
\end{definition}

\begin{theorem}[Invertibility and support identities]
    \label{thm:mpu_ambient_support_inverse_extension}
    \lean{Matrix.PosSemidef.supportInvExtension_mul_self,
        Matrix.PosSemidef.self_mul_supportInvExtension,
        Matrix.PosSemidef.supportInvExtension_mul_inverse,
        Matrix.PosSemidef.supportInvExtension_inverse_mul,
        Matrix.PosSemidef.supportInvExtensionUnit,
        Matrix.PosSemidef.isUnit_supportInvExtension,
        Matrix.PosSemidef.isUnit_det_supportInvExtension}
    \leanok
    \uses{def:mpu_ambient_support_inverse_extension}
    The matrices $X_A$ and $Z_A$ are mutually inverse, and
    \begin{align}
        X_AA=AX_A & =P.
    \end{align}
    Consequently, $X_A$ is a unit in the matrix algebra, with inverse $Z_A$,
    and $\det(X_A)$ is a unit in $\mathbb C$.

    In the notation of Proposition~IV.5 of \cite{Cirac2017MPU}, let $L$ and
    $R$ be the positive semidefinite outer factors with support projections
    $P$ and $Q$, respectively. The required witnesses are
    \begin{align}
        X & =X_L=L^++(\Id-P), \\
        Z & =X_R=R^++(\Id-Q).
    \end{align}
    Thus
    \begin{align}
        X_L L & =P, \\
        R X_R & =Q,
    \end{align}
    or equivalently $XL=P$ and $RZ=Q$. Here the paper's matrix $Z$ is $X_R$,
    not the inverse-extension matrix $Z_R$ from
    Definition~\ref{def:mpu_ambient_support_inverse_extension}.
    % This proves the existence step used in paper_v2.tex, lines 786--797.
\end{theorem}

\begin{proof}\leanok
    The support identities give
    \begin{align}
        A^+A=AA^+ & =P.
    \end{align}
    Their complementary forms are
    \begin{align}
        A^+(\Id-P)=(\Id-P)A^+ & =0,
    \end{align}
    \begin{align}
        A(\Id-P)=(\Id-P)A & =0,
    \end{align}
    and idempotence of $P$ gives
    \begin{align}
        (\Id-P)^2 & =\Id-P.
    \end{align}
    Therefore
    \begin{align}
        X_AZ_A & =P+(\Id-P)=\Id.
    \end{align}
    Likewise,
    \begin{align}
        Z_AX_A & =P+(\Id-P)=\Id,
    \end{align}
    while multiplication by $A$ gives $X_AA=AX_A=P$.
\end{proof}

\begin{proposition}[Constancy of the source ranks]
    \label{prop:continuity-index}
    \notready
    \discussion{6022}
    \uses{def:mpu_canonical_form, ThmFund1, thm:mpu_common_pow_four_simple_blocking, thm:mpu_two_projection_angle_defect_sum, thm:mpu_two_projection_reduced_projection, thm:mpu_two_projection_reduced_projection_right_factor, thm:mpu_reduced_representative, thm:mpu_rank_lower_semicontinuous, thm:mpu_fixed_product_preconnected, thm:mpu_factor_free_sandwich_continuous, thm:mpu_virtual_sandwich_source_ranks, thm:mpu_ambient_support_inverse_extension}
    Let $[0,1]\ni x\mapsto\mathcal W(x)$ be a continuous path of tensors
    generating MPUs, without requiring the path to be in canonical form. For
    each $x$, take the canonical form and the two source-cut ranks $r(x)$ and
    $\ell(x)$ of its associated simple tensor. Then $r(x)$ and $\ell(x)$ are
    constant on $[0,1]$.
    % Source lines: paper_v2.tex 733--818.
\end{proposition}

\begin{proposition}[Admissible constancy of the source ranks]
    \label{prop:mpu_admissible_continuity_index}
    \notready
    \discussion{6022}
    % **Scope restriction (pathwise reduced source data):** every associated
    % simple tensor must carry the datum required by
    % Theorem~\ref{thm:mpu_admissible_simple_tensor_equivalence}.
    % Documented in docs/paper-gaps/mpu_canonical_form_full_support.tex.
    \uses{thm:mpu_admissible_simple_tensor_equivalence, thm:mpu_bounded_simple_blocking, def:mpu_reduced_full_support_source_datum, def:mpu_admissible_equivalence_datum}
    Let $[0,1]\ni x\mapsto\mathcal W(x)$ be a continuous path of tensors
    generating MPUs, without requiring the path to be in canonical form.
    Suppose that the path carries an admissible equivalence datum with blocking
    length $m$. Then the ranks $r(x)$ and $\ell(x)$ of the actual blocked path
    tensors $\mathcal W(x)_m$ are constant on $[0,1]$.
    % Source lines: paper_v2.tex 733--818.
\end{proposition}

\begin{theorem}[Index Theorem]
    \label{IndexTh}
    \notready
    \discussion{6023}
    \uses{def:index, index-well-defined, lem:mpu_index_tensor, lem:mpu_index_composition, prop:continuity-index, def:equivalent-tensors}
    For MPU tensors under the source's standing canonical-form convention, take
    the index defined from any simple blocking. This index has the following
    properties.
    \begin{enumerate}
        \item It does not change under blocking.
        \item It is additive under tensor products and composition.
        \item It is constant under continuous deformations through MPU tensors.
        \item Two MPU tensors have the same index if and only if they are
              equivalent.
    \end{enumerate}
    This is the unrestricted statement of
    \cite[Theorem~IV.6]{Cirac2017MPU}.
    % Source lines: paper_v2.tex 824--865.
\end{theorem}

\begin{theorem}[Admissible Index Theorem]
    \label{thm:mpu_admissible_index}
    \notready
    \discussion{6023}
    % **Scope restriction (admissible tensors and operations):** every simple
    % block and standard form used below carries reduced full-support source
    % data. Documented in docs/paper-gaps/mpu_canonical_form_full_support.tex.
    \uses{prop:mpu_admissible_index_well_defined, lem:mpu_admissible_index_tensor, lem:mpu_admissible_index_composition, def:equivalent-tensors, def:mpu_admissible_equivalence_datum, prop:mpu_admissible_continuity_index, def:mpu_admissible_standard_form, thm:mpu_admissible_fundamental, def:mpu_reduced_full_support_source_datum}
    Assume the following operation-specific admissibility hypotheses. Every pair
    of simple blockings compared for blocking independence carries reduced
    full-support source data. For a tensor product $\mathcal U\otimes\mathcal V$,
    there is one positive length $k$ such that $\mathcal U_k$, $\mathcal V_k$,
    and $(\mathcal U\otimes\mathcal V)_k$ are simple and carry their own data.
    For a composition tensor $\mathcal W$ of $\mathcal U$ and $\mathcal U'$, there
    is one positive length $k$ such that
    \begin{align}
        \mathcal U_k,\quad \mathcal U'_k,\quad \mathcal W_k,\quad
        \mathcal U_{2k},\quad \mathcal U'_{2k},\quad \mathcal W_{4k}
    \end{align}
    are simple and carry their own data. Continuous-deformation and equivalence
    claims use the path data of
    Definition~\ref{def:mpu_admissible_equivalence_datum}. Under these
    hypotheses, the MPU index has the following properties.
    \begin{enumerate}
        \item It does not change under blocking.
        \item It is additive under tensor products and composition.
        \item It is constant under continuous deformations through MPU tensors.
        \item Admissibly equivalent MPU tensors have the same index. Conversely,
              tensors with the same index are admissibly equivalent whenever the
              standard-form interpolation is supplied with an admissible path datum.
    \end{enumerate}
    These are the conclusions of \cite[Theorem~IV.6]{Cirac2017MPU}, restricted
    to the explicit admissibility hypotheses above.
    % Source lines: paper_v2.tex 824--865.
\end{theorem}

\begin{corollary}[Continuous choice of standard form]
    \label{coruvSF}
    \notready
    \discussion{6024}
    \uses{prop:continuity-index, IndexTh, SF}
    For every continuous path $p\mapsto\mathcal U(p)$ of MPU tensors, not
    necessarily in canonical form, there are an integer $k_0\leq D^4$ and
    continuous paths of unitaries $u(p)$ and $v(p)$ such that
    $U^{(2k_0N)}(p)$ has standard form given by $u(p)$ and $v(p)$.
    % Source lines: paper_v2.tex 867--910.
\end{corollary}

\begin{corollary}[Admissible continuous choice of standard form]
    \label{cor:mpu_admissible_continuous_standard_form}
    \notready
    \discussion{6024}
    % **Scope restriction (continuous admissible standard forms):** the chosen
    % simple blocks carry reduced full-support source data varying continuously
    % along the path. Documented in
    % docs/paper-gaps/mpu_canonical_form_full_support.tex.
    \uses{prop:mpu_admissible_continuity_index, thm:mpu_admissible_index, def:mpu_admissible_standard_form, thm:mpu_admissible_simple_tensor_equivalence, def:mpu_reduced_full_support_source_datum, def:mpu_admissible_equivalence_datum}
    Let $k_0$ be the common blocking length in an admissible equivalence
    datum for a continuous path $\mathcal U(p)$ of MPU tensors, and suppose that
    this supplied length satisfies $k_0\leq D^4$. Then there are continuous
    paths of unitaries $u(p)$ and $v(p)$ such that the MPU
    $U^{(2k_0N)}(p)$ has standard form given by $u(p)$ and $v(p)$ for every
    $p\in[0,1]$. This is the reduced-source version of
    \cite[Corollary~IV.7]{Cirac2017MPU}.
    % Source lines: paper_v2.tex 867--910.
\end{corollary}

\subsection{Local characterization of symmetries}

\begin{proposition}[Conjugation symmetry]
    \label{prop1conj}
    \notready
    \discussion{5713}
    \uses{FundamentalMPU, SF}
    An MPU in standard form is invariant under complex conjugation if and only
    if there are unitaries $x$ and $y$ such that either
    $x=x^T$ and $y=y^T$, or $x=-x^T$ and $y=-y^T$, and
    \begin{align}
        \overline u & =(x\otimes y)u,                 &
        \overline v & =v(y^\dagger\otimes x^\dagger).
    \end{align}
    % Source lines: paper_v2.tex 1008--1039.
\end{proposition}

\begin{proposition}[Admissible conjugation symmetry]
    \label{prop:mpu_admissible_conjugation_symmetry}
    \notready
    \discussion{5713}
    % **Scope restriction (conjugate comparison tensor):** both standard-form
    % tensors compared by Theorem~\ref{thm:mpu_admissible_fundamental} carry
    % reduced full-support
    % source data. Documented in
    % docs/paper-gaps/mpu_canonical_form_full_support.tex.
    \uses{thm:mpu_admissible_fundamental, def:mpu_admissible_standard_form, def:mpu_reduced_full_support_source_datum}
    Let a tensor and its conjugate comparison tensor both be equipped with
    reduced full-support source data and written in admissible standard form.
    The generated
    MPU is invariant under complex conjugation if and only if there are
    unitaries $x$ and $y$ such that either
    $x=x^T$ and $y=y^T$, or $x=-x^T$ and $y=-y^T$, and
    \begin{align}
        \overline u & =(x\otimes y)u,                 &
        \overline v & =v(y^\dagger\otimes x^\dagger).
    \end{align}
    This is the reduced-source version of the conjugation characterization in
    \cite[Proposition~V.1]{Cirac2017MPU}.
    % Source lines: paper_v2.tex 1008--1039.
\end{proposition}

\begin{lemma}[Symmetric branch of the unitary congruence normal form]
    \label{lem:mpu_symmetric_unitary_congruence}
    \lean{Matrix.exists_symmetric_unitary_squareRoot}
    \leanok
    Let $x$ be a symmetric unitary matrix. There is a symmetric unitary matrix
    $S$ such that
    \begin{align}
        x & = S^T S.
    \end{align}
    This is the symmetric branch of
    \cite[Lemma~\texttt{lemma:conjclass-normalform-continuous}]{Cirac2017MPU}.
    % Source lines: paper_v2.tex 1054--1064.
\end{lemma}

\begin{proof}\leanok
    Write the finite spectral decomposition of $x$ as
    \begin{align}
        x & = \sum_{\lambda\in\sigma(x)} \lambda P_\lambda
    \end{align}
    and set
    \begin{align}
        S & = \sum_{\lambda\in\sigma(x)} \sqrt{\lambda}\,P_\lambda.
    \end{align}
    Each spectral value has modulus one, so its square root has modulus one,
    and hence $S$ is unitary and $S^2=x$. On the finite spectrum, interpolate
    the square-root function by a polynomial $q$. Then $S=q(x)$, and
    $x^T=x$ implies $S^T=q(x^T)=q(x)=S$. Therefore $x=S^2=S^T S$.
\end{proof}

\begin{lemma}[Skew congruence from paired spectral projectors]
    \label{lem:mpu_skew_unitary_paired_projectors}
    \lean{Matrix.exists_skew_unitary_congruence_of_paired_projectors}
    \leanok
    Let $K$ and $I$ be finite sets, let $E_k\in\mathbb R$ for $k\in K$, and
    let $P_k$ be Hermitian projectors on $\mathbb C^I$. Suppose that the family
    $\{P_k,P_k^T\}_{k\in K}$ consists of mutually orthogonal projectors and
    sums to the identity. Define
    \begin{align}
        x & = \sum_{k\in K} e^{-iE_k}(P_k-P_k^T).
    \end{align}
    Then the matrices
    \begin{align}
        S & = e^{-i\pi/4}\sum_{k\in K} e^{-iE_k/2}(P_k+P_k^T),
    \end{align}
    and
    \begin{align}
        \widetilde\Lambda & = i\sum_{k\in K}(P_k-P_k^T)
    \end{align}
    are unitary, $S^T=S$, $\overline{\widetilde\Lambda}=\widetilde\Lambda$,
    $\widetilde\Lambda^T=-\widetilde\Lambda$, and
    \begin{align}
        x & = S^T\widetilde\Lambda S.
    \end{align}
    This is the paired-projector construction in the skew branch of
    \cite[Lemma~\texttt{lemma:conjclass-normalform-continuous}]{Cirac2017MPU}.
    % Source lines: paper_v2.tex 1066--1082.
\end{lemma}

\begin{proof}\leanok
    Write the two projectors in each pair as separate members of one complete
    orthogonal family. A weighted sum of this family is unitary whenever every
    weight has modulus one. This applies to $S$, whose two weights at $k$
    both equal $e^{-i\pi/4}e^{-iE_k/2}$, and to
    $\widetilde\Lambda$, whose weights are $i$ and $-i$.
    Transposition exchanges the two members of each pair, proving the symmetry
    of $S$ and the skew symmetry of $\widetilde\Lambda$. Hermiticity gives
    $\overline{P_k}=P_k^T$, hence
    $\overline{\widetilde\Lambda}=\widetilde\Lambda$. Finally, orthogonality
    reduces the product to its weights, and
    \begin{align}
        \left(e^{-i\pi/4}e^{-iE_k/2}\right)^2 i & = e^{-iE_k}.
    \end{align}
    The transposed member has the negative weight, which gives
    $x=S^T\widetilde\Lambda S$.
\end{proof}

\begin{lemma}[Normal form of symmetric and skew-symmetric unitaries]
    \label{lemma:conjclass-normalform-continuous}
    \notready
    \discussion{5713}
    \uses{lem:mpu_symmetric_unitary_congruence, lem:mpu_skew_unitary_paired_projectors}
    Let $x$ be unitary and satisfy $x=x^T$ or $x=-x^T$. There are a symmetric
    unitary $S$ and a matrix $\widetilde\Lambda$ such that
    \begin{align}
        x & =S^T\widetilde\Lambda S, & S & =S^T.
    \end{align}
    In the symmetric case, $\widetilde\Lambda$ is real and symmetric. In the
    skew-symmetric case, it is real and skew-symmetric. In both cases,
    $\det(\widetilde\Lambda)=1$.
    % Source lines: paper_v2.tex 1041--1084.
\end{lemma}

\begin{lemma}[Odd-dimensional skew-symmetric determinant]
    \label{lem:mpu_skew_symmetric_det_zero}
    \lean{Matrix.det_eq_zero_of_transpose_eq_neg_of_odd_card}
    \leanok
    Let $I$ be a finite set of odd cardinality and let
    $A\in\operatorname{Mat}_I(\mathbb C)$ satisfy $A^T=-A$. Then
    \begin{align}
        \det A & =0.
    \end{align}
\end{lemma}

\begin{proof}\leanok
    Transpose invariance, skew symmetry, and the determinant of a negated
    matrix give
    \begin{align}
        \det A
         & =\det(A^T)
        =\det(-A)
        =(-1)^{|I|}\det A
        =-\det A.
    \end{align}
    Since the complex numbers have characteristic zero, $\det A=0$.
\end{proof}

\begin{lemma}[Parity of a nonsingular skew-symmetric matrix]
    \label{lem:mpu_skew_symmetric_even_card}
    \lean{Matrix.even_card_of_transpose_eq_neg_of_det_ne_zero}
    \leanok
    Let $I$ be finite and let $A\in\operatorname{Mat}_I(\mathbb C)$ satisfy
    $A^T=-A$ and $\det A\ne0$. Then $|I|$ is even.
\end{lemma}

\begin{proof}\leanok
    \uses{lem:mpu_skew_symmetric_det_zero}
    If $|I|$ were odd, Lemma~\ref{lem:mpu_skew_symmetric_det_zero} would give
    $\det A=0$. Hence $|I|$ is not odd and is therefore even.
\end{proof}

\begin{theorem}[Even dimension of a skew-symmetric unitary]
    \label{thm:mpu_skew_unitary_even_dimension}
    \lean{Matrix.UnitaryGroup.even_dimension_of_transpose_eq_neg}
    \leanok
    Let $U\in\operatorname{Mat}_d(\mathbb C)$ be unitary and satisfy
    $U^T=-U$. Then $d$ is even. This is the parity obstruction in
    \cite[Proposition~V.2, Case~II]{Cirac2017MPU}.
    % Source lines: paper_v2.tex 1098--1126, especially line 1125.
\end{theorem}

\begin{proof}\leanok
    \uses{lem:mpu_skew_symmetric_even_card}
    The determinant of a unitary matrix is a unit, so it is nonzero. Apply
    Lemma~\ref{lem:mpu_skew_symmetric_even_card} to $U$.
\end{proof}

\begin{proposition}[Orthogonal and symplectic standard forms]
    \label{corconj}
    \notready
    \discussion{5713}
    \uses{prop1conj, lemma:conjclass-normalform-continuous, FundamentalMPU, thm:mpu_skew_unitary_even_dimension}
    Let an MPU in standard form be invariant under complex conjugation. Its
    standard-form unitaries can be chosen in exactly one of the following two
    cases. In Case I,
    \begin{align}
        \overline{u'} & =u', & \overline{v'} & =v',
    \end{align}
    uniquely up to local orthogonal transformations. In Case II,
    \begin{align}
        \overline{u'} & =(\Sigma_\ell\otimes\Sigma_r)u', &
        \overline{v'} & =v'(\Sigma_r\otimes\Sigma_\ell),
    \end{align}
    uniquely up to local symplectic transformations, where
    $\Sigma_x=\Id_{x/2}\otimes\left(\begin{smallmatrix}0&1\\-1&0\end{smallmatrix}\right)$.
    Case II can occur only when $r$ and $\ell$ are even.
    % Source lines: paper_v2.tex 1094--1163.
\end{proposition}

\begin{proposition}[Admissible orthogonal and symplectic standard forms]
    \label{prop:mpu_admissible_conjugation_standard_forms}
    \notready
    \discussion{5713}
    % **Scope restriction (standard-form comparisons):** every alternative
    % standard-form tensor compared by the fundamental theorem carries its own
    % reduced full-support source datum. Documented in
    % docs/paper-gaps/mpu_canonical_form_full_support.tex.
    \uses{prop:mpu_admissible_conjugation_symmetry, lemma:conjclass-normalform-continuous, thm:mpu_admissible_fundamental, def:mpu_reduced_full_support_source_datum}
    Let an MPU in admissible standard form be invariant under complex
    conjugation, and
    suppose every standard-form tensor compared below carries reduced
    full-support source data. Its
    standard-form unitaries can be chosen in exactly one of the following two
    cases. In Case I,
    \begin{align}
        \overline{u'} & =u', & \overline{v'} & =v',
    \end{align}
    and the choice is unique up to local orthogonal transformations. In Case II,
    \begin{align}
        \overline{u'} & =(\Sigma_\ell\otimes\Sigma_r)u', &
        \overline{v'} & =v'(\Sigma_r\otimes\Sigma_\ell),
    \end{align}
    where
    \begin{align}
        \Sigma_x & =\Id_{x/2}\otimes Y,                    &
        Y        & =\begin{pmatrix}0&1\\-1&0\end{pmatrix}.
    \end{align}
    The choice is unique up to local transformations $S_x$ satisfying
    $S_x^T\Sigma_xS_x=\Sigma_x$. Case II can occur only when $r$ and $\ell$
    are even.
    % Source lines: paper_v2.tex 1094--1163.
\end{proposition}

\begin{corollary}[Reduction of the symplectic case to orthogonal unitaries]
    \label{cor:sympl-to-orthogonal}
    \notready
    \discussion{5713}
    \uses{corconj}
    In Case II of Proposition~\ref{corconj}, one can write
    \begin{align}
        u' & =(\Id_{r\ell/4}\otimes R)\widetilde u, &
        v' & =\widetilde v(\Id_{r\ell/4}\otimes R),
    \end{align}
    where $\widetilde u$ and $\widetilde v$ are orthogonal and
    \begin{align}
        R & =\frac{1-i}{2}[\Id_4+i(Y\otimes Y)].
    \end{align}
    % Source lines: paper_v2.tex 1165--1185.
\end{corollary}

\begin{corollary}[Admissible reduction of the symplectic case to orthogonal unitaries]
    \label{cor:mpu_admissible_symplectic_to_orthogonal}
    \notready
    \discussion{5713}
    % **Scope restriction (admissible standard form):** Case II refers to the
    % source-data-restricted standard form of Proposition~\ref{prop:mpu_admissible_conjugation_standard_forms}.
    % Documented in docs/paper-gaps/mpu_canonical_form_full_support.tex.
    \uses{prop:mpu_admissible_conjugation_standard_forms, def:mpu_reduced_full_support_source_datum}
    In the admissible Case II of Proposition~\ref{prop:mpu_admissible_conjugation_standard_forms}, one can write
    \begin{align}
        u' & =(\Id_{r\ell/4}\otimes R)\widetilde u, &
        v' & =\widetilde v(\Id_{r\ell/4}\otimes R),
    \end{align}
    where $\widetilde u$ and $\widetilde v$ are orthogonal and
    \begin{align}
        R & =\frac{1-i}{2}[\Id_4+i(Y\otimes Y)].
    \end{align}
    % Source lines: paper_v2.tex 1165--1185.
\end{corollary}

\begin{proposition}[Local characterization of time reversal]
    \label{PropChiral}
    \notready
    \discussion{5714}
    \uses{FundamentalMPU, SF}
    Let an MPU be in standard form, with one-site physical space
    $\mathbb C^{d_0}\otimes\mathbb C^{d_0}$. It satisfies
    $U^{(N)}=U^{(N)\dagger}$ for every even $N$ if and only if there is a
    unitary $x$ such that
    \begin{align}
        (\Id_1\otimes v_{23}\otimes\Id_4)(u_{12}\otimes u_{34})
         & =\pm(x_{12}\otimes x_{34}^\dagger)
        (\Id_1\otimes v_{23}^\dagger\otimes\Id_4).
    \end{align}
    % Source lines: paper_v2.tex 1201--1252.
\end{proposition}

\begin{proposition}[Admissible local characterization of time reversal]
    \label{prop:mpu_admissible_local_time_reversal}
    \notready
    \discussion{5714}
    % **Scope restriction (two-site blocked comparison):** the tensor obtained by
    % blocking two sites, and the adjoint comparison tensor to which the
    % fundamental theorem is applied, carry reduced full-support source data.
    % No preservation under blocking is asserted. Documented in
    % docs/paper-gaps/mpu_canonical_form_full_support.tex.
    \uses{thm:mpu_admissible_fundamental, prop:mpu_admissible_index_well_defined, def:mpu_admissible_index, def:mpu_admissible_standard_form, def:mpu_reduced_full_support_source_datum}
    Let an MPU be in admissible standard form. Suppose its one-site tensor
    $\mathcal U$, its actual two-site block $\mathcal U_2$, and the
    physical-adjoint comparison tensor $\mathcal U_2^\sharp$ are simple and each
    carries its own reduced full-support source datum. Write the one-site
    physical space as $\mathbb C^{d_0}\otimes\mathbb C^{d_0}$. Time-reversal
    symmetry then forces $\operatorname{ind}(\mathcal U)=0$. Under these
    hypotheses, $U^{(N)}=U^{(N)\dagger}$ for every even $N$ if and only if there is a
    unitary $x$ such that
    \begin{align}
        (\Id_1\otimes v_{23}\otimes\Id_4)(u_{12}\otimes u_{34})
         & =\pm(x_{12}\otimes x_{34}^\dagger)
        (\Id_1\otimes v_{23}^\dagger\otimes\Id_4).
    \end{align}
    % Source lines: paper_v2.tex 1201--1252.
\end{proposition}

\begin{proof}
    \uses{thm:mpu_admissible_fundamental, prop:mpu_admissible_index_well_defined, def:mpu_admissible_index}
    Under time reversal, $\mathcal U_2$ and $\mathcal U_2^\sharp$ generate the
    same MPU. The admissible fundamental theorem identifies their source spaces,
    so their corresponding source-cut ranks agree. Physical adjunction exchanges
    the two cuts. Hence
    \begin{align}
        r(\mathcal U_2)
         & =r(\mathcal U_2^\sharp)
        =\ell(\mathcal U_2).
    \end{align}
    Definition~\ref{def:mpu_admissible_index} therefore gives
    $\operatorname{ind}(\mathcal U_2)=0$, and blocking independence gives
    $\operatorname{ind}(\mathcal U)=0$. Apply the admissible fundamental theorem
    to the supplied standard forms of $\mathcal U_2$ and
    $\mathcal U_2^\sharp$ and follow the source proof. The converse follows by
    contracting the displayed local relation around every even periodic chain.
    % Source lines: paper_v2.tex 1201--1252.
\end{proof}

\begin{proposition}[Local characterization of transposition]
    \label{prop:local-char-transposition}
    \notready
    \discussion{5714}
    \uses{FundamentalMPU, SF}
    Let an MPU be in standard form, with one-site physical space
    $\mathbb C^{d_0}\otimes\mathbb C^{d_0}$. It satisfies
    $U^{(N)}=U^{(N)T}$ for every even $N$ if and only if there are a unitary
    $x$ and a phase $e^{i\phi}$ such that
    \begin{align}
        (\Id_1\otimes v_{23}\otimes\Id_4)(u_{12}\otimes u_{34})
         & =e^{i\phi}(x_{12}\otimes x_{34}^T)
        (\Id_1\otimes v_{23}^T\otimes\Id_4).
    \end{align}
    % Source lines: paper_v2.tex 1257--1288.
\end{proposition}

\begin{proposition}[Admissible local characterization of transposition]
    \label{prop:mpu_admissible_local_transposition}
    \notready
    \discussion{5714}
    % **Scope restriction (two-site blocked comparison):** the tensor obtained by
    % blocking two sites, and the transposed comparison tensor to which the
    % fundamental theorem is applied, carry reduced full-support source data.
    % No preservation under blocking is asserted. Documented in
    % docs/paper-gaps/mpu_canonical_form_full_support.tex.
    \uses{thm:mpu_admissible_fundamental, prop:mpu_admissible_index_well_defined, def:mpu_admissible_index, def:mpu_admissible_standard_form, def:mpu_reduced_full_support_source_datum}
    Let an MPU be in admissible standard form. Suppose its one-site tensor
    $\mathcal U$, its actual two-site block $\mathcal U_2$, and the transposed
    comparison tensor $\mathcal U_2^{\mathsf T}$ are simple and each carries its
    own reduced full-support source datum. Write the one-site physical space as
    $\mathbb C^{d_0}\otimes\mathbb C^{d_0}$. Transposition symmetry then forces
    $\operatorname{ind}(\mathcal U)=0$. Under these hypotheses,
    $U^{(N)}=U^{(N)T}$ for every even $N$ if and only if there are a unitary
    $x$ and a phase $e^{i\phi}$ such that
    \begin{align}
        (\Id_1\otimes v_{23}\otimes\Id_4)(u_{12}\otimes u_{34})
         & =e^{i\phi}(x_{12}\otimes x_{34}^T)
        (\Id_1\otimes v_{23}^T\otimes\Id_4).
    \end{align}
    % Source lines: paper_v2.tex 1257--1288.
\end{proposition}

\begin{proof}
    \uses{thm:mpu_admissible_fundamental, prop:mpu_admissible_index_well_defined, def:mpu_admissible_index}
    Under transposition, $\mathcal U_2$ and $\mathcal U_2^{\mathsf T}$ generate
    the same MPU. The admissible fundamental theorem identifies their source
    spaces, so their corresponding source-cut ranks agree. Transposition
    exchanges the two cuts. Hence
    \begin{align}
        r(\mathcal U_2)
         & =r(\mathcal U_2^{\mathsf T})
        =\ell(\mathcal U_2).
    \end{align}
    Thus $\operatorname{ind}(\mathcal U_2)=0$ by
    Definition~\ref{def:mpu_admissible_index}, and blocking independence gives
    $\operatorname{ind}(\mathcal U)=0$. Apply the admissible fundamental theorem
    to the supplied standard forms of $\mathcal U_2$ and
    $\mathcal U_2^{\mathsf T}$ and repeat the time-reversal argument, retaining
    the unrestricted phase $e^{i\phi}$. The converse follows by contracting the
    local relation around every even periodic chain.
    % Source lines: paper_v2.tex 1257--1288.
\end{proof}

\subsection{Equivalence under symmetries}

\begin{definition}[Finite-chain operator symmetry]
    \label{def:mpu_finite_chain_operator_symmetry}
    \lean{MPOTensor.FiniteChainOperatorSymmetry,
        MPOTensor.IsInvariantUnderSymmetry}
    \leanok
    At physical dimension $d$, a finite-chain operator symmetry $\mathcal S$
    specifies a set $A$ of applicable chain lengths and, for every length $N$,
    a map
    \begin{align}
        \mathcal S_{d,N}\colon
        \operatorname{Mat}_{d^N}(\mathbb C)
         & \longrightarrow \operatorname{Mat}_{d^N}(\mathbb C).
    \end{align}
    An MPO tensor $\mathcal W$ is invariant under $\mathcal S$ if
    \begin{align}
        \mathcal S_{d,N}[W^{(N)}] & =W^{(N)}
    \end{align}
    for every $N\in A$.
    % Source lines: paper_v2.tex 1345--1354.
\end{definition}

\begin{definition}[Coherent symmetry family under identity ancillas]
    \label{def:mpu_coherent_symmetry_family}
    \lean{MPOTensor.tensorPhysicalIdOperator,
        MPOTensor.FiniteChainOperatorSymmetryFamily,
        MPOTensor.FiniteChainOperatorSymmetryFamily.at}
    \leanok
    \uses{def:mpu_finite_chain_operator_symmetry, def:equivalent-tensors}
    A coherent symmetry family supplies an action $\mathcal S_{d,N}$ for every
    physical dimension $d$, with one common set $A$ of applicable lengths. For
    an ancilla dimension $x$, let
    \begin{align}
        \iota_x(X)
         & = R_{d,x,N}\bigl(X\otimes I_{x^N}\bigr)R_{d,x,N}^{-1},
    \end{align}
    where $R_{d,x,N}$ is the canonical sitewise identification of the
    $d$-dimensional and $x$-dimensional configurations. Coherence means
    \begin{align}
        \iota_x\bigl(\mathcal S_{d,N}(X)\bigr)
         & = \mathcal S_{dx,N}\bigl(\iota_x(X)\bigr).
    \end{align}
    Thus the action on the full enlarged operator algebra is supplied data; it
    is not obtained by extending its values on operators $X\otimes I$.
    % Source lines: paper_v2.tex 1356--1364.
\end{definition}

\begin{theorem}[Identity-ancilla transport of symmetry invariance]
    \label{thm:mpu_symmetry_identity_ancilla_invariant}
    \lean{MPOTensor.IsInvariantUnderSymmetry.tensorPhysicalId}
    \leanok
    \uses{def:mpu_coherent_symmetry_family, def:equivalent-tensors}
    If $\mathcal U$ is invariant under the $d$-dimensional component of a
    coherent symmetry family, then $\mathcal U^{(x)}$ is invariant under its
    $dx$-dimensional component.
\end{theorem}

\begin{proof}\leanok
    The finite-chain operator of $\mathcal U^{(x)}$ is $\iota_x(U^{(N)})$.
    Semiconjugacy and invariance of $U^{(N)}$ give
    \begin{align}
        \mathcal S_{dx,N}\bigl(\iota_x(U^{(N)})\bigr)
         & = \iota_x\bigl(\mathcal S_{d,N}(U^{(N)})\bigr)
        = \iota_x(U^{(N)}).
    \end{align}
\end{proof}

\begin{definition}[Blocked finite-chain symmetry]
    \label{def:mpu_blocked_finite_chain_symmetry}
    \lean{MPOTensor.FiniteChainOperatorSymmetry.block,
        MPOTensor.FiniteChainOperatorSymmetry.block_action}
    \leanok
    \uses{def:mpu_finite_chain_operator_symmetry, lem:eval_word_block}
    For a blocking length $k$, the blocked symmetry has applicable lengths
    \begin{align}
        A_k & = \{N\mid Nk\in A\}.
    \end{align}
    Its action is obtained by conjugating $\mathcal S_{d,Nk}$ along the
    canonical identification
    $(\mathbb C^{d^k})^{\otimes N}\cong(\mathbb C^d)^{\otimes Nk}$.
    % Source line: paper_v2.tex 1366.
\end{definition}

\begin{theorem}[Blocking transport of symmetry invariance]
    \label{thm:mpu_symmetry_blocking_invariant}
    \lean{MPOTensor.IsInvariantUnderSymmetry.blockTensor}
    \leanok
    \uses{def:mpu_blocked_finite_chain_symmetry, thm:mpdo_closed_chain_physical_blocking}
    If $\mathcal U$ is invariant under $\mathcal S$, then its $k$-site blocking
    is invariant under the blocked symmetry.
\end{theorem}

\begin{proof}\leanok
    Let $R_{d,N,k}$ be the canonical identification from blocked length-$N$
    configurations to unblocked length-$Nk$ configurations. For $N\in A_k$,
    one has $Nk\in A$ and
    \begin{align}
        \mathcal S_{k,N}\bigl[U_k^{(N)}\bigr]
         & = R_{d,N,k}^{-1}\!\left(
        \mathcal S_{d,Nk}\!\left[
                R_{d,N,k}\bigl(U_k^{(N)}\bigr)
        \right]\right)                               \\
         & = R_{d,N,k}^{-1}\!\left(
        \mathcal S_{d,Nk}\bigl[U^{(Nk)}\bigr]\right) \\
         & = R_{d,N,k}^{-1}\bigl(U^{(Nk)}\bigr)      \\
         & = U_k^{(N)}.
    \end{align}
    The second and final equalities are the inverse reindexing identities, and
    the third is invariance of $U^{(Nk)}$ because $Nk\in A$.
\end{proof}

\begin{definition}[Strict equivalence under a symmetry]
    \label{def:strictly-equivalent-symmetry}
    \lean{MPOTensor.StrictlyEquivalentUnderSymmetry}
    \leanok
    \discussion{6948}
    \uses{def:strictly-equivalent-tensors, def:mpu_finite_chain_operator_symmetry}
    % **Local fix:** paper_v2.tex 1345--1353 repeats the symmetry-transformed
    % operator on the right-hand side. Invariance requires the original MPU.
    % Documented in docs/paper-gaps/mpu_canonical_form_full_support.tex.
    Fix a virtual bond dimension $D$. Two canonical-form tensors $\mathcal U$
    and $\mathcal V$ are strictly equivalent under a symmetry $\mathcal S$ if
    they have the same physical dimension and there is a continuous path
    $\mathcal W(p)$ of fixed-bond MPU tensors from $\mathcal U$ to $\mathcal V$,
    not necessarily in canonical form, such that
    \begin{align}
        \mathcal S[W^{(N)}(p)] & =W^{(N)}(p)
    \end{align}
    for every $p\in[0,1]$ and every applicable chain length $N$.
    % Source lines: paper_v2.tex 1340--1354.
\end{definition}

\begin{lemma}[Forgetting symmetry from strict equivalence]
    \label{lem:mpu_strict_symmetry_equivalence_forget}
    \lean{MPOTensor.StrictlyEquivalentUnderSymmetry.toStrictlyEquivalent}
    \leanok
    \uses{def:strictly-equivalent-symmetry, def:strictly-equivalent-tensors}
    Strict equivalence under a symmetry implies strict equivalence.
\end{lemma}

\begin{proof}\leanok
    The same path lies in the larger locus obtained by forgetting invariance
    under $\mathcal S$.
\end{proof}

\begin{definition}[Admissible strict equivalence under a symmetry]
    \label{def:mpu_admissible_strict_symmetry_equivalence}
    \notready
    \discussion{6017}
    \uses{def:strictly-equivalent-tensors, def:mpu_admissible_equivalence_datum}
    Two tensors $\mathcal U$ and $\mathcal V$ are strictly equivalent under a
    symmetry $\mathcal S$ if they have the same physical dimension and there is
    a continuous path $\mathcal W(p)$ of MPU tensors from $\mathcal U$ to
    $\mathcal V$ such that
    \begin{align}
        \mathcal S[W^{(N)}(p)] & =W^{(N)}(p)
    \end{align}
    for every $p\in[0,1]$ and every chain length $N$ under consideration. The
    path need not be in canonical form. The equivalence is admissible if its
    witnessing path carries an admissible equivalence datum for the actual
    blocked path tensor and the fixed conjugate, physical-adjoint, transposed,
    and $\mathcal S$-image comparison paths, and if the common blocking length
    $m$ of that datum satisfies $m\leq D^4$.
    % Source lines: paper_v2.tex 1340--1354.
\end{definition}

\begin{definition}[Equivalence under a symmetry]
    \label{def:equivalent-symmetry}
    \notready
    \discussion{7422}
    \uses{def:strictly-equivalent-symmetry, def:equivalent-tensors}
    Let $\mathcal U$ and $\mathcal V$ have physical dimensions $d_a$ and
    $d_b$. They are equivalent under a symmetry $\mathcal S$ if there are a
    positive blocking length $k$ and coprime positive ancilla dimensions
    $p_a,p_b$ with $p_ad_a=p_bd_b$ such that $(\mathcal U^{(p_a)})_k$ and
    $(\mathcal V^{(p_b)})_k$ are strictly equivalent under $\mathcal S$.
    % Source lines: paper_v2.tex 1356--1366.
\end{definition}

\begin{definition}[Fixed-bond coherent-family equivalence under a symmetry]
    \label{def:mpu_coherent_family_equivalent_symmetry}
    \lean{MPOTensor.EquivalentUnderSymmetry}
    \leanok
    \discussion{7015}
    \uses{def:strictly-equivalent-symmetry, def:equivalent-tensors,
        def:mpu_coherent_symmetry_family, def:mpu_blocked_finite_chain_symmetry}
    % **Scope restriction (fixed bond and coherent family):** the symmetry is
    % supplied at every physical dimension and semiconjugates the canonical
    % identity-ancilla embeddings. Documented in
    % docs/paper-gaps/mpu_symmetry_ancilla_transport.tex.
    Fix a virtual bond dimension $D$, and let $\mathcal U$ and $\mathcal V$
    have physical dimensions $d_a$ and $d_b$. They are coherently equivalent
    under a dimension-indexed symmetry family $\mathcal S$ if they are MPUs and
    there are a positive blocking length $k$ and coprime positive ancilla
    dimensions $p_a,p_b$ with $p_ad_a=p_bd_b$ such that
    $(\mathcal U^{(p_a)})_k$ and
    $(\mathcal V^{(p_b)})_k$ are strictly equivalent under the $k$-site
    blocking of the supplied enlarged-dimension symmetry. The family obeys the
    identity-ancilla semiconjugacy law, and its applicable blocked lengths are
    precisely the lengths $N$ for which $Nk\in A$. The ancillas are attached
    before the common blocking.
    % Source lines: paper_v2.tex 1356--1366, under the stated restrictions.
\end{definition}

\begin{theorem}[Forgetting symmetry from coherent stabilized equivalence]
    \label{thm:mpu_symmetry_equivalence_forget}
    \lean{MPOTensor.EquivalentUnderSymmetry.toEquivalent}
    \leanok
    \uses{def:mpu_coherent_family_equivalent_symmetry, def:equivalent-tensors,
        lem:mpu_strict_symmetry_equivalence_forget}
    Fixed-bond equivalence under a coherent symmetry family implies ordinary
    equivalence after the same ancillas and blocking.
\end{theorem}

\begin{proof}\leanok
    Retain the MPU hypotheses, positive coprime ancilla dimensions, common
    positive blocking length, and physical-dimension equality from the
    symmetry-preserving witness. Forgetting invariance from its final strict
    path gives the required ordinary stabilized equivalence.
\end{proof}

\begin{definition}[Admissible equivalence under a symmetry]
    \label{def:mpu_admissible_symmetry_equivalence}
    \notready
    \discussion{6017}
    \uses{def:mpu_admissible_strict_symmetry_equivalence, def:equivalent-tensors, def:mpu_admissible_equivalence_datum}
    Let $\mathcal U$ and $\mathcal V$ have physical dimensions $d_a$ and
    $d_b$. They are equivalent under a symmetry $\mathcal S$ if there are a
    positive blocking length $k$ and coprime positive ancilla dimensions
    $p_a,p_b$ with $p_ad_a=p_bd_b$ such that $(\mathcal U^{(p_a)})_k$ and
    $(\mathcal V^{(p_b)})_k$ are strictly equivalent under $\mathcal S$. Thus
    the ancillas are attached before the common blocking. The equivalence is
    admissible if this strict-equivalence witness carries an admissible
    equivalence datum for the actual blocked path and the fixed conjugate,
    physical-adjoint, transposed, and $\mathcal S$-image comparison paths, with
    common blocking length at most $D^4$.
    % Source lines: paper_v2.tex 1356--1364.
\end{definition}

\begin{lemma}[Standard-form path criterion]
    \label{lem:equivalent-symmetry}
    \notready
    \discussion{6017}
    \uses{def:strictly-equivalent-symmetry, coruvSF, FundamentalMPU}
    Two MPUs described by simple tensors are strictly equivalent under
    $\mathcal S$ if and only if their standard-form unitaries are joined by
    continuous families $\widetilde u(p)$ and $\widetilde v(p)$ whose generated
    MPU is invariant under $\mathcal S$ for every $p\in[0,1]$.
    % Source lines: paper_v2.tex 1366--1380.
\end{lemma}

\begin{lemma}[Admissible standard-form path criterion]
    \label{lem:mpu_admissible_symmetry_path_criterion}
    \notready
    \discussion{6017}
    % **Scope restriction (admissible symmetry path):** the strict-equivalence
    % witness carries the path datum of
    % Definition~\ref{def:mpu_admissible_equivalence_datum}. Documented in
    % docs/paper-gaps/mpu_canonical_form_full_support.tex.
    \uses{def:mpu_admissible_strict_symmetry_equivalence, cor:mpu_admissible_continuous_standard_form, thm:mpu_admissible_fundamental, def:mpu_admissible_equivalence_datum}
    Fix a positive integer $m\leq D^4$. The actual $m$-site blocked endpoint
    MPUs $U_1^{(mN)}$ and $U_2^{(mN)}$ are admissibly strictly equivalent under
    $\mathcal S$ if and only if there are continuous families of unitaries
    $\widetilde u(p)$ and $\widetilde v(p)$ whose endpoints generate
    $U_1^{(mN)}$ and $U_2^{(mN)}$, whose generated MPU is invariant under
    $\mathcal S$ for every $p\in[0,1]$, and whose tensor path and fixed
    conjugate, physical-adjoint, transposed, and $\mathcal S$-image comparison
    paths carry continuously varying reduced full-support source data. These
    families arise from the actual blocked path $\mathcal W(p)_m$ in an
    admissible equivalence datum of common length $m$; the blocked path itself is
    the length-one admissible witness for the displayed endpoint families. When
    $m>1$, this criterion does not assert strict equivalence of the unblocked
    families $U_1^{(N)}$ and $U_2^{(N)}$. This is the reduced-source form of
    \cite[Lemma following the symmetry-equivalence definitions]{Cirac2017MPU}.
    % Source lines: paper_v2.tex 1366--1380.
\end{lemma}

\begin{theorem}[Conjugation symmetry when one source rank is odd]
    \label{conjphs}
    \notready
    \discussion{5713}
    \uses{prop1conj, corconj, lem:equivalent-symmetry}
    Suppose that either $r$ or $\ell$ is odd. Two conjugation-invariant MPUs
    $\mathcal U_1$ and $\mathcal U_2$ are strictly equivalent under
    conjugation if and only if
    \begin{align}
        \det(u_1v_1) & =\det(u_2v_2).
    \end{align}
    % Source lines: paper_v2.tex 1401--1446.
\end{theorem}

\begin{theorem}[Admissible conjugation symmetry when one source rank is odd]
    \label{thm:mpu_admissible_conjugation_odd_rank}
    \notready
    \discussion{5713}
    % **Scope restriction (endpoints and admissible symmetry path):** each
    % original endpoint and its conjugate comparison carry source data, and
    % equivalence means admissible strict equivalence in
    % Definition~\ref{def:mpu_admissible_equivalence_datum}. Documented in
    % docs/paper-gaps/mpu_canonical_form_full_support.tex.
    \uses{prop:mpu_admissible_conjugation_standard_forms, lem:mpu_admissible_symmetry_path_criterion, def:mpu_admissible_equivalence_datum, def:mpu_reduced_full_support_source_datum}
    For $i=1,2$, suppose the original one-site endpoint tensor
    $\mathcal U_i$ and its conjugate comparison tensor are simple and each
    carries reduced full-support source data. Suppose that either $r$ or $\ell$
    is odd and that the standard-form interpolation determined below carries an
    admissible path datum of blocking length one. Then $\mathcal U_1$ and
    $\mathcal U_2$ are admissibly strictly equivalent
    under conjugation symmetry if and only if
    \begin{align}
        \det(u_1v_1) & =\det(u_2v_2).
    \end{align}
    % Source lines: paper_v2.tex 1401--1446.
\end{theorem}

\begin{corollary}[Equivalence of conjugation-symmetric MPUs after blocking]
    \label{cor:all-equivalent-under-conjugation}
    \notready
    \uses{conjphs, def:equivalent-symmetry, lem:mpu_blocking}
    If either $r$ or $\ell$ is odd, all conjugation-invariant MPUs with these
    source ranks are equivalent under conjugation symmetry.
    % Source lines: paper_v2.tex 1450--1460.
\end{corollary}

\begin{corollary}[Admissible equivalence of conjugation-symmetric MPUs after blocking]
    \label{cor:mpu_admissible_all_conjugation_equivalent}
    \notready
    \discussion{5713}
    % **Scope restriction (endpoints, blocking, and admissible path):** the
    % one-site and two-site endpoints and their conjugate comparisons carry
    % source data, and the resulting symmetry path is admissible. No
    % preservation under blocking is asserted. Documented in
    % docs/paper-gaps/mpu_canonical_form_full_support.tex.
    \uses{prop:mpu_admissible_conjugation_standard_forms, cor:mpu_admissible_equiv_under_conjugation, def:mpu_admissible_symmetry_equivalence, def:mpu_admissible_equivalence_datum, def:mpu_reduced_full_support_source_datum}
    For $i=1,2$, suppose the endpoint tensor $\mathcal U_i$, its conjugate
    comparison $\overline{\mathcal U_i}$, the actual two-site block
    $\mathcal U_{i,2}$, and its conjugate comparison
    $\overline{\mathcal U_{i,2}}$ are simple and each carries reduced
    full-support source data. Suppose that either $r$ or $\ell$ is odd and that
    the resulting standard-form interpolation carries an admissible path datum
    of blocking length one. Then
    $\mathcal U_1$ and $\mathcal U_2$ are admissibly equivalent under
    conjugation symmetry.
    % Source lines: paper_v2.tex 1450--1460.
\end{corollary}

\begin{proof}
    \uses{prop:mpu_admissible_conjugation_standard_forms, cor:mpu_admissible_equiv_under_conjugation}
    If either $r$ or $\ell$ is odd, both endpoints lie in Case I of
    Proposition~\ref{prop:mpu_admissible_conjugation_standard_forms}. The full
    admissible classification in
    Corollary~\ref{cor:mpu_admissible_equiv_under_conjugation} handles both
    possible parities of the source ranks after blocking and gives the claimed
    admissible equivalence.
\end{proof}

\begin{theorem}[Strict conjugation equivalence for even source ranks]
    \label{prop:conjsym-strict-equiv}
    \notready
    \discussion{5713}
    \uses{prop1conj, corconj, cor:sympl-to-orthogonal, lem:equivalent-symmetry, lemma:conjclass-normalform-continuous}
    Suppose $r$ and $\ell$ are even. Two conjugation-invariant MPUs in the
    standard forms of Proposition~\ref{corconj} are strictly equivalent under
    conjugation if and only if they belong to the same case and
    \begin{align}
        \det(u_1) & =\det(u_2), &
        \det(v_1) & =\det(v_2).
    \end{align}
    % Source lines: paper_v2.tex 1464--1546.
\end{theorem}

\begin{theorem}[Admissible strict conjugation equivalence for even source ranks]
    \label{thm:mpu_admissible_conjugation_even_rank}
    \notready
    \discussion{5713}
    % **Scope restriction (admissible symmetry path):** equivalence means
    % admissible strict equivalence in
    % Definition~\ref{def:mpu_admissible_equivalence_datum}. Documented in
    % docs/paper-gaps/mpu_canonical_form_full_support.tex.
    \uses{prop:mpu_admissible_conjugation_standard_forms, cor:mpu_admissible_symplectic_to_orthogonal, lem:mpu_admissible_symmetry_path_criterion, lemma:conjclass-normalform-continuous, def:mpu_admissible_equivalence_datum}
    Suppose that $r$ and $\ell$ are both even and that the standard-form
    interpolation determined below carries an admissible path datum of blocking
    length one. Choose admissible standard forms of two conjugation-invariant
    MPUs as in
    Proposition~\ref{prop:mpu_admissible_conjugation_standard_forms}. They are admissibly strictly equivalent under
    conjugation symmetry if and only if they belong
    to the same case of Proposition~\ref{prop:mpu_admissible_conjugation_standard_forms} and
    \begin{align}
        \det(u_1) & =\det(u_2), &
        \det(v_1) & =\det(v_2).
    \end{align}
    % Source lines: paper_v2.tex 1464--1546.
\end{theorem}

\begin{corollary}[Equivalence under conjugation]
    \label{cor:equiv-under-conjugation}
    \notready
    \discussion{5713}
    \uses{corconj, prop:conjsym-strict-equiv, def:equivalent-symmetry}
    Two conjugation-invariant MPUs are equivalent under conjugation symmetry if
    and only if they belong to the same case of Proposition~\ref{corconj}.
    % Source lines: paper_v2.tex 1548--1567.
\end{corollary}

\begin{corollary}[Admissible equivalence under conjugation]
    \label{cor:mpu_admissible_equiv_under_conjugation}
    \notready
    \discussion{5713}
    % **Scope restriction (endpoints, blocking, and admissible path):** the
    % original one-site endpoints and their conjugate comparisons carry source
    % data, all blocked endpoints used below carry their own data, and the
    % symmetry path is admissible. Documented in
    % docs/paper-gaps/mpu_canonical_form_full_support.tex.
    \uses{thm:mpu_admissible_conjugation_odd_rank, thm:mpu_admissible_conjugation_even_rank, prop:mpu_admissible_conjugation_standard_forms, prop:mpu_admissible_index_well_defined, def:mpu_admissible_symmetry_equivalence, def:mpu_admissible_equivalence_datum, def:mpu_reduced_full_support_source_datum}
    For $i=1,2$, suppose the original one-site endpoint tensor
    $\mathcal U_i$, its conjugate comparison $\overline{\mathcal U_i}$, the
    actual two-site block $\mathcal U_{i,2}$, and its conjugate comparison
    $\overline{\mathcal U_{i,2}}$ are simple and each carries reduced
    full-support source data. Suppose also that the resulting standard-form
    interpolation carries an admissible path datum of blocking length one. Then
    $\mathcal U_1$ and $\mathcal U_2$ are admissibly equivalent under
    conjugation symmetry if and only if they belong to the same case of
    Proposition~\ref{prop:mpu_admissible_conjugation_standard_forms}.
    % Source lines: paper_v2.tex 1548--1567.
\end{corollary}

\begin{proof}
    \uses{thm:mpu_admissible_conjugation_odd_rank, thm:mpu_admissible_conjugation_even_rank, prop:mpu_admissible_conjugation_standard_forms, prop:mpu_admissible_index_well_defined, def:mpu_admissible_symmetry_equivalence}
    Suppose first that either $r$ or $\ell$ is odd. Then both original
    endpoints lie in Case I. After one blocking, write $r^{(2)}$ and
    $\ell^{(2)}$ for the blocked source ranks. Admissible blocking independence,
    applied to the supplied one-site and two-site endpoint data, gives
    \begin{align}
        r^{(2)} & =dr,
    \end{align}
    and
    \begin{align}
        \ell^{(2)} & =d\ell.
    \end{align}
    For each endpoint,
    \begin{align}
        \det(u_i^{(2)}v_i^{(2)}) & =\det(u_iv_i)^{2r\ell}=1.
    \end{align}
    Blocking preserves Case I. If either $r^{(2)}$ or $\ell^{(2)}$ is odd, the
    odd-rank theorem applied to the blocked endpoints gives admissible strict
    equivalence. If both blocked ranks are even, then $d$ is even because one
    original source rank is odd. Hence $r\ell=d^2$ is even.
    % **Local fix:** paper_v2.tex 1554--1559 states the two separate determinant
    % identities before checking parity. They are used here only when the
    % blocked source ranks are both even, which forces the required even power.
    % Documented in docs/paper-gaps/mpu_canonical_form_full_support.tex.
    Consequently,
    \begin{align}
        \det(u_i^{(2)}) & =(\det(u_i)^2\det(v_i))^{r\ell}=1,
    \end{align}
    and
    \begin{align}
        \det(v_i^{(2)}) & =\det(v_i)^{r\ell}=1.
    \end{align}
    The even-rank theorem now applies to the preserved Case I endpoints. Thus
    the odd-source-rank branch has a single admissible equivalence class without
    imposing a parity condition on the blocked ranks.

    If $r$ and $\ell$ are both even, admissible blocking independence shows that
    the blocked ranks remain even. Since $r\ell$ is even, the same calculations
    give
    \begin{align}
        \det(u_i^{(2)}) & =(\det(u_i)^2\det(v_i))^{r\ell}=1,
    \end{align}
    and
    \begin{align}
        \det(v_i^{(2)}) & =\det(v_i)^{r\ell}=1.
    \end{align}
    Blocking preserves Case I and Case II. The even-rank theorem therefore shows
    that the blocked endpoints are admissibly strictly equivalent exactly when
    they belong to the same case. This is precisely admissible equivalence of
    the original endpoints.
\end{proof}

\begin{lemma}[SWAP trace contraction]
    \label{lem:matrix_trace_swap_kronecker}
    \lean{Matrix.trace_swapMatrix_mul_kronecker,
        Matrix.trace_swapMatrix_mul_kronecker_conjTranspose_of_mem_unitaryGroup}
    \leanok
    For matrices $A,B\in\mathbb C^{n\times n}$ and the exchange matrix $\mathbb S$ on
    $\mathbb C^n\otimes\mathbb C^n$,
    \begin{align}
        \tr\!\left[\mathbb S(A\otimes B)\right] & =\tr(AB).
    \end{align}
    In particular, if $U$ is unitary, then
    \begin{align}
        \tr\!\left[\mathbb S(U\otimes U^\dagger)\right] & =n.
    \end{align}
    % Source identity: paper_v2.tex lines 1608--1610.
\end{lemma}

\begin{proof}\leanok
    In product coordinates, the exchange matrix has entries
    $\mathbb S_{(i,j),(k,l)}=\delta_{i,l}\delta_{j,k}$. Therefore
    \begin{align}
        \tr\!\left[\mathbb S(A\otimes B)\right]
         & =\sum_{i,j}A_{j,i}B_{i,j}
        =\sum_{i,j}A_{i,j}B_{j,i}
        =\tr(AB).
    \end{align}
    Taking $A=U$ and $B=U^\dagger$ gives
    $\tr(UU^\dagger)=\tr(\Id)=n$.
\end{proof}

\begin{lemma}[Trace formula for the time-reversal sign]
    \label{eq:sigma-from-trace}
    \notready
    \discussion{5714}
    \uses{PropChiral, lem:matrix_trace_swap_kronecker}
    For a time-reversal-invariant MPU in standard form, the sign $\sigma$ is
    \begin{align}
        \sigma & =\frac1{d^2}\tr\!\left[\mathbb S_{12,34}
            v_{23}u_{12}u_{34}v_{23}\right],
    \end{align}
    where $\mathbb S_{12,34}$ exchanges the pairs of sites $12$ and $34$.
    % Source lines: paper_v2.tex 1588--1611.
\end{lemma}

\begin{lemma}[Admissible trace formula for the time-reversal sign]
    \label{lem:mpu_admissible_sigma_trace}
    \notready
    \discussion{5714}
    % **Scope restriction (two-site blocked comparison):** the trace formula uses
    % the blocked source data required by Proposition~\ref{prop:mpu_admissible_local_time_reversal}.
    % Documented in docs/paper-gaps/mpu_canonical_form_full_support.tex.
    \uses{prop:mpu_admissible_local_time_reversal, def:mpu_reduced_full_support_source_datum}
    For a time-reversal-invariant MPU in admissible standard form, suppose its
    one-site tensor $\mathcal U$, its actual two-site block $\mathcal U_2$, and
    the physical-adjoint comparison tensor $\mathcal U_2^\sharp$ are simple and
    each carries its own reduced full-support source datum. Then the sign
    $\sigma$ in
    Proposition~\ref{prop:mpu_admissible_local_time_reversal} is
    \begin{align}
        \sigma & =\frac1{d^2}\tr\!\left[\mathbb S_{12,34}v_{23}u_{12}u_{34}v_{23}\right],
    \end{align}
    where $\mathbb S_{12,34}$ exchanges the pairs of sites $12$ and $34$.
    % Source lines: paper_v2.tex 1588--1611.
\end{lemma}

\begin{lemma}[Gauge invariance of the time-reversal sign]
    \label{lem:mpu_sigma_gauge_invariant}
    \notready
    \discussion{5714}
    \uses{eq:sigma-from-trace}
    Under $u'=(x\otimes y)u$ and
    $v'=v(y^\dagger\otimes x^\dagger)$, the trace formula in
    Lemma~\ref{eq:sigma-from-trace} has the same value for $(u',v')$ as for
    $(u,v)$. Thus $\sigma$ is independent of the standard-form gauge.
    % Source proof: paper_v2.tex 1624--1630.
\end{lemma}

\begin{lemma}[Admissible gauge invariance of the time-reversal sign]
    \label{lem:mpu_admissible_sigma_gauge_invariant}
    \notready
    \discussion{5714}
    \uses{lem:mpu_admissible_sigma_trace}
    % Source proof: paper_v2.tex 1624--1630.
    Under the hypotheses of Lemma~\ref{lem:mpu_admissible_sigma_trace}, let $x$ and $y$ be
    unitaries and change the standard-form gauge by
    \begin{align}
        u' & =(x\otimes y)u,                 &
        v' & =v(y^\dagger\otimes x^\dagger).
    \end{align}
    Then the trace formula has the same value for $(u',v')$ as for $(u,v)$.
    Consequently the sign $\sigma$ is independent of the standard-form gauge.
\end{lemma}

\begin{proof}
    \uses{lem:mpu_admissible_sigma_trace}
    Write $x_i$ and $y_i$ for the unitaries acting on site $i$. Then
    $u'_{12}=x_1y_2u_{12}$, $u'_{34}=x_3y_4u_{34}$, and
    $v'_{23}=v_{23}y_2^\dagger x_3^\dagger$.
    Substitution into Lemma~\ref{lem:mpu_admissible_sigma_trace}, cancellation on the
    middle legs, and
    $\mathbb S_{12,34}x_1y_4=x_3y_2\mathbb S_{12,34}$ give
    \begin{align}
        \tr\!\left[\mathbb S_{12,34}v'_{23}u'_{12}u'_{34}v'_{23}\right]
         & =\tr\!\left[x_3y_2\mathbb S_{12,34}
        v_{23}u_{12}u_{34}v_{23}y_2^\dagger x_3^\dagger\right]           \\
         & =\tr\!\left[\mathbb S_{12,34}v_{23}u_{12}u_{34}v_{23}\right].
    \end{align}
    The last equality follows from cyclicity of the trace and unitarity of
    $x$ and $y$.
\end{proof}

\begin{proposition}[Well-definedness of the time-reversal sign]
    \label{prop:sigma-well-defined}
    \notready
    \discussion{5714}
    \uses{eq:sigma-from-trace, lem:mpu_sigma_gauge_invariant}
    For a time-reversal-invariant MPU in standard form, the sign
    $\sigma=\pm1$ is independent of the gauge and is constant under continuous
    deformations that preserve time-reversal symmetry.
    % Source lines: paper_v2.tex 1613--1631.
\end{proposition}

\begin{proposition}[Admissible well-definedness of the time-reversal sign]
    \label{prop:mpu_admissible_sigma_well_defined}
    \notready
    \discussion{5714}
    % **Scope restriction (admissible deformation):** the deformation carries the
    % path datum of Definition~\ref{def:mpu_admissible_equivalence_datum}.
    % Documented in docs/paper-gaps/mpu_canonical_form_full_support.tex.
    \uses{lem:mpu_admissible_sigma_gauge_invariant, lem:mpu_admissible_symmetry_path_criterion, def:mpu_reduced_full_support_source_datum, def:mpu_admissible_equivalence_datum}
    Suppose a time-reversal-preserving path carries an admissible equivalence
    datum of blocking length one. At every path point, require the one-site
    tensor, its actual two-site block, and the physical-adjoint comparison of
    that two-site block to be simple and to carry continuously varying reduced
    full-support source data. Then the sign
    $\sigma=\pm1$ is independent of the admissible standard-form gauge and is
    constant along the path.
    % Source lines: paper_v2.tex 1613--1631.
\end{proposition}

\begin{proposition}[Blocking parity of the time-reversal sign]
    \label{prop:sigma1-and-sigma2-equal}
    \notready
    \discussion{5714}
    \uses{prop:sigma-well-defined, PropChiral, lem:mpu_blocking}
    Let $\mathcal U$ be a time-reversal-invariant MPU in standard form, and let
    $\sigma^{(k)}$ be the sign obtained from a standard form of its $k$-site
    block. Then
    \begin{align}
        \sigma^{(k)}
         & =\begin{cases}
                \sigma^{(1)}, & k\text{ odd},  \\
                \sigma^{(2)}, & k\text{ even}.
            \end{cases}
    \end{align}
    % Source lines: paper_v2.tex 1633--1733.
\end{proposition}

\begin{proposition}[Admissible blocking parity of the time-reversal sign]
    \label{prop:mpu_admissible_sigma_blocking_parity}
    \notready
    \discussion{5714}
    % **Scope restriction (blocked source data):** every blocked tensor used to
    % define $\sigma^{(k)}$, including its two-site comparison tensor, carries
    % reduced full-support source data. No preservation under blocking is
    % asserted. Documented in
    % docs/paper-gaps/mpu_canonical_form_full_support.tex.
    \uses{lem:mpu_admissible_sigma_gauge_invariant, prop:mpu_admissible_local_time_reversal, lem:mpu_blocking, def:mpu_reduced_full_support_source_datum}
    For every blocking length $k$ under consideration, suppose the endpoint
    tensor $\mathcal U_k$, its actual two-site block $\mathcal U_{2k}$, and the
    physical-adjoint comparison tensor $\mathcal U_{2k}^\sharp$ are simple and
    each carries its own reduced full-support source datum. Let $\sigma^{(k)}$
    be the sign of the resulting standard form. Then
    \begin{align}
        \sigma^{(k)}
         & =\begin{cases}
                \sigma^{(1)}, & k\text{ odd},  \\
                \sigma^{(2)}, & k\text{ even}.
            \end{cases}
    \end{align}
    % Source lines: paper_v2.tex 1633--1733.
\end{proposition}

\begin{proof}
    \uses{lem:mpu_admissible_sigma_gauge_invariant, prop:mpu_admissible_local_time_reversal, lem:mpu_blocking}
    Gauge invariance permits the standard blocked representatives used in the
    source proof. Choose
    \begin{align}
        u^{(2)} & =v_{23}u_{12}u_{34}, & v^{(2)} & =v_{45},
    \end{align}
    and put $\zeta=\sigma^{(2)}/\sigma^{(1)}$. The one- and two-site local
    time-reversal relations give the unwinding identity
    \begin{align}
        \zeta\big[(\Id\otimes v\otimes\Id)(u\otimes x)\big]^\dagger
         & =(\Id\otimes v\otimes\Id)(x^\dagger\otimes u).
    \end{align}
    Applying this identity successively through a block of length $k$ gives
    \begin{align}
        \sigma^{(k)}
         & =\sigma^{(1)}\zeta^{k-1}.
    \end{align}
    Since each sign is $\pm1$, this equals $\sigma^{(1)}$ for odd $k$ and
    $\sigma^{(2)}$ for even $k$.
    % Source lines: paper_v2.tex 1657--1732.
\end{proof}

\begin{corollary}[Necessary time-reversal invariants]
    \label{cor:sigma-equivalence-necessary}
    \notready
    \discussion{5714}
    \uses{prop:sigma1-and-sigma2-equal, def:strictly-equivalent-symmetry, def:equivalent-symmetry}
    Strictly equivalent time-reversal-invariant MPUs have the same
    $\sigma^{(1)}$ and $\sigma^{(2)}$. Equivalent time-reversal-invariant MPUs
    have the same $\sigma^{(2)}$.
    % Source lines: paper_v2.tex 1735--1743.
\end{corollary}

\begin{corollary}[Admissible time-reversal invariants of supplied witnesses]
    \label{cor:mpu_admissible_sigma_equivalence_necessary}
    \notready
    \discussion{5714}
    % **Scope restriction (admissible symmetry equivalence):** the endpoint
    % blockings and symmetry paths carry the data of
    % Definition~\ref{def:mpu_admissible_equivalence_datum}. The non-strict
    % conclusion concerns only the supplied ancilla-enlarged blocked endpoints.
    % Documented in docs/paper-gaps/mpu_canonical_form_full_support.tex.
    \uses{prop:mpu_admissible_sigma_blocking_parity, def:mpu_admissible_strict_symmetry_equivalence, def:mpu_admissible_symmetry_equivalence, def:mpu_reduced_full_support_source_datum, def:mpu_admissible_equivalence_datum}
    Let two time-reversal-invariant MPUs be compared by an admissible symmetry
    path $p\mapsto\mathcal W(p)$ whose path datum has blocking length one.
    Suppose that the five path families $\mathcal W(p)$,
    $\mathcal W(p)_2$, $\mathcal W(p)_4$,
    $\mathcal W(p)_2^\sharp$, and $\mathcal W(p)_4^\sharp$ are simple and carry
    continuously varying reduced full-support source data. Then admissible
    strict equivalence requires equality of both $\sigma^{(1)}$ and
    $\sigma^{(2)}$.

    For admissible equivalence after adjoining ancillas and applying the source
    blocking, let $p\mapsto\mathcal W_{\mathrm{anc}}(p)$ be the actual
    ancilla-enlarged blocked witness path, not a reduced or equivalent
    representative, and regard it as a length-one admissible path. Require
    continuous reduced full-support source data along
    $\mathcal W_{\mathrm{anc}}(p)$, its actual two- and four-site blocks, and
    the physical-adjoint comparisons of those two blocks. The conclusion is
    only
    \begin{align}
        \sigma^{(2)}(\mathcal W_{\mathrm{anc}}(0))
         & =\sigma^{(2)}(\mathcal W_{\mathrm{anc}}(1)).
    \end{align}
    No equality between $\sigma^{(2)}$ of the original two MPUs is asserted.
    Such a conclusion requires preservation under adjoining identity ancillas,
    which is not established here.
    % Tracked by issue #6138.
    % Source lines: paper_v2.tex 1735--1743.
\end{corollary}

\begin{proposition}[Conjugation class and the MPS sign]
    \label{prop:mps-class-eq-mpu-class-for-conjugation}
    \notready
    \discussion{5714}
    \uses{corconj, cor:equiv-under-conjugation}
    For a conjugation-invariant MPU, the MPS sign
    $G\overline G=\pm\Id$ agrees with Case I and Case II, respectively, of
    Proposition~\ref{corconj}.
    % Source lines: paper_v2.tex 1791--1805.
\end{proposition}

\begin{proposition}[Admissible conjugation class and the MPS sign]
    \label{prop:mpu_admissible_mps_conjugation_class}
    \notready
    \discussion{5714}
    % **Scope restriction (admissible conjugation class):** the standard forms,
    % blockings, and comparison paths carry the required source data. Documented
    % in docs/paper-gaps/mpu_canonical_form_full_support.tex.
    \uses{prop:mpu_admissible_conjugation_standard_forms, cor:mpu_admissible_equiv_under_conjugation, def:mpu_reduced_full_support_source_datum, def:mpu_admissible_equivalence_datum}
    Let $\mathcal U$ be conjugation invariant. Suppose $\mathcal U$ and its
    conjugate comparison tensor are simple and carry their own reduced
    full-support source data. Let $p\mapsto\mathcal W(p)$ be the actual
    standard-form interpolation from $\mathcal U$ to its Case I or Case II
    representative. Suppose this interpolation carries an admissible
    conjugation-symmetry path datum of blocking length one: the path
    $\mathcal W(p)$ and its fixed conjugate, physical-adjoint, transposed, and
    conjugation-image comparison paths are simple and carry continuously varying
    source data. Then the sign
    $G\overline G=\pm\Id$ in the MPS classification agrees with Case I and
    Case II, respectively, of
    Proposition~\ref{prop:mpu_admissible_conjugation_standard_forms}.
    % Source lines: paper_v2.tex 1791--1805.
\end{proposition}

\begin{proposition}[Time-reversal class and the MPS sign]
    \label{prop:dagger-mps-class-and-mpu-class-equal}
    \notready
    \discussion{5714}
    \uses{PropChiral, prop:sigma1-and-sigma2-equal}
    For a time-reversal-invariant MPU in standard form, the MPS sign
    $G\overline G=\pm\Id$ equals the sign $\sigma^{(2)}$ obtained after
    blocking two sites.
    % Source lines: paper_v2.tex 1807--1829.
\end{proposition}

\begin{proposition}[Admissible time-reversal class and the MPS sign]
    \label{prop:mpu_admissible_mps_time_reversal_class}
    \notready
    \discussion{5714}
    % **Scope restriction (one-, two-, and four-site source data):** the tensors
    % used to define both $\sigma^{(1)}$ and $\sigma^{(2)}$ carry their own data.
    % Documented in docs/paper-gaps/mpu_canonical_form_full_support.tex.
    \uses{prop:mpu_admissible_local_time_reversal, prop:mpu_admissible_sigma_blocking_parity, def:mpu_reduced_full_support_source_datum, def:mpu_admissible_equivalence_datum}
    Let an MPU in admissible standard form be invariant under time reversal.
    Suppose $\mathcal U$, $\mathcal U_2$, and $\mathcal U_4$, together with the
    physical-adjoint comparison tensors $\mathcal U_2^\sharp$ and
    $\mathcal U_4^\sharp$, are simple and each carries its own reduced
    full-support source datum. Then the MPS sign $G\overline G=\pm\Id$ equals
    the sign $\sigma^{(2)}$.
    % Source lines: paper_v2.tex 1807--1829.
\end{proposition}

\subsection{Symmetry examples}

\begin{proposition}[All phases under conjugation]
    \label{ex:all-phases-under-conjugation}
    \notready
    \discussion{5715}
    \uses{conjphs, prop:conjsym-strict-equiv, corconj, cor:sympl-to-orthogonal}
    Let $u_+=v_+=\Id_{r\ell}$ and
    $u_-=v_-=\operatorname{diag}(-1,1,\ldots,1)$. If at least one of $r$ and
    $\ell$ is odd, $(u_+,v_+)$ and $(u_+,v_-)$ realize the two classes of
    Theorem~\ref{conjphs}. If both are even, the four pairs $(u_\pm,v_\pm)$
    realize the four Case I classes. Multiplying by the matrix $R$ of
    Corollary~\ref{cor:sympl-to-orthogonal} realizes the four Case II classes.
    % Source lines: paper_v2.tex 1845--1876.
\end{proposition}

\begin{proposition}[Admissible phases under conjugation]
    \label{ex:mpu_admissible_all_conjugation_phases}
    \notready
    \discussion{5715}
    % **Scope restriction (realization by admissible standard forms):** each
    % displayed standard-form tensor carries reduced full-support source data.
    % Documented in docs/paper-gaps/mpu_canonical_form_full_support.tex.
    \uses{thm:mpu_admissible_conjugation_odd_rank, thm:mpu_admissible_conjugation_even_rank, prop:mpu_admissible_conjugation_standard_forms, cor:mpu_admissible_symplectic_to_orthogonal, def:mpu_reduced_full_support_source_datum, def:mpu_admissible_equivalence_datum}
    Assume the standard-form tensors displayed below carry reduced full-support
    source data and that the comparison paths used to classify them carry
    admissible path data of blocking length one. Let
    \begin{align}
        u_+ & =v_+=\Id_{r\ell},                        &
        u_- & =v_-=\operatorname{diag}(-1,1,\ldots,1).
    \end{align}
    If at least one of $r$ and $\ell$ is odd, the pairs $(u_+,v_+)$ and
    $(u_+,v_-)$ represent the two classes of Theorem~\ref{thm:mpu_admissible_conjugation_odd_rank}. If $r$
    and $\ell$ are both even, the four pairs $(u_\pm,v_\pm)$ represent the
    four Case I classes of Theorem~\ref{thm:mpu_admissible_conjugation_even_rank}. With
    \begin{align}
        u'_\pm & =(\Id_{r\ell/4}\otimes R)u_\pm, &
        v'_\pm & =v_\pm(\Id_{r\ell/4}\otimes R),
    \end{align}
    the four pairs $(u'_\pm,v'_\pm)$ represent the four Case II classes.
    % Source lines: paper_v2.tex 1845--1876.
\end{proposition}

\begin{proposition}[CZX]
    \label{example:czx}
    \notready
    \discussion{5715}
    \uses{prop:sigma1-and-sigma2-equal, prop:dagger-mps-class-and-mpu-class-equal}
    The CZX MPU is defined by
    \begin{align}
        u & =CZ(\sigma_x\otimes\sigma_x), & v & =CZ,
    \end{align}
    where $CZ=\operatorname{diag}(1,1,1,-1)$. Let
    \begin{align}
        H & =\frac1{\sqrt2}\begin{pmatrix}1&1\\1&-1\end{pmatrix}
    \end{align}
    be the Hadamard matrix. The corresponding tensor is
    \begin{align}
        \mathcal U_{\ell r}^{ud}
         & =(\sigma_x)_{ud}\delta_{u,\ell}H_{\ell,r}.
    \end{align}
    It has $\sigma^{(1)}=\sigma^{(2)}=-1$.
    % Source lines: paper_v2.tex 1906--1943.
\end{proposition}

\begin{proposition}[Admissible CZX data]
    \label{ex:mpu_admissible_czx}
    \notready
    \discussion{5715}
    % **Scope restriction (blocked CZX data):** the one-, two-, and four-site
    % tensors used to define the two signs carry reduced full-support source data.
    % Documented in docs/paper-gaps/mpu_canonical_form_full_support.tex.
    \uses{prop:mpu_admissible_sigma_blocking_parity, prop:mpu_admissible_mps_time_reversal_class, def:mpu_reduced_full_support_source_datum, def:mpu_admissible_equivalence_datum}
    Suppose the CZX tensors $\mathcal U$, $\mathcal U_2$, and $\mathcal U_4$,
    together with the physical-adjoint comparison tensors
    $\mathcal U_2^\sharp$ and $\mathcal U_4^\sharp$, are simple and each carries
    its own reduced full-support source datum. The CZX MPU is defined by
    \begin{align}
        u  & =CZ(\sigma_x\otimes\sigma_x),   & v & =CZ, \\
        CZ & =\operatorname{diag}(1,1,1,-1).
    \end{align}
    Its tensor is
    \begin{align}
        \mathcal U_{\ell r}^{ud}
          & =(\sigma_x)_{ud} \delta_{u,\ell}H_{\ell,r},           &
        H & =\frac1{\sqrt2}\begin{pmatrix}1&1\\1&-1\end{pmatrix}.
    \end{align}
    It has $\sigma^{(1)}=\sigma^{(2)}=-1$.
    % Source lines: paper_v2.tex 1906--1943.
\end{proposition}

\begin{proposition}[The other three time-reversal classes]
    \label{ex:dagger-all-but-czx}
    \notready
    \discussion{5715}
    \uses{example:czx, prop:sigma1-and-sigma2-equal}
    The standard forms
    \begin{align}
        u & =v=\Id,                                      \\
        u & =i\Id,                           & v & =\Id, \\
        u & =i\,CZ(\sigma_x\otimes\sigma_x), & v & =CZ
    \end{align}
    realize, respectively,
    \begin{align}
        (\sigma^{(1)},\sigma^{(2)})
         & =(1,1),\quad(-1,1),\quad(1,-1).
    \end{align}
    % Source lines: paper_v2.tex 1945--1962.
\end{proposition}

\begin{proposition}[Admissible remaining time-reversal classes]
    \label{ex:mpu_admissible_time_reversal_classes}
    \notready
    \discussion{5715}
    % **Scope restriction (blocked example data):** every standard form and
    % blocking used to define the signs carries reduced full-support source data.
    % Documented in docs/paper-gaps/mpu_canonical_form_full_support.tex.
    \uses{ex:mpu_admissible_czx, prop:mpu_admissible_sigma_blocking_parity, def:mpu_reduced_full_support_source_datum, def:mpu_admissible_equivalence_datum}
    For each displayed example, suppose the tensors $\mathcal U$,
    $\mathcal U_2$, and $\mathcal U_4$, together with the physical-adjoint
    comparison tensors $\mathcal U_2^\sharp$ and $\mathcal U_4^\sharp$, are
    simple and each carries its own reduced full-support source datum.
    The standard forms
    \begin{align}
        u & =v=\Id,                                      \\
        u & =i\Id,                           & v & =\Id, \\
        u & =i\,CZ(\sigma_x\otimes\sigma_x), & v & =CZ
    \end{align}
    realize, respectively,
    \begin{align}
        (\sigma^{(1)},\sigma^{(2)})
         & =(1,1),\quad(-1,1),\quad(1,-1).
    \end{align}
    % Source lines: paper_v2.tex 1945--1962.
\end{proposition}

\begin{definition}[Right- and left-shift tensors]
    \label{def:mpu_shift_tensors}
    \lean{MPOTensor.rightShiftTensor}
    \lean{MPOTensor.leftShiftTensor}
    \leanok
    \uses{def:mpdo_physical_adjoint}
    The right shift on $(\mathbb C^d)^{\otimes N}$ is
    \begin{align}
        T^{(N)}\lvert n_1,\ldots,n_N\rangle
         & =\lvert n_N,n_1,\ldots,n_{N-1}\rangle.
    \end{align}
    It has bond dimension $d$ and local matrices
    $A^{ij}=\lvert i\rangle\langle j\rvert$.
    The physical adjoint of $A$ defines the left-shift tensor.
    % Source lines: paper_v2.tex 1980--1987.
\end{definition}

\begin{theorem}[Shift tensor formulas]
    \label{thm:mpu_shift_tensor_formulas}
    \lean{MPOTensor.mpo_rightShiftTensor_apply}
    \lean{MPOTensor.mpo_rightShiftTensor}
    \lean{MPOTensor.mpo_leftShiftTensor}
    \lean{MPOTensor.rightShiftTensor_isMPU}
    \lean{MPOTensor.leftShiftTensor_isMPU}
    \leanok
    \uses{def:mpu_shift_tensors, def:mpu, thm:mpu_physical_adjoint_identities}
    For every $N>0$, the tensors in Definition~\ref{def:mpu_shift_tensors}
    generate $T^{(N)}$ and $T^{(N)\dagger}$, respectively.  Both tensors are
    matrix product unitaries.
\end{theorem}

\begin{proof}\leanok
    \uses{def:mpo_eval_word, def:mpo_rotate_config, lem:eval_word_of_fn_prod, thm:trace_eval_word_eq_sum_cyclic, def:mpu, thm:mpu_physical_adjoint_identities}
    Closing the virtual indices of $A^{ij}=\lvert i\rangle\langle j\rvert$
    leaves one nonzero virtual configuration.  If
    $\tau_n=\sigma_{n+1}$ for every $n$ modulo $N$, then
    \begin{align}
        \langle\sigma\vert T^{(N)}\vert\tau\rangle & =1.
    \end{align}
    Otherwise the same matrix entry is zero.  Hence the output configuration is
    $(\tau_{N-1},\tau_0,\ldots,\tau_{N-2})$, which fixes the permutation-matrix
    convention and gives the right shift.  Its permutation matrix is unitary.
    Physical adjunction gives the conjugate transpose, hence the left shift and
    its unitarity.
\end{proof}

\begin{theorem}[Source cuts and ranks of the cyclic shifts]
    \label{thm:mpu_shift_source_cuts_ranks}
    \lean{MPOTensor.sourceCutM₁_rightShiftTensor}
    \lean{MPOTensor.sourceCutM₂_rightShiftTensor}
    \lean{MPOTensor.rightRank_rightShiftTensor}
    \lean{MPOTensor.leftRank_rightShiftTensor}
    \lean{MPOTensor.sourceCutM₁_leftShiftTensor}
    \lean{MPOTensor.sourceCutM₂_leftShiftTensor}
    \lean{MPOTensor.rightRank_leftShiftTensor}
    \lean{MPOTensor.leftRank_leftShiftTensor}
    \leanok
    \uses{def:mpu_shift_tensors, def:mpu_source_matrix_one, def:mpu_source_matrix_two, def:mpu_right_rank, def:mpu_left_rank}
    Let $A$ be the right-shift tensor and let
    $z_{(\alpha,i)}=\delta_{\alpha i}$ be the vectorized identity. Its raw
    source cuts are
    \begin{align}
        M_1(A) & =\Id,            \\
        M_2(A) & =zz^{\mathsf T}.
    \end{align}
    For the left-shift tensor $A^\sharp$ they are
    \begin{align}
        M_1(A^\sharp) & =zz^{\mathsf T}, \\
        M_2(A^\sharp) & =\Id.
    \end{align}
    In every physical dimension, including $d=0$,
    \begin{align}
        r(A)           & =d^2, \\
        \ell(A^\sharp) & =d^2.
    \end{align}
    If $d>0$, the remaining two ranks are
    \begin{align}
        \ell(A)     & =1, \\
        r(A^\sharp) & =1.
    \end{align}
    These are the raw source-cut quantities from
    \cite[Section~III]{Cirac2017MPU}; no choice of source factors is involved.
\end{theorem}

\begin{proof}\leanok
    \uses{def:mpu_shift_tensors, def:mpu_source_matrix_one, def:mpu_source_matrix_two, def:mpu_right_rank, def:mpu_left_rank, thm:mpu_physical_adjoint_source_cuts, thm:mpu_physical_adjoint_source_ranks}
    The entry formula $A^{ij}_{\alpha\beta}
        =\delta_{i\alpha}\delta_{j\beta}$ gives $M_1(A)=\Id$ and
    $M_2(A)=zz^{\mathsf T}$. Hence the first rank is $d^2$. For $d>0$, the
    vector $z$ has a nonzero diagonal coordinate, so its outer product has
    rank one. Physical adjunction exchanges the two cuts and their ranks,
    which gives all four left-shift formulas without a second rank
    calculation.
\end{proof}

\begin{theorem}[Unitarity of the swap operator]
    \label{thm:mpu_swap_matrix_unitary}
    \lean{Matrix.swapMatrix_isUnitaryBetween}
    \leanok
    Let $d\in\N$. On $\C^d\otimes\C^d$, let $\mathbb S$ be the swap operator
    used in \cite[Section~VII]{Cirac2017MPU}, defined on the standard basis by
    $\mathbb S(e_k\otimes e_\ell)=e_\ell\otimes e_k$. With columns indexed by
    source coordinates $(k,\ell)$ and rows indexed by target coordinates
    $(i,j)$, its entries are
    \begin{align}
        \mathbb S_{(i,j),(k,\ell)} & =\delta_{i,\ell}\delta_{j,k}.
    \end{align}
    Then $\mathbb S$ is unitary between the two product-coordinate orders:
    \begin{align}
        \mathbb S^\dagger\mathbb S & =\Id, \\
        \mathbb S\mathbb S^\dagger & =\Id.
    \end{align}
    % Operator characterization: paper_v2.tex lines 1316--1320.
    % Section VII uses: lines 1999--2034, equations eq:SF_u1_u3,
    % eq:uv2_U2, and eq:uv2_U3; the unitary interpolation has endpoint
    % S at lines 2148--2150.
\end{theorem}

\begin{proof}\leanok
    The swap operator is self-adjoint and involutive, so
    $\mathbb S^\dagger=\mathbb S$ and $\mathbb S^2=\Id$. Substitution gives
    both identities.
\end{proof}

\begin{definition}[Three topological-insulator shift MPUs]
    \label{threeMPU}
    \lean{MPOTensor.identityMPUTensor}
    \lean{MPOTensor.shiftExampleU₁}
    \lean{MPOTensor.shiftExampleU₂}
    \lean{MPOTensor.shiftExampleU₃}
    \leanok
    \uses{def:mpu_shift_tensors, def:mpu_independent_tensor_product, def:mpo_stacked_layers}
    On two $d$-dimensional spins per site, define
    \begin{align}
        U_1^{(N)} & =(\Id\otimes\Id)^{\otimes N},   \\
        U_2^{(N)} & =T^{(N)\dagger}\otimes T^{(N)}, \\
        U_3^{(N)} & =T^{(N)}\otimes T^{(N)\dagger}.
    \end{align}
    % Source lines: paper_v2.tex 1988--1995.
\end{definition}

\begin{definition}[Trace-normalized shift weight]
    \label{def:mpu_shift_trace_weight}
    \lean{MPOTensor.shiftPaperWeight,
        MPOTensor.shiftPaperWeightSquared,
        MPOTensor.shiftPaperWeightSquared_eq}
    \leanok
    \uses{def:mpu_source_weight, def:mpu_shift_tensors, threeMPU}
    Canonical form II requires the positive right fixed point $\rho$ to have
    trace one. If the virtual weight on the bond space of a single shift is
    chosen to be a scalar identity, this normalization forces
    \begin{align}
        \rho_d & =d^{-1}\Id_d
    \end{align}
    whenever $d>0$. The same formula is defined at $d=0$, although it then has
    trace zero. The tensors $U_2$ and $U_3$ have bond dimension $d^2$, so the
    corresponding scalar trace-one matrix is
    \begin{align}
        \rho_{d^2} & =(d^2)^{-1}\Id_{d^2}.
    \end{align}
    These are the scalar identity choices compatible with the source weight in
    \cite[equation labelled~\texttt{Y1Y1X1X1}]{Cirac2017MPU}. No transfer-map
    fixed-point equation is asserted here.

    The supplied source-factor witnesses below instead use
    $\rho_{\mathrm{aux}}=\Id_d$ for each constituent shift and
    $\rho_{\mathrm{aux}}=\Id_{d^2}$ after tensoring. The paper gates computed
    from them carry reciprocal factors $d$ and $d^{-1}$ because these identity
    weights differ from $\rho_d$ and $\rho_{d^2}$ when $d>1$.
    % Source normalization: paper_v2.tex 269--280.
    % Source-factor weight: paper_v2.tex 487--495.
    % Shift and tensor-product dimensions: paper_v2.tex 1980--1994.
\end{definition}

\begin{theorem}[Positivity and normalization of the shift weight]
    \label{thm:mpu_shift_trace_weight_normalized}
    \lean{MPOTensor.shiftPaperWeight_posSemidef,
        MPOTensor.shiftPaperWeight_posDef,
        MPOTensor.shiftPaperWeight_trace,
        MPOTensor.shiftPaperWeightSquared_posSemidef,
        MPOTensor.shiftPaperWeightSquared_posDef,
        MPOTensor.shiftPaperWeightSquared_trace}
    \leanok
    \uses{def:mpu_shift_trace_weight}
    For every $d\in\mathbb N$, the matrix $\rho_d$ is positive semidefinite. If
    $d>0$, then $\rho_d$ is positive definite and
    \begin{align}
        \tr(\rho_d) & =1.
    \end{align}
    Under the same hypothesis, $\rho_{d^2}$ is positive definite with trace one;
    it is positive semidefinite for every $d$.
\end{theorem}

\begin{proof}\leanok
    \uses{def:mpu_shift_trace_weight}
    The scalar $d^{-1}$ is nonnegative, and it is strictly positive when
    $d>0$. Hence the claims follow from positivity of the identity matrix.
    For positive $d$,
    \begin{align}
        \tr(\rho_d) & =d^{-1}\tr(\Id_d)=d^{-1}d=1.
    \end{align}
    Replacing $d$ by $d^2$ proves the product-space statements.
\end{proof}

\begin{definition}[Trace-normalized source factors for the right shift]
    \label{def:mpu_right_shift_paper_source_factors}
    \lean{MPOTensor.rightShiftPaperSourceFactors}
    \leanok
    \uses{def:mpu_shift_trace_weight, def:threeMPU_supplied_source_factors}
    Assume $d>0$, and let $T$ be the right-shift tensor. Starting from the
    displayed source factors for $T$ with identity weight, rescale the first
    cut by
    \begin{align}
        X_1' & =\sqrt d\,X_1,
             & Y_1'           & =d^{-1/2}Y_1,
             & Z_1'           & =\sqrt d\,Z_1,
    \end{align}
    and leave $X_2,Y_2,Z_2$ unchanged. Then these are supplied source factors
    for the trace-normalized scalar weight $\rho_d=d^{-1}\Id_d$ prescribed by
    CFII. In particular,
    \begin{align}
        M_i & =X_i'Y_i',
            & (X_1')^\dagger(\Id_d\otimes\rho_d)X_1' & =\Id,
            & (X_2')^\dagger X_2'                    & =\Id,
            & Y_i'Z_i'                               & =\Id.
    \end{align}
    This is the normalization of
    \cite[equations~\texttt{Erightleft}, \texttt{X1X2b}, and
        \texttt{YZ=1}]{Cirac2017MPU} for the right shift. This definition
    concerns the source-factor metric; it does not assert a transfer
    fixed-point equation.
\end{definition}

\begin{theorem}[The paper source tensor $u$ for the right shift]
    \label{thm:mpu_right_shift_paper_source_u}
    \lean{MPOTensor.sourceU_rightShiftPaperSourceFactors_apply,
        MPOTensor.sourceU_rightShiftPaperSourceFactors_gram,
        MPOTensor.sourceU_rightShiftPaperSourceFactors_isIsometry}
    \leanok
    \uses{def:mpu_right_shift_paper_source_factors, def:mpu_source_uv,
        thm:mpu_shift_source_cuts_ranks}
    Assume $d>0$. Identify the unique left source coordinate with $0$ and the
    right source coordinate with $(a,b)$. For the trace-normalized source
    factors of the right shift,
    \begin{align}
        u_{(0,(a,b)),(i_1,i_2)}
         & =\delta_{a,i_1}\delta_{b,i_2}.
    \end{align}
    Hence
    \begin{align}
        \sum_{l,r}u_{(l,r),q}\overline{u_{(l,r),p}}
         & =\delta_{p,q},
         & u^\dagger u    & =\Id.
    \end{align}
    This fixes the order and normalization of the two triangles in
    \cite[equations~\texttt{uu} and~\texttt{uUnitary}]{Cirac2017MPU} for this
    basic MPU. It is the explicit normalization and orientation regression for
    Theorem~\ref{thm:mpu_source_u_retained_interior} and
    Lemma~\ref{lem:mpu_admissible_source_u_isometry}.
\end{theorem}

\begin{proof}\leanok
    \uses{def:mpu_right_shift_paper_source_factors, def:mpu_source_uv}
    For the identity-weight factors, direct substitution in
    $u=Y_2\mathbin{-}Y_1$ gives
    \begin{align}
        u^{(0)}_{(0,(a,b)),(i_1,i_2)}
         & =d^{1/2}\delta_{a,i_1}\delta_{b,i_2}.
    \end{align}
    The trace-normalized first cut replaces $Y_1$ by $d^{-1/2}Y_1$ and leaves
    $Y_2$ unchanged. Therefore
    \begin{align}
        u_{(0,(a,b)),(i_1,i_2)}
         & =d^{-1/2}u^{(0)}_{(0,(a,b)),(i_1,i_2)}
        =\delta_{a,i_1}\delta_{b,i_2}.
    \end{align}
    Summing the product of two such entries over $(a,b)$ gives
    $\delta_{p,q}$, which is the entrywise identity $u^\dagger u=\Id$.
\end{proof}

\begin{theorem}[Local physical-species swap for the shift MPUs]
    \label{thm:threeMPU_local_physical_swap}
    \lean{MPOTensor.shiftExampleU₃_physicalSwap_eq_swapMatrix_mul_shiftExampleU₂_mul_swapMatrix}
    \leanok
    \uses{threeMPU, def:mpu_shift_tensors, def:mpu_independent_tensor_product}
    Write $[U_a]^{(i_1,i_2),(j_1,j_2)}$ for the virtual matrix of the
    one-site tensor $U_a$ with physical output coordinate $(i_1,i_2)$ and
    input coordinate $(j_1,j_2)$. Let $\mathbb S$ be the SWAP matrix on the
    two virtual factors, defined by
    \begin{align}
        \mathbb S\lvert\alpha_1,\alpha_2\rangle
         & =\lvert\alpha_2,\alpha_1\rangle.
    \end{align}
    Then
    \begin{align}
        [U_3]^{(i_2,i_1),(j_2,j_1)}
         & =\mathbb S[U_2]^{(i_1,i_2),(j_1,j_2)}\mathbb S.
    \end{align}
    Thus exchanging the two physical species exchanges the local tensors
    $U_2$ and $U_3$, with the same exchange on their virtual factors.
    % Local refinement of the finite-chain statements at paper_v2.tex
    % 1999--2001 and 2240--2245.
\end{theorem}

\begin{proof}\leanok
    \uses{threeMPU, def:mpu_shift_tensors, def:mpu_independent_tensor_product}
    In virtual coordinates $(\alpha_1,\alpha_2)$ and
    $(\beta_1,\beta_2)$, multiplication by $\mathbb S$ on both sides selects
    the entry with row $(\alpha_2,\alpha_1)$ and column
    $(\beta_2,\beta_1)$. Consequently,
    \begin{align}
        [U_3]^{(i_2,i_1),(j_2,j_1)}_{(\alpha_1,\alpha_2),(\beta_1,\beta_2)}
         & =T^{i_2j_2}_{\alpha_1\beta_1}
        (T^\dagger)^{i_1j_1}_{\alpha_2\beta_2}                                   \\
         & =(T^\dagger)^{i_1j_1}_{\alpha_2\beta_2}
        T^{i_2j_2}_{\alpha_1\beta_1}                                             \\
         & =[U_2]^{(i_1,i_2),(j_1,j_2)}_{(\alpha_2,\alpha_1),(\beta_2,\beta_1)}.
    \end{align}
    This is exactly the entrywise form of conjugation by $\mathbb S$.
\end{proof}

\begin{theorem}[Finite-chain formulas for the three shift MPUs]
    \label{thm:threeMPU_chain_formulas}
    \lean{MPOTensor.mpo_identityMPUTensor}
    \lean{MPOTensor.identityMPUTensor_isMPU}
    \lean{MPOTensor.mpo_shiftExampleU₁}
    \lean{MPOTensor.mpo_shiftExampleU₂}
    \lean{MPOTensor.mpo_shiftExampleU₃}
    \lean{MPOTensor.shiftExampleU₁_isMPU}
    \lean{MPOTensor.shiftExampleU₂_isMPU}
    \lean{MPOTensor.shiftExampleU₃_isMPU}
    \leanok
    \uses{threeMPU, thm:mpu_shift_tensor_formulas, thm:mpu_finite_tensor_product, thm:mpu_tensor_product, def:mpu}
    The three tensors generate the operators displayed in
    Definition~\ref{threeMPU} on every nonempty chain, and each tensor is a
    matrix product unitary.
\end{theorem}

\begin{proof}\leanok
    \uses{def:mpo_stacked_layers, thm:mpu_finite_tensor_product, thm:mpu_shift_tensor_formulas, thm:mpu_tensor_product, def:mpu}
    The bond-one identity tensor generates the identity at every length.  The
    finite-chain operator of an independent tensor product is the Kronecker
    product of the two finite-chain operators, after the sitewise identification
    $(\mathbb C^d\otimes\mathbb C^d)^{\otimes N}
        \cong(\mathbb C^d)^{\otimes N}\otimes(\mathbb C^d)^{\otimes N}$.  Applying
    this identity to the right and left shifts gives the three displayed
    formulas.  The identity tensor is unitary, and tensor products preserve
    unitarity, so all three tensors are matrix product unitaries.
\end{proof}

\begin{theorem}[Specified-tensor source-index values of the shifts]
    \label{thm:mpu_shift_specified_source_index_values}
    \lean{MPOTensor.sourceIndexValue_rightShiftTensor,
        MPOTensor.sourceIndexValue_leftShiftTensor,
        MPOTensor.sourceIndexValue_identityMPUTensor,
        MPOTensor.sourceIndexValue_shiftExampleU₁,
        MPOTensor.sourceIndexValue_shiftExampleU₂,
        MPOTensor.sourceIndexValue_shiftExampleU₃}
    \leanok
    \uses{def:mpu_specified_source_index_value,
        thm:mpu_shift_source_cuts_ranks,
        thm:mpu_tensor_product_source_ranks, threeMPU}
    Assume $d>0$. For the displayed right- and left-shift tensors $T$ and
    $T^\sharp$, the specified-tensor values are
    \begin{align}
        \operatorname{ind}_{\mathrm{src}}(T) & =\log_2 d.
    \end{align}
    Likewise,
    \begin{align}
        \operatorname{ind}_{\mathrm{src}}(T^\sharp) & =-\log_2 d.
    \end{align}
    The bond-one identity tensor and the three tensors $U_1,U_2,U_3$ satisfy
    \begin{align}
        \operatorname{ind}_{\mathrm{src}}(\Id)
        =\operatorname{ind}_{\mathrm{src}}(U_1)
        =\operatorname{ind}_{\mathrm{src}}(U_2)
        =\operatorname{ind}_{\mathrm{src}}(U_3) & =0.
    \end{align}
    Here $\operatorname{ind}_{\mathrm{src}}$ is the specified-tensor quantity
    from Definition~\ref{def:mpu_specified_source_index_value}, evaluated on the
    displayed unblocked tensors. It is not the public MPU index
    $\operatorname{ind}$ of Definition~\ref{def:index}, whose definition uses a
    simple blocking and requires independence of that choice.
    % **Scope restriction (specified tensors):** the source states the public
    % blocking-independent MPU index, while this theorem computes only the raw
    % value of each displayed tensor. Documented in
    % docs/paper-gaps/mpu_shift_specified_tensor_index_scope.tex; elimination is
    % tracked by issues #7011 and #5856.
    % Source: paper_v2.tex lines 2037--2041.
\end{theorem}

\begin{proof}\leanok
    \uses{def:mpu_specified_source_index_value,
        thm:mpu_shift_source_cuts_ranks,
        thm:mpu_tensor_product_source_ranks, threeMPU}
    The right shift has source ranks $(d^2,1)$, while the left shift has
    $(1,d^2)$. Substitution into the source-index expression gives the first
    two formulas. The identity tensor has equal source ranks $(d,d)$. Source
    ranks multiply under the independent tensor product, so each of
    $U_1,U_2,U_3$ has equal right and left rank $d^2$. Their logarithmic
    differences therefore vanish.
\end{proof}

\begin{definition}[Supplied source factors for the three shift MPUs]
    \label{def:threeMPU_supplied_source_factors}
    \lean{MPOTensor.identitySourceFactors,
        MPOTensor.rightShiftSourceFactors,
        MPOTensor.leftShiftSourceFactors,
        MPOTensor.SourceFactors.independentTensorProductOfIdentityWeight,
        MPOTensor.identityTensorIdentityMatrix,
        MPOTensor.identityTensorIdentityMatrix_apply,
        MPOTensor.swapTensorSwapMatrix_mul_identitySwapIdentityMatrix_apply,
        MPOTensor.identitySwapIdentityMatrix_mul_swapTensorSwapMatrix_apply,
        MPOTensor.shiftExampleU₁LeftRankEquiv,
        MPOTensor.shiftExampleU₁RightRankEquiv,
        MPOTensor.shiftExampleU₂LeftRankEquiv,
        MPOTensor.shiftExampleU₂RightRankEquiv,
        MPOTensor.shiftExampleU₃RightRankEquiv_apply,
        MPOTensor.shiftExampleU₁SourceFactors,
        MPOTensor.shiftExampleU₂SourceFactors,
        MPOTensor.shiftExampleU₃SourceFactors}
    \leanok
    \uses{threeMPU, def:mpu_weighted_source_factors, thm:mpu_shift_source_cuts_ranks}
    Assume $d>0$. Equip $U_1$ with the tensor product of the two bond-one
    identity source factorizations.  Equip $U_2$ with the tensor product of the
    left-shift and right-shift source factorizations, in that order, and equip
    $U_3$ with the tensor product in the reverse order.  In all three cases the
    virtual weight is the identity.  These are explicit supplied source-factor
    witnesses; no identification with an arbitrary compact singular-value
    decomposition is asserted.
\end{definition}

\begin{theorem}[Unitarity of the four-spin shift matrices]
    \label{thm:threeMPU_four_spin_matrix_unitarity}
    \lean{MPOTensor.identitySwapIdentityMatrix_isUnitaryBetween,
        MPOTensor.swapTensorSwapMatrix_isUnitaryBetween,
        MPOTensor.swapTensorSwapMatrix_mul_identitySwapIdentityMatrix_isUnitaryBetween,
        MPOTensor.identitySwapIdentityMatrix_mul_swapTensorSwapMatrix_isUnitaryBetween}
    \leanok
    Let $\mathbb S$ exchange two $d$-dimensional spins. The two swap factors
    appearing in the blocked forms of $U_2$ and $U_3$, and both ordered
    products, are unitary:
    \begin{align}
        \Id\otimes\mathbb S\otimes\Id,\quad
        \mathbb S\otimes\mathbb S,\quad
        (\mathbb S\otimes\mathbb S)(\Id\otimes\mathbb S\otimes\Id),\quad
        (\Id\otimes\mathbb S\otimes\Id)(\mathbb S\otimes\mathbb S).
    \end{align}
    % Source lines: paper_v2.tex 2018--2034.
\end{theorem}

\begin{proof}\leanok
    The swap satisfies $\mathbb S^\dagger=\mathbb S$ and
    $\mathbb S^2=\Id$. Kronecker products preserve isometries, and a square
    isometry is unitary. Thus both factors are unitary. Closure of the unitary
    matrices under multiplication gives the two ordered products.
\end{proof}

\begin{theorem}[Supplied paper gates for $U_1$]
    \label{thm:threeMPU_supplied_source_gates_one}
    \lean{MPOTensor.shiftExampleU₁SourceURowEquiv,
        MPOTensor.shiftExampleU₁SourceVColumnEquiv,
        MPOTensor.shiftExampleU₁_sourceU_fourSpin_apply,
        MPOTensor.shiftExampleU₁_sourceV_fourSpin_apply}
    \leanok
    \uses{def:threeMPU_supplied_source_factors, def:mpu_source_uv}
    In the four-spin coordinates of the supplied source factorization,
    \begin{align}
        u_1 & =\Id\otimes\Id,
            & v_1             & =\Id\otimes\Id.
    \end{align}
\end{theorem}

\begin{proof}\leanok
    \uses{def:threeMPU_supplied_source_factors, def:mpu_source_uv}
    The tensor-product formulas for $u=Y_2\mathbin{-}Y_1$ and
    $v=X_1\mathbin{-}X_2$ reduce each entry to the product of two primitive
    identity entries.
\end{proof}

\begin{theorem}[Paper gates in the two-site factorizations of $U_2$]
    \label{thm:threeMPU_U2_supplied_source_gates_two}
    \lean{MPOTensor.shiftExampleU₂SourceURowEquiv,
        MPOTensor.shiftExampleU₂SourceVColumnEquiv,
        MPOTensor.shiftExampleU₂_sourceU_fourSpin_apply,
        MPOTensor.shiftExampleU₂_sourceV_fourSpin_apply,
        MPOTensor.shiftExampleU₂_blockTwo_apply_eq_sum_X₂_mul_sourceU_mul_X₁,
        MPOTensor.shiftExampleU₂_blockTwo_apply_eq_sum_sourceV_mul_Y₁_mul_Y₂}
    \leanok
    \uses{def:threeMPU_supplied_source_factors, def:mpu_source_uv,
        thm:mpu_source_v_open_two_site, thm:mpu_source_u_open_two_site}
    Assume $d>0$. In the four-spin coordinates of the paper, the supplied
    source gates satisfy
    \begin{align}
        u_2^{(2)}    & =d\,(\Id\otimes\mathbb S\otimes\Id), \\
        d\,v_2^{(2)} & =(\mathbb S\otimes\mathbb S)
        (\Id\otimes\mathbb S\otimes\Id).
    \end{align}
    The blocked tensor factors through $X_2,u_2^{(2)},X_1$ and, equivalently,
    through $v_2^{(2)},Y_1,Y_2$. The scalar factors record the normalization of
    the supplied source-factor witnesses.
\end{theorem}

\begin{proof}\leanok
    \uses{def:threeMPU_supplied_source_factors, def:mpu_source_uv,
        thm:mpu_source_v_open_two_site, thm:mpu_source_u_open_two_site}
    The tensor-product formulas for the paper gates reduce their entries to
    products of the left- and right-shift source gates. The Kronecker deltas
    give the two displayed swap matrices, with the factors $d$ and $d^{-1}$
    fixed by the identity-weight normalization. Substitute these evaluations
    into the two exact open source-gate factorizations.
\end{proof}

\begin{theorem}[Paper gates in the two-site factorizations of $U_3$]
    \label{thm:threeMPU_U3_supplied_source_gates_two}
    \lean{MPOTensor.shiftExampleU₃SourceURowEquiv,
        MPOTensor.shiftExampleU₃SourceVColumnEquiv,
        MPOTensor.shiftExampleU₃_sourceU_fourSpin_apply,
        MPOTensor.shiftExampleU₃_sourceV_fourSpin_apply,
        MPOTensor.shiftExampleU₃_blockTwo_apply_eq_sum_X₂_mul_sourceU_mul_X₁,
        MPOTensor.shiftExampleU₃_blockTwo_apply_eq_sum_sourceV_mul_Y₁_mul_Y₂}
    \leanok
    \uses{def:threeMPU_supplied_source_factors, def:mpu_source_uv,
        thm:mpu_source_v_open_two_site, thm:mpu_source_u_open_two_site}
    Assume $d>0$. In the four-spin coordinates of the paper, the supplied
    source gates satisfy
    \begin{align}
        u_3^{(2)}    & =d\,(\Id\otimes\mathbb S\otimes\Id)
        (\mathbb S\otimes\mathbb S),                       \\
        d\,v_3^{(2)} & =\Id\otimes\mathbb S\otimes\Id.
    \end{align}
    The blocked tensor factors through $X_2,u_3^{(2)},X_1$ and, equivalently,
    through $v_3^{(2)},Y_1,Y_2$. The scalar factors record the normalization of
    the supplied source-factor witnesses.
\end{theorem}

\begin{proof}\leanok
    \uses{def:threeMPU_supplied_source_factors, def:mpu_source_uv,
        thm:mpu_source_v_open_two_site, thm:mpu_source_u_open_two_site}
    Reverse the two primitive shift factors in the tensor-product formulas for
    $u$ and $v$. Their Kronecker deltas give the displayed ordered swap products,
    with the factors $d$ and $d^{-1}$ fixed by the identity-weight normalization.
    Substitute these evaluations into the two exact open source-gate
    factorizations.
\end{proof}

\begin{proposition}[Standard forms of the three shift MPUs]
    \label{prop:threeMPU_standard_forms}
    \notready
    \discussion{5715}
    \uses{threeMPU, def:mpu_source_uv}
    Assume $d>0$. For the fixed standard-form source factors, the corresponding
    unitaries satisfy
    \begin{align}
        u_1                          & =\Id\otimes\Id,                  & v_1 & =\Id\otimes\Id, \\
        u_3                          & =\mathbb S,                      & v_3 & =\mathbb S,     \\
        u_2^{(2)}                    & =\Id\otimes\mathbb S\otimes\Id,  &
        v_2^{(2)}                    & =(\mathbb S\otimes\mathbb S)
        (\Id\otimes\mathbb S\otimes\Id),                                                        \\
        u_3^{(2)}                    & =(\Id\otimes\mathbb S\otimes\Id)
        (\mathbb S\otimes\mathbb S), &
        v_3^{(2)}                    & =\Id\otimes\mathbb S\otimes\Id.
    \end{align}
    These statements concern the paper gates $u=Y_2\mathbin{-}Y_1$ and
    $v=X_1\mathbin{-}X_2$ for trace-normalized standard-form factors. The
    supplied identity-weight computations above give the same gate matrices
    only after transferring reciprocal scalar factors between $u$ and $v$.
    Constructing the trace-normalized factors and proving the displayed
    equalities remains.
    % Source lines: paper_v2.tex 2002--2042.
\end{proposition}

\begin{definition}[Transformation with swaps]
    \label{threeMPU2}
    \lean{MPOTensor.shiftPhysicalSwap}
    \lean{MPOTensor.shiftPhysicalSwap_apply}
    \lean{MPOTensor.shiftExampleTildeU₁}
    \lean{MPOTensor.shiftExampleTildeU₂}
    \lean{MPOTensor.shiftExampleTildeU₃}
    \leanok
    \uses{threeMPU, def:mpu_physical_pair_swap}
    Let $\mathbb S$ exchange the two $d$-dimensional spins at one site, with
    matrix entries
    \begin{align}
        \mathbb S_{(p,q),(r,s)}=\delta_{p,s}\delta_{q,r}.
    \end{align}
    With the row index as the output physical index and the column index as the
    input physical index, define
    \begin{align}
        \widetilde{\mathcal U}_a^{\,IJ}
        =\sum_K\mathbb S_{IK}\mathcal U_a^{\,KJ},
        \qquad a\in\{1,2,3\}.
    \end{align}
    Thus the swap acts on the output leg, and the corresponding closed operator
    is left-multiplied by $\mathbb S^{\otimes N}$.
    % Source lines: paper_v2.tex 2045--2062.
\end{definition}

\begin{theorem}[Finite-chain formulas for the swap-transformed shifts]
    \label{thm:threeMPU2_chain_formulas}
    \lean{MPOTensor.mpo_ketLeftMul}
    \lean{MPOTensor.IsMPU.ketLeftMul}
    \lean{MPOTensor.shiftPhysicalSwap_mem_unitaryGroup}
    \lean{MPOTensor.mpo_shiftExampleTildeU₁}
    \lean{MPOTensor.mpo_shiftExampleTildeU₂}
    \lean{MPOTensor.mpo_shiftExampleTildeU₃}
    \lean{MPOTensor.shiftExampleTildeU₁_isMPU}
    \lean{MPOTensor.shiftExampleTildeU₂_isMPU}
    \lean{MPOTensor.shiftExampleTildeU₃_isMPU}
    \leanok
    \uses{threeMPU2, thm:threeMPU_chain_formulas,
        def:mpdo_sitewise_physical_matrix, def:mpu}
    On every nonempty chain,
    \begin{align}
        \widetilde U_1^{(N)} & =\mathbb S^{\otimes N}, \\
        \widetilde U_2^{(N)} & =\mathbb S^{\otimes N}
        (T^{(N)\dagger}\otimes T^{(N)}),               \\
        \widetilde U_3^{(N)} & =\mathbb S^{\otimes N}
        (T^{(N)}\otimes T^{(N)\dagger}).
    \end{align}
    Each transformed tensor is a matrix product unitary.
    % Source lines: paper_v2.tex 2052--2062.
\end{theorem}

\begin{proof}\leanok
    \uses{threeMPU2, thm:threeMPU_chain_formulas,
        def:mpdo_sitewise_physical_matrix, def:mpu}
    Acting by a one-site matrix $P$ on every output physical index gives
    \begin{align}
        \operatorname{MPO}_N(P\mathcal M)
        =P^{\otimes N}\operatorname{MPO}_N(\mathcal M).
    \end{align}
    Substituting $P=\mathbb S$ and the three shift formulas gives the stated
    identities. The swap is a permutation matrix and hence unitary. If
    $\mathcal M$ is an MPU, then
    \begin{align}
        \bigl(P^{\otimes N}\operatorname{MPO}_N(\mathcal M)\bigr)^\dagger
        \bigl(P^{\otimes N}\operatorname{MPO}_N(\mathcal M)\bigr)
         & =\operatorname{MPO}_N(\mathcal M)^\dagger
        (P^{\otimes N})^\dagger P^{\otimes N}
        \operatorname{MPO}_N(\mathcal M)             \\
         & =\Id.
    \end{align}
    Thus the sitewise left action preserves the MPU property.
\end{proof}

\begin{theorem}[Conjugacy of the third swap-transformed shift]
    \label{thm:tildeU3}
    \lean{MPOTensor.mpo_shiftExampleTildeU₃_eq_swap_mul_tildeU₂_mul_swap}
    \leanok
    \uses{threeMPU2, thm:threeMPU2_chain_formulas}
    On every nonempty chain,
    \begin{align}
        \widetilde U_3^{(N)}
        =\mathbb S^{\otimes N}\widetilde U_2^{(N)}\mathbb S^{\otimes N}.
    \end{align}
    This is the conjugacy identity in
    \cite[lines~2059--2062]{Cirac2017MPU}.
    % Source label: tildeU3.
\end{theorem}

\begin{proof}\leanok
    \uses{thm:threeMPU2_chain_formulas}
    Split each physical configuration into its two spin configurations.  The
    three operators are permutation matrices.  In these coordinates,
    $\mathbb S^{\otimes N}$ exchanges the two configurations,
    $U_2^{(N)}$ acts by $(T^{(N)\dagger},T^{(N)})$, and $U_3^{(N)}$ acts by
    $(T^{(N)},T^{(N)\dagger})$.  Exchanging the two factors on both sides of
    the second action gives the third, which proves the identity.
\end{proof}

\begin{theorem}[Standard forms of the swap-transformed shifts]
    \label{thm:threeMPU2_remaining_formulas}
    \notready
    \discussion{5715}
    \uses{threeMPU2, def:mpu_source_uv}
    The first two standard forms are
    \begin{align}
        \widetilde u_1 & =\mathbb S,     &
        \widetilde v_1 & =\Id\otimes\Id,   \\
        \widetilde u_2 & =\Id\otimes\Id, &
        \widetilde v_2 & =\mathbb S.
    \end{align}
    % Source lines: paper_v2.tex 2087--2099.
\end{theorem}

\begin{lemma}[Transformation with swaps and symmetry equivalence]
    \label{lemma:sym-trafo-swap}
    \notready
    \discussion{5715}
    \uses{threeMPU, threeMPU2, def:equivalent-symmetry}
    Two MPUs $U_i^{(N)}$ and $U_j^{(N)}$ are in the same phase under time
    reversal combined with the local swap, transposition combined with the
    local swap, or conjugation if and only if the corresponding
    $\widetilde U_i^{(N)}$ and $\widetilde U_j^{(N)}$ are in the same phase
    under time reversal, transposition, or conjugation, respectively.
    % Source lines: paper_v2.tex 2045--2085.
\end{lemma}

\begin{lemma}[Admissible transformation with swaps and symmetry equivalence]
    \label{lem:mpu_admissible_swap_symmetry_equivalence}
    \notready
    \discussion{5715}
    % **Scope restriction (admissible phase comparison):** every blocked endpoint
    % and comparison path carries the datum of
    % Definition~\ref{def:mpu_admissible_equivalence_datum}. Documented in
    % docs/paper-gaps/mpu_canonical_form_full_support.tex.
    \uses{threeMPU, threeMPU2, def:mpu_admissible_strict_symmetry_equivalence, def:mpu_admissible_symmetry_equivalence, def:mpu_reduced_full_support_source_datum, def:mpu_admissible_equivalence_datum}
    Fix one of the three symmetry comparisons and let
    $p\mapsto\mathcal W(p)$ be the actual source-constructed interpolation used
    in that comparison. Suppose it carries an admissible symmetry-path datum of
    blocking length one. Thus the path $\mathcal W(p)$ and the four fixed
    comparison paths $\overline{\mathcal W(p)}$, $\mathcal W(p)^\sharp$,
    $\mathcal W(p)^{\mathsf T}$, and $\mathcal S[\mathcal W(p)]$ are simple and
    carry continuously varying reduced full-support source data. Suppose
    also that every one- and two-site endpoint tensor used in the local
    comparison carries its own datum. Then $U_i^{(N)}$ and $U_j^{(N)}$ are in
    the same admissible phase under time reversal combined with the local swap,
    transposition combined with the local swap, or conjugation if and only if
    the corresponding $\widetilde U_i^{(N)}$ and $\widetilde U_j^{(N)}$ are in
    the same admissible phase under time reversal, transposition, or
    conjugation, respectively.
    % Source lines: paper_v2.tex 2045--2085.
\end{lemma}

\begin{proposition}[Conjugation phases of $U_2$ and $U_3$]
    \label{prop:U2-U3-trivial-conjugation}
    \notready
    \discussion{5715}
    \uses{threeMPU, conjphs, prop:conjsym-strict-equiv}
    After blocking two sites, $U_2^{(N)}$ and $U_3^{(N)}$ are both in the
    trivial phase under conjugation.
    % Source lines: paper_v2.tex 2110--2134.
\end{proposition}

\begin{proposition}[Admissible conjugation phases of $U_2$ and $U_3$]
    \label{prop:mpu_admissible_u2_u3_conjugation}
    \notready
    \discussion{5715}
    % **Scope restriction (two-site blocked data):** both blocked tensors carry
    % reduced full-support source data. No preservation under blocking is
    % asserted. Documented in
    % docs/paper-gaps/mpu_canonical_form_full_support.tex.
    \uses{threeMPU, thm:mpu_admissible_conjugation_odd_rank, thm:mpu_admissible_conjugation_even_rank, def:mpu_reduced_full_support_source_datum, def:mpu_admissible_equivalence_datum, def:mpu_admissible_symmetry_equivalence}
    Suppose the two-site blocked tensors carry reduced full-support source data
    and the standard-form comparison paths carry admissible path data of
    blocking length one. Then the actual two-site blocked tensors of $U_2$ and
    $U_3$ are both in the trivial admissible strict phase under conjugation.
    Consequently $U_2$ and $U_3$ lie in the same admissible equivalence class.
    % Source lines: paper_v2.tex 2110--2134.
\end{proposition}

\begin{theorem}[Continuous unitary path to the factor swap]
    \label{thm:mpu-continuous-unitary-swap-path}
    \lean{Matrix.swapSymmetricProjection,
        Matrix.swapAntisymmetricProjection,
        Matrix.swapPath,
        Matrix.exists_continuous_unitary_swapPath}
    \leanok
    Let $\mathbb S$ denote the swap operator on
    $\mathbb C^d\otimes\mathbb C^d$. There is a continuous map
    $S\colon\mathbb R\to \mathrm{U}(d^2)$ such that
    \begin{align}
        \bigl(S(0),S(1)\bigr) & = \bigl(\Id,\mathbb S\bigr).
    \end{align}
    One concrete choice is
    \begin{align}
        S(x) & = P_{\mathrm{sym}} + e^{\mathrm{i}\pi x}P_{\mathrm{anti}},
    \end{align}
    where
    \begin{align}
        P_{\mathrm{sym}} & = \frac{\Id+\mathbb S}{2},
    \end{align}
    and
    \begin{align}
        P_{\mathrm{anti}} & = \frac{\Id-\mathbb S}{2}.
    \end{align}
    % Source lines: paper_v2.tex 2148--2151 and 2229--2234.
\end{theorem}

\begin{proof}
    \leanok
    The swap is a self-adjoint involution:
    \begin{align}
        \mathbb S^2 & = \Id,
    \end{align}
    and
    \begin{align}
        \mathbb S^\dagger & = \mathbb S.
    \end{align}
    Hence the two projections satisfy
    \begin{align}
        \bigl(P_{\mathrm{sym}}^2,P_{\mathrm{anti}}^2\bigr)
         & = \bigl(P_{\mathrm{sym}},P_{\mathrm{anti}}\bigr),
    \end{align}
    \begin{align}
        \bigl(P_{\mathrm{sym}}P_{\mathrm{anti}},
        P_{\mathrm{anti}}P_{\mathrm{sym}}\bigr) & = (0,0),
    \end{align}
    \begin{align}
        \bigl(P_{\mathrm{sym}}^\dagger,P_{\mathrm{anti}}^\dagger\bigr)
         & = \bigl(P_{\mathrm{sym}},P_{\mathrm{anti}}\bigr),
    \end{align}
    and
    \begin{align}
        P_{\mathrm{sym}}+P_{\mathrm{anti}} & = \Id.
    \end{align}
    Since $\lvert e^{\mathrm{i}\pi x}\rvert=1$, these identities imply
    \begin{align}
        S(x)S(x)^\dagger
         & = P_{\mathrm{sym}}+P_{\mathrm{anti}}=\Id.
    \end{align}
    Continuity follows from continuity of the exponential, while $e^0=1$ and
    $e^{\mathrm{i}\pi}=-1$ give the endpoint identities.
\end{proof}

\begin{proposition}[Time-reversal and transposition phases of $U_2$ and $U_3$]
    \label{prop:U2-U3-trivial-tr-transpose}
    \notready
    \discussion{5715}
    \uses{threeMPU2, lemma:sym-trafo-swap, thm:mpu-continuous-unitary-swap-path, def:strictly-equivalent-symmetry}
    The MPUs $\widetilde U_2^{(N)}$ and $\widetilde U_3^{(N)}$ are in the same
    phase under both time reversal and transposition.
    % Source lines: paper_v2.tex 2136--2155.
\end{proposition}

\begin{proposition}[Admissible time-reversal and transposition phases of $U_2$ and $U_3$]
    \label{prop:mpu_admissible_u2_u3_time_reversal_transposition}
    \notready
    \discussion{5715}
    % **Scope restriction (admissible symmetry path):** the comparison path
    % carries the datum of Definition~\ref{def:mpu_admissible_equivalence_datum}.
    % Documented in docs/paper-gaps/mpu_canonical_form_full_support.tex.
    \uses{threeMPU2, lem:mpu_admissible_swap_symmetry_equivalence, def:mpu_admissible_strict_symmetry_equivalence, def:mpu_reduced_full_support_source_datum, def:mpu_admissible_equivalence_datum}
    For each of time reversal and transposition, let
    $p\mapsto\mathcal W_{\mathcal S}(p)$ be the actual source-constructed
    interpolation from $\widetilde U_2$ to $\widetilde U_3$. Suppose it carries
    an admissible symmetry-path datum of blocking length one. Thus the path
    $\mathcal W_{\mathcal S}(p)$ and its fixed conjugate, physical-adjoint,
    transposed, and $\mathcal S$-image comparison paths are simple and carry
    continuously varying reduced full-support source data. Suppose also that
    the one- and two-site endpoint tensors used in the local comparison carry
    their own data. Then $\widetilde U_2^{(N)}$ and
    $\widetilde U_3^{(N)}$ are in the same admissible phase under both time
    reversal and transposition.
    % Source lines: paper_v2.tex 2136--2155.
\end{proposition}

\begin{proposition}[Conjugation phases of $U_1$ and $U_2$]
    \label{prop:eqiv-U1-U2-conjugate}
    \notready
    \discussion{5715}
    \uses{threeMPU2, lemma:sym-trafo-swap, conjphs, prop:conjsym-strict-equiv, cor:equiv-under-conjugation}
    The MPUs $U_1^{(N)}$ and $U_2^{(N)}$ lie in different strict conjugation
    phases exactly when $d=4k+2$ for a nonnegative integer $k$. Otherwise they
    lie in the same strict phase. All these phases are in Case I, so the two
    MPUs are equivalent after blocking and adjoining ancillas.
    % Source lines: paper_v2.tex 2158--2200.
\end{proposition}

\begin{proposition}[Admissible conjugation phases of $U_1$ and $U_2$]
    \label{prop:mpu_admissible_u1_u2_conjugation}
    \notready
    \discussion{5715}
    % **Scope restriction (admissible blockings and paths):** every standard form,
    % blocked endpoint, and comparison path carries the required source data.
    % Documented in docs/paper-gaps/mpu_canonical_form_full_support.tex.
    \uses{threeMPU2, lem:mpu_admissible_swap_symmetry_equivalence, thm:mpu_admissible_conjugation_odd_rank, thm:mpu_admissible_conjugation_even_rank, cor:mpu_admissible_equiv_under_conjugation, def:mpu_reduced_full_support_source_datum, def:mpu_admissible_equivalence_datum}
    For each endpoint tensor $\mathcal A$ among
    $U_1$, $U_2$, $\widetilde U_1$, and $\widetilde U_2$, suppose that
    $\mathcal A$, its conjugate comparison $\overline{\mathcal A}$, its actual
    two-site block $\mathcal A_2$, and the conjugate comparison
    $\overline{\mathcal A_2}$ are simple and carry their own reduced
    full-support source data. For every strict comparison used below, let
    $p\mapsto\mathcal W(p)$ be the actual
    source-constructed interpolation between the corresponding endpoint
    tensors. Suppose it carries an admissible conjugation-symmetry path datum of
    blocking length one: the path and its fixed conjugate, physical-adjoint,
    transposed, and conjugation-image comparison paths are simple and carry
    continuously varying source data. Require a length-one admissible path datum
    also for the actual blocked interpolation used in the final equivalence
    claim. Under these hypotheses,
    $U_1^{(N)}$ and $U_2^{(N)}$ are in different admissible strict phases under
    conjugation exactly when $k$ is a nonnegative integer and
    \begin{align}
        d & =4k+2.
    \end{align}
    Otherwise they are in the same admissible strict phase. All these standard
    forms lie in Case I, so $U_1^{(N)}$ and $U_2^{(N)}$ are admissibly equivalent
    after blocking and adding ancillas.
    % Source lines: paper_v2.tex 2158--2200.
\end{proposition}

\begin{definition}[Two-species ancilla configurations]
    \label{def:mpu_shift_ancilla_config}
    \lean{MPOTensor.ShiftAncillaConfig}
    \leanok
    \discussion{7470}
    For a periodic chain of length $N$, the configuration space consists of
    triples $((\sigma_1,\sigma_2),a)$ of two physical configurations and one
    $d$-dimensional ancilla configuration.
\end{definition}

\begin{definition}[Sitewise physical swap]
    \label{def:mpu_shift_ancilla_swap_12}
    \lean{MPOTensor.shiftAncillaSwap₁₂}
    \leanok
    \discussion{7470}
    \uses{def:mpu_shift_ancilla_config}
    The sitewise physical swap is
    \begin{align}
        S^{1,2}((\sigma_1,\sigma_2),a)
         & =((\sigma_2,\sigma_1),a).
    \end{align}
    % Source: paper_v2.tex 2217--2229 and Figure fig:TR-ancilla.
\end{definition}

\begin{definition}[Sitewise second-species--ancilla swap]
    \label{def:mpu_shift_ancilla_swap_2a}
    \lean{MPOTensor.shiftAncillaSwap₂a}
    \leanok
    \discussion{7470}
    \uses{def:mpu_shift_ancilla_config}
    The sitewise swap of the second species with the ancilla is
    \begin{align}
        S^{2,a}((\sigma_1,\sigma_2),a)
         & =((\sigma_1,a),\sigma_2).
    \end{align}
    % Source: paper_v2.tex 2217--2229 and Figure fig:TR-ancilla.
\end{definition}

\begin{definition}[Crossed first-species--ancilla swap]
    \label{def:mpu_shift_ancilla_swap_1a_next}
    \lean{MPOTensor.shiftAncillaSwap₁aNext,
        MPOTensor.shiftAncillaSwap₁aNext_apply}
    \leanok
    \discussion{7470}
    \uses{def:mpu_shift_ancilla_config, def:mpo_rotate_config}
    Let $R$ cyclically rotate configurations on the chain. The crossed swap is
    \begin{align}
        S^{1,a}_{+}((\sigma_1,\sigma_2),a)
         & =((Ra,\sigma_2),R^{-1}\sigma_1).
    \end{align}
    % Source: paper_v2.tex 2217--2229 and Figure fig:TR-ancilla.
\end{definition}

\begin{definition}[Literal three-swap circuit]
    \label{def:mpu_shift_ancilla_three_swap}
    \lean{MPOTensor.shiftAncillaThreeSwap}
    \leanok
    \discussion{7470}
    \uses{def:mpu_shift_ancilla_swap_12,
        def:mpu_shift_ancilla_swap_2a,
        def:mpu_shift_ancilla_swap_1a_next}
    Let $C$ be the permutation obtained by applying $S^{1,a}_{+}$ first, then
    $S^{1,2}$, and finally $S^{2,a}$.
    % Source: paper_v2.tex 2217--2229 and Figure fig:TR-ancilla.
\end{definition}

\begin{definition}[Counter-shift permutation]
    \label{def:mpu_shift_ancilla_counter_shift}
    \lean{MPOTensor.shiftAncillaCounterShift}
    \leanok
    \discussion{7470}
    \uses{def:mpu_shift_ancilla_config, def:mpo_rotate_config}
    Define
    \begin{align}
        K((\sigma_1,\sigma_2),a)
         & =((R\sigma_2,R^{-1}\sigma_1),a).
    \end{align}
    % Source: paper_v2.tex 2217--2229 and Figure fig:TR-ancilla.
\end{definition}

\begin{theorem}[Literal three-swap permutation identity]
    \label{thm:mpu_shift_ancilla_three_swap_endpoint}
    \lean{MPOTensor.shiftAncillaThreeSwap_conj_swap₁₂,
        MPOTensor.shiftAncillaThreeSwap_conj_swap₁₂_apply,
        MPOTensor.permMatrix_shiftAncillaThreeSwap_endpoint,
        MPOTensor.permMatrix_shiftAncillaThreeSwap_conj_swap₁₂}
    \leanok
    \discussion{7470}
    \uses{threeMPU2,
        def:mpu_shift_ancilla_swap_12,
        def:mpu_shift_ancilla_three_swap,
        def:mpu_shift_ancilla_counter_shift}
    The literal three-swap circuit satisfies
    \begin{align}
        C^{-1}S^{1,2}C & =K.
    \end{align}
    Writing $P_f$ for the permutation matrix of $f$, one also has
    \begin{align}
        P_C P_{S^{1,2}} P_{C^{-1}} & =P_K.
    \end{align}
    At the coordinate level, $K$ is the configuration map of
    $\widetilde U_3$, whereas $\widetilde U_2$ has configuration map
    $((\sigma_1,\sigma_2),a)\mapsto
        ((R^{-1}\sigma_2,R\sigma_1),a)$. Thus the literal source labels do not
    produce the claimed $\widetilde U_2$ endpoint on a nondegenerate periodic
    chain. The theorem does not identify $P_K$ with an ancilla-enlarged MPU and
    asserts neither a continuous path nor a symmetry action on the enlarged
    operator algebra.
    % Source: paper_v2.tex 2217--2229 and Figure fig:TR-ancilla.
\end{theorem}

\begin{proof}\leanok
    \uses{def:mpu_shift_ancilla_swap_12,
        def:mpu_shift_ancilla_three_swap,
        def:mpu_shift_ancilla_counter_shift}
    Evaluating the three swaps gives
    \begin{align}
        C^{-1}S^{1,2}C((\sigma_1,\sigma_2),a)
         & =((R\sigma_2,R^{-1}\sigma_1),a).
    \end{align}
    Extensionality gives the permutation equality. The permutation-matrix
    identity follows by functoriality.
\end{proof}

\begin{proposition}[Ancilla equivalence of $U_1$ and $U_2$]
    \label{prop:U1-U2-equiv-ancillatrick}
    \notready
    \discussion{5715}
    \uses{threeMPU2,
        lemma:sym-trafo-swap,
        thm:mpu-continuous-unitary-swap-path,
        def:strictly-equivalent-symmetry}
    After adding one ancilla of dimension $d$ per site,
    $\widetilde U_1^{(N)}$ and $\widetilde U_2^{(N)}$ are strictly equivalent
    under both time reversal and transposition.
    % Source lines: paper_v2.tex 2203--2235.
\end{proposition}

\begin{proposition}[Admissible ancilla equivalence of $U_1$ and $U_2$]
    \label{prop:mpu_admissible_u1_u2_ancilla}
    \notready
    \discussion{5715}
    % **Scope restriction (ancilla endpoint and admissible path data):** the
    % ancilla-enlarged tensors carry source data and the comparison path is
    % admissible. No preservation under adjoining ancillas is asserted.
    % Documented in docs/paper-gaps/mpu_canonical_form_full_support.tex.
    \uses{threeMPU2,
        lem:mpu_admissible_swap_symmetry_equivalence,
        def:mpu_admissible_strict_symmetry_equivalence,
        def:mpu_admissible_symmetry_equivalence,
        def:mpu_reduced_full_support_source_datum,
        def:mpu_admissible_equivalence_datum}
    After adding one ancilla of dimension $d$ per site and performing the source
    blocking, denote the actual enlarged blocked endpoints by
    $\widehat{\mathcal U}_1$ and $\widehat{\mathcal U}_2$. For either time
    reversal or transposition, let $p\mapsto\mathcal W_{\mathcal S}(p)$ be the
    source-constructed interpolation from $\widehat{\mathcal U}_1$ to
    $\widehat{\mathcal U}_2$. Suppose this interpolation carries an admissible
    symmetry-path datum of blocking length one. Thus the path
    $\mathcal W_{\mathcal S}(p)$ and its fixed conjugate, physical-adjoint,
    transposed, and $\mathcal S$-image comparison paths are simple and carry
    continuously varying reduced full-support source data. Suppose also that
    the one- and two-site enlarged endpoint tensors used in the local comparison
    carry their own data. Then $\widehat{\mathcal U}_1$ and
    $\widehat{\mathcal U}_2$ are admissibly strictly equivalent under both time
    reversal and transposition. No equivalence of the original ancilla-free
    tensors is asserted, because the equal ancilla dimensions used here need not
    satisfy the coprimality condition in
    Definition~\ref{def:mpu_admissible_symmetry_equivalence}.
    % Source lines: paper_v2.tex 2203--2235.
\end{proposition}

\section{Finite-region observable algebras}

\begin{definition}[Finite-region local algebra]
    \label{def:qca_local_algebra}
    \lean{SpinChain.Config,SpinChain.LocalAlgebra}
    \leanok
    Fix a nonnegative integer $d$.  For a finite region
    $\Lambda\subset\mathbb Z$, the configuration space and local observable
    algebra are
    \begin{align}
        \operatorname{Config}_d(\Lambda)
         & =\{0,\ldots,d-1\}^{\Lambda},                     \\
        \mathcal A_\Lambda
         & =M_{\operatorname{Config}_d(\Lambda)}(\mathbb C)
        \cong M_d^{\otimes\Lambda}.
    \end{align}
    The algebra $\mathcal A_\Lambda$ carries its operator norm.  This is the
    finite-region construction used in the QCA appendix of
    \cite{Cirac2017MPU}; no completion is taken here.
\end{definition}

\begin{definition}[Canonical inclusion of finite-region observables]
    \label{def:qca_local_inclusion}
    \lean{SpinChain.localInclusion}
    \leanok
    \uses{def:qca_local_algebra}
    If $\Lambda\subseteq\Gamma$ are finite regions, define
    \begin{align}
        \iota_{\Lambda,\Gamma} & \colon\mathcal A_\Lambda\longrightarrow
        \mathcal A_\Gamma,                                                      \\
        \iota_{\Lambda,\Gamma}(A)
                               & = A\otimes \mathbf 1_{\Gamma\setminus\Lambda}.
    \end{align}
    Here product configurations are identified with configurations on
    $\Gamma$.  This is a unital star-algebra homomorphism.
\end{definition}

\begin{lemma}[Injectivity of the finite-region inclusion]
    \label{lem:qca_local_inclusion_injective}
    \lean{SpinChain.localInclusion_injective}
    \leanok
    \uses{def:qca_local_inclusion}
    For $d>0$, the map $\iota_{\Lambda,\Gamma}$ is injective.
\end{lemma}

\begin{proof}\leanok
    Since $d>0$, choose a configuration $c$ on
    $\Gamma\setminus\Lambda$.  For configurations $x,y$ on $\Lambda$,
    \begin{align}
        \iota_{\Lambda,\Gamma}(A)((x,c),(y,c))
         & =A(x,y)\cdot\mathbf 1_{\Gamma\setminus\Lambda}(c,c)
        =A(x,y).
    \end{align}
    Hence $\iota_{\Lambda,\Gamma}(A)=\iota_{\Lambda,\Gamma}(B)$ implies
    $A=B$.
\end{proof}

\begin{lemma}[Operator-norm preservation]
    \label{lem:qca_local_inclusion_norm}
    \lean{SpinChain.localInclusion_norm}
    \leanok
    \uses{def:qca_local_inclusion}
    If $d>0$, then for every $A\in\mathcal A_\Lambda$,
    \begin{align}
        \lVert\iota_{\Lambda,\Gamma}(A)\rVert & =\lVert A\rVert.
    \end{align}
\end{lemma}

\begin{proof}\leanok
    \uses{lem:qca_local_inclusion_injective}
    An injective star-algebra homomorphism between complex $C^*$-algebras is
    isometric.
\end{proof}

\begin{lemma}[Isometry of the finite-region inclusion]
    \label{lem:qca_local_inclusion_isometry}
    \lean{SpinChain.localInclusion_isometry}
    \leanok
    \uses{def:qca_local_inclusion}
    For $d>0$, the map $\iota_{\Lambda,\Gamma}$ is an operator-norm isometry.
\end{lemma}

\begin{proof}\leanok
    \uses{lem:qca_local_inclusion_injective}
    Use injectivity and the general theorem that an injective star-algebra
    homomorphism between complex $C^*$-algebras is an isometry.
\end{proof}

\begin{lemma}[Identity inclusion]
    \label{lem:qca_local_inclusion_refl}
    \lean{SpinChain.localInclusion_refl}
    \leanok
    \uses{def:qca_local_inclusion}
    For every finite region $\Lambda$,
    $\iota_{\Lambda,\Lambda}$ is the identity on $\mathcal A_\Lambda$.
\end{lemma}

\begin{proof}\leanok
    When the two regions agree, their complement is empty.  Hence
    \begin{align}
        \iota_{\Lambda,\Lambda}(A)
         & =A\otimes\mathbf 1_{\varnothing}
        =A\otimes 1
        =A.
    \end{align}
\end{proof}

\begin{lemma}[Composition of finite-region inclusions]
    \label{lem:qca_local_inclusion_trans}
    \lean{SpinChain.localInclusion_trans}
    \leanok
    \uses{def:qca_local_inclusion}
    If $\Lambda\subseteq\Gamma\subseteq\Delta$, then
    \begin{align}
        \iota_{\Gamma,\Delta}\circ\iota_{\Lambda,\Gamma}
         & =\iota_{\Lambda,\Delta}.
    \end{align}
\end{lemma}

\begin{proof}\leanok
    The complement decomposes as the disjoint union
    \begin{align}
        \Delta\setminus\Lambda
         & =(\Gamma\setminus\Lambda)\sqcup(\Delta\setminus\Gamma), \\
        \mathbf 1_{\Delta\setminus\Lambda}
         & =\mathbf 1_{\Gamma\setminus\Lambda}
        \otimes\mathbf 1_{\Delta\setminus\Gamma}.
    \end{align}
    Therefore
    \begin{align}
        \iota_{\Gamma,\Delta}(\iota_{\Lambda,\Gamma}(A))
         & =(A\otimes\mathbf 1_{\Gamma\setminus\Lambda})
        \otimes\mathbf 1_{\Delta\setminus\Gamma}         \\
         & =A\otimes\mathbf 1_{\Delta\setminus\Lambda}   \\
         & =\iota_{\Lambda,\Delta}(A).
    \end{align}
\end{proof}

\begin{theorem}[Disjoint finite-region observables commute]
    \label{thm:qca_disjoint_local_inclusion_commute}
    \lean{SpinChain.localInclusion_commute_of_disjoint}
    \leanok
    \uses{def:qca_local_inclusion}
    If $\Lambda\cap\Gamma=\varnothing$, $A\in\mathcal A_\Lambda$, and
    $B\in\mathcal A_\Gamma$, then their images in the algebra of the union
    satisfy
    \begin{align}
        \iota_{\Lambda,\Lambda\cup\Gamma}(A)
        \iota_{\Gamma,\Lambda\cup\Gamma}(B)
         & =\iota_{\Gamma,\Lambda\cup\Gamma}(B)
        \iota_{\Lambda,\Lambda\cup\Gamma}(A).
    \end{align}
    This tensor-factor identity is local algebra infrastructure for the
    quasi-local algebra used in \cite[lines~2292--2300]{Cirac2017MPU}; the
    appendix does not state it as a separate theorem.
    % Source context: paper_v2.tex 2292--2300.
\end{theorem}

\begin{proof}\leanok
    For configurations $x,y$ on $\Lambda\cup\Gamma$, split each intermediate
    configuration $z$ into $(z|_\Lambda,z|_\Gamma)$.  The matrix entry of the
    left-hand side is
    \begin{align}
        \sum_{z_\Lambda,z_\Gamma}
        A(x|_\Lambda,z_\Lambda)
        \delta_{x|_\Gamma,z_\Gamma}
        B(z_\Gamma,y|_\Gamma)
        \delta_{z_\Lambda,y|_\Lambda}
         & =A(x|_\Lambda,y|_\Lambda)
        B(x|_\Gamma,y|_\Gamma).
    \end{align}
    Reversing the two factors gives the same scalar product.
\end{proof}

\begin{theorem}[Finite relative commutant]
    \label{thm:qca_finite_relative_commutant}
    \lean{SpinChain.mem_range_localInclusion_iff_commute_complement}
    \leanok
    \uses{def:qca_local_inclusion}
    Let $\Lambda\subseteq\Gamma$ be finite regions and
    $A\in\mathcal A_\Gamma$. Then $A$ acts trivially outside $\Lambda$ if
    and only if it commutes with the full observable algebra on the complement:
    \begin{align}
        A\in\iota_{\Lambda,\Gamma}(\mathcal A_\Lambda)
         & \iff
        \forall B\in\mathcal A_{\Gamma\setminus\Lambda},\quad
        A\iota_{\Gamma\setminus\Lambda,\Gamma}(B)
        =\iota_{\Gamma\setminus\Lambda,\Gamma}(B)A.
    \end{align}
    This is a finite-dimensional tensor-factor identity. It is not a statement
    of Schumacher--Werner or the QCA appendix of \cite{Cirac2017MPU}.
\end{theorem}

\begin{proof}\leanok
    \uses{lem:amplified_center_scalar}
    Splitting configurations gives
    $\mathcal A_\Gamma\cong
        \mathcal A_\Lambda\otimes\mathcal A_{\Gamma\setminus\Lambda}$. Under this
    identification, the two inclusions are $C\mapsto C\otimes\mathbf 1$ and
    $B\mapsto\mathbf 1\otimes B$. Write an arbitrary matrix in blocks indexed
    by configurations on $\Lambda$. If it commutes with every
    $\mathbf 1\otimes B$, each block commutes with the full matrix algebra on
    $\Gamma\setminus\Lambda$ and is therefore scalar. Thus the matrix is
    $C\otimes\mathbf 1$ for some $C\in\mathcal A_\Lambda$. The converse follows
    from
    \begin{align}
        (C\otimes\mathbf 1)(\mathbf 1\otimes B)
         & =C\otimes B
        =(\mathbf 1\otimes B)(C\otimes\mathbf 1).
    \end{align}
\end{proof}

\subsection{Bipartite support algebras}

Let $I$ and $J$ be finite sets.  For
$X\in M_{I\times J}(\mathbb C)$, the existing left coefficient slice and
right operator block are the matrices
\begin{align}
    X^{\mathrm L}_{rs}(i,j) & =X((i,r),(j,s)), &
    X^{\mathrm R}_{ij}(r,s) & =X((i,r),(j,s)).
\end{align}
For a star-subalgebra $\mathcal D\subseteq M_K(\mathbb C)$, write
$\mathcal D'=\{Y\in M_K(\mathbb C)\mid [Y,Z]=0\text{ for every }Z\in\mathcal D\}$
for its commutant.  The construction below is the full complex matrix-algebra
specialization of the coefficient construction in
\cite[Section~4.3]{SchumacherWerner2004QCA} and of the leastness and commutant
statements in \cite[Lemma~7]{Gross2012IndexTheory}.

\textbf{Local fix (star-closed input):}
The printed statement of Lemma~7 says only that $\mathcal A$ is a
subalgebra.  Its coefficient-generated support and leastness clauses remain
valid for arbitrary algebra input, but its commutant characterization fails
unless $\mathcal A$ is closed under the involution.  For instance, the algebra
$\operatorname{span}\{\mathbf 1,E_{12}\}\subset M_2(\mathbb C)$ gives a
counterexample to the printed commutant clause if involution closure is
omitted.  Schumacher and Werner first define the ordinary algebra generated
by the coefficient support and then derive its adjoint closure from adjoint
closure of the observable-algebra input
\cite[lines~1161--1171]{SchumacherWerner2004QCA}.  Accordingly, the statements
below use a star-subalgebra.  This is a necessary correction to the printed
hypothesis.  In the QCA application the input is the image of a local
observable algebra under a star-automorphism
\cite[Section~7.2]{Gross2012IndexTheory}.  The MPU application likewise takes
the image of a local star-algebra under a star-automorphism
\cite[lines~2300--2306]{Cirac2017MPU}.  This correction, including the
counterexample without involution closure and the exact application ranges,
is recorded in \cite{gap:gnvw12_support_algebra_star_closure}.

\textbf{Scope restriction (full matrix factors):}
After this correction, Lemma~7 permits arbitrary finite-dimensional
$C^*$-algebras $\mathcal B_1$ and $\mathcal B_2$, whereas the statements below
take $\mathcal B_1=M_I(\mathbb C)$ and $\mathcal B_2=M_J(\mathbb C)$.
Full matrix factors suffice for the MPU spin chain because every one-site
algebra is isomorphic to $M_d(\mathbb C)$
\cite[lines~2320--2334]{Cirac2017MPU}.  The general-factor result is not
claimed here; see \cite{gap:gnvw12_support_algebra_full_matrix_scope}.

No full-matrix structure theorem, dimension formula, Margolus gate, depth-two
circuit, MPU standard form, or index comparison is asserted here.

\begin{theorem}[Adjoint of a right coefficient block]
    \label{thm:qca_operator_block_adjoint}
    \lean{Matrix.operatorBlock_conjTranspose}
    \leanok
    For every $X\in M_{I\times J}(\mathbb C)$ and $i,j\in I$,
    \begin{align}
        (X^\dagger)^{\mathrm R}_{ij}
         & =\bigl(X^{\mathrm R}_{ji}\bigr)^\dagger.
    \end{align}
\end{theorem}

\begin{proof}\leanok
    Compare the $(r,s)$ entries.  Both sides equal
    $\overline{X((j,s),(i,r))}$.
\end{proof}

\begin{definition}[Left coefficient set]
    \label{def:qca_left_coefficient_set}
    \lean{Matrix.leftCoefficientSet}
    \leanok
    Let $\mathcal A\subseteq M_{I\times J}(\mathbb C)$ be a star-subalgebra.
    Its set of left coefficients is
    \begin{align}
        \operatorname{Coeff}_{\mathrm L}(\mathcal A)
         & =\bigl\{X^{\mathrm L}_{rs}\ \bigm|\
        X\in\mathcal A,\ r,s\in J\bigr\}
        \subseteq M_I(\mathbb C).
    \end{align}
\end{definition}

\begin{definition}[Right coefficient set]
    \label{def:qca_right_coefficient_set}
    \lean{Matrix.rightCoefficientSet}
    \leanok
    Let $\mathcal A\subseteq M_{I\times J}(\mathbb C)$ be a star-subalgebra.
    Its set of right coefficients is
    \begin{align}
        \operatorname{Coeff}_{\mathrm R}(\mathcal A)
         & =\bigl\{X^{\mathrm R}_{ij}\ \bigm|\
        X\in\mathcal A,\ i,j\in I\bigr\}
        \subseteq M_J(\mathbb C).
    \end{align}
\end{definition}

\begin{definition}[Left support algebra]
    \label{def:qca_left_support_algebra}
    \lean{Matrix.leftSupportAlgebra}
    \leanok
    \uses{def:qca_left_coefficient_set}
    The left support algebra of $\mathcal A$ is the star-subalgebra
    \begin{align}
        \operatorname{Spp}_{\mathrm L}(\mathcal A)
         & =\operatorname{alg}^*
        \bigl(\operatorname{Coeff}_{\mathrm L}(\mathcal A)\bigr)
        \subseteq M_I(\mathbb C)
    \end{align}
    generated by all left coefficients.
\end{definition}

\begin{definition}[Right support algebra]
    \label{def:qca_right_support_algebra}
    \lean{Matrix.rightSupportAlgebra}
    \leanok
    \uses{def:qca_right_coefficient_set}
    The right support algebra of $\mathcal A$ is the star-subalgebra
    \begin{align}
        \operatorname{Spp}_{\mathrm R}(\mathcal A)
         & =\operatorname{alg}^*
        \bigl(\operatorname{Coeff}_{\mathrm R}(\mathcal A)\bigr)
        \subseteq M_J(\mathbb C)
    \end{align}
    generated by all right coefficients.
\end{definition}

\begin{theorem}[Left coefficient generators belong to the support algebra]
    \label{thm:qca_left_coefficient_mem_support}
    \lean{Matrix.bipartite_slice_mem_left_support_algebra}
    \leanok
    \uses{def:qca_left_support_algebra}
    If $X\in\mathcal A$ and $r,s\in J$, then
    \begin{align}
        X^{\mathrm L}_{rs}
         & \in\operatorname{Spp}_{\mathrm L}(\mathcal A).
    \end{align}
\end{theorem}

\begin{proof}\leanok
    \uses{def:qca_left_coefficient_set, def:qca_left_support_algebra}
    The matrix $X^{\mathrm L}_{rs}$ is one of the generators of the defining
    star-algebra closure.
\end{proof}

\begin{theorem}[Right coefficient generators belong to the support algebra]
    \label{thm:qca_right_coefficient_mem_support}
    \lean{Matrix.operator_block_mem_right_support_algebra}
    \leanok
    \uses{def:qca_right_support_algebra}
    If $X\in\mathcal A$ and $i,j\in I$, then
    \begin{align}
        X^{\mathrm R}_{ij}
         & \in\operatorname{Spp}_{\mathrm R}(\mathcal A).
    \end{align}
\end{theorem}

\begin{proof}\leanok
    \uses{def:qca_right_coefficient_set, def:qca_right_support_algebra}
    The matrix $X^{\mathrm R}_{ij}$ is one of the generators of the defining
    star-algebra closure.
\end{proof}

\begin{theorem}[Reconstruction from left coefficients]
    \label{thm:qca_left_coefficient_reconstruction}
    \lean{Matrix.eq_sum_left_coefficient_kronecker_single}
    \leanok
    For every $X\in M_{I\times J}(\mathbb C)$,
    \begin{align}
        X & =\sum_{r,s\in J}X^{\mathrm L}_{rs}\otimes E_{rs}.
    \end{align}
\end{theorem}

\begin{proof}\leanok
    At the entry $((i,r),(j,s))$, only the matrix unit $E_{rs}$ contributes,
    and its coefficient is $X^{\mathrm L}_{rs}(i,j)=X((i,r),(j,s))$.
\end{proof}

\begin{theorem}[Reconstruction from right coefficients]
    \label{thm:qca_right_coefficient_reconstruction}
    \lean{Matrix.eq_sum_single_kronecker_right_coefficient}
    \leanok
    For every $X\in M_{I\times J}(\mathbb C)$,
    \begin{align}
        X & =\sum_{i,j\in I}E_{ij}\otimes X^{\mathrm R}_{ij}.
    \end{align}
\end{theorem}

\begin{proof}\leanok
    At the entry $((i,r),(j,s))$, only the matrix unit $E_{ij}$ contributes,
    and its coefficient is $X^{\mathrm R}_{ij}(r,s)=X((i,r),(j,s))$.
\end{proof}

\begin{theorem}[Minimality of the left support algebra]
    \label{thm:qca_left_support_minimality}
    \lean{Matrix.left_support_algebra_le_iff}
    \leanok
    \uses{def:qca_left_support_algebra}
    For every star-subalgebra $\mathcal B\subseteq M_I(\mathbb C)$,
    \begin{align}
        \operatorname{Spp}_{\mathrm L}(\mathcal A)\subseteq\mathcal B
        \quad\Longleftrightarrow\quad
         & X^{\mathrm L}_{rs}\in\mathcal B
        \quad\text{for every }X\in\mathcal A\text{ and }r,s\in J.
    \end{align}
\end{theorem}

\begin{proof}\leanok
    \uses{def:qca_left_coefficient_set, def:qca_left_support_algebra}
    Use the universal property of the star-subalgebra generated by
    $\operatorname{Coeff}_{\mathrm L}(\mathcal A)$.
\end{proof}

\begin{theorem}[Minimality of the right support algebra]
    \label{thm:qca_right_support_minimality}
    \lean{Matrix.right_support_algebra_le_iff}
    \leanok
    \uses{def:qca_right_support_algebra}
    For every star-subalgebra $\mathcal C\subseteq M_J(\mathbb C)$,
    \begin{align}
        \operatorname{Spp}_{\mathrm R}(\mathcal A)\subseteq\mathcal C
        \quad\Longleftrightarrow\quad
         & X^{\mathrm R}_{ij}\in\mathcal C
        \quad\text{for every }X\in\mathcal A\text{ and }i,j\in I.
    \end{align}
\end{theorem}

\begin{proof}\leanok
    \uses{def:qca_right_coefficient_set, def:qca_right_support_algebra}
    Use the universal property of the star-subalgebra generated by
    $\operatorname{Coeff}_{\mathrm R}(\mathcal A)$.
\end{proof}

\begin{definition}[Kronecker-product submodule]
    \label{def:qca_kronecker_submodule}
    \lean{Matrix.kroneckerSubmodule}
    \leanok
    For submodules $S\subseteq M_I(\mathbb C)$ and
    $T\subseteq M_J(\mathbb C)$, define
    \begin{align}
        S\boxtimes T
         & =\operatorname{span}_{\mathbb C}
        \{A\otimes B\mid A\in S,\ B\in T\}
        \subseteq M_{I\times J}(\mathbb C).
    \end{align}
    This is the tensor-product subspace used below, rather than a
    conjunction of two one-sided containment conditions.
\end{definition}

\begin{theorem}[Coefficient criterion for the Kronecker-product submodule]
    \label{thm:qca_mem_kronecker_submodule_iff}
    \lean{Matrix.mem_kroneckerSubmodule_iff}
    \leanok
    \uses{def:qca_kronecker_submodule}
    For $X\in M_{I\times J}(\mathbb C)$,
    \begin{align}
        X\in S\boxtimes T
        \quad\Longleftrightarrow\quad
         & X^{\mathrm L}_{rs}\in S
        \quad\text{for every }r,s\in J, \\
         & \text{and}\quad
        X^{\mathrm R}_{ij}\in T
        \quad\text{for every }i,j\in I.
    \end{align}
\end{theorem}

\begin{proof}\leanok
    \uses{def:qca_kronecker_submodule}
    For $A\in S$ and $B\in T$, direct evaluation gives
    \begin{align}
        (A\otimes B)^{\mathrm L}_{rs} & =B_{rs}A, &
        (A\otimes B)^{\mathrm R}_{ij} & =A_{ij}B.
    \end{align}
    Thus the forward implication extends by linearity.
    Conversely, expand the right blocks in a basis of $T$.  Applying the dual
    coordinate functionals to those blocks produces matrices in $S$ from the
    left-slice hypothesis, and the resulting finite sum reconstructs $X$ as a
    sum of Kronecker products from $S$ and $T$.
\end{proof}

\begin{theorem}[Kronecker-submodule criterion for left support]
    \label{thm:qca_left_support_tensor_submodule}
    \lean{Matrix.left_support_algebra_le_iff_le_kronecker_submodule}
    \leanok
    \uses{def:qca_left_support_algebra, def:qca_kronecker_submodule}
    For every star-subalgebra $\mathcal B\subseteq M_I(\mathbb C)$,
    \begin{align}
        \operatorname{Spp}_{\mathrm L}(\mathcal A)\subseteq\mathcal B
        \quad\Longleftrightarrow\quad
        \mathcal A & \subseteq
        \mathcal B\boxtimes M_J(\mathbb C).
    \end{align}
\end{theorem}

\begin{proof}\leanok
    \uses{thm:qca_left_support_minimality, thm:qca_mem_kronecker_submodule_iff}
    By Theorem~\ref{thm:qca_mem_kronecker_submodule_iff} with the right
    submodule the full matrix space, membership in the Kronecker span is
    exactly membership of every left coefficient in $\mathcal B$.  Apply
    Theorem~\ref{thm:qca_left_support_minimality}.
\end{proof}

\begin{theorem}[Kronecker-submodule criterion for right support]
    \label{thm:qca_right_support_kronecker_submodule}
    \lean{Matrix.right_support_algebra_le_iff_le_kronecker_submodule}
    \leanok
    \uses{def:qca_right_support_algebra, def:qca_kronecker_submodule}
    For every star-subalgebra $\mathcal C\subseteq M_J(\mathbb C)$,
    \begin{align}
        \operatorname{Spp}_{\mathrm R}(\mathcal A)\subseteq\mathcal C
        \quad\Longleftrightarrow\quad
        \mathcal A & \subseteq M_I(\mathbb C)\boxtimes\mathcal C.
    \end{align}
\end{theorem}

\begin{proof}\leanok
    \uses{thm:qca_right_support_minimality, thm:qca_mem_kronecker_submodule_iff}
    The coefficient criterion reduces the right-hand containment to membership
    of every right block in $\mathcal C$.  This is equivalent to the left-hand
    inclusion by Theorem~\ref{thm:qca_right_support_minimality}.
\end{proof}

\begin{theorem}[Two-sided support-algebra containment]
    \label{thm:qca_two_sided_support_containment}
    \lean{Matrix.le_support_kroneckerSubmodule}
    \leanok
    \uses{def:qca_left_support_algebra, def:qca_right_support_algebra, def:qca_kronecker_submodule}
    Every bipartite star-subalgebra satisfies the literal two-sided containment
    \begin{align}
        \mathcal A
         & \subseteq
        \operatorname{Spp}_{\mathrm L}(\mathcal A)
        \boxtimes
        \operatorname{Spp}_{\mathrm R}(\mathcal A).
    \end{align}
    This is the full complex matrix-algebra specialization of the
    support-algebra containment in Schumacher--Werner Section~4.3 and of
    Gross--Nesme--Vogts--Werner equation~(22).  The latter follows after
    applying Lemma~7 to both tensor factors.  The displayed assertion is an
    inclusion in an actual tensor-product span, not merely the conjunction of
    the two one-sided inclusions.  The MPU appendix applies this construction
    to $\mathcal R_{2x}$ in
    \cite[lines~2320--2334]{Cirac2017MPU}.
\end{theorem}

\begin{proof}\leanok
    \uses{thm:qca_mem_kronecker_submodule_iff, thm:qca_left_coefficient_mem_support, thm:qca_right_coefficient_mem_support}
    If $X\in\mathcal A$, every left coefficient of $X$ lies in the left support
    algebra and every right coefficient lies in the right support algebra.
    Theorem~\ref{thm:qca_mem_kronecker_submodule_iff} places $X$ in their
    Kronecker-product submodule.
\end{proof}

\begin{theorem}[Left Kronecker commutation is coefficientwise]
    \label{thm:qca_left_kronecker_commute_iff}
    \lean{Matrix.commute_left_kronecker_embed_iff}
    \leanok
    For $B\in M_I(\mathbb C)$ and $X\in M_{I\times J}(\mathbb C)$,
    \begin{align}
        [B\otimes\mathbf 1_J,X]=0
        \quad\Longleftrightarrow\quad
        [B,X^{\mathrm L}_{rs}] & =0
        \quad\text{for every }r,s\in J.
    \end{align}
\end{theorem}

\begin{proof}\leanok
    For $i,j\in I$ and $r,s\in J$, the two product entries are
    \begin{align}
        [(B\otimes\mathbf 1_J)X]_{(i,r),(j,s)}
         & =\sum_{k\in I}B_{ik}X_{(k,r),(j,s)}
        =(BX^{\mathrm L}_{rs})_{ij},           \\
        [X(B\otimes\mathbf 1_J)]_{(i,r),(j,s)}
         & =\sum_{k\in I}X_{(i,r),(k,s)}B_{kj}
        =(X^{\mathrm L}_{rs}B)_{ij}.
    \end{align}
    Equality of the two products is therefore equivalent to
    $[B,X^{\mathrm L}_{rs}]=0$ for every $r,s\in J$.
\end{proof}

\begin{theorem}[Right Kronecker commutation is coefficientwise]
    \label{thm:qca_right_kronecker_commute_iff}
    \lean{Matrix.commute_right_kronecker_embed_iff}
    \leanok
    For $C\in M_J(\mathbb C)$ and $X\in M_{I\times J}(\mathbb C)$,
    \begin{align}
        [\mathbf 1_I\otimes C,X]=0
        \quad\Longleftrightarrow\quad
        [C,X^{\mathrm R}_{ij}] & =0
        \quad\text{for every }i,j\in I.
    \end{align}
\end{theorem}

\begin{proof}\leanok
    For $i,j\in I$ and $r,s\in J$, the two product entries are
    \begin{align}
        [(\mathbf 1_I\otimes C)X]_{(i,r),(j,s)}
         & =\sum_{t\in J}C_{rt}X_{(i,t),(j,s)}
        =(CX^{\mathrm R}_{ij})_{rs},           \\
        [X(\mathbf 1_I\otimes C)]_{(i,r),(j,s)}
         & =\sum_{t\in J}X_{(i,r),(j,t)}C_{ts}
        =(X^{\mathrm R}_{ij}C)_{rs}.
    \end{align}
    Equality of the two products is therefore equivalent to
    $[C,X^{\mathrm R}_{ij}]=0$ for every $i,j\in I$.
\end{proof}

\begin{theorem}[Left support-algebra centralizer]
    \label{thm:qca_left_support_centralizer}
    \lean{Matrix.left_kronecker_embed_mem_centralizer_iff}
    \leanok
    \uses{def:qca_left_support_algebra}
    For $B\in M_I(\mathbb C)$,
    \begin{align}
        B\otimes\mathbf 1_J\in\mathcal A'
        \quad\Longleftrightarrow\quad
        B & \in\operatorname{Spp}_{\mathrm L}(\mathcal A)'.
    \end{align}
    Equivalently,
    \begin{align}
        \operatorname{Spp}_{\mathrm L}(\mathcal A)'
         & =\{B\in M_I(\mathbb C)\mid B\otimes\mathbf 1_J\in\mathcal A'\}.
    \end{align}
\end{theorem}

\begin{proof}\leanok
    \uses{thm:qca_left_coefficient_mem_support, thm:qca_left_kronecker_commute_iff, def:qca_left_coefficient_set, def:qca_left_support_algebra}
    By Theorem~\ref{thm:qca_left_kronecker_commute_iff}, commutation of
    $B\otimes\mathbf 1_J$ with every $X\in\mathcal A$ is equivalent to
    commutation of $B$ with every left coefficient.  The coefficient set is
    closed under adjoints because $\mathcal A$ is a star-subalgebra.  The
    universal property of the generated star-subalgebra then extends this
    commutation relation to all of
    $\operatorname{Spp}_{\mathrm L}(\mathcal A)$; the converse follows from
    generator membership.
\end{proof}

\begin{theorem}[Right support-algebra centralizer]
    \label{thm:qca_right_support_centralizer}
    \lean{Matrix.right_kronecker_embed_mem_centralizer_iff}
    \leanok
    \uses{def:qca_right_support_algebra}
    For $C\in M_J(\mathbb C)$,
    \begin{align}
        \mathbf 1_I\otimes C\in\mathcal A'
        \quad\Longleftrightarrow\quad
        C & \in\operatorname{Spp}_{\mathrm R}(\mathcal A)'.
    \end{align}
    Equivalently,
    \begin{align}
        \operatorname{Spp}_{\mathrm R}(\mathcal A)'
         & =\{C\in M_J(\mathbb C)\mid \mathbf 1_I\otimes C\in\mathcal A'\}.
    \end{align}
\end{theorem}

\begin{proof}\leanok
    \uses{thm:qca_operator_block_adjoint, thm:qca_right_coefficient_mem_support, thm:qca_right_kronecker_commute_iff, def:qca_right_coefficient_set, def:qca_right_support_algebra}
    By Theorem~\ref{thm:qca_right_kronecker_commute_iff}, commutation of
    $\mathbf 1_I\otimes C$ with every $X\in\mathcal A$ is equivalent to
    commutation of $C$ with every right coefficient.  Theorem~\ref{thm:qca_operator_block_adjoint}
    shows that the right coefficient set is closed under adjoints.  The
    universal property of the generated star-subalgebra extends commutation to
    all of $\operatorname{Spp}_{\mathrm R}(\mathcal A)$; the converse follows
    from generator membership.
\end{proof}

\begin{theorem}[Commutation of overlapping support algebras]
    \label{thm:qca_overlapping_support_algebras_commute}
    \lean{Matrix.supportAlgebras_commute_of_overlappingLifts_commute}
    \leanok
    \uses{def:qca_left_support_algebra, def:qca_right_support_algebra}
    Let $I$, $J$, and $K$ be finite sets, and let
    \begin{align}
        \mathcal A_1 & \subseteq M_{I\times J}(\mathbb C), &
        \mathcal A_2 & \subseteq M_{J\times K}(\mathbb C)
    \end{align}
    be star-subalgebras.  Suppose that, for every
    $X\in\mathcal A_1$ and $Y\in\mathcal A_2$,
    \begin{align}
        (X\otimes\mathbf 1_K)(\mathbf 1_I\otimes Y)
         & =(\mathbf 1_I\otimes Y)(X\otimes\mathbf 1_K).
        \label{eq:qca_overlapping_support_lifts}
    \end{align}
    Then the two support algebras on the common factor commute:
    \begin{align}
        RL & =LR
           &     & \text{for every }
        R\in\operatorname{Spp}_{\mathrm R}(\mathcal A_1),\quad
        L\in\operatorname{Spp}_{\mathrm L}(\mathcal A_2).
        \label{eq:qca_overlapping_support_conclusion}
    \end{align}
    This is the full complex matrix-algebra specialization of
    Schumacher--Werner Lemma \texttt{sppcomm}
    \cite[lines~1174--1194]{SchumacherWerner2004QCA} and GNVW Lemma
    \texttt{sppcomm}
    \cite[Section~7.1, lines~1221--1246]{Gross2012IndexTheory}.

    \textbf{Local fixes.}
    For the generated $C^*$-algebras, the input must be closed under the
    involution; see \cite{gap:gnvw12_support_algebra_star_closure}.  In both
    printed coefficient proofs, the basis element $e'_\nu$ belongs to the
    third factor, rather than the second factor as printed; see
    \cite{gap:gnvw12_support_algebra_outer_factor_basis}.

    \textbf{Scope restriction (full matrix factors).}
    The cited lemmas permit arbitrary finite-dimensional $C^*$-algebra
    factors, whereas the statement above treats full complex matrix algebras.
    This restriction suffices for the MPU application in
    \cite[lines~2300--2306 and 2313--2334]{Cirac2017MPU}; the general-factor
    statement remains open as recorded in
    \cite{gap:gnvw12_support_algebra_full_matrix_scope}.
\end{theorem}

\begin{proof}\leanok
    \uses{def:qca_left_coefficient_set, def:qca_right_coefficient_set, def:qca_left_support_algebra, def:qca_right_support_algebra}
    For $X\in\mathcal A_1$, $Y\in\mathcal A_2$, $i,i'\in I$, and
    $k,k'\in K$, applying
    Theorem~\cite[``Middle-factor coefficients of commuting overlapping operators'']{QICLean2026Blueprint} to
    (\ref{eq:qca_overlapping_support_lifts}) gives
    \begin{align}
        X^{\mathrm R}_{ii'}Y^{\mathrm L}_{kk'}
         & =Y^{\mathrm L}_{kk'}X^{\mathrm R}_{ii'}.
        \label{eq:qca_overlapping_support_coefficients}
    \end{align}
    Applying the same identity to $Y^\dagger\in\mathcal A_2$, with $k$ and
    $k'$ exchanged, also gives commutation with
    $(Y^{\mathrm L}_{kk'})^\dagger$.  Hence every generator of
    $\operatorname{Spp}_{\mathrm R}(\mathcal A_1)$ lies in the
    star-subalgebra commuting with all left coefficients of $\mathcal A_2$.
    Passing to the generated star-subalgebra yields
    \begin{align}
        LR          & =RL,        &
        L^\dagger R & =RL^\dagger
        \label{eq:qca_overlapping_support_first_closure}
    \end{align}
    for every $R\in\operatorname{Spp}_{\mathrm R}(\mathcal A_1)$ and every
    left coefficient $L$ of $\mathcal A_2$.

    Since $\operatorname{Spp}_{\mathrm R}(\mathcal A_1)$ is closed under
    adjoints, applying the first equality in
    (\ref{eq:qca_overlapping_support_first_closure}) to $R^\dagger$ shows that
    each left coefficient commutes with both $R$ and $R^\dagger$.  Passing to
    the star-subalgebra generated by the left coefficients gives
    (\ref{eq:qca_overlapping_support_conclusion}).
\end{proof}

\subsection{Products of commuting matrix factors}

Let $\mathcal C$ be a nonzero complex star-algebra and let
$\mathcal A,\mathcal B\subseteq\mathcal C$ be unital star-subalgebras such
that $ab=ba$ for every $a\in\mathcal A$ and $b\in\mathcal B$.

\begin{definition}[Multiplication of commuting factors]
    \label{def:qca_commuting_factor_multiplication}
    \lean{StarSubalgebra.commutingMulMap,
        StarSubalgebra.commutingMulMap_tmul}
    \leanok
    Multiplication defines a unital star-algebra homomorphism
    \begin{align}
        \mu_{\mathcal A,\mathcal B}
        \colon \mathcal A\otimes_{\mathbb C}\mathcal B
         & \longrightarrow \mathcal C, \notag \\
        \mu_{\mathcal A,\mathcal B}(a\otimes b)
         & =ab.
    \end{align}
\end{definition}

\begin{theorem}[Algebraic product of commuting factors]
    \label{thm:qca_commuting_factor_product}
    \lean{StarSubalgebra.commutingMulMap_range,
        StarSubalgebra.sup_toSubmodule_eq_mul}
    \leanok
    \uses{def:qca_commuting_factor_multiplication}
    The image of $\mu_{\mathcal A,\mathcal B}$ is the algebraic join
    $\mathcal A\vee\mathcal B$.  Its underlying complex vector space is
    \begin{align}
        \mathcal A\mathcal B
         & =\operatorname{span}_{\mathbb C}
        \{ab\mid a\in\mathcal A,\ b\in\mathcal B\}.
        \label{eq:qca_commuting_factor_product_span}
    \end{align}
    Thus every element of the join is a finite complex-linear sum of
    products $ab$.
\end{theorem}

\begin{proof}\leanok
    \uses{def:qca_commuting_factor_multiplication}
    The involution reverses products.  Since $a^*\in\mathcal A$ and
    $b^*\in\mathcal B$ commute,
    \begin{align}
        \mu_{\mathcal A,\mathcal B}\bigl((a\otimes b)^*\bigr)
         & =a^*b^*=b^*a^*=(ab)^*.
    \end{align}
    The image contains both factors because $a=a\mathbf 1$ and
    $b=\mathbf 1b$.  Conversely, every pure tensor maps to an element of
    the join, and linearity gives the reverse inclusion.  The description
    of the underlying vector space is the range formula for the linear
    multiplication map.
\end{proof}

\begin{theorem}[Product of two full matrix factors]
    \label{thm:qca_commuting_full_matrix_product}
    \lean{StarSubalgebra.commutingMulMap_injective_of_equiv_matrix,
        StarSubalgebra.commutingSupEquivMatrix,
        StarSubalgebra.nonempty_equiv_sup_matrix,
        StarSubalgebra.finrank_sup_of_equiv_matrix}
    \leanok
    \uses{def:qca_commuting_factor_multiplication}
    Suppose that $\mathcal C=M_I(\mathbb C)$ for a finite nonempty set
    $I$, that $p,q>0$, and that
    \begin{align}
        \mathcal A & \cong M_p(\mathbb C), \notag \\
        \mathcal B & \cong M_q(\mathbb C)
    \end{align}
    as star-algebras.  Then $\mu_{\mathcal A,\mathcal B}$ is injective and
    \begin{align}
        \mathcal A\vee\mathcal B
         & \cong M_{pq}(\mathbb C), \notag \\
        \dim_{\mathbb C}(\mathcal A\vee\mathcal B)
         & =(pq)^2.
        \label{eq:qca_commuting_full_matrix_product}
    \end{align}
    \textbf{Scope restriction (full matrix, single block):}
    This is the unital single-block specialization of Schumacher--Werner
    Proposition~\texttt{Cscom}
    \cite[lines~2119--2138]{SchumacherWerner2004QCA}, used for the adjacent
    support algebras in
    \cite[lines~1270--1308]{Gross2012IndexTheory}.  It does not assert the
    direct-sum block-pair theorem or the multiplicity equation of
    Proposition~\texttt{Cscom}, nor that the displayed matrix algebra fills
    the ambient algebra.  The omitted statement and its elimination plan are
    recorded in \cite{gap:gnvw12_cscom_single_block_scope}.
\end{theorem}

\begin{proof}\leanok
    \uses{def:qca_commuting_factor_multiplication,
        thm:qca_commuting_factor_product}
    The source of the multiplication homomorphism is
    \begin{align}
        M_p(\mathbb C)\otimes_{\mathbb C}M_q(\mathbb C)
         & \cong M_{pq}(\mathbb C).
    \end{align}
    This is a simple ring.  The kernel of a unital homomorphism from a
    simple ring to a nonzero ring is not the whole ring, and is therefore
    zero.  Hence multiplication is injective.  Its image is the join by
    Theorem~\ref{thm:qca_commuting_factor_product}; the matrix presentation
    and dimension formula in
    Equation~(\ref{eq:qca_commuting_full_matrix_product}) follow.
\end{proof}

\begin{theorem}[Canonical Kronecker coordinates]
    \label{thm:qca_commuting_factor_kronecker_coordinates}
    \lean{Matrix.leftKroneckerEmbed_sup_rightKroneckerEmbed_toSubmodule}
    \leanok
    \uses{def:qca_kronecker_submodule}
    Let $I,J$ be finite sets, and let
    $\mathcal S\subseteq M_I(\mathbb C)$ and
    $\mathcal T\subseteq M_J(\mathbb C)$ be unital star-subalgebras.  In
    the canonical bipartite coordinates, the underlying complex vector
    space of the join of $\mathcal S\otimes\mathbf 1_J$ and
    $\mathbf 1_I\otimes\mathcal T$ is exactly
    \begin{align}
        \mathcal S\boxtimes\mathcal T
         & =\operatorname{span}_{\mathbb C}
        \{S\otimes T\mid S\in\mathcal S,\ T\in\mathcal T\}.
        \label{eq:qca_commuting_factor_kronecker_coordinates}
    \end{align}
\end{theorem}

\begin{proof}\leanok
    \uses{thm:qca_commuting_factor_product}
    If $I$ or $J$ is empty, the ambient matrix space is zero, so both
    subspaces in the claimed equality are zero.  Assume that both sets are
    nonempty.  The two canonical factors commute, and their products satisfy
    \begin{align}
        (S\otimes\mathbf 1_J)(\mathbf 1_I\otimes T)
         & =S\otimes T.
    \end{align}
    Equation~(\ref{eq:qca_commuting_factor_product_span}) therefore becomes
    exactly
    (\ref{eq:qca_commuting_factor_kronecker_coordinates}).
\end{proof}

\begin{definition}[Translation of a finite region]
    \label{def:qca_translate_region}
    \lean{SpinChain.translateRegion,SpinChain.mem_translateRegion,
        SpinChain.mem_translateRegion_add,SpinChain.translateRegion_zero,
        SpinChain.translateRegion_add,
        SpinChain.translateRegion_mono}
    \leanok
    For a finite region $\Lambda\subset\mathbb Z$ and $a\in\mathbb Z$, define
    \begin{align}
        \Lambda+a & =\{i+a\mid i\in\Lambda\}.
    \end{align}
    Thus
    \begin{align}
        j\in\Lambda+a & \iff j-a\in\Lambda, \\
        \Lambda+0     & =\Lambda,           \\
        (\Lambda+a)+b & =\Lambda+(a+b).
    \end{align}
    If $\Lambda\subseteq\Gamma$, then $\Lambda+a\subseteq\Gamma+a$.
\end{definition}

\begin{definition}[Sumset of finite regions]
    \label{def:qca_region_sumset}
    \lean{SpinChain.regionSumset,SpinChain.mem_regionSumset,
        SpinChain.regionSumset_singleton_right,SpinChain.regionSumset_mono,
        SpinChain.regionSumset_mono_left,SpinChain.regionSumset_mono_right,
        SpinChain.regionSumset_zero_left,SpinChain.regionSumset_zero_right,
        SpinChain.regionSumset_assoc}
    \leanok
    \uses{def:qca_translate_region}
    For finite regions $\Lambda,\mathcal N\subset\mathbb Z$, define their
    sumset by
    \begin{align}
        \Lambda+\mathcal N
         & =\{i+a\mid i\in\Lambda,\ a\in\mathcal N\}.
    \end{align}
    Hence $j\in\Lambda+\mathcal N$ exactly when $j=i+a$ for some
    $i\in\Lambda$ and $a\in\mathcal N$. The sumset is monotone in both
    arguments, and the singleton $\{0\}$ is an identity. A singleton
    displacement recovers translation,
    \begin{align}
        \Lambda+\{a\} & =\Lambda+a.
    \end{align}
    Sumsets are associative:
    \begin{align}
        (\Lambda+\mathcal N)+\mathcal M & =\Lambda+(\mathcal N+\mathcal M).
    \end{align}
    These identities support composition of propagation bounds. This is the
    finite region $\Lambda+\mathcal N$ in the propagation condition
    of the QCA appendix of \cite{Cirac2017MPU}.
\end{definition}

\begin{definition}[Translation of sites and configurations]
    \label{def:qca_finite_configuration_translation}
    \lean{SpinChain.siteTranslation_apply_val,
        SpinChain.siteTranslation_symm_apply_val,SpinChain.Config.translation,
        SpinChain.Config.translation_apply,SpinChain.Config.translation_symm_apply}
    \leanok
    \uses{def:qca_local_algebra, def:qca_translate_region}
    Translation gives bijections
    \begin{align}
        t_{a,\Lambda} & \colon\Lambda\longrightarrow\Lambda+a,
                      & t_{a,\Lambda}(i)                       & =i+a,    \\
        T_{a,\Lambda} & \colon\operatorname{Config}_d(\Lambda)
        \longrightarrow\operatorname{Config}_d(\Lambda+a),
                      & (T_{a,\Lambda}x)(j)                    & =x(j-a).
    \end{align}
    Their inverses translate sites and configurations by $-a$.
\end{definition}

\begin{definition}[Finite-region translation of local observables]
    \label{def:qca_finite_local_translation}
    \lean{SpinChain.localTranslation,SpinChain.localTranslation_apply}
    \leanok
    \uses{def:qca_local_algebra, def:qca_finite_configuration_translation}
    Relabelling configurations defines a star-algebra equivalence
    \begin{align}
        \tau_{a,\Lambda} & \colon\mathcal A_\Lambda
        \longrightarrow\mathcal A_{\Lambda+a},                          \\
        \tau_{a,\Lambda}(A)(x,y)
                         & =A(T_{a,\Lambda}^{-1}x,T_{a,\Lambda}^{-1}y).
    \end{align}
    This is the finite-region translation associated with the lattice
    translation $i\mapsto i+a$.
\end{definition}

\begin{lemma}[Norm preservation of finite-region translation]
    \label{lem:qca_finite_local_translation_norm}
    \lean{SpinChain.localTranslation_norm,
        SpinChain.localTranslation_isometry}
    \leanok
    \uses{def:qca_finite_local_translation}
    Let $d$ be a nonnegative integer, let $\Lambda\subset\mathbb Z$ be finite,
    and let $a\in\mathbb Z$. For every $A\in\mathcal A_\Lambda$,
    \begin{align}
        \lVert\tau_{a,\Lambda}(A)\rVert & =\lVert A\rVert.
    \end{align}
    Consequently, $\tau_{a,\Lambda}$ is an isometry from
    $\mathcal A_\Lambda$ onto $\mathcal A_{\Lambda+a}$.
\end{lemma}

\begin{proof}\leanok
    \uses{def:qca_finite_local_translation}
    Translation bijectively relabels the configuration basis. The resulting
    star-algebra equivalence between the two finite matrix algebras preserves
    the operator norm. Applying this identity to $A-B$ and using linearity gives
    \begin{align}
        \lVert\tau_{a,\Lambda}(A)-\tau_{a,\Lambda}(B)\rVert
         & =\lVert\tau_{a,\Lambda}(A-B)\rVert
        =\lVert A-B\rVert,
    \end{align}
    so distances are preserved.
\end{proof}

\begin{lemma}[Finite-region translation action laws]
    \label{lem:qca_finite_translation_action}
    \lean{SpinChain.siteTranslation_zero,SpinChain.siteTranslation_add,
        SpinChain.Config.translation_apply_siteTranslation,
        SpinChain.Config.translation_zero,SpinChain.Config.translation_add,
        SpinChain.localTranslation_apply_translation,
        SpinChain.localTranslation_zero,SpinChain.localTranslation_add}
    \leanok
    \uses{def:qca_finite_configuration_translation, def:qca_finite_local_translation}
    Translation by zero acts trivially on sites, configurations, and
    finite-region observables. Successive translations by $a$ and $b$ agree
    with translation by $a+b$. Pointwise, for $i\in\Lambda$,
    $x,y\in\operatorname{Config}_d(\Lambda)$, and
    $A\in\mathcal A_\Lambda$,
    \begin{align}
        t_{0,\Lambda}(i)                   & =i,                                      \\
        t_{b,\Lambda+a}(t_{a,\Lambda}(i))
                                           & =t_{a+b,\Lambda}(i),                     \\
        (T_{a,\Lambda}x)(t_{a,\Lambda}(i)) & =x(i),                                   \\
        (T_{b,\Lambda+a}T_{a,\Lambda}x)
        (t_{b,\Lambda+a}(t_{a,\Lambda}(i)))
                                           & =(T_{a+b,\Lambda}x)(t_{a+b,\Lambda}(i)), \\
        (\tau_{a,\Lambda}(A))(T_{a,\Lambda}x,T_{a,\Lambda}y)
                                           & =A(x,y),                                 \\
        (\tau_{b,\Lambda+a}\tau_{a,\Lambda}(A))
        (T_{b,\Lambda+a}T_{a,\Lambda}x,
        T_{b,\Lambda+a}T_{a,\Lambda}y)
                                           & =(\tau_{a+b,\Lambda}(A))
        (T_{a+b,\Lambda}x,T_{a+b,\Lambda}y).
    \end{align}
    These identities express the additive action without identifying the
    distinct configuration spaces before applying their canonical translation
    bijections.
\end{lemma}

\begin{proof}\leanok
    The site identities follow from
    \begin{align}
        (i+a)+b & =i+(a+b).
    \end{align}
    Relabelling a configuration is precomposition by the inverse site
    translation, so evaluation at the corresponding translated site gives
    \begin{align}
        (T_{a,\Lambda}x)(t_{a,\Lambda}(i)) & =x(i).
    \end{align}
    Applying this identity successively proves the configuration action laws.
    Finally, the local observable translation relabels both matrix indices by
    $T_{a,\Lambda}^{-1}$; evaluating on translated configurations therefore
    gives $A(x,y)$, and applying this entrywise identity successively proves the
    local observable action laws.
\end{proof}

\begin{lemma}[Naturality of finite-region translation]
    \label{lem:qca_finite_local_translation_naturality}
    \lean{SpinChain.localTranslation_localInclusion}
    \leanok
    \uses{def:qca_local_inclusion, def:qca_translate_region, def:qca_finite_local_translation}
    If $\Lambda\subseteq\Gamma$, then for every
    $A\in\mathcal A_\Lambda$,
    \begin{align}
        \tau_{a,\Gamma}(\iota_{\Lambda,\Gamma}(A))
         & =\iota_{\Lambda+a,\Gamma+a}(\tau_{a,\Lambda}(A)).
    \end{align}
\end{lemma}

\begin{proof}\leanok
    The lattice translation identifies $\Gamma\setminus\Lambda$ with
    $(\Gamma+a)\setminus(\Lambda+a)$. Hence relabelling configurations
    carries the identity on the first complement to the identity on the
    second complement, and
    \begin{align}
        \tau_{a,\Gamma}(A\otimes\mathbf 1_{\Gamma\setminus\Lambda})
         & =\tau_{a,\Lambda}(A)
        \otimes\mathbf 1_{(\Gamma+a)\setminus(\Lambda+a)}.
    \end{align}
\end{proof}

\begin{definition}[Algebraic local algebra]
    \label{def:qca_algebraic_local_algebra}
    \lean{SpinChain.instDirectedSystemLocalAlgebra,
        SpinChain.AlgebraicLocalAlgebra,SpinChain.localObservable,
        SpinChain.localObservable_localInclusion}
    \leanok
    \uses{def:qca_local_algebra, def:qca_local_inclusion, lem:qca_local_inclusion_refl, lem:qca_local_inclusion_trans}
    The finite-region local algebras form a directed system under the maps
    $\iota_{\Lambda,\Gamma}$.  Its algebraic direct limit is the algebraic
    local algebra
    \begin{align}
        \mathcal A_{\mathrm{loc}}
         & =\varinjlim_{\Lambda\Subset\mathbb Z}\mathcal A_\Lambda.
    \end{align}
    For every finite region $\Lambda$, write
    $\iota_\Lambda\colon\mathcal A_\Lambda\to\mathcal A_{\mathrm{loc}}$
    for the canonical inclusion.  Whenever $\Lambda\subseteq\Gamma$, these
    inclusions satisfy
    \begin{align}
        \iota_\Gamma\circ\iota_{\Lambda,\Gamma} & =\iota_\Lambda.
    \end{align}
    This is the union
    $\bigcup_{\Lambda\Subset\mathbb Z}\mathcal A_\Lambda$ in the Appendix of
    \cite{Cirac2017MPU}.  No norm completion is taken.
\end{definition}

\begin{definition}[Finite-chain realization of a periodic MPU]
    \label{def:mpu_finite_chain_realization}
    \lean{SpinChain.finiteChainRegion,SpinChain.card_finiteChainRegion,
        SpinChain.finiteChainSiteEquiv,SpinChain.finiteChainConfigEquiv,
        SpinChain.finiteChainMPOMatrix}
    \leanok
    \uses{def:mpo_family, def:qca_local_algebra}
    For $a\in\mathbb Z$ and $N\in\mathbb N$, let
    $I_{a,N}=[a,a+N)\cap\mathbb Z$.  The increasing bijection
    $\{0,\ldots,N-1\}\simeq I_{a,N}$ identifies the periodic configuration
    basis with $\operatorname{Config}_d(I_{a,N})$.  In this basis, write
    $U_{a}^{(N)}$ for the matrix obtained by simultaneously reindexing the rows
    and columns of $U^{(N)}$.
\end{definition}

\begin{definition}[Finite-chain MPU conjugation]
    \label{def:mpu_finite_chain_conjugation}
    \lean{SpinChain.finiteChainMPUUnitary,
        SpinChain.finiteChainConjugation,
        SpinChain.finiteChainConjugation_apply}
    \leanok
    \uses{def:mpu, def:mpu_finite_chain_realization}
    If $\mathcal U$ is an MPU and $N>1$, then $U_a^{(N)}$ is unitary and
    defines the star-algebra automorphism
    \begin{align}
        \operatorname{Ad}_{U_a^{(N)}}(B)
         & =U_a^{(N)}B\,U_a^{(N)\dagger}.
    \end{align}
    The factor order is the one in
    \cite[lines~2300--2306]{Cirac2017MPU}.
\end{definition}

\begin{proof}\leanok
    \uses{def:mpu, thm:mpu_reindex_unitary}
    Simultaneous reindexing preserves both unitarity equations for $U^{(N)}$.
    Conjugation by the resulting unitary preserves multiplication, the unit,
    complex scalars, and the adjoint.
\end{proof}

\begin{definition}[Finite-chain term on a local observable]
    \label{def:mpu_finite_chain_local_observable}
    \lean{SpinChain.finiteChainLocalObservable,
        SpinChain.finiteChainLocalObservable_apply}
    \leanok
    \uses{def:qca_local_inclusion, def:qca_algebraic_local_algebra, def:mpu_finite_chain_conjugation}
    Let $\Lambda\subseteq I_{a,N}$ and $A\in\mathcal A_\Lambda$.  Set
    $A_L=\iota_{\Lambda,I_{a,N}}(A)$.  For $N>1$, the finite-chain term in
    \cite[lines~2300--2306]{Cirac2017MPU} is the star-algebra map
    \begin{align}
        \mathcal A_\Lambda & \longrightarrow\mathcal A_{\mathrm{loc}}, \\
        A                  & \longmapsto
        \iota_{I_{a,N}}\!\left(U_a^{(N)}A_LU_a^{(N)\dagger}\right).
    \end{align}
    No stabilization in $N$ is asserted.
\end{definition}

\begin{definition}[Operator norm on the local algebra]
    \label{def:qca_algebraic_local_operator_norm}
    \lean{SpinChain.algebraicLocalOperatorNorm,
        SpinChain.algebraicLocalOperatorNorm_localObservable,
        SpinChain.instNormAlgebraicLocalAlgebra,
        SpinChain.norm_localObservable}
    \leanok
    \uses{def:qca_algebraic_local_algebra, lem:qca_local_inclusion_norm}
    Assume $d>0$.  If a local observable $A\in\mathcal A_{\mathrm{loc}}$
    is represented by $A_\Lambda\in\mathcal A_\Lambda$ on a finite region
    $\Lambda$, define its operator norm by
    \begin{align}
        \lVert A\rVert & =\lVert A_\Lambda\rVert.
    \end{align}
    This does not depend on the finite region or representative.  It is the
    norm whose topology is completed in the Appendix of
    \cite{Cirac2017MPU} to obtain the quasi-local algebra.
\end{definition}

\begin{proof}\leanok
    If $A_\Lambda$ and $A_\Gamma$ represent the same local observable, then
    the defining relation for the inductive limit gives a finite region
    $\Delta\supseteq\Lambda\cup\Gamma$ such that
    \begin{align}
        \iota_{\Lambda,\Delta}(A_\Lambda)
         & =\iota_{\Gamma,\Delta}(A_\Gamma).
    \end{align}
    Both inclusions preserve the operator norm, so
    $\lVert A_\Lambda\rVert=\lVert A_\Gamma\rVert$.
\end{proof}

\begin{definition}[Normed star algebra of local observables]
    \label{def:qca_algebraic_local_normed_star_algebra}
    \lean{SpinChain.instNormedAddCommGroupAlgebraicLocalAlgebra,
        SpinChain.instNormedRingAlgebraicLocalAlgebra,
        SpinChain.instNormedAlgebraAlgebraicLocalAlgebra}
    \leanok
    \uses{def:qca_algebraic_local_operator_norm}
    For $d>0$, the local algebra with the operator norm is a normed complex
    star algebra.  No completeness assertion is made.
\end{definition}

\begin{proof}\leanok
    Choose finite-region representatives of the local observables and enlarge
    their regions to a common finite region.  The norm axioms, the product
    inequality, and compatibility with complex scalar multiplication then
    follow from the corresponding properties of the finite-region matrix
    algebra.
\end{proof}

\begin{lemma}[C-star identity for local observables]
    \label{lem:qca_algebraic_local_norm_star_mul_self}
    \lean{SpinChain.norm_star_mul_self}
    \leanok
    \uses{def:qca_algebraic_local_normed_star_algebra}
    For $d>0$ and every local observable $A$,
    \begin{align}
        \lVert A^*A\rVert & =\lVert A\rVert^2.
    \end{align}
\end{lemma}

\begin{proof}\leanok
    Choose a finite region $\Lambda$ and a representative
    $A_\Lambda\in\mathcal A_\Lambda$.  Then
    \begin{align}
        \lVert A^*A\rVert
         & =\lVert A_\Lambda^*A_\Lambda\rVert
        =\lVert A_\Lambda\rVert^2
        =\lVert A\rVert^2.
    \end{align}
\end{proof}

\begin{definition}[C-star axiom structure on local observables]
    \label{def:qca_algebraic_local_cstar_ring}
    \lean{SpinChain.instCStarRingAlgebraicLocalAlgebra}
    \leanok
    \uses{lem:qca_algebraic_local_norm_star_mul_self}
    The normed star algebra of local observables satisfies the C-star axiom.
    This structure does not include completeness.
\end{definition}

\begin{proof}\leanok
    The C-star identity gives the required norm inequality in the reverse
    direction; submultiplicativity gives the other direction.
\end{proof}

\begin{lemma}[Local observables embed isometrically]
    \label{lem:qca_local_observable_isometry}
    \lean{SpinChain.localObservable_isometry}
    \leanok
    \uses{def:qca_algebraic_local_operator_norm}
    For $d>0$ and every finite region $\Lambda$, the canonical map
    \begin{align}
        \iota_\Lambda\colon\mathcal A_\Lambda
         & \longrightarrow\mathcal A_{\mathrm{loc}}
    \end{align}
    is an operator-norm isometry.
\end{lemma}

\begin{proof}\leanok
    \uses{def:qca_algebraic_local_operator_norm}
    The norm of $\iota_\Lambda(A)$ is defined to be the operator norm of its
    representative $A\in\mathcal A_\Lambda$.  Hence
    $\lVert\iota_\Lambda(A)\rVert=\lVert A\rVert$; applying this equality to
    differences proves preservation of distance.
\end{proof}

\begin{theorem}[Star-algebra universal property]
    \label{thm:qca_algebraic_local_algebra_universal}
    \lean{SpinChain.algebraicLocalAlgebraLift,
        SpinChain.algebraicLocalAlgebraLift_localObservable,
        SpinChain.algebraicLocalAlgebra_starAlgHom_ext,
        SpinChain.algebraicLocalAlgebraLift_unique}
    \leanok
    \uses{def:qca_algebraic_local_algebra}
    Let $B$ be a complex star algebra.  Suppose that for every finite region
    $\Lambda$ there is a star-algebra homomorphism
    $g_\Lambda\colon\mathcal A_\Lambda\to B$.  Whenever
    $\Lambda\subseteq\Gamma$, suppose that
    \begin{align}
        g_\Gamma\circ\iota_{\Lambda,\Gamma} & =g_\Lambda.
    \end{align}
    Then there is a unique star-algebra homomorphism
    $g\colon\mathcal A_{\mathrm{loc}}\to B$ satisfying
    $g\circ\iota_\Lambda=g_\Lambda$ for every finite region $\Lambda$.
\end{theorem}

\begin{proof}\leanok
    Define $g$ on a local observable represented by
    $A\in\mathcal A_\Lambda$ as $g_\Lambda(A)$.  Compatibility makes this
    independent of the representing region.  Since each $g_\Lambda$ preserves
    addition, multiplication, the unit, scalar multiplication, and the star
    operation, so does $g$.  Every element of the direct limit has a
    finite-region representative, which also proves uniqueness.
\end{proof}

\begin{definition}[Support of a local observable]
    \label{def:qca_supported_in}
    \lean{SpinChain.SupportedIn}
    \leanok
    \uses{def:qca_algebraic_local_algebra}
    A local observable $A\in\mathcal A_{\mathrm{loc}}$ is supported in the
    finite region $\Lambda$ if it lies in the image of
    $\iota_\Lambda\colon\mathcal A_\Lambda\to\mathcal A_{\mathrm{loc}}$.
\end{definition}

\begin{theorem}[Finite supports of local observables]
    \label{thm:qca_local_observable_finite_support}
    \lean{SpinChain.localObservable_injective,SpinChain.SupportedIn.mono,
        SpinChain.SupportedIn.union_left,SpinChain.SupportedIn.union_right,
        SpinChain.SupportedIn.add,SpinChain.SupportedIn.mul,
        SpinChain.SupportedIn.star,SpinChain.SupportedIn.zero,
        SpinChain.SupportedIn.one,SpinChain.exists_supportedIn}
    \leanok
    \uses{def:qca_supported_in, lem:qca_local_inclusion_injective}
    Every local observable is supported in some finite region.  If $A$ is
    supported in $\Lambda$ and $\Lambda\subseteq\Gamma$, then $A$ is supported
    in $\Gamma$.  If $A$ and $B$ are supported in $\Lambda$ and $\Gamma$,
    respectively, then $A+B$ and $AB$ are supported in $\Lambda\cup\Gamma$.
    The adjoint $A^*$ is supported in $\Lambda$, and both $0$ and $1$ are
    supported in every finite region.  For positive on-site dimension, each
    map $\iota_\Lambda$ is injective.
\end{theorem}

\begin{proof}\leanok
    \uses{def:qca_algebraic_local_algebra, lem:qca_local_inclusion_injective}
    Enlarge representatives by tensoring with the identity on the added sites.
    Two representatives can therefore be moved to the finite union of their
    regions, where sums and products are formed.  The canonical inclusions are
    unital star-algebra homomorphisms, so they preserve the adjoint, zero, and
    the unit.  For $d>0$, every transition map between finite-region algebras is
    injective; the canonical map from each member of such a directed system to
    its direct limit is therefore injective.  Finally, every element of an
    algebraic direct limit has a representative in one member of the directed
    system.
\end{proof}

\begin{theorem}[Disjointly supported local observables commute]
    \label{thm:qca_disjoint_supported_local_commute}
    \lean{SpinChain.localObservable_commute_of_disjoint,
        SpinChain.SupportedIn.commute_of_disjoint}
    \leanok
    \uses{def:qca_supported_in}
    If $X,Y\in\mathcal A_{\mathrm{loc}}$ are supported in finite regions
    $\Lambda$ and $\Gamma$ with $\Lambda\cap\Gamma=\varnothing$, then
    \begin{align}
        XY & =YX.
    \end{align}
    In particular, for $A\in\mathcal A_\Lambda$ and
    $B\in\mathcal A_\Gamma$,
    \begin{align}
        \iota_\Lambda(A)\iota_\Gamma(B)
         & =\iota_\Gamma(B)\iota_\Lambda(A).
    \end{align}
    This is the algebraic-local form of the tensor-factor identity underlying
    the construction in \cite[lines~2292--2300]{Cirac2017MPU}, not a named
    theorem of the appendix.
    % Source context: paper_v2.tex 2292--2300.
\end{theorem}

\begin{proof}\leanok
    \uses{thm:qca_disjoint_local_inclusion_commute}
    Choose representatives $X=\iota_\Lambda(A)$ and
    $Y=\iota_\Gamma(B)$.  Embed both into
    $\mathcal A_{\Lambda\cup\Gamma}$.  The preceding finite-region identity
    gives
    \begin{align}
        XY
         & =\iota_{\Lambda\cup\Gamma}
        \bigl(\iota_{\Lambda,\Lambda\cup\Gamma}(A)
        \iota_{\Gamma,\Lambda\cup\Gamma}(B)\bigr) \\
         & =\iota_{\Lambda\cup\Gamma}
        \bigl(\iota_{\Gamma,\Lambda\cup\Gamma}(B)
        \iota_{\Lambda,\Lambda\cup\Gamma}(A)\bigr)
        =YX.
    \end{align}
\end{proof}

\begin{theorem}[Algebraic-local relative commutant]
    \label{thm:qca_algebraic_local_relative_commutant}
    \lean{SpinChain.supportedIn_iff_commute_disjoint}
    \leanok
    \uses{def:qca_supported_in}
    Assume $d>0$. A local observable $X\in\mathcal A_{\mathrm{loc}}$ is
    supported in a finite region $\Lambda$ if and only if it commutes with every
    local observable having a finite support disjoint from $\Lambda$:
    \begin{align}
        X\in\iota_\Lambda(\mathcal A_\Lambda)
         & \iff
        \forall \Gamma,\ \forall Y,\quad
        \bigl(\Gamma\cap\Lambda=\varnothing\ \land\
        Y\in\iota_\Gamma(\mathcal A_\Gamma)\bigr)\Longrightarrow XY=YX.
    \end{align}
    This is an algebraic-local statement. It does not assert the corresponding
    characterization in the norm completion and is not stated by
    Schumacher--Werner. Their Theorem~6 gives the generalized Margolus structure,
    and Corollary~7 says that the inverse of a nearest-neighbor QCA exists and is a
    nearest-neighbor QCA.
\end{theorem}

\begin{proof}\leanok
    \uses{thm:qca_finite_relative_commutant, thm:qca_disjoint_supported_local_commute, thm:qca_local_observable_finite_support}
    The forward implication is commutation of disjointly supported observables.
    Conversely, choose a finite support $\Delta$ of $X$ and represent $X$ by
    $A\in\mathcal A_{\Lambda\cup\Delta}$. The hypothesis says that $A$ commutes
    with every observable on $(\Lambda\cup\Delta)\setminus\Lambda$. Since
    $d>0$, the canonical map from $\mathcal A_{\Lambda\cup\Delta}$ into
    $\mathcal A_{\mathrm{loc}}$ is injective, so this commutation already holds
    in the finite-region algebra. The finite relative-commutant theorem gives
    $A=\iota_{\Lambda,\Lambda\cup\Delta}(A_\Lambda)$ for some
    $A_\Lambda\in\mathcal A_\Lambda$. Hence
    $X=\iota_\Lambda(A_\Lambda)$.
\end{proof}

\begin{definition}[Algebraic-local translation homomorphism]
    \label{def:qca_algebraic_local_translation_hom}
    \lean{SpinChain.algebraicLocalTranslationHom,
        SpinChain.algebraicLocalTranslationHom_localObservable}
    \leanok
    \uses{def:qca_finite_local_translation, def:qca_algebraic_local_algebra}
    For $a\in\mathbb Z$, finite-region translations induce a star-algebra
    homomorphism
    \begin{align}
        \tau_a\colon\mathcal A_{\mathrm{loc}} & \longrightarrow
        \mathcal A_{\mathrm{loc}}
    \end{align}
    determined on every finite-region representative by
    \begin{align}
        \tau_a(\iota_\Lambda(A))
         & =\iota_{\Lambda+a}(\tau_{a,\Lambda}(A)).
    \end{align}
\end{definition}

\begin{proof}\leanok
    \uses{def:qca_translate_region, def:qca_algebraic_local_algebra, lem:qca_finite_local_translation_naturality, thm:qca_algebraic_local_algebra_universal}
    Naturality gives, for $\Lambda\subseteq\Gamma$,
    \begin{align}
        \iota_{\Gamma+a}\!\left(\tau_{a,\Gamma}(\iota_{\Lambda,\Gamma}(A))\right)
         & =\iota_{\Gamma+a}\!\left(\iota_{\Lambda+a,\Gamma+a}(\tau_{a,\Lambda}(A))\right) \\
         & =\iota_{\Lambda+a}(\tau_{a,\Lambda}(A)).
    \end{align}
    Hence the finite-region maps form a compatible family and induce the stated
    homomorphism on the algebraic direct limit.
\end{proof}

\begin{lemma}[Action laws for algebraic-local translation homomorphisms]
    \label{lem:qca_algebraic_local_translation_hom_action}
    \lean{SpinChain.algebraicLocalTranslationHom_zero,
        SpinChain.algebraicLocalTranslationHom_add}
    \leanok
    \uses{def:qca_algebraic_local_translation_hom}
    The homomorphisms satisfy
    \begin{align}
        \tau_0            & =\operatorname{id}, \\
        \tau_b\circ\tau_a & =\tau_{a+b}.
    \end{align}
\end{lemma}

\begin{proof}\leanok
    \uses{def:qca_translate_region, def:qca_algebraic_local_algebra, lem:qca_finite_translation_action, thm:qca_algebraic_local_algebra_universal}
    On a representative $A\in\mathcal A_\Lambda$, the two sides of the
    composition law are represented on the equal finite regions
    $(\Lambda+a)+b$ and $\Lambda+(a+b)$. The finite-region action identity
    therefore gives
    \begin{align}
        \tau_b(\tau_a(\iota_\Lambda(A)))
         & =\iota_{\Lambda+(a+b)}
        (\tau_{a+b,\Lambda}(A))
        =\tau_{a+b}(\iota_\Lambda(A)).
    \end{align}
    The zero law is proved in the same way. Equality on all finite-region
    representatives determines equality on the algebraic direct limit.
\end{proof}

\begin{definition}[Translation of algebraic-local observables]
    \label{def:qca_algebraic_local_translation}
    \lean{SpinChain.algebraicLocalTranslation,
        SpinChain.algebraicLocalTranslation_apply,
        SpinChain.algebraicLocalTranslation_localObservable}
    \leanok
    \uses{def:qca_algebraic_local_translation_hom}
    For every $a\in\mathbb Z$, the homomorphism $\tau_a$ is a star-algebra
    automorphism of $\mathcal A_{\mathrm{loc}}$, and on finite-region
    representatives
    \begin{align}
        \tau_a(\iota_\Lambda(A))
         & =\iota_{\Lambda+a}(\tau_{a,\Lambda}(A)).
    \end{align}
\end{definition}

\begin{proof}\leanok
    \uses{def:qca_algebraic_local_translation_hom, lem:qca_algebraic_local_translation_hom_action}
    The homomorphism action laws give
    \begin{align}
        \tau_{-a}\circ\tau_a & =\tau_0=\operatorname{id}, \\
        \tau_a\circ\tau_{-a} & =\tau_0=\operatorname{id}.
    \end{align}
    Thus $\tau_a$ is bijective with inverse $\tau_{-a}$. The formula on
    finite-region representatives is the defining formula for the homomorphism.
\end{proof}

\begin{lemma}[Algebraic-local translation action laws]
    \label{lem:qca_algebraic_local_translation_action}
    \lean{SpinChain.algebraicLocalTranslation_zero,
        SpinChain.algebraicLocalTranslation_add,
        SpinChain.algebraicLocalTranslation_symm}
    \leanok
    \uses{def:qca_algebraic_local_translation}
    The automorphisms form an additive action of $\mathbb Z$:
    \begin{align}
        \tau_0            & =\operatorname{id}, \\
        \tau_b\circ\tau_a & =\tau_{a+b},        \\
        \tau_a^{-1}       & =\tau_{-a}.
    \end{align}
\end{lemma}

\begin{proof}\leanok
    \uses{def:qca_algebraic_local_translation, lem:qca_algebraic_local_translation_hom_action}
    The first two identities are the homomorphism action laws expressed as
    identities of star-algebra automorphisms. Taking $b=-a$ in the composition
    law gives
    \begin{align}
        \tau_{-a}\circ\tau_a & =\tau_0=\operatorname{id}, \\
        \tau_a\circ\tau_{-a} & =\tau_0=\operatorname{id},
    \end{align}
    which is the inverse identity.
\end{proof}

\begin{lemma}[Exact translation of algebraic-local supports]
    \label{lem:qca_algebraic_local_translation_support}
    \lean{SpinChain.algebraicLocalTranslation_supportedIn,
        SpinChain.algebraicLocalTranslation_supportedIn_iff}
    \leanok
    \uses{def:qca_supported_in, def:qca_algebraic_local_translation}
    A local observable $A$ is supported in a finite region $\Lambda$ if and only
    if $\tau_a(A)$ is supported in $\Lambda+a$.
\end{lemma}

\begin{proof}\leanok
    \uses{def:qca_translate_region, def:qca_algebraic_local_translation, lem:qca_algebraic_local_translation_hom_action}
    If $A=\iota_\Lambda(A_\Lambda)$, then
    \begin{align}
        \tau_a(A)
         & =\iota_{\Lambda+a}(\tau_{a,\Lambda}(A_\Lambda)),
    \end{align}
    which proves the forward implication. Apply the same argument to
    $\tau_a(A)$ with displacement $-a$, and use
    $\tau_{-a}\tau_a=\operatorname{id}$ and
    $(\Lambda+a)-a=\Lambda$, for the reverse implication.
\end{proof}

\begin{lemma}[Norm isometry of algebraic-local translation]
    \label{lem:qca_algebraic_local_translation_isometry}
    \lean{SpinChain.norm_algebraicLocalTranslation,
        SpinChain.algebraicLocalTranslation_isometry}
    \leanok
    \uses{def:qca_algebraic_local_operator_norm, def:qca_algebraic_local_translation}
    If $d>0$, then every translation preserves the operator norm and distance:
    \begin{align}
        \lVert\tau_a(A)\rVert           & =\lVert A\rVert,   \\
        \lVert\tau_a(A)-\tau_a(B)\rVert & =\lVert A-B\rVert.
    \end{align}
\end{lemma}

\begin{proof}\leanok
    \uses{def:qca_algebraic_local_operator_norm, def:qca_algebraic_local_translation, lem:qca_finite_local_translation_norm, thm:qca_local_observable_finite_support}
    Choose $A_\Lambda\in\mathcal A_\Lambda$ with
    $A=\iota_\Lambda(A_\Lambda)$. Finite-region translation is a star-algebra
    equivalence of matrix algebras, hence preserves the operator norm. Therefore
    \begin{align}
        \lVert\tau_a(A)\rVert
         & =\lVert\tau_{a,\Lambda}(A_\Lambda)\rVert
        =\lVert A_\Lambda\rVert
        =\lVert A\rVert.
    \end{align}
    Applying this equality to $A-B$ gives distance preservation.
\end{proof}

\begin{definition}[C-star structure on a completion]
    \label{def:completion_cstar_algebra}
    \lean{UniformSpace.Completion.toComplAlgHom,
        UniformSpace.Completion.toComplStarAlgHom,
        UniformSpace.Completion.coe_toComplStarAlgHom,
        UniformSpace.Completion.instStar,UniformSpace.Completion.coe_star,
        UniformSpace.Completion.instContinuousStar,
        UniformSpace.Completion.instStarRing,
        UniformSpace.Completion.instCStarRing,
        UniformSpace.Completion.instStarModule,
        UniformSpace.Completion.instCStarAlgebra}
    \leanok
    The norm completion of a complex normed star algebra satisfying the C-star
    identity carries the continuously extended involution and a C-star algebra
    structure. The canonical map into the completion is a star-algebra
    homomorphism.
\end{definition}

\begin{proof}\leanok
    Extend the isometric involution continuously. The algebraic identities and
    the C-star identity hold on the dense canonical image and hence on the
    completion.
\end{proof}

\begin{definition}[Quasi-local observable algebra]
    \label{def:qca_quasi_local_algebra}
    \lean{SpinChain.QuasiLocalAlgebra,SpinChain.algebraicToQuasiLocal}
    \leanok
    \uses{def:qca_algebraic_local_cstar_ring, def:qca_algebraic_local_normed_star_algebra, def:completion_cstar_algebra}
    For $d>0$, the quasi-local observable algebra is the operator-norm
    completion
    \begin{align}
        \mathcal A
         & =\overline{\mathcal A_{\mathrm{loc}}}^{\lVert\cdot\rVert}.
    \end{align}
    The involution on $\mathcal A$ is the unique continuous extension of the
    involution on $\mathcal A_{\mathrm{loc}}$.  It makes $\mathcal A$ a complex
    C-star algebra, and the canonical star-algebra homomorphism
    \begin{align}
        j\colon\mathcal A_{\mathrm{loc}} & \longrightarrow\mathcal A
    \end{align}
    is the completion map.
\end{definition}

\begin{proof}\leanok
    The involution is isometric on $\mathcal A_{\mathrm{loc}}$, hence extends
    continuously to the completion.  The identities
    $(A+B)^*=A^*+B^*$, $(AB)^*=B^*A^*$, and $(A^*)^*=A$
    hold on the dense subalgebra $j(\mathcal A_{\mathrm{loc}})$ and therefore
    on $\mathcal A$.  The conjugate-linear identity
    $(cA)^*=\overline cA^*$ extends by the same argument.  Density also extends
    \begin{align}
        \lVert A^*A\rVert & =\lVert A\rVert^2,
    \end{align}
    so the completed algebra satisfies the C-star identity.
\end{proof}

\begin{definition}[Finite-region quasi-local observable]
    \label{def:qca_quasi_local_observable}
    \lean{SpinChain.quasiLocalObservable}
    \leanok
    \uses{def:qca_quasi_local_algebra, def:qca_algebraic_local_algebra}
    For every finite region $\Lambda$, define
    \begin{align}
        j_\Lambda\colon\mathcal A_\Lambda & \longrightarrow\mathcal A,
                                          & j_\Lambda                  & =j\circ\iota_\Lambda.
    \end{align}
\end{definition}

\begin{lemma}[Finite-region observables in the quasi-local algebra]
    \label{lem:qca_quasi_local_observable}
    \lean{SpinChain.norm_algebraicToQuasiLocal,
        SpinChain.norm_quasiLocalObservable,
        SpinChain.quasiLocalObservable_isometry,
        SpinChain.quasiLocalObservable_injective,
        SpinChain.quasiLocalObservable_localInclusion}
    \leanok
    \uses{def:qca_quasi_local_observable, def:qca_algebraic_local_operator_norm, def:qca_algebraic_local_algebra}
    For every finite region $\Lambda$, the map $j_\Lambda$ is an injective
    isometry.  If $\Lambda\subseteq\Gamma$, then
    \begin{align}
        j_\Gamma\circ\iota_{\Lambda,\Gamma} & =j_\Lambda.
    \end{align}
\end{lemma}

\begin{proof}\leanok
    \uses{def:qca_quasi_local_observable, def:qca_algebraic_local_operator_norm, def:qca_algebraic_local_algebra}
    The algebraic inclusion $\iota_\Lambda$ and the completion map $j$ both
    preserve the norm.  Their composite is therefore an isometry and hence
    injective.  Compatibility follows from
    $\iota_\Gamma\circ\iota_{\Lambda,\Gamma}=\iota_\Lambda$.
\end{proof}

\begin{theorem}[Density of the canonical completion map]
    \label{thm:completion_map_dense_range}
    \lean{UniformSpace.Completion.denseRange_toComplStarAlgHom}
    \leanok
    \uses{def:completion_cstar_algebra}
    The canonical star-algebra homomorphism from a normed star algebra into its
    completion has dense range.
\end{theorem}

\begin{proof}\leanok
    This is the defining density property of the completion.
\end{proof}

\begin{theorem}[Density of finite-region observables]
    \label{thm:qca_quasi_local_density}
    \lean{SpinChain.denseRange_algebraicToQuasiLocal,
        SpinChain.dense_iUnion_range_quasiLocalObservable}
    \leanok
    \uses{def:qca_quasi_local_algebra, def:qca_quasi_local_observable, thm:qca_local_observable_finite_support, thm:completion_map_dense_range}
    Algebraic local observables are dense in the quasi-local algebra.  Equivalently,
    \begin{align}
        \overline{j(\mathcal A_{\mathrm{loc}})}
         & =\overline{\bigcup_{\Lambda\Subset\mathbb Z}
            j_\Lambda(\mathcal A_\Lambda)}
        =\mathcal A.
    \end{align}
\end{theorem}

\begin{proof}\leanok
    \uses{thm:qca_local_observable_finite_support}
    Density of $j(\mathcal A_{\mathrm{loc}})$ is the defining density property
    of the completion.  Every algebraic local observable has a representative
    in some finite region, so
    \begin{align}
        j(\mathcal A_{\mathrm{loc}})
         & =\bigcup_{\Lambda\Subset\mathbb Z}j_\Lambda(\mathcal A_\Lambda).
    \end{align}
    Taking closures proves the second equality.  This theorem concerns the union
    over all finite regions and makes no claim about the image of a fixed region.
\end{proof}

The appendix introduces the quasi-local algebra and then invokes lattice
translation in the definition of a QCA and in its covariance condition
\cite[lines~2292--2306]{Cirac2017MPU}.  The following completion result and
translation laws make this background structure explicit; they are not stated
as separate results in the appendix.

\begin{definition}[Extension of star-algebra equivalences to completions]
    \label{def:completion_star_alg_equiv}
    \lean{UniformSpace.Completion.mapStarAlgEquiv,
        UniformSpace.Completion.mapStarAlgEquiv_apply,
        UniformSpace.Completion.mapStarAlgEquiv_coe}
    \leanok
    Let $A$ and $B$ be complex normed star algebras, and let
    $f\colon A\to B$ be a star-algebra equivalence.  If $f$ and $f^{-1}$ are
    continuous, then $f$ extends to a star-algebra equivalence
    \begin{align}
        \widehat f\colon\widehat A & \longrightarrow\widehat B
    \end{align}
    between their norm completions.  Its underlying map is the continuous
    extension of $f$.  Writing $j_A$ and $j_B$ for the canonical maps into the
    completions, the extension satisfies
    \begin{align}
        \widehat f(j_A(x)) & =j_B(f(x)).
    \end{align}
\end{definition}

\begin{proof}\leanok
    Apply the completion functor to $f$ and $f^{-1}$.  On the dense canonical
    images their composites satisfy
    \begin{align}
        \widehat{f^{-1}}\bigl(\widehat f(j_A(x))\bigr) & =j_A(x), &
        \widehat f\bigl(\widehat{f^{-1}}(j_B(y))\bigr) & =j_B(y).
    \end{align}
    Continuity gives the two inverse identities on the completions.
    Multiplication, addition, scalar multiplication, and the star operation
    commute with the extension because the corresponding identities hold on
    the canonical images and both sides are continuous.
\end{proof}

\begin{definition}[Translation of quasi-local observables]
    \label{def:qca_quasi_local_translation}
    \lean{SpinChain.quasiLocalTranslation,
        SpinChain.quasiLocalTranslation_apply,
        SpinChain.quasiLocalTranslation_algebraicToQuasiLocal,
        SpinChain.quasiLocalTranslation_quasiLocalObservable}
    \leanok
    \uses{def:completion_star_alg_equiv, def:qca_algebraic_local_translation, def:qca_quasi_local_algebra, def:qca_quasi_local_observable, lem:qca_algebraic_local_translation_isometry}
    For $d>0$ and every $a\in\mathbb Z$, algebraic-local translation extends to
    a star-algebra automorphism
    \begin{align}
        \tau_a\colon\mathcal A & \longrightarrow\mathcal A
    \end{align}
    of the quasi-local algebra.  On the dense algebraic-local subalgebra and on
    finite-region observables it satisfies
    \begin{align}
        \tau_a(j(A)) & =j(\tau_a(A)),                               \\
        \tau_a(j_\Lambda(A_\Lambda))
                     & =j_{\Lambda+a}(\tau_{a,\Lambda}(A_\Lambda)).
    \end{align}
\end{definition}

\begin{proof}\leanok
    \uses{def:completion_star_alg_equiv, def:qca_algebraic_local_translation, def:qca_quasi_local_algebra, def:qca_quasi_local_observable, lem:qca_algebraic_local_translation_action, lem:qca_algebraic_local_translation_isometry}
    Algebraic-local translation and its inverse translation by $-a$ are
    operator-norm isometries, hence continuous.  Extend this equivalence to the
    norm completion.  The extension identity gives
    \begin{align}
        \tau_a(j(A)) & =j(\tau_a(A)).
    \end{align}
    Substituting $A=\iota_\Lambda(A_\Lambda)$ and using the finite-region
    translation formula gives the second identity.
\end{proof}

\begin{lemma}[Quasi-local translation action laws]
    \label{lem:qca_quasi_local_translation_action}
    \lean{SpinChain.quasiLocalTranslation_zero,
        SpinChain.quasiLocalTranslation_add,
        SpinChain.quasiLocalTranslation_symm}
    \leanok
    \uses{def:qca_quasi_local_translation}
    For $d>0$, the quasi-local automorphisms form an additive action of
    $\mathbb Z$:
    \begin{align}
        \tau_0            & =\operatorname{id}, \\
        \tau_b\circ\tau_a & =\tau_{a+b},        \\
        \tau_a^{-1}       & =\tau_{-a}.
    \end{align}
\end{lemma}

\begin{proof}\leanok
    \uses{def:qca_quasi_local_translation, lem:qca_algebraic_local_translation_action}
    On the canonical image of the algebraic-local algebra,
    \begin{align}
        \tau_b\bigl(\tau_a(j(A))\bigr)
                     & =j\bigl(\tau_b(\tau_a(A))\bigr)
        =j(\tau_{a+b}(A)),                             \\
        \tau_0(j(A)) & =j(A).
    \end{align}
    Continuity extends these identities to the completion.  Taking
    $(a,b)=(a,-a)$ and $(-a,a)$ identifies the inverse of $\tau_a$ with
    $\tau_{-a}$.
\end{proof}

\begin{definition}[Translation covariance]
    \label{def:qca_translation_covariant}
    \lean{SpinChain.TranslationCovariant}
    \leanok
    \uses{def:qca_quasi_local_translation, def:qca_quasi_local_algebra}
    For $d>0$, a star-algebra automorphism $\omega$ of the quasi-local algebra
    $\mathcal A$ is translation covariant if
    \begin{align}
        \omega\circ\tau_a=\tau_a\circ\omega
    \end{align}
    for every $a\in\mathbb Z$.  This is the covariance condition in the
    definition of a one-dimensional QCA \cite[line~2298]{Cirac2017MPU}.
\end{definition}

\begin{theorem}[Pointwise characterization of translation covariance]
    \label{thm:qca_translation_covariant_iff}
    \lean{SpinChain.translationCovariant_iff}
    \leanok
    \uses{def:qca_translation_covariant}
    For $d>0$, a star-algebra automorphism $\omega$ of $\mathcal A$ is
    translation covariant if and only if
    \begin{align}
        \omega\bigl(\tau_a(A)\bigr)=\tau_a\bigl(\omega(A)\bigr)
    \end{align}
    for every $a\in\mathbb Z$ and every $A\in\mathcal A$.
\end{theorem}

\begin{proof}\leanok
    \uses{def:qca_translation_covariant}
    Two automorphisms are equal if and only if they agree on every observable.
    Apply this observation to the two composites in the definition.
\end{proof}

\begin{theorem}[Translation covariance from unit translation]
    \label{thm:qca_translation_covariant_iff_commute_one}
    \lean{SpinChain.translationCovariant_iff_commute_one}
    \leanok
    \uses{def:qca_translation_covariant}
    For $d>0$, a star-algebra automorphism $\omega$ of $\mathcal A$ commutes with
    every integer translation if and only if it commutes with unit translation:
    \begin{align}
        \omega\circ\tau_1=\tau_1\circ\omega.
    \end{align}
\end{theorem}

\begin{proof}\leanok
    \uses{def:qca_translation_covariant, lem:qca_quasi_local_translation_action}
    The forward implication follows by taking the translation step to be $1$.
    Conversely, use two-sided integer induction.  The zero case follows from
    $\tau_0=\operatorname{id}$.  For the successor step, the additive action law
    gives $\tau_{i+1}=\tau_1\circ\tau_i$, and
    \begin{align}
        \omega\circ\tau_{i+1}
         & =\omega\circ\tau_1\circ\tau_i
        =\tau_1\circ\omega\circ\tau_i    \\
         & =\tau_1\circ\tau_i\circ\omega
        =\tau_{i+1}\circ\omega.
    \end{align}
    For the predecessor step, $\tau_{-1}=\tau_1^{-1}$ transfers commutation
    to $\tau_{-1}$, while $\tau_{i-1}=\tau_{-1}\circ\tau_i$, so
    \begin{align}
        \omega\circ\tau_{i-1}
         & =\omega\circ\tau_{-1}\circ\tau_i
        =\tau_{-1}\circ\omega\circ\tau_i    \\
         & =\tau_{-1}\circ\tau_i\circ\omega
        =\tau_{i-1}\circ\omega.
    \end{align}
\end{proof}

\begin{theorem}[Identity and composition of translation-covariant automorphisms]
    \label{thm:qca_translation_covariant_identity_composition}
    \lean{SpinChain.translationCovariant_refl,
        SpinChain.TranslationCovariant.trans}
    \leanok
    \uses{def:qca_translation_covariant}
    For $d>0$, the identity automorphism of $\mathcal A$ is translation
    covariant.  If $\omega$ and $\eta$ are translation covariant, then so is
    the composite
    \begin{align}
        A\longmapsto\eta\bigl(\omega(A)\bigr),
    \end{align}
    which first applies $\omega$ and then $\eta$.
\end{theorem}

\begin{proof}\leanok
    \uses{def:qca_translation_covariant}
    The identity commutes pointwise with every translation.  For the composite,
    translation covariance gives
    \begin{align}
        \eta\bigl(\omega(\tau_a(A))\bigr)
         & =\eta\bigl(\tau_a(\omega(A))\bigr)
        =\tau_a\bigl(\eta(\omega(A))\bigr).
    \end{align}
\end{proof}

\begin{lemma}[Norm isometry of quasi-local translation]
    \label{lem:qca_quasi_local_translation_isometry}
    \lean{SpinChain.norm_quasiLocalTranslation,
        SpinChain.quasiLocalTranslation_isometry}
    \leanok
    \uses{def:qca_quasi_local_translation, def:qca_quasi_local_algebra}
    For $d>0$, every quasi-local translation preserves the operator norm and
    distance:
    \begin{align}
        \lVert\tau_a(A)\rVert           & =\lVert A\rVert,   \\
        \lVert\tau_a(A)-\tau_a(B)\rVert & =\lVert A-B\rVert.
    \end{align}
\end{lemma}

\begin{proof}\leanok
    \uses{def:qca_quasi_local_translation, def:qca_quasi_local_algebra}
    A star-algebra equivalence between complex C-star algebras preserves the
    C-star norm.  Since $\tau_a$ is linear,
    \begin{align}
        \lVert\tau_a(A)-\tau_a(B)\rVert
         & =\lVert\tau_a(A-B)\rVert
        =\lVert A-B\rVert.
    \end{align}
\end{proof}

\begin{definition}[Finite-region support in the quasi-local algebra]
    \label{def:qca_quasi_local_supported_in}
    \lean{SpinChain.QuasiLocalSupportedIn,
        SpinChain.QuasiLocalSupportedIn.quasiLocalObservable,
        SpinChain.QuasiLocalSupportedIn.iff_exists}
    \leanok
    \uses{def:qca_quasi_local_observable}
    For $d>0$, a quasi-local observable $x\in\mathcal A$ is supported in a
    finite region $\Lambda\subset\mathbb Z$ if it belongs to the image of the
    canonical embedding $j_\Lambda\colon\mathcal A_\Lambda\to\mathcal A$.
    Equivalently,
    \begin{align}
        x\text{ is supported in }\Lambda
         & \iff x\in j_\Lambda(\mathcal A_\Lambda)         \\
         & \iff \text{there exists }A\in\mathcal A_\Lambda
        \text{ such that }j_\Lambda(A)=x.
    \end{align}
    In particular, $j_\Lambda(A)$ is supported in $\Lambda$ for every
    $A\in\mathcal A_\Lambda$.
\end{definition}

\begin{proof}\leanok
    \uses{def:qca_quasi_local_observable}
    These assertions follow directly from membership in the range of
    $j_\Lambda$.
\end{proof}

\begin{lemma}[Monotonicity of quasi-local support]
    \label{lem:qca_quasi_local_support_mono}
    \lean{SpinChain.QuasiLocalSupportedIn.mono}
    \leanok
    \uses{def:qca_quasi_local_supported_in}
    Let $d>0$.  If $x\in\mathcal A$ is supported in $\Lambda$ and
    $\Lambda\subseteq\Gamma$, then $x$ is supported in $\Gamma$.
\end{lemma}

\begin{proof}\leanok
    \uses{def:qca_quasi_local_supported_in, def:qca_local_inclusion, lem:qca_quasi_local_observable}
    Write $x=j_\Lambda(A)$.  Enlarge $A$ to $\mathcal A_\Gamma$ by tensoring
    with the identity on $\Gamma\setminus\Lambda$.  Compatibility of the
    canonical embeddings identifies the image of this enlargement with $x$.
\end{proof}

\begin{theorem}[Disjointly supported quasi-local observables commute]
    \label{thm:qca_disjoint_quasi_local_commute}
    \lean{SpinChain.quasiLocalObservable_commute_of_disjoint,
        SpinChain.QuasiLocalSupportedIn.commute_of_disjoint}
    \leanok
    \uses{def:qca_quasi_local_supported_in}
    Let $d>0$.  If $x,y\in\mathcal A$ are supported in finite regions
    $\Lambda$ and $\Gamma$ with $\Lambda\cap\Gamma=\varnothing$, then
    \begin{align}
        xy & =yx.
    \end{align}
    Equivalently, for $A\in\mathcal A_\Lambda$ and
    $B\in\mathcal A_\Gamma$,
    \begin{align}
        j_\Lambda(A)j_\Gamma(B) & =j_\Gamma(B)j_\Lambda(A).
    \end{align}
    This is finite-support infrastructure for the quasi-local algebra in
    \cite[lines~2292--2300]{Cirac2017MPU}; it is not stated there as a separate
    theorem.
    % Source context: paper_v2.tex 2292--2300.
\end{theorem}

\begin{proof}\leanok
    \uses{thm:qca_disjoint_supported_local_commute}
    The completion map is a star-algebra homomorphism.  Applying it to the
    algebraic-local identity gives
    \begin{align}
        j(XY)=j(YX),
    \end{align}
    and finite-region support provides representatives
    $x=j_\Lambda(A)$ and $y=j_\Gamma(B)$.
\end{proof}

The following finite-dimensional averaging fact is standard and is not a
claim from \cite{Cirac2017MPU}.  It will be used to characterize exact local
support by commutation with observables on the complementary tensor factor.

\begin{definition}[Complementary-factor Weyl twirl]
    \label{def:qca_complementary_weyl_twirl}
    \lean{Matrix.reindexedWeyl,Matrix.complementaryWeylTwirl}
    \leanok
    Let $S$ and $C$ be finite sets, let $n=|C|>0$, and choose an equivalence
    $e\colon C\simeq\mathbb Z/n\mathbb Z$ and a primitive $n$-th root of unity
    $\zeta$.  If $W_{a,b}$ is the corresponding reindexed Weyl operator on
    $\mathbb C^C$, define
    \begin{align}
        \mathcal E_C(M)
         & = \frac{1}{n^2}\sum_{a,b\in\mathbb Z/n\mathbb Z}
        (\mathbf 1_S\otimes W_{a,b})M
        (\mathbf 1_S\otimes W_{a,b})^\dagger.
    \end{align}
\end{definition}

\begin{lemma}[Complementary-factor Weyl average]
    \label{lem:qca_complementary_weyl_twirl_eq}
    \lean{Matrix.sum_kronecker_one_reindexedWeyl_conj,
        Matrix.complementaryWeylTwirl_eq_partialTraceRight_kronecker,
        Matrix.complementaryWeylTwirl_eq_smul_partialTraceRight_kronecker_one}
    \leanok
    \uses{def:qca_complementary_weyl_twirl}
    For every $M\in\operatorname{Mat}_{S\times C}(\mathbb C)$,
    \begin{align}
        \mathcal E_C(M)
         & = \operatorname{tr}_C(M)\otimes\frac{\mathbf 1_C}{n}                \\
         & = \left(\frac{1}{n}\operatorname{tr}_C(M)\right)\otimes\mathbf 1_C.
    \end{align}
\end{lemma}

\begin{proof}\leanok
    Split $M$ into $C\times C$ blocks indexed by $S$.  The Weyl average sends
    each block $B$ to $(\operatorname{tr}B/n)\mathbf 1_C$, so its block trace is
    the corresponding entry of $\operatorname{tr}_C(M)$.
\end{proof}

\begin{lemma}[Norm estimate for the complementary-factor twirl]
    \label{lem:qca_complementary_weyl_twirl_norm}
    \lean{Matrix.norm_complementaryWeylTwirl_le,
        Matrix.norm_complementaryWeylTwirl_sub_le}
    \leanok
    \uses{def:qca_complementary_weyl_twirl}
    Let $S$ be a nonempty finite set.  The average is contractive in the operator
    norm.  Moreover, if every complementary Weyl conjugation fixes $X$, then
    \begin{align}
        \lVert\mathcal E_C(M)-X\rVert
         & \le \lVert M-X\rVert.
    \end{align}
\end{lemma}

\begin{proof}\leanok
    Unitary conjugation preserves the operator norm, and the coefficients of the
    finite average sum to one.  Apply this contraction to $M-X$.
\end{proof}

\begin{theorem}[Exact support from outside commutation]
    \label{thm:qca_quasi_local_supported_iff_commute_disjoint}
    \lean{SpinChain.quasiLocalSupportedIn_iff_commute_disjoint}
    \leanok
    \uses{def:qca_quasi_local_supported_in}
    Let $d>0$, let $\Lambda\subset\mathbb Z$ be finite, and let
    $x\in\mathcal A$.  Then
    \begin{align}
        x\in\mathcal A_\Lambda
        \quad\Longleftrightarrow\quad
        [x,y]=0
        \text{ for every }y\in\mathcal A_\Gamma
        \text{ and every finite }\Gamma\cap\Lambda=\varnothing.
    \end{align}
    This is standard infrastructure for the UHF algebra used in TNLean.  It is
    not a theorem stated in \cite{Cirac2017MPU} or by Schumacher--Werner,
    quant-ph/0405174.
\end{theorem}

\begin{proof}\leanok
    \uses{thm:qca_disjoint_quasi_local_commute, lem:qca_quasi_local_observable, thm:qca_quasi_local_density, lem:qca_local_inclusion_norm, thm:mpu_reindex_unitary, def:qca_complementary_weyl_twirl, lem:qca_complementary_weyl_twirl_eq}
    The forward implication follows from commutation of disjoint tensor factors.
    Conversely, approximate $x$ by $A\in\mathcal A_\Delta$, where
    $\Lambda\subseteq\Delta$.  Split
    $\Delta=\Lambda\sqcup(\Delta\setminus\Lambda)$ and average $A$ over Weyl
    unitaries on the second factor.  Every such conjugation fixes $x$ by the
    commutation hypothesis, so unitary invariance and convexity give
    \begin{align}
        \lVert\mathcal E_{\Delta\setminus\Lambda}(A)-x\rVert
         & \le \lVert A-x\rVert.
    \end{align}
    The Weyl identity writes the average as $B\otimes\mathbf 1$ for some
    $B\in\mathcal A_\Lambda$.  Thus elements of $\mathcal A_\Lambda$ approximate
    $x$ arbitrarily well.  The embedding of $\mathcal A_\Lambda$ is isometric,
    hence has closed range, and therefore $x\in\mathcal A_\Lambda$.
\end{proof}

The following theorem gives the scalar center of the homogeneous UHF algebra.  GNVW
invoke this conclusion for the quasi-local algebra at
\cite[lines~1276--1282]{Gross2012IndexTheory}, specifically line~1279, but do
not prove it there.  The homogeneous construction in
\cite[lines~2292--2298]{Cirac2017MPU} likewise defines the norm completion
without stating this conclusion separately.  Schumacher--Werner give the
corresponding homogeneous norm-completed construction in
\cite[lines~263--285]{SchumacherWerner2004QCA}, again without a separate
scalar-center theorem.

\begin{theorem}[Scalar center of the homogeneous quasi-local algebra]
    \label{thm:qca_quasi_local_center_scalar}
    \lean{SpinChain.existsUnique_eq_smul_one_of_commute_quasiLocalObservable,
        SpinChain.existsUnique_eq_smul_one_of_commute,
        SpinChain.instIsCentralQuasiLocalAlgebra,
        SpinChain.quasiLocalAlgebra_center_eq_bot}
    \leanok
    \uses{def:qca_quasi_local_algebra, def:qca_quasi_local_observable}
    Let $d>0$ and let $z\in\mathcal A$.  If
    \begin{align}
        \forall\Gamma\subset\mathbb Z\text{ finite},\quad
        \forall A\in\mathcal A_\Gamma,\qquad
        z j_\Gamma(A) & =j_\Gamma(A)z,
    \end{align}
    then there is a unique $c\in\mathbb C$ such that
    \begin{align}
        z & =c\mathbf 1.
    \end{align}
    Consequently,
    \begin{align}
        \{z\in\mathcal A\mid za=az\text{ for every }a\in\mathcal A\}
         & =\mathbb C\mathbf 1.
    \end{align}
    This statement concerns the fixed on-site dimension $d$; it does not assert
    the site-dependent extension used in GNVW.
\end{theorem}

\begin{proof}\leanok
    \uses{thm:qca_quasi_local_supported_iff_commute_disjoint, def:qca_quasi_local_supported_in, lem:qca_quasi_local_observable}
    Apply Theorem~\ref{thm:qca_quasi_local_supported_iff_commute_disjoint}
    with $\Lambda=\varnothing$.  The commutation hypothesis gives
    $z\in j_\varnothing(\mathcal A_\varnothing)$, so there is an
    $A\in\mathcal A_\varnothing$ with $z=j_\varnothing(A)$.  Let
    $e\colon\mathbb C\xrightarrow{\ \sim\ }\mathcal A_\varnothing$ be the
    canonical isomorphism and write $A=e(c)$.  If also
    $z=j_\varnothing(e(c'))$, then injectivity of $j_\varnothing$ and of $e$
    give
    \begin{align}
        j_\varnothing(e(c))=j_\varnothing(e(c'))
         & \Longrightarrow e(c)=e(c')
        \Longrightarrow c=c'.
    \end{align}
    Thus $c$ is the unique scalar for which $z=c\mathbf 1$.

    Write $Z(\mathcal A)$ for the center.  Since every
    $j_\Gamma(B)$ belongs to $\mathcal A$, the two inclusions are
    \begin{align}
        z\in Z(\mathcal A)
         & \Longrightarrow
        \bigl(\forall \Gamma\subset\mathbb Z\text{ finite},\
        \forall B\in\mathcal A_\Gamma,\
        z j_\Gamma(B)=j_\Gamma(B)z\bigr)
        \Longrightarrow z\in\mathbb C\mathbf 1, \\
        z=c\mathbf 1\quad(c\in\mathbb C)
         & \Longrightarrow
        \bigl(\forall a\in\mathcal A,\
        (c\mathbf 1)a=a(c\mathbf 1)\bigr)
        \Longrightarrow z\in Z(\mathcal A).
    \end{align}
    Hence $Z(\mathcal A)=\mathbb C\mathbf 1$.
\end{proof}

The appendix requires that ``there exists a finite subset
$\mathcal N\subset\mathbb Z$'' such that, for every finite
$\Lambda\subset\mathbb Z$,
\begin{align}
    \omega(\mathcal A_\Lambda) & \subset
    \mathcal A_{\Lambda+\mathcal N}.
\end{align}
This is the propagation condition stated at line~2298 of
\cite{Cirac2017MPU}.  The definitions below isolate this forward inclusion.
They do not impose translation covariance or a propagation bound for
$\omega^{-1}$.

\begin{definition}[Propagation within a fixed neighborhood]
    \label{def:qca_propagates_within}
    \lean{SpinChain.PropagatesWithin}
    \leanok
    \uses{def:qca_quasi_local_supported_in, def:qca_region_sumset}
    Let $d>0$, let $\omega\colon\mathcal A\to\mathcal A$ be a star-algebra
    automorphism, and let $\mathcal N\subset\mathbb Z$ be finite.  The
    automorphism $\omega$ propagates within $\mathcal N$ if, for every finite
    region $\Lambda$ and every $x\in\mathcal A$,
    \begin{align}
        x\text{ supported in }\Lambda
         & \implies
        \omega(x)\text{ supported in }\Lambda+\mathcal N.
    \end{align}
\end{definition}

\begin{definition}[Finite forward propagation]
    \label{def:qca_finite_propagation}
    \lean{SpinChain.HasFinitePropagation}
    \leanok
    \uses{def:qca_propagates_within}
    Let $d>0$ and let $\omega\colon\mathcal A\to\mathcal A$ be a star-algebra
    automorphism.  The automorphism $\omega$ has finite forward propagation if
    there exists a finite set $\mathcal N\subset\mathbb Z$ such that $\omega$
    propagates within $\mathcal N$.
\end{definition}

\begin{lemma}[Equivalent formulations of forward propagation]
    \label{lem:qca_propagates_within_iff}
    \lean{SpinChain.propagatesWithin_iff_quasiLocalObservable,
        SpinChain.propagatesWithin_iff_range_subset,
        SpinChain.propagatesWithin_iff_exists_local}
    \leanok
    \uses{def:qca_quasi_local_supported_in, def:qca_propagates_within, def:qca_quasi_local_observable, def:qca_region_sumset}
    Let $d>0$, let $\omega\colon\mathcal A\to\mathcal A$ be a star-algebra
    automorphism, and let $\mathcal N\subset\mathbb Z$ be finite.  The following
    conditions are equivalent:
    \begin{align}
         & \omega\text{ propagates within }\mathcal N; \tag{a}              \\
         & \forall\Lambda\Subset\mathbb Z,\ \forall A\in\mathcal A_\Lambda,
        \quad \omega(j_\Lambda(A))\text{ is supported in }
        \Lambda+\mathcal N; \tag{b}                                         \\
         & \forall\Lambda\Subset\mathbb Z,\quad
        \omega\bigl(j_\Lambda(\mathcal A_\Lambda)\bigr)
        \subseteq j_{\Lambda+\mathcal N}
        (\mathcal A_{\Lambda+\mathcal N}); \tag{c}                          \\
         & \forall\Lambda\Subset\mathbb Z,\ \forall A\in\mathcal A_\Lambda,
        \ \exists B\in\mathcal A_{\Lambda+\mathcal N},\quad
        \omega(j_\Lambda(A))=j_{\Lambda+\mathcal N}(B). \tag{d}
    \end{align}
    Thus condition~\textup{(c)} is precisely the range-inclusion formula from
    line~2298 of \cite{Cirac2017MPU}, expressed inside the quasi-local algebra,
    and condition~\textup{(d)} is its finite-region witness form.
\end{lemma}

\begin{proof}\leanok
    \uses{def:qca_quasi_local_supported_in, def:qca_propagates_within, def:qca_quasi_local_observable}
    Every observable supported in $\Lambda$ has the form $j_\Lambda(A)$, so
    \textup{(a)} and \textup{(b)} are equivalent.  Replacing support by
    membership in the range of the canonical embedding gives \textup{(c)}.
    Finally, membership in
    $j_{\Lambda+\mathcal N}(\mathcal A_{\Lambda+\mathcal N})$ is equivalent to
    the existence of a representative $B$ in that finite-region algebra, which
    gives \textup{(d)}.
\end{proof}

\begin{lemma}[Enlargement of propagation neighborhoods]
    \label{lem:qca_propagation_neighborhood_mono}
    \lean{SpinChain.PropagatesWithin.mono,
        SpinChain.HasFinitePropagation.exists_superset}
    \leanok
    \uses{def:qca_propagates_within, def:qca_finite_propagation}
    Let $d>0$ and let $\omega\colon\mathcal A\to\mathcal A$ be a star-algebra
    automorphism.  If $\mathcal N\subseteq\mathcal M$ and $\omega$ propagates
    within $\mathcal N$, then $\omega$ propagates within $\mathcal M$.  More
    explicitly, for every finite region $\Lambda$,
    \begin{align}
        \Lambda+\mathcal N & \subseteq\Lambda+\mathcal M,                              \\
        x\text{ supported in }\Lambda
                           & \implies \omega(x)\text{ supported in }\Lambda+\mathcal N
        \implies \omega(x)\text{ supported in }\Lambda+\mathcal M.
    \end{align}
    Consequently, if $\omega$ has finite forward propagation, then for every
    prescribed finite set $\mathcal N\subset\mathbb Z$ there is a finite set
    $\mathcal M\subset\mathbb Z$ such that
    \begin{align}
        \mathcal N & \subseteq\mathcal M,                 \\
                   & \forall\Lambda\Subset\mathbb Z,\quad
        \omega(\mathcal A_\Lambda)\subseteq
        \mathcal A_{\Lambda+\mathcal M}.
    \end{align}
    This is the neighborhood-enlargement consequence of the propagation
    condition in Appendix line~2298 of \cite{Cirac2017MPU}.
\end{lemma}

\begin{proof}\leanok
    \uses{def:qca_region_sumset, lem:qca_quasi_local_support_mono, def:qca_propagates_within, def:qca_finite_propagation}
    Monotonicity of the sumset in its second argument gives
    $\Lambda+\mathcal N\subseteq\Lambda+\mathcal M$.  Monotonicity of
    quasi-local support then enlarges the support of $\omega(x)$ from
    $\Lambda+\mathcal N$ to $\Lambda+\mathcal M$.  For the final assertion,
    choose a finite propagation neighborhood $\mathcal K$ and take
    $\mathcal M=\mathcal K\cup\mathcal N$.
\end{proof}

\begin{lemma}[Composition of forward propagation bounds]
    \label{lem:qca_forward_propagation_comp}
    \lean{SpinChain.PropagatesWithin.trans,
        SpinChain.HasFinitePropagation.trans}
    \leanok
    \uses{def:qca_region_sumset, def:qca_propagates_within, def:qca_finite_propagation}
    Let $d>0$, let $\omega,\eta\colon\mathcal A\to\mathcal A$ be star-algebra
    automorphisms, and let $\mathcal N,\mathcal M\subset\mathbb Z$ be finite.
    Suppose that $\omega$ propagates within $\mathcal N$ and $\eta$ propagates
    within $\mathcal M$.  The composite $\eta\circ\omega$ first applies
    $\omega$ and then $\eta$, so
    \begin{align}
        (\eta\circ\omega)(x) & =\eta(\omega(x)),                                         \\
        x\text{ supported in }\Lambda
                             & \implies \omega(x)\text{ supported in }\Lambda+\mathcal N \\
                             & \implies \eta(\omega(x))\text{ supported in }
        (\Lambda+\mathcal N)+\mathcal M
        =\Lambda+(\mathcal N+\mathcal M).
    \end{align}
    Hence $\eta\circ\omega$ propagates within $\mathcal N+\mathcal M$.
    In particular, if $\omega$ and $\eta$ have finite forward propagation, then
    $\eta\circ\omega$ has finite forward propagation.  This is the composition
    closure of the finite-neighborhood condition in Appendix line~2298 of
    \cite{Cirac2017MPU}.
\end{lemma}

\begin{proof}\leanok
    \uses{def:qca_region_sumset, def:qca_propagates_within, def:qca_finite_propagation}
    Apply the two propagation bounds successively and use associativity of finite
    region sumsets.  For the existential assertion, choose one finite propagation
    neighborhood for each automorphism and use their sumset.
\end{proof}

\begin{lemma}[Symmetric-interval normalization of finite propagation]
    \label{lem:qca_finite_propagation_symmetric_interval}
    \lean{SpinChain.exists_subset_symmetric_Icc,
        SpinChain.PropagatesWithin.exists_symmetric_Icc,
        SpinChain.HasFinitePropagation.exists_symmetric_Icc}
    \leanok
    \uses{def:qca_propagates_within, def:qca_finite_propagation}
    Every finite set $\mathcal N\subset\mathbb Z$ is contained in
    $[-R,R]\cap\mathbb Z$ for some $R\in\mathbb N$.  This includes
    $\mathcal N=\varnothing$, for which one may take $R=0$.  If $\omega$
    propagates within a fixed set $\mathcal N$, then $R$ may be chosen so that
    \begin{align}
        \mathcal N & \subseteq[-R,R]\cap\mathbb Z,                  \\
        \omega     & \text{ propagates within }[-R,R]\cap\mathbb Z.
    \end{align}
    Consequently, finite forward propagation admits a symmetric-interval
    witness:
    \begin{align}
        \exists R\in\mathbb N,\qquad
        \omega & \text{ propagates within }[-R,R]\cap\mathbb Z.
    \end{align}
\end{lemma}

\begin{proof}\leanok
    \uses{def:qca_finite_propagation, lem:qca_propagation_neighborhood_mono}
    Let $R$ be the maximum of $|n|$ over $n\in\mathcal N$, with value $0$
    when $\mathcal N$ is empty.  Then $-R\leq n\leq R$ for every
    $n\in\mathcal N$, which gives
    $\mathcal N\subseteq[-R,R]\cap\mathbb Z$.  Enlargement of propagation
    neighborhoods gives the fixed-neighborhood assertion.  Finally, choose a
    finite propagation neighborhood and apply this enlargement to its symmetric
    interval.
\end{proof}

\begin{lemma}[Disjointness for reflected neighborhoods]
    \label{lem:qca_disjoint_region_sumset_neg_iff}
    \lean{SpinChain.disjoint_regionSumset_neg_iff}
    \leanok
    \uses{def:qca_region_sumset}
    Let $\Gamma,\Lambda,\mathcal N\subset\mathbb Z$ be finite.  Then
    \begin{align}
        \Gamma\cap(\Lambda+(-\mathcal N))=\varnothing
        \quad\Longleftrightarrow\quad
        (\Gamma+\mathcal N)\cap\Lambda=\varnothing.
    \end{align}
\end{lemma}

\begin{proof}\leanok
    If $g+n=\lambda$ for $g\in\Gamma$, $n\in\mathcal N$, and
    $\lambda\in\Lambda$, then
    $g=\lambda+(-n)\in\Lambda+(-\mathcal N)$, proving the forward implication.
    Conversely, if $g\in\Gamma$ and $g=\lambda+(-n)$ with
    $\lambda\in\Lambda$ and $n\in\mathcal N$, then
    $\lambda=g+n\in\Gamma+\mathcal N$.
\end{proof}

\begin{theorem}[Inverse propagation and inverse covariance]
    \label{thm:qca_inverse_propagation}
    \lean{SpinChain.PropagatesWithin.symm,
        SpinChain.HasFinitePropagation.symm,
        SpinChain.TranslationCovariant.symm}
    \leanok
    \uses{def:qca_propagates_within, def:qca_finite_propagation, def:qca_translation_covariant}
    Let $d>0$, let $\omega$ be a star-algebra automorphism of $\mathcal A$, and
    let $\mathcal N\subset\mathbb Z$ be finite.  If $\omega$ propagates within
    $\mathcal N$, then $\omega^{-1}$ propagates within $-\mathcal N$.  Therefore
    finite forward propagation of $\omega$ implies finite forward propagation of
    $\omega^{-1}$.  Moreover, if $\omega$ is translation covariant, then so is
    $\omega^{-1}$.

    This is a consequence of the forward-only locality definition, not part of
    that definition.  Schumacher--Werner, quant-ph/0405174, derive inverse
    locality through Lemma~5, Theorem~6, and Corollary~7 using the generalized
    Margolus structure.  The proof here instead uses TNLean's exact
    outside-commutant infrastructure.
\end{theorem}

\begin{proof}\leanok
    \uses{lem:qca_disjoint_region_sumset_neg_iff, thm:qca_quasi_local_supported_iff_commute_disjoint}
    Let $x\in\mathcal A_\Lambda$ and let
    $y\in\mathcal A_\Gamma$, where
    $\Gamma\cap(\Lambda+(-\mathcal N))=\varnothing$.  Forward locality gives
    $\omega(y)\in\mathcal A_{\Gamma+\mathcal N}$.  The regions
    $\Gamma+\mathcal N$ and $\Lambda$ are disjoint, so $x$ commutes with
    $\omega(y)$.  Applying $\omega^{-1}$ shows that $\omega^{-1}(x)$ commutes
    with $y$.  The exact outside-commutant characterization now gives
    \begin{align}
        \omega^{-1}(x) & \in\mathcal A_{\Lambda+(-\mathcal N)}.
    \end{align}
    Choosing the reflected finite neighborhood proves the existential statement.
    For each lattice translation $\tau_a$, multiplying the covariance identity
    on the left and right by $\omega^{-1}$ gives
    \begin{align}
        \omega\circ\tau_a=\tau_a\circ\omega
         & \implies
        \omega^{-1}\circ\tau_a=\tau_a\circ\omega^{-1}.
    \end{align}
\end{proof}

\begin{definition}[One-dimensional quantum cellular automaton]
    \label{def:qca}
    \lean{SpinChain.IsQCA}
    \leanok
    \uses{def:qca_translation_covariant, def:qca_finite_propagation}
    Let $d>0$ and let $\omega\colon\mathcal A\to\mathcal A$ be a star-algebra
    automorphism.  The automorphism $\omega$ is a one-dimensional quantum
    cellular automaton if it is translation covariant and has finite forward
    propagation.  Thus there exists a finite set $\mathcal N\subset\mathbb Z$
    such that, for every finite $\Lambda\subset\mathbb Z$,
    \begin{align}
        \omega(\mathcal A_\Lambda) & \subseteq
        \mathcal A_{\Lambda+\mathcal N}.
    \end{align}
    This is the definition at line~2298 of \cite{Cirac2017MPU}.  It imposes only
    the displayed forward propagation condition and makes no locality claim for
    $\omega^{-1}$.
\end{definition}

\begin{theorem}[Exact interface for quantum cellular automata]
    \label{thm:qca_iff_mk_projections}
    \lean{SpinChain.isQCA_iff,SpinChain.IsQCA.mk,
        SpinChain.IsQCA.translationCovariant,
        SpinChain.IsQCA.hasFinitePropagation}
    \leanok
    \uses{def:qca}
    For $d>0$, a star-algebra automorphism $\omega$ of $\mathcal A$ is a
    quantum cellular automaton if and only if it is translation covariant and
    has finite forward propagation.  These two properties therefore construct an
    automaton, and every quantum cellular automaton satisfies each of them.
\end{theorem}

\begin{proof}\leanok
    \uses{def:qca}
    This is the conjunction defining a quantum cellular automaton, together with
    its introduction and elimination rules.
\end{proof}

\begin{theorem}[Identity quantum cellular automaton]
    \label{thm:qca_identity}
    \lean{SpinChain.isQCA_refl}
    \leanok
    \uses{def:qca}
    For $d>0$, the identity automorphism of $\mathcal A$ is a quantum cellular
    automaton.
\end{theorem}

\begin{proof}\leanok
    \uses{def:qca_finite_propagation, def:qca_propagates_within, def:qca_region_sumset, thm:qca_iff_mk_projections, thm:qca_translation_covariant_identity_composition}
    The identity is translation covariant.  For finite forward propagation,
    choose the neighborhood $\{0\}$.  An observable supported in $\Lambda$
    remains unchanged, and $\Lambda+\{0\}=\Lambda$.
\end{proof}

\begin{theorem}[Composition of quantum cellular automata]
    \label{thm:qca_composition}
    \lean{SpinChain.IsQCA.trans}
    \leanok
    \uses{def:qca}
    Let $d>0$.  If $\omega$ and $\eta$ are quantum cellular automata on
    $\mathcal A$, then $\eta\circ\omega$ is a quantum cellular automaton.  This
    composite first applies $\omega$ and then $\eta$.  No propagation condition
    for either inverse is required.
\end{theorem}

\begin{proof}\leanok
    \uses{thm:qca_iff_mk_projections, thm:qca_translation_covariant_identity_composition, lem:qca_forward_propagation_comp}
    Both defining properties are closed under composition in the stated order.
    Translation covariance of $\omega$ and $\eta$ gives translation covariance
    of $\eta\circ\omega$.  If $\mathcal N$ and $\mathcal M$ are forward
    propagation neighborhoods for $\omega$ and $\eta$, respectively, then
    $\mathcal N+\mathcal M$ is a forward propagation neighborhood for
    $\eta\circ\omega$.  Hence the composite is a quantum cellular automaton.
\end{proof}

\begin{theorem}[Inverse quantum cellular automaton]
    \label{thm:qca_inverse}
    \lean{SpinChain.IsQCA.symm}
    \leanok
    \uses{def:qca}
    Let $d>0$.  If $\omega$ is a quantum cellular automaton on $\mathcal A$, then
    $\omega^{-1}$ is also a quantum cellular automaton.
\end{theorem}

\begin{proof}\leanok
    \uses{thm:qca_iff_mk_projections, thm:qca_inverse_propagation}
    Translation covariance and finite forward propagation are both preserved by
    inversion.  These are the two defining properties of a quantum cellular
    automaton.
\end{proof}

\begin{definition}[Expansion of a blocked region]
    \label{def:qca_expanded_region}
    \lean{SpinChain.expandedRegion,SpinChain.blockSiteEquiv}
    \leanok
    Let $L>0$ be a block length and let $\Lambda\subset\mathbb Z$ be finite.
    The expansion of $\Lambda$ is
    \begin{align}
        \widehat\Lambda_L
         & =\{Lx+r:x\in\Lambda,\ 0\leq r<L\}.
    \end{align}
    Every $z\in\widehat\Lambda_L$ has unique block coordinates $(x,r)$ with
    $x\in\Lambda$ and $0\leq r<L$.  Increasing $r$ gives the consecutive-site
    order within the block.  This is the finite-region site grouping used in
    \cite[lines~2308 and 2313--2320]{Cirac2017MPU}; it does not assert the
    QCA-to-MPU converse stated at line~2308.
\end{definition}

\begin{theorem}[Coordinates and monotonicity of expanded regions]
    \label{thm:qca_expanded_region_coordinates}
    \lean{SpinChain.mem_expandedRegion,SpinChain.expandedRegion_mono,
        SpinChain.mem_expandedRegion_blockSite,
        SpinChain.blockSiteEquiv_apply_fst_val,
        SpinChain.blockSiteEquiv_apply_snd,
        SpinChain.blockSiteEquiv_symm_apply_val}
    \leanok
    \uses{def:qca_expanded_region}
    An integer $z$ belongs to $\widehat\Lambda_L$ exactly when its quotient by
    $L$ belongs to $\Lambda$.  The inverse block-coordinate map sends $(x,r)$
    to $Lx+r$, and $\Lambda\subseteq\Gamma$ implies
    $\widehat\Lambda_L\subseteq\widehat\Gamma_L$.
\end{theorem}

\begin{proof}\leanok
    \uses{def:qca_expanded_region}
    Apply Euclidean division by the positive integer $L$.  Its quotient and
    residue are unique, the residue lies in $\{0,\ldots,L-1\}$, and taking the
    inverse coordinates gives $Lx+r$.  Monotonicity follows by retaining the
    same quotient and residue.
\end{proof}

\begin{definition}[Blocking equivalence for finite configurations]
    \label{def:qca_configuration_blocking}
    \lean{SpinChain.Config.blocking}
    \leanok
    \uses{def:qca_expanded_region}
    For on-site dimension $d$, configurations of $d^L$-level blocked spins on
    $\Lambda$ are equivalent to configurations of $d$-level spins on
    $\widehat\Lambda_L$.  The blocked label at $x$ is decoded into its ordered
    length-$L$ word, and its $r$th letter is placed at $Lx+r$.
\end{definition}

\begin{theorem}[Evaluation and restriction of blocked configurations]
    \label{thm:qca_configuration_blocking_formulas}
    \lean{SpinChain.Config.blocking_apply,
        SpinChain.Config.blocking_symm_apply,
        SpinChain.Config.restrict_blocking}
    \leanok
    \uses{def:qca_configuration_blocking, thm:qca_expanded_region_coordinates}
    Evaluation at $Lx+r$ reads the $r$th letter of the blocked spin at $x$.
    Conversely, encoding collects these letters in increasing residue order.
    At the configuration level, decoding commutes with restriction from
    $\Gamma$ to $\Lambda\subseteq\Gamma$.
\end{theorem}

\begin{proof}\leanok
    \uses{def:qca_configuration_blocking, thm:qca_expanded_region_coordinates}
    The two evaluation formulas are the inverse laws for the ordered word
    encoding.  Restriction preserves the quotient block coordinate and residue,
    so both sides have the same value at every expanded site.
\end{proof}

\begin{definition}[Finite local-algebra blocking equivalence]
    \label{def:qca_local_blocking}
    \lean{SpinChain.localBlocking}
    \leanok
    \uses{def:qca_local_algebra, def:qca_configuration_blocking}
    Reindexing matrix rows and columns by the configuration blocking
    equivalence gives a star-algebra equivalence
    \begin{align}
        \mathcal A^{(d^L)}_\Lambda
         & \simeq \mathcal A^{(d)}_{\widehat\Lambda_L}.
    \end{align}
    This is the finite-region equivalence underlying the algebraic-local
    blocking construction below.  No quasi-local completion, compatibility
    with translations, propagation statement, circuit representation, or
    transport of quantum cellular automata is asserted here.
\end{definition}

\begin{theorem}[Finite local-algebra blocking formulas]
    \label{thm:qca_local_blocking_formulas}
    \lean{SpinChain.localBlocking_apply,SpinChain.localBlocking_norm,
        SpinChain.localBlocking_isometry}
    \leanok
    \uses{def:qca_local_blocking}
    Matrix entries are transported by applying the inverse configuration
    equivalence to both indices.  The blocking equivalence preserves the
    operator norm and is an isometry.
\end{theorem}

\begin{proof}\leanok
    \uses{def:qca_local_blocking}
    The entry formula is matrix reindexing.  A star-algebra equivalence between
    finite matrix algebras preserves the C-star norm and hence distances.
\end{proof}

\begin{definition}[Finite block hull]
    \label{def:qca_block_hull}
    \lean{SpinChain.blockHull}
    \leanok
    For $L>0$ and a finite original region $\Delta\subset\mathbb Z$, define
    its block hull by
    \begin{align}
        \operatorname{Hull}_L(\Delta)
         & =\{\lfloor z/L\rfloor:z\in\Delta\}.
    \end{align}
    This is the finite set of length-$L$ blocks that meet $\Delta$.
\end{definition}

\begin{lemma}[Block-hull containment and monotonicity]
    \label{lem:qca_block_hull_properties}
    \lean{SpinChain.subset_expandedRegion_blockHull,SpinChain.blockHull_mono,
        SpinChain.blockHull_expandedRegion}
    \leanok
    \uses{def:qca_block_hull, def:qca_expanded_region}
    Let $L>0$.  Every finite original region is contained in the expansion of
    its block hull,
    \begin{align}
        \Delta\subseteq
        \widehat{\operatorname{Hull}_L(\Delta)}_L.
    \end{align}
    If $\Delta\subseteq\Theta$, then
    $\operatorname{Hull}_L(\Delta)\subseteq
        \operatorname{Hull}_L(\Theta)$.  Expansion and block hull are exact in the
    blocked direction:
    \begin{align}
        \operatorname{Hull}_L(\widehat\Lambda_L) & =\Lambda.
    \end{align}
\end{lemma}

\begin{proof}\leanok
    \uses{def:qca_block_hull, thm:qca_expanded_region_coordinates}
    The quotient coordinate of each $z\in\Delta$ belongs to the image defining
    the hull, which proves containment after expansion.  Monotonicity follows
    by retaining the same witness $z$ in the larger region.  Finally, every
    expanded site has quotient in $\Lambda$, while for each $x\in\Lambda$ the
    site $Lx$ belongs to $\widehat\Lambda_L$ and has quotient $x$.  Hence
    $\operatorname{Hull}_L(\widehat\Lambda_L)=\Lambda$.
\end{proof}

\begin{lemma}[Blocking commutes with finite local inclusions]
    \label{lem:qca_local_blocking_inclusion}
    \lean{SpinChain.localBlocking_localInclusion}
    \leanok
    \uses{def:qca_local_blocking, thm:qca_configuration_blocking_formulas, def:qca_local_inclusion}
    Let $L>0$.  For finite blocked regions $\Lambda\subseteq\Gamma$, the
    square formed by finite blocking and the two canonical local inclusions
    commutes:
    \begin{align}
        \mathcal B_{L,\Gamma}\,\iota_{\Lambda,\Gamma}
         & =\iota_{\widehat\Lambda_L,\widehat\Gamma_L}\,
        \mathcal B_{L,\Lambda}.
    \end{align}
\end{lemma}

\begin{proof}\leanok
    \uses{thm:qca_local_blocking_formulas, thm:qca_configuration_blocking_formulas}
    For configurations $x,y$ on $\widehat\Gamma_L$, let $\widetilde x$ and
    $\widetilde y$ be their block encodings on $\Gamma$.  The two matrix
    entries in the commuting square are
    \begin{align}
        \bigl[\mathcal B_{L,\Gamma}
            (\iota_{\Lambda,\Gamma}(A))\bigr]_{x,y}
         & =A_{\widetilde x|_\Lambda,\widetilde y|_\Lambda}
        \mathbf 1_{\{\widetilde x|_{\Gamma\setminus\Lambda}
        =\widetilde y|_{\Gamma\setminus\Lambda}\}},         \\
        \bigl[\iota_{\widehat\Lambda_L,\widehat\Gamma_L}
            (\mathcal B_{L,\Lambda}(A))\bigr]_{x,y}
         & =A_{\widetilde x|_\Lambda,\widetilde y|_\Lambda}
        \mathbf 1_{\{x|_{\widehat\Gamma_L\setminus\widehat\Lambda_L}
                =y|_{\widehat\Gamma_L\setminus\widehat\Lambda_L}\}}.
    \end{align}
    Restriction commutes with block encoding, and decoding identifies
    $\Gamma\setminus\Lambda$ with the original sites in
    $\widehat\Gamma_L\setminus\widehat\Lambda_L$.  Hence the two indicator
    conditions are equivalent and the entries agree.
\end{proof}

\begin{definition}[Forward and reverse algebraic-local blocking maps]
    \label{def:qca_algebraic_blocking_maps}
    \lean{SpinChain.algebraicLocalBlockingHom,
        SpinChain.algebraicLocalUnblockingHom}
    \leanok
    \uses{def:qca_algebraic_local_algebra, def:qca_local_blocking, def:qca_block_hull}
    Let $L>0$.  The compatible finite blocking maps induce a star-algebra
    homomorphism
    \begin{align}
        \mathcal B_L:\mathcal A_{\mathrm{loc}}^{(d^L)}
        \longrightarrow\mathcal A_{\mathrm{loc}}^{(d)}.
    \end{align}
    In the reverse direction there is a star-algebra homomorphism
    \begin{align}
        \mathcal U_L:\mathcal A_{\mathrm{loc}}^{(d)}
        \longrightarrow\mathcal A_{\mathrm{loc}}^{(d^L)}.
    \end{align}
    A representative on an arbitrary original region $\Delta$ is first
    included into $\widehat{\operatorname{Hull}_L(\Delta)}_L$ and then finite
    blocking is inverted.  Thus the reverse construction does not require
    $\Delta$ itself to be an expanded region.
\end{definition}

\begin{proof}\leanok
    \uses{lem:qca_block_hull_properties, lem:qca_local_blocking_inclusion, thm:qca_algebraic_local_algebra_universal}
    For $\Lambda\subseteq\Gamma$, the finite commuting square gives
    \begin{align}
        \iota_{\widehat\Gamma_L}
        \mathcal B_{L,\Gamma}(\iota_{\Lambda,\Gamma}(A))
         & =\iota_{\widehat\Lambda_L}\mathcal B_{L,\Lambda}(A),
    \end{align}
    so the forward finite-region maps form a compatible family.  For the
    reverse family, block-hull monotonicity sends $\Delta\subseteq\Theta$ to
    $\operatorname{Hull}_L(\Delta)\subseteq
        \operatorname{Hull}_L(\Theta)$.  Transitivity of the canonical inclusions
    and the inverse commuting square then give
    \begin{align}
        \iota_{\operatorname{Hull}_L(\Theta)}
        \mathcal U_{L,\Theta}(\iota_{\Delta,\Theta}(A))
         & =\iota_{\operatorname{Hull}_L(\Delta)}
        \mathcal U_{L,\Delta}(A).
    \end{align}
    The universal property of the algebraic direct limit therefore induces
    both star-algebra homomorphisms.
\end{proof}

\begin{lemma}[Evaluation of the algebraic-local blocking maps]
    \label{lem:qca_algebraic_blocking_maps_evaluation}
    \lean{SpinChain.algebraicLocalBlockingHom_localObservable,
        SpinChain.algebraicLocalUnblockingHom_localObservable,
        SpinChain.algebraicLocalBlockingHom_comp_unblockingHom}
    \leanok
    \uses{def:qca_algebraic_blocking_maps, lem:qca_block_hull_properties}
    Let $L>0$.  The forward map applies finite blocking to a representative on
    $\Lambda$ and regards the result as supported on $\widehat\Lambda_L$.  The
    reverse map enlarges a representative on $\Delta$ to the expansion of its
    block hull, applies the inverse finite blocking equivalence, and regards
    the result as supported on $\operatorname{Hull}_L(\Delta)$.  The forward
    composite is the identity,
    \begin{align}
        \mathcal B_L\mathcal U_L
         & =\operatorname{id}_{\mathcal A_{\mathrm{loc}}^{(d)}}.
    \end{align}
    This holds for every on-site dimension $d$.
\end{lemma}

\begin{proof}\leanok
    \uses{lem:qca_local_blocking_inclusion, lem:qca_block_hull_properties}
    These are the evaluation formulas for the two compatible-family lifts.
    Forward compatibility is the commuting square.  Reverse compatibility
    follows by monotonicity of block hulls, composition of local inclusions,
    and the same square applied in the inverse direction.  Forward blocking
    after the reverse map reduces to the canonical inclusion of $\Delta$ into
    its expanded block hull.
\end{proof}

\begin{definition}[Algebraic-local site-blocking equivalence]
    \label{def:qca_algebraic_blocking}
    \lean{SpinChain.algebraicLocalBlocking}
    \leanok
    \uses{def:qca_algebraic_blocking_maps}
    For $d>0$ and $L>0$, forward blocking gives a star-algebra equivalence
    \begin{align}
        \mathcal A_{\mathrm{loc}}^{(d^L)}
         & \simeq \mathcal A_{\mathrm{loc}}^{(d)}.
    \end{align}
    This formalizes only the algebraic-local change of site grouping used at
    lines~2308 and 2313--2320 of \cite{Cirac2017MPU}.  It does not assert the
    QCA-to-MPU converse at line~2308, extend blocking to the quasi-local
    completion, intertwine translations, provide a propagation estimate or a
    circuit representation, or transport a quantum cellular automaton.
\end{definition}

\begin{proof}\leanok
    \uses{lem:qca_algebraic_blocking_maps_evaluation, thm:qca_local_observable_finite_support, lem:qca_local_blocking_inclusion}
    The identity $\mathcal B_L\mathcal U_L=\operatorname{id}$ makes
    $\mathcal B_L$ surjective.  For injectivity, represent two blocked
    observables on finite regions $\Lambda$ and $\Gamma$, enlarge both to
    $\Omega=\Lambda\cup\Gamma$, and suppose their images under
    $\mathcal B_L$ agree.  Injectivity of the canonical map from the finite
    algebra on $\widehat\Omega_L$, followed by injectivity of
    $\mathcal B_{L,\Omega}$ and the finite commuting square, gives equality of
    the two representatives in the algebra on $\Omega$.  Thus
    $\mathcal B_L$ is bijective, with inverse $\mathcal U_L$, and defines the
    stated star-algebra equivalence.
\end{proof}

\begin{lemma}[Algebraic-local blocking formulas and norm]
    \label{lem:qca_algebraic_blocking_formulas}
    \lean{SpinChain.algebraicLocalBlocking_apply,
        SpinChain.algebraicLocalBlocking_localObservable,
        SpinChain.algebraicLocalBlocking_symm_apply,
        SpinChain.algebraicLocalBlocking_symm_localObservable,
        SpinChain.algebraicLocalBlocking_symm_localObservable_expandedRegion,
        SpinChain.algebraicLocalBlocking_norm,
        SpinChain.algebraicLocalBlocking_isometry,
        SpinChain.algebraicLocalBlocking_symm_isometry}
    \leanok
    \uses{def:qca_algebraic_blocking, lem:qca_algebraic_blocking_maps_evaluation}
    Let $d>0$ and $L>0$.  The equivalence agrees with the forward lift, and
    its inverse agrees with the block-hull unblocking map.  For
    $A\in\mathcal A^{(d^L)}_\Lambda$ and
    $C\in\mathcal A^{(d)}_\Delta$,
    \begin{align}
        \mathcal B_L\bigl(\iota_\Lambda(A)\bigr)
         & =\iota_{\widehat\Lambda_L}
        \bigl(\mathcal B_{L,\Lambda}(A)\bigr),    \\
        \mathcal B_L^{-1}\bigl(\iota_\Delta(C)\bigr)
         & =\iota_{\operatorname{Hull}_L(\Delta)}
        \left(\mathcal B^{-1}_{L,\operatorname{Hull}_L(\Delta)}
        \bigl(\iota_{\Delta,
                    \widehat{\operatorname{Hull}_L(\Delta)}_L}(C)\bigr)\right).
    \end{align}
    In particular, if $C\in\mathcal A^{(d)}_{\widehat\Lambda_L}$, then
    \begin{align}
        \mathcal B_L^{-1}\bigl(\iota_{\widehat\Lambda_L}(C)\bigr)
         & =\iota_\Lambda\bigl(\mathcal B_{L,\Lambda}^{-1}(C)\bigr).
    \end{align}
    The forward and inverse maps are operator-norm isometries.  In particular,
    each preserves the norm.
\end{lemma}

\begin{proof}\leanok
    \uses{lem:qca_algebraic_blocking_maps_evaluation, lem:qca_block_hull_properties, thm:qca_local_blocking_formulas, def:qca_algebraic_local_operator_norm}
    The first two formulas are the evaluation identities for the forward and
    reverse maps.  For an expanded region, the block-hull identity removes the
    enlargement and leaves the inverse finite blocking equivalence.  On a
    finite representative, forward norm preservation reduces to norm
    preservation by finite blocking and by the canonical local inclusion;
    distance preservation follows by applying the norm identity to
    differences.  The inverse isometry follows from the forward isometry and
    the right-inverse identity, rather than repeating the norm calculation.
\end{proof}

\begin{definition}[Quasi-local site-blocking equivalence]
    \label{def:qca_quasi_local_blocking}
    \lean{SpinChain.quasiLocalBlocking}
    \leanok
    \uses{def:qca_quasi_local_algebra, def:qca_algebraic_blocking, lem:qca_algebraic_blocking_formulas}
    For $d>0$ and $L>0$, site blocking extends uniquely by continuity to a
    star-algebra equivalence
    \begin{align}
        \overline{\mathcal B}_L:
        \mathcal A^{(d^L)} & \simeq\mathcal A^{(d)}.
    \end{align}
    This is only the completion of the change of site grouping used at
    \cite[lines~2308 and 2313--2320]{Cirac2017MPU}.  This slice does not define
    scaled translations or carry-aware neighborhoods, prove propagation
    estimates, construct blocked automorphisms, transport the property of being
    a quantum cellular automaton, realize a circuit, or reconstruct an MPU from
    a quantum cellular automaton.
\end{definition}

\begin{proof}\leanok
    \uses{lem:qca_algebraic_blocking_formulas, def:completion_star_alg_equiv}
    The algebraic blocking equivalence and its inverse are isometries, hence
    continuous.  Extending both maps to the operator-norm completions gives
    mutually inverse star-algebra homomorphisms.
\end{proof}

\begin{lemma}[Quasi-local blocking evaluation and norm]
    \label{lem:qca_quasi_local_blocking_formulas}
    \lean{SpinChain.quasiLocalBlocking_apply,
        SpinChain.quasiLocalBlocking_algebraicToQuasiLocal,
        SpinChain.quasiLocalBlocking_symm_algebraicToQuasiLocal,
        SpinChain.quasiLocalBlocking_quasiLocalObservable,
        SpinChain.quasiLocalBlocking_symm_quasiLocalObservable,
        SpinChain.quasiLocalBlocking_symm_quasiLocalObservable_expandedRegion,
        SpinChain.norm_quasiLocalBlocking,
        SpinChain.quasiLocalBlocking_isometry}
    \leanok
    \uses{def:qca_quasi_local_blocking, lem:qca_algebraic_blocking_formulas, lem:qca_quasi_local_observable}
    Let $d>0$ and $L>0$.  Write $j^{(d)}$ for the canonical map from the
    algebraic-local algebra into its completion.  The underlying map of
    $\overline{\mathcal B}_L$ is the continuous extension of $\mathcal B_L$,
    and on the two dense algebraic-local subalgebras,
    \begin{align}
        \overline{\mathcal B}_L\bigl(j^{(d^L)}(A)\bigr)
         & =j^{(d)}\bigl(\mathcal B_L(A)\bigr),        \\
        \overline{\mathcal B}_L^{-1}\bigl(j^{(d)}(C)\bigr)
         & =j^{(d^L)}\bigl(\mathcal B_L^{-1}(C)\bigr).
    \end{align}
    For $A\in\mathcal A^{(d^L)}_\Lambda$ and
    $C\in\mathcal A^{(d)}_\Delta$, these identities become
    \begin{align}
        \overline{\mathcal B}_L\bigl(j^{(d^L)}_\Lambda(A)\bigr)
         & =j^{(d)}_{\widehat\Lambda_L}
        \bigl(\mathcal B_{L,\Lambda}(A)\bigr),        \\
        \overline{\mathcal B}_L^{-1}\bigl(j^{(d)}_\Delta(C)\bigr)
         & =j^{(d^L)}_{\operatorname{Hull}_L(\Delta)}
        \left(\mathcal B^{-1}_{L,\operatorname{Hull}_L(\Delta)}
        \bigl(\iota_{\Delta,
                    \widehat{\operatorname{Hull}_L(\Delta)}_L}(C)\bigr)\right).
    \end{align}
    If $C\in\mathcal A^{(d)}_{\widehat\Lambda_L}$, the inverse formula reduces
    exactly to
    \begin{align}
        \overline{\mathcal B}_L^{-1}
        \bigl(j^{(d)}_{\widehat\Lambda_L}(C)\bigr)
         & =j^{(d^L)}_\Lambda\bigl(\mathcal B_{L,\Lambda}^{-1}(C)\bigr).
    \end{align}
    Finally, for all quasi-local observables $X$ and $Y$ on the blocked chain,
    \begin{align}
        \lVert\overline{\mathcal B}_L(X)\rVert & =\lVert X\rVert,   \\
        \lVert\overline{\mathcal B}_L(X)-
        \overline{\mathcal B}_L(Y)\rVert       & =\lVert X-Y\rVert.
    \end{align}
\end{lemma}

\begin{proof}\leanok
    \uses{def:completion_star_alg_equiv, def:qca_quasi_local_observable, lem:qca_algebraic_blocking_formulas, lem:qca_block_hull_properties}
    The first identity is the defining completion map, and the two dense
    formulas are its canonical-image formulas for the algebraic equivalence and
    its inverse.  Composing them with the finite-region canonical maps gives the
    two finite formulas.  The exact expanded-region formula uses
    $\operatorname{Hull}_L(\widehat\Lambda_L)=\Lambda$.  A star-algebra
    equivalence of C-star algebras preserves the norm; applying this to $X-Y$
    gives distance preservation.
\end{proof}

\begin{lemma}[Support transport under site blocking]
    \label{lem:qca_quasi_local_blocking_support}
    \lean{SpinChain.quasiLocalBlocking_supportedIn,
        SpinChain.quasiLocalBlocking_symm_supportedIn,
        SpinChain.quasiLocalBlocking_supportedIn_expandedRegion_iff}
    \leanok
    \uses{lem:qca_quasi_local_blocking_formulas, def:qca_quasi_local_supported_in, lem:qca_block_hull_properties}
    Let $d>0$ and $L>0$.  If a blocked-chain observable $X$ is supported in
    $\Lambda$, then $\overline{\mathcal B}_L(X)$ is supported in
    $\widehat\Lambda_L$.  If an original-chain observable $Y$ is supported in
    $\Delta$, then $\overline{\mathcal B}_L^{-1}(Y)$ is supported in
    $\operatorname{Hull}_L(\Delta)$.  In particular, support transport is exact
    on expanded regions:
    \begin{align}
        \overline{\mathcal B}_L(X)
        \text{ is supported in }\widehat\Lambda_L
        \quad\Longleftrightarrow\quad
        X\text{ is supported in }\Lambda.
    \end{align}
\end{lemma}

\begin{proof}\leanok
    \uses{lem:qca_quasi_local_blocking_formulas, lem:qca_block_hull_properties}
    Choose a finite-region representative and apply the corresponding forward
    or inverse finite evaluation formula.  For the reverse implication in the
    exact statement, apply inverse blocking and simplify the inverse composite
    and the block hull of the expanded region.
\end{proof}

\begin{lemma}[Expansion of a translated blocked region]
    \label{lem:qca_expanded_region_translation}
    \lean{SpinChain.expandedRegion_translateRegion}
    \leanok
    \uses{def:qca_translate_region, def:qca_expanded_region}
    Let $L>0$, let $a\in\mathbb Z$, and let $\Lambda\subset\mathbb Z$ be a
    finite blocked region.  Then
    \begin{align}
        \widehat{\Lambda+a}_L
         & =\widehat\Lambda_L+aL.
    \end{align}
    This is the translation rule for the site grouping used in
    \cite[lines~2308 and 2313--2320]{Cirac2017MPU}.
\end{lemma}

\begin{proof}\leanok
    \uses{thm:qca_expanded_region_coordinates}
    Write each original site uniquely as $xL+r$, with $0\leq r<L$.
    Translation of the blocked coordinate by $a$ sends this site to
    $(x+a)L+r=xL+r+aL$.  The argument is unchanged when $x$ or $a$ is
    negative.
\end{proof}

\begin{definition}[Carry-aware blocked neighborhood]
    \label{def:qca_blocked_neighborhood}
    \lean{SpinChain.blockedNeighborhood}
    \leanok
    \uses{def:qca_expanded_region, def:qca_block_hull, def:qca_region_sumset}
    Let $L>0$ and let $\mathcal N\subset\mathbb Z$ be a finite original-lattice
    neighborhood.  Write $\widehat{\{0\}}_L=\{0,\ldots,L-1\}$ for the sites in
    the reference block.  Define the induced blocked-lattice neighborhood by
    \begin{align}
        \mathcal C_L(\mathcal N)
         & =\operatorname{Hull}_L
        \bigl(\widehat{\{0\}}_L+\mathcal N\bigr).
    \end{align}
    The full reference block is included because the block reached after a
    displacement can depend on the starting residue.  This is the finite-region
    geometry implicit in the grouping of sites in
    \cite[lines~2308 and 2313--2320]{Cirac2017MPU}.
\end{definition}

\begin{lemma}[Residue description of the blocked neighborhood]
    \label{lem:qca_blocked_neighborhood_membership}
    \lean{SpinChain.mem_blockedNeighborhood,
        SpinChain.mem_blockedNeighborhood_iff_exists_eq}
    \leanok
    \uses{def:qca_blocked_neighborhood, thm:qca_expanded_region_coordinates}
    Let $L>0$, let $\mathcal N\subset\mathbb Z$ be finite, and let
    $b\in\mathbb Z$.  Then
    \begin{align}
        b\in\mathcal C_L(\mathcal N)
        \quad\Longleftrightarrow\quad
         & \exists r\in\{0,\ldots,L-1\},\ \exists n\in\mathcal N,
        \quad \left\lfloor\frac{r+n}{L}\right\rfloor=b,             \\
        \quad\Longleftrightarrow\quad
         & \exists r,s\in\{0,\ldots,L-1\},\ \exists n\in\mathcal N,
        \quad r+n=bL+s.
    \end{align}
    The quotient is the Euclidean quotient by the positive integer $L$.
    Consequently these formulas include negative displacements and both
    positive and negative block carries.  In general,
    $\operatorname{Hull}_L(\mathcal N)$ alone omits some carries caused by the
    starting residue.  This makes explicit the boundary geometry implicit in
    \cite[lines~2313--2320]{Cirac2017MPU}.
\end{lemma}

\begin{proof}\leanok
    \uses{def:qca_blocked_neighborhood, def:qca_block_hull, thm:qca_expanded_region_coordinates}
    An element of $\widehat{\{0\}}_L+\mathcal N$ has the form $r+n$ with
    $0\leq r<L$ and $n\in\mathcal N$.  Taking its block coordinate gives the
    first equivalence.  Euclidean division writes $r+n=bL+s$ with the unique
    remainder $0\leq s<L$, which gives the second equivalence without any sign
    restriction on $n$.
\end{proof}

\begin{lemma}[Exact block hull after adding a neighborhood]
    \label{lem:qca_block_hull_region_sumset}
    \lean{SpinChain.blockHull_regionSumset_expandedRegion}
    \leanok
    \uses{def:qca_blocked_neighborhood, def:qca_region_sumset, def:qca_expanded_region, def:qca_block_hull}
    Let $L>0$.  For every finite blocked region $\Lambda$ and every finite
    original-lattice neighborhood $\mathcal N$,
    \begin{align}
        \operatorname{Hull}_L
        \bigl(\widehat\Lambda_L+\mathcal N\bigr)
         & =\Lambda+\mathcal C_L(\mathcal N).
    \end{align}
    This is the exact finite-region carry identity associated with the site
    grouping in \cite[lines~2313--2320]{Cirac2017MPU}.
\end{lemma}

\begin{proof}\leanok
    \uses{lem:qca_blocked_neighborhood_membership, thm:qca_expanded_region_coordinates}
    Write an expanded site as $xL+r$, where $x\in\Lambda$ and $0\leq r<L$.
    For $n\in\mathcal N$, Euclidean division gives
    \begin{align}
        \left\lfloor\frac{xL+r+n}{L}\right\rfloor
         & =x+\left\lfloor\frac{r+n}{L}\right\rfloor.
    \end{align}
    The second summand ranges over exactly $\mathcal C_L(\mathcal N)$ by
    Lemma~\ref{lem:qca_blocked_neighborhood_membership}.  This proves both
    inclusions, including for negative $x$ and $n$.
\end{proof}

\begin{lemma}[Carry-aware expansion containment]
    \label{lem:qca_region_sumset_expanded_region_subset}
    \lean{SpinChain.regionSumset_expandedRegion_subset}
    \leanok
    \uses{lem:qca_block_hull_region_sumset, lem:qca_block_hull_properties}
    Let $L>0$.  For every finite blocked region $\Lambda$ and every finite
    original-lattice neighborhood $\mathcal N$,
    \begin{align}
        \widehat\Lambda_L+\mathcal N
         & \subseteq
        \widehat{\Lambda+\mathcal C_L(\mathcal N)}_L.
    \end{align}
\end{lemma}

\begin{proof}\leanok
    \uses{lem:qca_block_hull_region_sumset, lem:qca_block_hull_properties}
    Every finite original region is contained in the expansion of its block
    hull.  Apply this containment to
    $\widehat\Lambda_L+\mathcal N$ and use
    Lemma~\ref{lem:qca_block_hull_region_sumset} to identify that hull.
\end{proof}

\begin{lemma}[Algebraic-local blocking intertwines scaled translations]
    \label{lem:qca_algebraic_blocking_translation}
    \lean{SpinChain.algebraicLocalBlocking_translation}
    \leanok
    \uses{lem:qca_expanded_region_translation, def:qca_algebraic_blocking, def:qca_algebraic_local_translation}
    Let $d,L>0$.  If $\mathcal B_L$ denotes the algebraic-local blocking
    equivalence, then for every $a\in\mathbb Z$,
    \begin{align}
        \mathcal B_L\,\tau^{(d^L)}_a
         & =\tau^{(d)}_{aL}\,\mathcal B_L.
    \end{align}
    This is the translation rule induced by the blocking of sites used in
    \cite[lines~2308 and 2313--2320]{Cirac2017MPU}.
\end{lemma}

\begin{proof}\leanok
    \uses{lem:qca_expanded_region_translation, lem:qca_algebraic_blocking_formulas, def:qca_algebraic_local_translation}
    On a representative supported in a finite blocked region $\Lambda$, both
    sides apply the same matrix reindexing.  Their supports agree because
    \begin{align}
        \widehat{\Lambda+a}_L
         & =\widehat\Lambda_L+aL.
    \end{align}
    Equality on all finite-region representatives proves the identity on the
    algebraic direct limit.
\end{proof}

\begin{lemma}[Quasi-local blocking intertwines scaled translations]
    \label{lem:qca_quasi_local_blocking_translation}
    \lean{SpinChain.quasiLocalBlocking_translation}
    \leanok
    \uses{lem:qca_algebraic_blocking_translation, def:qca_quasi_local_blocking, def:qca_quasi_local_translation}
    Let $d,L>0$.  If $\overline{\mathcal B}_L$ denotes the quasi-local blocking
    equivalence, then for every $a\in\mathbb Z$,
    \begin{align}
        \overline{\mathcal B}_L\,\overline\tau^{(d^L)}_a
         & =\overline\tau^{(d)}_{aL}\,\overline{\mathcal B}_L.
    \end{align}
    This is the completion of the translation rule induced by the site grouping
    in \cite[lines~2308 and 2313--2320]{Cirac2017MPU}.
\end{lemma}

\begin{proof}\leanok
    \uses{lem:qca_algebraic_blocking_translation, lem:qca_quasi_local_blocking_formulas, def:qca_quasi_local_translation}
    The two sides are continuous maps on the quasi-local algebra.  On the dense
    algebraic-local subalgebra their values agree by
    Lemma~\ref{lem:qca_algebraic_blocking_translation}.  They therefore agree
    on the completion.
\end{proof}

Line~2308 of \cite{Cirac2017MPU} invokes a converse representation theorem:
after blocking finitely many sites, the restriction of every one-dimensional
QCA to the local algebra is represented by a depth-two circuit of
nearest-neighbor unitaries, and hence by an MPU in standard form.  The results
below concern the preceding change of grouping.  A
given QCA remains a QCA under any prescribed finite blocking.  This one-way
transport alone does not produce the circuit or MPU representation asserted by
the converse theorem.

\begin{definition}[Blocked automorphism]
    \label{def:qca_blocked_automorphism}
    \lean{SpinChain.blockedAutomorphism}
    \leanok
    \uses{def:qca_quasi_local_blocking}
    Let $d,L>0$, let
    $\overline{\mathcal B}_L\colon\mathcal A^{(d^L)}\simeq\mathcal A^{(d)}$
    be the quasi-local blocking equivalence, and let $\omega$ be a star-algebra
    automorphism of $\mathcal A^{(d)}$.  The induced automorphism of the blocked
    chain is
    \begin{align}
        \omega^{[L]}
         & =\overline{\mathcal B}_L^{-1}\circ\omega\circ
        \overline{\mathcal B}_L.
    \end{align}
\end{definition}

\begin{lemma}[Evaluation of the blocked automorphism]
    \label{lem:qca_blocked_automorphism_apply}
    \lean{SpinChain.blockedAutomorphism_apply}
    \leanok
    \uses{def:qca_blocked_automorphism}
    For every $X\in\mathcal A^{(d^L)}$,
    \begin{align}
        \omega^{[L]}(X)
         & =\overline{\mathcal B}_L^{-1}
        \bigl(\omega(\overline{\mathcal B}_L(X))\bigr).
    \end{align}
\end{lemma}

\begin{proof}\leanok
    \uses{def:qca_blocked_automorphism}
    This is the evaluation formula for the defining conjugation.
\end{proof}

\begin{lemma}[Inverse of the blocked automorphism]
    \label{lem:qca_blocked_automorphism_symm}
    \lean{SpinChain.blockedAutomorphism_symm}
    \leanok
    \uses{def:qca_blocked_automorphism}
    The inverse of the blocked automorphism is the blocked inverse:
    \begin{align}
        (\omega^{[L]})^{-1}
         & =(\omega^{-1})^{[L]}.
    \end{align}
\end{lemma}

\begin{proof}\leanok
    \uses{lem:qca_blocked_automorphism_apply}
    Evaluate both sides on a blocked-chain observable.  The blocking equivalence
    and its inverse cancel on the two sides of $\omega^{-1}$.
\end{proof}

\begin{lemma}[Propagation neighborhood after site blocking]
    \label{lem:qca_propagates_within_blocked}
    \lean{SpinChain.PropagatesWithin.blocked}
    \leanok
    \uses{def:qca_blocked_automorphism, def:qca_blocked_neighborhood, def:qca_propagates_within}
    Let $\mathcal N\subset\mathbb Z$ be finite.  If $\omega$ propagates within
    $\mathcal N$, then $\omega^{[L]}$ propagates within the carry-aware blocked
    neighborhood $\mathcal C_L(\mathcal N)$.  Explicitly, for every finite
    blocked region $\Lambda$ and every $X\in\mathcal A^{(d^L)}$ supported in
    $\Lambda$, the observable $\omega^{[L]}(X)$ is supported in
    $\Lambda+\mathcal C_L(\mathcal N)$.
\end{lemma}

\begin{proof}\leanok
    \uses{lem:qca_blocked_automorphism_apply, lem:qca_quasi_local_blocking_support, lem:qca_block_hull_region_sumset}
    Forward blocking sends support in $\Lambda$ to support in
    $\widehat\Lambda_L$.  Apply the propagation bound for $\omega$, then apply
    inverse blocking.  The resulting support is
    \begin{align}
        \operatorname{Hull}_L
        \bigl(\widehat\Lambda_L+\mathcal N\bigr)
         & =\Lambda+\mathcal C_L(\mathcal N)
    \end{align}
    by the exact carry-aware block-hull identity.
\end{proof}

\begin{lemma}[Finite propagation after site blocking]
    \label{lem:qca_finite_propagation_blocked}
    \lean{SpinChain.HasFinitePropagation.blocked}
    \leanok
    \uses{def:qca_blocked_automorphism, def:qca_finite_propagation, lem:qca_propagates_within_blocked}
    If $\omega$ has finite forward propagation, then $\omega^{[L]}$ has finite
    forward propagation.
\end{lemma}

\begin{proof}\leanok
    \uses{lem:qca_propagates_within_blocked}
    Choose a finite propagation neighborhood $\mathcal N$ for $\omega$.
    Lemma~\ref{lem:qca_propagates_within_blocked} gives the finite neighborhood
    $\mathcal C_L(\mathcal N)$ for $\omega^{[L]}$.
\end{proof}

\begin{theorem}[Translation covariance of the blocked automorphism]
    \label{thm:qca_blocked_translation_covariant_iff}
    \lean{SpinChain.translationCovariant_blockedAutomorphism_iff}
    \leanok
    \uses{def:qca_blocked_automorphism, def:qca_translation_covariant, def:qca_quasi_local_translation}
    The blocked automorphism is translation covariant if and only if the
    original automorphism commutes with translation by one whole block:
    \begin{align}
        \omega^{[L]}\text{ is translation covariant}
        \quad\Longleftrightarrow\quad
        \omega\circ\overline\tau^{(d)}_L
         & =\overline\tau^{(d)}_L\circ\omega.
    \end{align}
    Thus blocked unit-translation covariance records commutation of $\omega$
    with translation by $L$, not necessarily commutation with original unit
    translation.  For $L\neq1$, the latter condition can be strictly stronger.
\end{theorem}

\begin{proof}\leanok
    \uses{lem:qca_blocked_automorphism_apply, thm:qca_translation_covariant_iff_commute_one, lem:qca_quasi_local_blocking_translation}
    By Theorem~\ref{thm:qca_translation_covariant_iff_commute_one}, translation
    covariance on the blocked chain is equivalent to commutation with blocked
    unit translation.  Under $\overline{\mathcal B}_L$, blocked unit translation
    corresponds to original translation by $L$.  Conjugating the commutation
    identity by the blocking equivalence proves both implications.
\end{proof}

\begin{lemma}[Translation covariance after site blocking]
    \label{lem:qca_translation_covariant_blocked}
    \lean{SpinChain.TranslationCovariant.blocked}
    \leanok
    \uses{def:qca_blocked_automorphism, def:qca_translation_covariant, thm:qca_blocked_translation_covariant_iff}
    If $\omega$ is translation covariant, then $\omega^{[L]}$ is translation
    covariant.
\end{lemma}

\begin{proof}\leanok
    \uses{thm:qca_blocked_translation_covariant_iff}
    Translation covariance of $\omega$ gives commutation with
    $\overline\tau^{(d)}_L$.  Apply
    Theorem~\ref{thm:qca_blocked_translation_covariant_iff}.
\end{proof}

\begin{lemma}[QCA transport through site blocking]
    \label{lem:qca_blocked}
    \lean{SpinChain.IsQCA.blocked}
    \leanok
    \uses{def:qca, def:qca_blocked_automorphism, lem:qca_finite_propagation_blocked, lem:qca_translation_covariant_blocked}
    Let $d,L>0$.  If $\omega$ is a one-dimensional QCA on
    $\mathcal A^{(d)}$, then
    \begin{align}
        \omega^{[L]}
         & =\overline{\mathcal B}_L^{-1}\circ\omega\circ
        \overline{\mathcal B}_L
    \end{align}
    is a one-dimensional QCA on $\mathcal A^{(d^L)}$.
\end{lemma}

\begin{proof}\leanok
    \uses{lem:qca_finite_propagation_blocked, lem:qca_translation_covariant_blocked}
    The two defining properties of a one-dimensional QCA are finite forward
    propagation and translation covariance.  They are preserved by
    Lemma~\ref{lem:qca_finite_propagation_blocked} and
    Lemma~\ref{lem:qca_translation_covariant_blocked}, respectively.
\end{proof}

Line~2308 of \cite{Cirac2017MPU} invokes the external
Schumacher--Werner structure theorem to obtain support algebras, a depth-two
nearest-neighbor circuit, and an MPU in standard form.  The next statements
establish only the elementary normalization of a finite propagation
neighborhood to $[-1,1]\cap\mathbb Z$ by regrouping sites.  They do not prove
any part of that external structure theorem beyond this change of grouping.

\begin{lemma}[Nearest-neighbor carries from a strict displacement bound]
    \label{lem:qca_blocked_neighborhood_subset_nearest_neighbor}
    \lean{SpinChain.blockedNeighborhood_subset_Icc_one}
    \leanok
    \uses{def:qca_blocked_neighborhood}
    Let $L>0$ and let $\mathcal N\subset\mathbb Z$ be finite.  If
    $|n|<L$ for every $n\in\mathcal N$, then
    \begin{align}
        \mathcal C_L(\mathcal N) & \subseteq[-1,1]\cap\mathbb Z.
    \end{align}
\end{lemma}

\begin{proof}\leanok
    \uses{lem:qca_blocked_neighborhood_membership}
    If $\mathcal N=\varnothing$, membership in
    $\mathcal C_L(\mathcal N)$ is impossible.  In general, a point
    $b\in\mathcal C_L(\mathcal N)$ supplies residues $0\leq r,s<L$ and
    $n\in\mathcal N$ such that
    \begin{align}
        r+n & =bL+s.
    \end{align}
    Since $-L<n<L$, the equality is incompatible with $b\leq-2$ and with
    $b\geq2$.  Thus $-1\leq b\leq1$, including for negative displacements.
\end{proof}

\begin{lemma}[Supremum criterion for nearest-neighbor carries]
    \label{lem:qca_blocked_neighborhood_subset_nearest_neighbor_sup}
    \lean{SpinChain.blockedNeighborhood_subset_Icc_one_of_sup}
    \leanok
    \uses{def:qca_blocked_neighborhood}
    Let $L>0$ and let $\mathcal N\subset\mathbb Z$ be finite.  If
    \begin{align}
        \sup_{n\in\mathcal N}|n| & <L,
    \end{align}
    then
    \begin{align}
        \mathcal C_L(\mathcal N) & \subseteq[-1,1]\cap\mathbb Z.
    \end{align}
\end{lemma}

\begin{proof}\leanok
    \uses{lem:qca_blocked_neighborhood_subset_nearest_neighbor}
    Every $|n|$ is bounded above by the finite supremum, so the strict
    supremum bound supplies the pointwise hypothesis.  When
    $\mathcal N=\varnothing$, the supremum is $0$ and the conclusion is
    vacuous.
\end{proof}

\begin{lemma}[Canonical nearest-neighbor block length]
    \label{lem:qca_blocked_neighborhood_canonical_nearest_neighbor}
    \lean{SpinChain.blockedNeighborhood_subset_Icc_one_sup_add_one}
    \leanok
    \uses{def:qca_blocked_neighborhood}
    For every finite $\mathcal N\subset\mathbb Z$, define
    \begin{align}
        L_{\mathcal N} & =1+\sup_{n\in\mathcal N}|n|.
    \end{align}
    Then $L_{\mathcal N}>0$ and
    \begin{align}
        \mathcal C_{L_{\mathcal N}}(\mathcal N)
         & \subseteq[-1,1]\cap\mathbb Z.
    \end{align}
    With the finite-supremum convention, $L_{\varnothing}=1$.
\end{lemma}

\begin{proof}\leanok
    \uses{lem:qca_blocked_neighborhood_subset_nearest_neighbor_sup}
    The finite supremum is strictly smaller than its successor, which is
    positive even when $\mathcal N$ is empty.
\end{proof}

\begin{lemma}[Nearest-neighbor propagation after a bounded blocking]
    \label{lem:qca_propagates_within_blocked_nearest_neighbor}
    \lean{SpinChain.PropagatesWithin.blocked_nearest_neighbor}
    \leanok
    \uses{def:qca_propagates_within, def:qca_blocked_automorphism}
    Let $d,L>0$.  Suppose that $\omega$ propagates within the finite
    neighborhood $\mathcal N$ and that $|n|<L$ for every $n\in\mathcal N$.
    For every finite blocked region $\Lambda$ and every observable
    $X\in\mathcal A^{(d^L)}$ supported in $\Lambda$, the observable
    $\omega^{[L]}(X)$ is supported in
    \begin{align}
        \Lambda+([-1,1]\cap\mathbb Z).
    \end{align}
\end{lemma}

\begin{proof}\leanok
    \uses{lem:qca_propagates_within_blocked, lem:qca_blocked_neighborhood_subset_nearest_neighbor, lem:qca_propagation_neighborhood_mono}
    Blocking first gives propagation within $\mathcal C_L(\mathcal N)$.
    The carry bound places this neighborhood inside
    $[-1,1]\cap\mathbb Z$, and enlargement of propagation neighborhoods gives
    the claim.
\end{proof}

\begin{theorem}[Finite propagation becomes nearest-neighbor]
    \label{thm:qca_finite_propagation_exists_blocked_nearest_neighbor}
    \lean{SpinChain.HasFinitePropagation.exists_blocked_nearest_neighbor}
    \leanok
    \uses{def:qca_finite_propagation, def:qca_propagates_within, def:qca_blocked_automorphism}
    Let $d>0$.  If $\omega$ has finite forward propagation, then there exists
    a positive integer $L$ such that, for every finite blocked region
    $\Lambda$ and every observable $X\in\mathcal A^{(d^L)}$ supported in
    $\Lambda$, the observable $\omega^{[L]}(X)$ is supported in
    \begin{align}
        \Lambda+([-1,1]\cap\mathbb Z).
    \end{align}
\end{theorem}

\begin{proof}\leanok
    \uses{lem:qca_propagates_within_blocked, lem:qca_blocked_neighborhood_canonical_nearest_neighbor, lem:qca_propagation_neighborhood_mono}
    Choose a finite propagation neighborhood $\mathcal N$ and take
    $L=1+\sup_{n\in\mathcal N}|n|$.  This $L$ is positive, including when
    $\mathcal N$ is empty.  Blocking gives propagation within
    $\mathcal C_L(\mathcal N)$; the canonical carry inclusion and propagation
    monotonicity then give the nearest-neighbor bound.
\end{proof}

\begin{theorem}[A QCA becomes nearest-neighbor after blocking]
    \label{thm:qca_exists_blocked_nearest_neighbor}
    \lean{SpinChain.IsQCA.exists_blocked_nearest_neighbor}
    \leanok
    \uses{def:qca, def:qca_blocked_automorphism, def:qca_propagates_within}
    Let $d>0$ and let $\omega$ be a one-dimensional QCA on
    $\mathcal A^{(d)}$.  There exists a positive integer $L$ such that
    $\omega^{[L]}$ is a QCA on $\mathcal A^{(d^L)}$ and, for every finite
    blocked region $\Lambda$ and every observable
    $X\in\mathcal A^{(d^L)}$ supported in $\Lambda$, the observable
    $\omega^{[L]}(X)$ is supported in
    \begin{align}
        \Lambda+([-1,1]\cap\mathbb Z).
    \end{align}
\end{theorem}

\begin{proof}\leanok
    \uses{lem:qca_blocked, thm:qca_finite_propagation_exists_blocked_nearest_neighbor}
    Apply the preceding theorem to the finite propagation of $\omega$, using
    the canonical positive block length.  At the same length, site blocking
    preserves both finite propagation and translation covariance, hence the
    QCA property.
\end{proof}

Lines~2292--2298 of \cite{Cirac2017MPU} specify the finite local
algebras, their canonical inclusions, translation covariance, and the locality
inclusion
\begin{align}
    \omega(\mathcal A_\Lambda)
     & \subseteq \mathcal A_{\Lambda+\mathcal N}.
\end{align}
The following statements record the finite-dimensional maps determined by
this inclusion.  They are consequences and coordinate formulations of the
cited locality relation, not separately named theorems of the paper.  The
left/right order below agrees with the adjacent-pair convention in
\cite{Gross2012IndexTheory}.

\begin{definition}[Finite-region restriction]
    \label{def:qca_finite_region_restriction}
    \lean{SpinChain.PropagatesWithin.localRestriction}
    \leanok
    \uses{def:qca_local_algebra, def:qca_local_inclusion, def:qca_region_sumset, def:qca_quasi_local_observable, def:qca_propagates_within}
    Let $d>0$, let $\omega$ be a star-algebra automorphism of
    $\mathcal A^{(d)}$, and suppose that $\omega$ propagates within the finite
    neighborhood $\mathcal N$.  For every finite region $\Lambda$, define
    the finite-region restriction
    \begin{align}
        \omega_{\Lambda,\mathcal N}\colon\mathcal A_\Lambda
         & \longrightarrow\mathcal A_{\Lambda+\mathcal N}
    \end{align}
    to be the unique star-algebra homomorphism satisfying
    \begin{align}
        \iota_{\Lambda+\mathcal N}
        \bigl(\omega_{\Lambda,\mathcal N}(A)\bigr)
         & =\omega\bigl(\iota_\Lambda(A)\bigr)
        \qquad(A\in\mathcal A_\Lambda).
    \end{align}
\end{definition}

\begin{theorem}[Finite-region restriction is an embedding]
    \label{thm:qca_finite_region_restriction_embedding}
    \lean{SpinChain.PropagatesWithin.quasiLocalObservable_localRestriction,
        SpinChain.PropagatesWithin.localRestriction_injective}
    \leanok
    \uses{def:qca_finite_region_restriction}
    The finite-region restriction is injective, and for every
    $A\in\mathcal A_\Lambda$ one has
    \begin{align}
        \iota_{\Lambda+\mathcal N}
        \bigl(\omega_{\Lambda,\mathcal N}(A)\bigr)
         & =\omega\bigl(\iota_\Lambda(A)\bigr).
    \end{align}
\end{theorem}

\begin{proof}\leanok
    \uses{def:qca_finite_region_restriction}
    The defining equality follows from the locality inclusion and the
    injectivity of $\iota_{\Lambda+\mathcal N}$.  If two local observables
    have the same image under $\omega_{\Lambda,\mathcal N}$, their canonical
    images become equal after applying $\omega$.  Injectivity of $\omega$ and
    of $\iota_\Lambda$ then gives equality of the two local observables.
\end{proof}

\begin{theorem}[Enlargement of finite-region restrictions]
    \label{thm:qca_finite_region_restriction_enlargement}
    \lean{SpinChain.PropagatesWithin.localRestriction_localInclusion}
    \leanok
    \uses{def:qca_finite_region_restriction}
    If $\Lambda\subseteq\Gamma$, then for every
    $A\in\mathcal A_\Lambda$,
    \begin{align}
        \omega_{\Gamma,\mathcal N}
        \bigl(\iota_{\Lambda,\Gamma}(A)\bigr)
         & =\iota_{\Lambda+\mathcal N,\Gamma+\mathcal N}
        \bigl(\omega_{\Lambda,\mathcal N}(A)\bigr).
    \end{align}
\end{theorem}

\begin{proof}\leanok
    \uses{thm:qca_finite_region_restriction_embedding}
    Apply the canonical inclusion into the quasi-local algebra to both sides.
    Both expressions become $\omega(\iota_\Lambda(A))$; injectivity of the
    target inclusion gives the equality.
\end{proof}

\begin{theorem}[Translation of finite-region restrictions]
    \label{thm:qca_finite_region_restriction_translation}
    \lean{SpinChain.PropagatesWithin.localRestriction_translation}
    \leanok
    \uses{def:qca_translation_covariant, def:qca_finite_local_translation, def:qca_finite_region_restriction}
    Suppose in addition that $\omega$ is translation covariant.  For
    $a\in\mathbb Z$ and $A\in\mathcal A_\Lambda$,
    \begin{align}
        \iota_{(\Lambda+a)+\mathcal N}
        \bigl(\omega_{\Lambda+a,\mathcal N}
            (\tau_{a,\Lambda}(A))\bigr)
         & =\iota_{(\Lambda+\mathcal N)+a}
        \bigl(\tau_{a,\Lambda+\mathcal N}
            (\omega_{\Lambda,\mathcal N}(A))\bigr).
    \end{align}
    Here $(\Lambda+a)+\mathcal N=(\Lambda+\mathcal N)+a$.
\end{theorem}

\begin{proof}\leanok
    \uses{thm:qca_finite_region_restriction_embedding, lem:qca_quasi_local_translation_action}
    The left-hand side is
    $\omega(\tau_a(\iota_\Lambda(A)))$.  Translation covariance moves
    $\tau_a$ past $\omega$, and the finite-region translation formula gives
    the right-hand side.
\end{proof}

\begin{definition}[Bipartite coordinates for disjoint local algebras]
    \label{def:qca_bipartite_local_algebra_coordinates}
    \lean{SpinChain.Config.disjointUnionEquiv,
        SpinChain.bipartiteLocalAlgebraEquiv}
    \leanok
    \uses{def:qca_local_algebra}
    Let $\Gamma,\Delta\subset\mathbb Z$ be disjoint finite regions.  Writing
    $\mathcal C_\Gamma$ and $\mathcal C_\Delta$ for their configuration
    spaces, restriction to the two regions gives
    \begin{align}
        \mathcal C_{\Gamma\cup\Delta}
         & \simeq \mathcal C_\Gamma\times\mathcal C_\Delta,
    \end{align}
    with the $\Gamma$-configuration first.  Reindexing matrix rows and
    columns gives a star-algebra equivalence
    \begin{align}
        \beta_{\Gamma,\Delta}\colon
        \mathcal A_{\Gamma\cup\Delta}
         & \simeq
        M_{\mathcal C_\Gamma\times\mathcal C_\Delta}(\mathbb C).
    \end{align}
\end{definition}

\begin{theorem}[Canonical inclusions in bipartite coordinates]
    \label{thm:qca_bipartite_local_algebra_inclusions}
    \lean{SpinChain.bipartiteLocalAlgebraEquiv_localInclusion_left,
        SpinChain.bipartiteLocalAlgebraEquiv_localInclusion_right}
    \leanok
    \uses{def:qca_bipartite_local_algebra_coordinates, def:qca_local_inclusion}
    For $A\in\mathcal A_\Gamma$ and $B\in\mathcal A_\Delta$,
    \begin{align}
        \beta_{\Gamma,\Delta}(\iota_{\Gamma,\Gamma\cup\Delta}(A))
         & =A\otimes\mathbf 1,  \\
        \beta_{\Gamma,\Delta}(\iota_{\Delta,\Gamma\cup\Delta}(B))
         & =\mathbf 1\otimes B.
    \end{align}
\end{theorem}

\begin{proof}\leanok
    \uses{def:qca_bipartite_local_algebra_coordinates}
    Evaluate both matrices on pairs of configurations.  For the first local
    inclusion, restriction to $\Gamma$ gives the matrix coefficient of $A$,
    while equality of the complementary $\Delta$-configurations gives the
    matrix coefficient of the identity.  This is $A\otimes\mathbf 1$.
    Interchanging the two regions gives $\mathbf 1\otimes B$.
\end{proof}

\begin{definition}[Finite local image in bipartite coordinates]
    \label{def:qca_bipartite_finite_local_image}
    \lean{SpinChain.PropagatesWithin.localRestrictionRange,
        SpinChain.PropagatesWithin.localRestrictionRangeEquiv,
        SpinChain.StarSubalgebra.bipartiteReindex,
        SpinChain.PropagatesWithin.bipartiteLocalRestrictionRange}
    \leanok
    \uses{def:qca_finite_region_restriction, thm:qca_finite_region_restriction_embedding, def:qca_bipartite_local_algebra_coordinates}
    Define the finite local image by
    \begin{align}
        \mathcal R_{\Lambda,\mathcal N}
         & =\omega_{\Lambda,\mathcal N}(\mathcal A_\Lambda)
        \subseteq\mathcal A_{\Lambda+\mathcal N}.
    \end{align}
    It is a star-subalgebra star-isomorphic to $\mathcal A_\Lambda$.  If
    $\Lambda+\mathcal N=\Gamma\cup\Delta$ with
    $\Gamma\cap\Delta=\varnothing$, then
    \begin{align}
        \beta_{\Gamma,\Delta}
        (\mathcal R_{\Lambda,\mathcal N})
         & \subseteq
        M_{\mathcal C_\Gamma\times\mathcal C_\Delta}(\mathbb C)
    \end{align}
    is the bipartite matrix star-subalgebra to which the left and right
    support-algebra constructions apply.  No support algebra is defined in
    this statement.
\end{definition}

\begin{definition}[Nearest-neighbour pair regions]
    \label{def:qca_nearest_neighbor_pair_regions}
    \lean{SpinChain.evenPair,SpinChain.leftPair,SpinChain.rightPair}
    \leanok
    For $x\in\mathbb Z$, set
    \begin{align}
        E_x & =\{2x,2x+1\},   &
        L_x & =\{2x-1,2x\},   &
        P_x & =\{2x+1,2x+2\}.
        \label{eq:qca_pair_regions}
    \end{align}
    The order $L_x,P_x$ is the left--right factor order in equation
    \texttt{RR2x} of
    \cite[lines~1251--1266]{Gross2012IndexTheory}.  The same four-site
    localization occurs in
    \cite[lines~2313--2318]{Cirac2017MPU}.
\end{definition}

\begin{lemma}[Geometry of the nearest-neighbour pair image]
    \label{lem:qca_nearest_neighbor_pair_geometry}
    \lean{SpinChain.regionSumset_evenPair,
        SpinChain.disjoint_leftPair_rightPair,
        SpinChain.rightPair_eq_leftPair_add_one}
    \leanok
    \uses{def:qca_nearest_neighbor_pair_regions}
    For the neighbourhood $\mathcal N=\{-1,0,1\}$,
    \begin{align}
        E_x+\mathcal N & =L_x\cup P_x, &
        L_x\cap P_x    & =\varnothing, &
        P_x            & =L_{x+1}.
        \label{eq:qca_pair_geometry}
    \end{align}
\end{lemma}

\begin{proof}\leanok
    Substitute (\ref{eq:qca_pair_regions}).  Adding $-1,0,1$ to the two
    points of $E_x$ gives
    \begin{align}
        E_x+\{-1,0,1\}
         & =\{2x-1,2x,2x+1,2x+2\}=L_x\cup P_x.
    \end{align}
    The largest point of $L_x$ is $2x$, while the smallest point of $P_x$
    is $2x+1$, so the two sets are disjoint.  Finally,
    $L_{x+1}=\{2x+1,2x+2\}=P_x$.
\end{proof}

\begin{definition}[Site-indexed support algebras]
    \label{def:qca_site_indexed_support_algebras}
    \lean{SpinChain.matrixToQuasiLocalObservable,
        SpinChain.PropagatesWithin.evenPairLocalImage,
        SpinChain.PropagatesWithin.evenSupportAlgebra,
        SpinChain.PropagatesWithin.oddSupportAlgebra,
        SpinChain.PropagatesWithin.embeddedEvenSupportAlgebra,
        SpinChain.PropagatesWithin.embeddedOddSupportAlgebra,
        SpinChain.PropagatesWithin.embeddedSupportAlgebra,
        SpinChain.PropagatesWithin.embeddedSupportAlgebra_even,
        SpinChain.PropagatesWithin.embeddedSupportAlgebra_odd}
    \leanok
    \uses{def:qca_bipartite_finite_local_image, def:qca_quasi_local_observable, def:qca_left_support_algebra, def:qca_right_support_algebra, def:qca_nearest_neighbor_pair_regions, lem:qca_nearest_neighbor_pair_geometry}
    Let $\omega$ be a star-algebra automorphism of the homogeneous
    quasi-local algebra which propagates within $\{-1,0,1\}$.  Write the
    finite image of $\mathcal A_{E_x}$ in the ordered $L_x\times P_x$
    coordinates as
    \begin{align}
        S_x
         & =\beta_{L_x,P_x}\bigl(\omega(\mathcal A_{E_x})\bigr)
        \subseteq
        M_{\mathcal C_{L_x}\times\mathcal C_{P_x}}(\mathbb C).
        \label{eq:qca_site_image}
    \end{align}
    Define
    \begin{align}
        \mathcal R_{2x}
         & =\operatorname{Spp}_{\mathrm L}(S_x)
        \subseteq M_{\mathcal C_{L_x}}(\mathbb C), \\
        \mathcal R_{2x+1}
         & =\operatorname{Spp}_{\mathrm R}(S_x)
        \subseteq M_{\mathcal C_{P_x}}(\mathbb C).
        \label{eq:qca_RR2x}
    \end{align}
    After the canonical finite-region embeddings, these form one family of
    star-subalgebras of the quasi-local algebra.  Here the matrix embedding is
    the composition
    \begin{align}
        M_{\mathcal C_\Lambda}(\mathbb C)
        \cong \mathcal A_\Lambda
        \longrightarrow \mathcal A.
    \end{align}
    Under this identification, the canonical local inclusion in
    \cite[lines~2292--2298]{Cirac2017MPU} is exactly the displayed
    composition.
    Formula (\ref{eq:qca_RR2x}) is equation \texttt{RR2x} in
    \cite[lines~1261--1266]{Gross2012IndexTheory}.  Cirac et al. use the
    homogeneous even factor in
    \cite[lines~2322--2328]{Cirac2017MPU}; the odd factor and unified family
    are supplied by the GNVW construction.

    \textbf{Scope restriction (homogeneous chain):}
    GNVW allow the cell size $d(y)$ to depend on $y$, whereas the present
    definition fixes one positive size $d$.  The site-dependent statement is
    not claimed here; see
    \cite{gap:gnvw12_site_dependent_adjacent_generation_scope}.
\end{definition}

\begin{theorem}[Adjacent-pair support containment]
    \label{thm:qca_adjacent_pair_support_containment}
    \lean{SpinChain.PropagatesWithin.evenPairLocalImage_le_support_kroneckerSubmodule}
    \leanok
    \uses{def:qca_site_indexed_support_algebras, thm:qca_two_sided_support_containment}
    For every $x\in\mathbb Z$, the finite image satisfies
    \begin{align}
        S_x
         & \subseteq
        \mathcal R_{2x}\boxtimes\mathcal R_{2x+1}
        \subseteq
        M_{\mathcal C_{L_x}\times\mathcal C_{P_x}}(\mathbb C).
        \label{eq:qca_adjacent_pair_containment}
    \end{align}
    This is the homogeneous, full-matrix instance used in the two-sided
    containment at
    \cite[lines~1276--1278]{Gross2012IndexTheory}; see
    \cite{gap:gnvw12_site_dependent_adjacent_generation_scope}.
    Cirac et al. state only the weaker inclusion in the left support algebra
    tensored with the unrestricted right physical factor at
    \cite[lines~2322--2328]{Cirac2017MPU}.
\end{theorem}

\begin{proof}\leanok
    \uses{thm:qca_two_sided_support_containment}
    Apply Theorem~\ref{thm:qca_two_sided_support_containment} to $S_x$ in
    the ordered coordinates of (\ref{eq:qca_site_image}).  Its left support
    is $\mathcal R_{2x}$ and its right support is
    $\mathcal R_{2x+1}$ by (\ref{eq:qca_RR2x}), which gives
    (\ref{eq:qca_adjacent_pair_containment}).
\end{proof}

\begin{theorem}[Physical support of the site-indexed support algebras]
    \label{thm:qca_site_indexed_support_regions}
    \lean{SpinChain.matrixToQuasiLocalObservable_supportedIn,
        SpinChain.PropagatesWithin.embeddedSupportAlgebra_even_supportedIn,
        SpinChain.PropagatesWithin.embeddedSupportAlgebra_odd_supportedIn}
    \leanok
    \uses{def:qca_site_indexed_support_algebras, def:qca_quasi_local_supported_in}
    If $\widehat{\mathcal R}_y$ denotes the canonical quasi-local copy of
    $\mathcal R_y$, then for every $x\in\mathbb Z$,
    \begin{align}
        \widehat{\mathcal R}_{2x}   & \subseteq\mathcal A_{L_x}, &
        \widehat{\mathcal R}_{2x+1} & \subseteq\mathcal A_{P_x}.
    \end{align}
    More generally, the canonical image of any matrix on a finite region
    $\Lambda$ is supported in $\Lambda$.
    This is the homogeneous specialization of the localization accompanying
    equation \texttt{RR2x} in
    \cite[lines~1261--1274]{Gross2012IndexTheory}; see
    \cite{gap:gnvw12_site_dependent_adjacent_generation_scope}.
\end{theorem}

\begin{proof}\leanok
    An element of $\widehat{\mathcal R}_{2x}$ is the canonical image of a
    matrix in $\mathcal R_{2x}\subseteq\mathcal A_{L_x}$ and is therefore
    supported in $L_x$.  The same argument places
    $\widehat{\mathcal R}_{2x+1}$ in $\mathcal A_{P_x}$.
\end{proof}

\begin{theorem}[Generation by two disjoint finite regions]
    \label{thm:qca_disjoint_region_generation}
    \lean{SpinChain.range_localObservable_mono,
        SpinChain.range_localObservable_union}
    \leanok
    \uses{def:qca_algebraic_local_algebra, def:qca_local_inclusion,
        def:qca_bipartite_local_algebra_coordinates,
        thm:qca_bipartite_local_algebra_inclusions,
        thm:qca_commuting_factor_kronecker_coordinates}
    If $\Lambda\subseteq\Gamma$ are finite regions, then
    $\iota_\Lambda(\mathcal A_\Lambda)
        \subseteq\iota_\Gamma(\mathcal A_\Gamma)$.  If the finite regions
    $\Lambda$ and $\Gamma$ are disjoint, then the two canonical images
    generate the image of their union,
    \begin{align}
        \iota_{\Lambda\cup\Gamma}(\mathcal A_{\Lambda\cup\Gamma})
         & =\iota_\Lambda(\mathcal A_\Lambda)
        \vee\iota_\Gamma(\mathcal A_\Gamma).
        \label{eq:qca_disjoint_region_generation}
    \end{align}
    This is the expansion of an observable on several cells into a product of
    observables on disjoint groups of cells used in
    \cite[lines~592--596]{Gross2012IndexTheory} and proved in
    \cite[lines~341--364]{SchumacherWerner2004QCA}.
\end{theorem}

\begin{proof}\leanok
    \uses{thm:qca_bipartite_local_algebra_inclusions,
        thm:qca_commuting_factor_kronecker_coordinates,
        thm:qca_mem_kronecker_submodule_iff}
    Monotonicity follows from
    $\iota_\Gamma\circ\iota_{\Lambda,\Gamma}=\iota_\Lambda$, and it gives the
    inclusion $\supseteq$ in
    (\ref{eq:qca_disjoint_region_generation}).  For the remaining inclusion,
    let $\mathcal V\subseteq\mathcal A_{\Lambda\cup\Gamma}$ be the
    star-subalgebra of those $A$ whose canonical image lies in the join on the
    right-hand side.  In the bipartite coordinates
    $\beta_{\Lambda,\Gamma}$, the canonical inclusion of
    $\mathcal A_\Lambda$ is $X\mapsto X\otimes\mathbf 1$ and that of
    $\mathcal A_\Gamma$ is $Y\mapsto\mathbf 1\otimes Y$, so
    $\beta_{\Lambda,\Gamma}(\mathcal V)$ contains
    $M_{\mathcal C_\Lambda}(\mathbb C)\otimes\mathbf 1$ and
    $\mathbf 1\otimes M_{\mathcal C_\Gamma}(\mathbb C)$.  The join of these two
    subalgebras has underlying vector space
    $M_{\mathcal C_\Lambda}(\mathbb C)\boxtimes
        M_{\mathcal C_\Gamma}(\mathbb C)$, and every matrix on
    $\mathcal C_\Lambda\times\mathcal C_\Gamma$ lies in that span.  Hence
    $\beta_{\Lambda,\Gamma}(\mathcal V)$ is everything, and
    $\beta_{\Lambda,\Gamma}$ is bijective, so
    $\mathcal V=\mathcal A_{\Lambda\cup\Gamma}$.
\end{proof}

\begin{lemma}[Even pairs cover the chain]
    \label{lem:qca_even_pair_cover}
    \lean{SpinChain.disjoint_evenPair_of_ne,
        SpinChain.subset_biUnion_evenPair}
    \leanok
    \uses{def:qca_nearest_neighbor_pair_regions}
    For $x\ne y$ the pairs $E_x$ and $E_y$ are disjoint, and every finite
    region $\Lambda$ satisfies
    \begin{align}
        \Lambda & \subseteq\bigcup_{x\in X_\Lambda}E_x,
        \label{eq:qca_even_pair_cover}
    \end{align}
    where $X_\Lambda=\{\lfloor i/2\rfloor\mid i\in\Lambda\}$.
\end{lemma}

\begin{proof}\leanok
    \uses{def:qca_nearest_neighbor_pair_regions}
    Both statements are integer arithmetic: the four sites
    $2x,2x+1,2y,2y+1$ are distinct for $x\ne y$, and every $i$ satisfies
    $i=2\lfloor i/2\rfloor$ or $i=2\lfloor i/2\rfloor+1$.
\end{proof}

\begin{theorem}[Generation of the algebraic local algebra by the pair algebras]
    \label{thm:qca_even_pair_algebraic_generation}
    \lean{SpinChain.range_localObservable_biUnion_evenPair_le,
        SpinChain.iSup_range_localObservable_evenPair}
    \leanok
    \uses{def:qca_nearest_neighbor_pair_regions,
        def:qca_algebraic_local_algebra,
        thm:qca_disjoint_region_generation, lem:qca_even_pair_cover,
        thm:qca_local_observable_finite_support}
    The pair algebras generate the algebraic local algebra,
    \begin{align}
        \bigvee_{x\in\mathbb Z}\iota_{E_x}(\mathcal A_{E_x})
         & =\mathcal A_{\mathrm{loc}}.
        \label{eq:qca_even_pair_algebraic_generation}
    \end{align}
    The completion is not involved: this is an equality of algebras of local
    observables.
\end{theorem}

\begin{proof}\leanok
    \uses{thm:qca_disjoint_region_generation, lem:qca_even_pair_cover,
        thm:qca_local_observable_finite_support}
    For a finite set $X$ of pair labels, induction on $X$ gives
    \begin{align}
        \iota_{\bigcup_{x\in X}E_x}
        \Bigl(\mathcal A_{\bigcup_{x\in X}E_x}\Bigr)
         & \subseteq\bigvee_{x\in\mathbb Z}
        \iota_{E_x}(\mathcal A_{E_x}).
        \label{eq:qca_even_pair_generation_induction}
    \end{align}
    For $X=\varnothing$ the region is empty and contained in $E_0$, so
    monotonicity applies.  For $X\cup\{x_0\}$ with $x_0\notin X$, the pair
    $E_{x_0}$ is disjoint from $\bigcup_{x\in X}E_x$, so
    (\ref{eq:qca_disjoint_region_generation}) splits the left-hand side into
    the join of $\iota_{E_{x_0}}(\mathcal A_{E_{x_0}})$ with the algebra
    handled by the induction hypothesis.  Every local observable is supported
    in some finite region $\Lambda$, which is covered by
    (\ref{eq:qca_even_pair_cover}); monotonicity and
    (\ref{eq:qca_even_pair_generation_induction}) place it in the join.
\end{proof}

\begin{definition}[Join of the site-indexed support algebras]
    \label{def:qca_support_algebra_join}
    \lean{SpinChain.PropagatesWithin.supportAlgebraJoin,
        SpinChain.PropagatesWithin.embeddedEvenSupportAlgebra_le_supportAlgebraJoin,
        SpinChain.PropagatesWithin.embeddedOddSupportAlgebra_le_supportAlgebraJoin}
    \leanok
    \uses{def:qca_site_indexed_support_algebras}
    For a star-algebra automorphism $\omega$ of the quasi-local algebra
    propagating within $\{-1,0,1\}$, define
    \begin{align}
        G & =\bigvee_{y\in\mathbb Z}\widehat{\mathcal R}_y
        \subseteq\mathcal A.
        \label{eq:qca_support_algebra_join}
    \end{align}
    Each $\widehat{\mathcal R}_y$ is contained in $G$.  This is the algebraic
    join associated with the algebra generated by all $\mathcal R_y$ in
    \cite[lines~1276--1279]{Gross2012IndexTheory}.  No claim that $G$ is
    proper or that it fails to be norm closed is made here.
\end{definition}

\begin{theorem}[Images of the pair algebras lie in the join]
    \label{thm:qca_evenPair_image_in_support_algebra_join}
    \lean{SpinChain.PropagatesWithin.quasiLocalObservable_evenPair_mem_supportAlgebraJoin,
        SpinChain.PropagatesWithin.algebraicToQuasiLocal_mem_supportAlgebraJoin}
    \leanok
    \uses{def:qca_support_algebra_join,
        def:qca_bipartite_finite_local_image,
        thm:qca_adjacent_pair_support_containment,
        thm:qca_even_pair_algebraic_generation}
    For every $x\in\mathbb Z$,
    \begin{align}
        \omega\bigl(j_{E_x}(\mathcal A_{E_x})\bigr) & \subseteq G,
        \label{eq:qca_evenPair_image_in_support_algebra_join}
    \end{align}
    and consequently
    $\omega\bigl(j(\mathcal A_{\mathrm{loc}})\bigr)\subseteq G$.
    Inclusion (\ref{eq:qca_evenPair_image_in_support_algebra_join}) is the
    containment
    $\alpha(\mathcal A_{2x}\otimes\mathcal A_{2x+1})\subseteq
        \mathcal R_{2x}\otimes\mathcal R_{2x+1}$
    at \cite[lines~1276--1278]{Gross2012IndexTheory}, in the homogeneous
    fixed-$d$ setting; see
    \cite{gap:gnvw12_site_dependent_adjacent_generation_scope}.
\end{theorem}

\begin{proof}\leanok
    \uses{thm:qca_adjacent_pair_support_containment,
        thm:qca_commuting_factor_kronecker_coordinates,
        thm:qca_bipartite_local_algebra_inclusions,
        thm:qca_even_pair_algebraic_generation}
    Write $S_x$ for the finite image of $\mathcal A_{E_x}$ in the ordered
    $L_x\times P_x$ coordinates.  By
    (\ref{eq:qca_adjacent_pair_containment}), $S_x$ is contained in
    $\mathcal R_{2x}\boxtimes\mathcal R_{2x+1}$, which is the underlying vector
    space of the join of $\mathcal R_{2x}\otimes\mathbf 1$ and
    $\mathbf 1\otimes\mathcal R_{2x+1}$.  The canonical quasi-local map on
    $L_x\cup P_x$ sends $X\otimes\mathbf 1$ to the canonical embedding of
    $X\in\mathcal R_{2x}$ and $\mathbf 1\otimes Y$ to the canonical embedding
    of $Y\in\mathcal R_{2x+1}$, so it maps that join into $G$.  It sends the
    image of $A\in\mathcal A_{E_x}$ to $\omega(j_{E_x}(A))$, which proves
    (\ref{eq:qca_evenPair_image_in_support_algebra_join}).  By
    (\ref{eq:qca_even_pair_algebraic_generation}), every local observable lies
    in the algebra generated by the $\iota_{E_x}(\mathcal A_{E_x})$, and
    $\omega\circ j$ is a star-algebra homomorphism, so its image lies in $G$.
\end{proof}

\begin{theorem}[Norm-closed generation by the support algebras]
    \label{thm:qca_support_algebra_topological_generation}
    \lean{SpinChain.PropagatesWithin.supportAlgebraJoin_topologicalClosure_eq_top}
    \leanok
    \uses{def:qca_support_algebra_join,
        thm:qca_evenPair_image_in_support_algebra_join,
        thm:qca_quasi_local_density}
    The site-indexed support algebras generate the quasi-local algebra in the
    norm-closed sense,
    \begin{align}
        \overline{G}^{\lVert\cdot\rVert} & =\mathcal A.
        \label{eq:qca_support_algebra_topological_generation}
    \end{align}
    In \cite[lines~1276--1279]{Gross2012IndexTheory} the algebras
    $\mathcal R_y$ are said to generate an algebra containing
    $\alpha(\mathcal A(\mathbb Z))$, which equals $\mathcal A(\mathbb Z)$
    because $\alpha$ is an automorphism.  Here this generation statement is
    made precise by the norm-closure equality
    (\ref{eq:qca_support_algebra_topological_generation}), without asserting
    that the algebraic join itself is proper or non-closed.  This is reusable
    homogeneous fixed-$d$ infrastructure and does not assert the
    site-dependent statement of the source; see
    \cite{gap:gnvw12_site_dependent_adjacent_generation_scope}.
\end{theorem}

\begin{proof}\leanok
    \uses{thm:qca_evenPair_image_in_support_algebra_join,
        thm:qca_quasi_local_density}
    An injective star-algebra homomorphism of complex C-star algebras is
    isometric, so $\omega$ is continuous and
    $\omega^{-1}(\overline G)$ is closed.  By
    (\ref{eq:qca_evenPair_image_in_support_algebra_join}) it contains
    $j(\mathcal A_{\mathrm{loc}})$, which is dense, so
    $\omega^{-1}(\overline G)=\mathcal A$.  Since $\omega$ is surjective,
    every $z\in\mathcal A$ equals $\omega(\omega^{-1}(z))\in\overline G$.
\end{proof}

\begin{corollary}[Commutant of the site-indexed support algebras]
    \label{cor:qca_support_algebra_commutant}
    \lean{SpinChain.PropagatesWithin.commute_of_forall_commute_embeddedSupportAlgebra}
    \leanok
    \uses{def:qca_support_algebra_join,
        thm:qca_support_algebra_topological_generation}
    If $z\in\mathcal A$ satisfies $[z,R]=0$ for every $y\in\mathbb Z$ and
    every $R\in\widehat{\mathcal R}_y$, then $[z,A]=0$ for every
    $A\in\mathcal A$.  This is the step from commutation with all
    $\mathcal R_y$ to commutation in the quasi-local algebra at
    \cite[lines~1279--1282]{Gross2012IndexTheory}; see
    \cite{gap:gnvw12_site_dependent_adjacent_generation_scope}.
\end{corollary}

\begin{proof}\leanok
    \uses{def:qca_support_algebra_join,
        thm:qca_support_algebra_topological_generation}
    The elements commuting with $z$ and with $z^*$ form a norm-closed
    star-subalgebra.  Each $\widehat{\mathcal R}_y$ is star closed, so the
    hypothesis applied to $R^*$ gives $[z^*,R]=0$ as well; hence that
    star-subalgebra contains every $\widehat{\mathcal R}_y$, therefore
    contains $G$, and being closed it contains
    $\overline G=\mathcal A$ by
    (\ref{eq:qca_support_algebra_topological_generation}).
\end{proof}

\begin{definition}[Tripartite finite-region coordinates]
    \label{def:qca_tripartite_local_coordinates}
    \lean{SpinChain.Config.disjointTripleEquiv,
        SpinChain.tripartiteLocalAlgebraEquiv}
    \leanok
    \uses{def:qca_bipartite_local_algebra_coordinates}
    For pairwise disjoint finite regions $\Lambda,\Gamma,\Delta$, successive
    restriction in left--middle--right order gives
    \begin{align}
        \mathcal C_{(\Lambda\cup\Gamma)\cup\Delta}
         & \cong(\mathcal C_\Lambda\times\mathcal C_\Gamma)
        \times\mathcal C_\Delta,                                \\
        \mathcal A_{(\Lambda\cup\Gamma)\cup\Delta}
         & \cong M_{(\mathcal C_\Lambda\times\mathcal C_\Gamma)
                \times\mathcal C_\Delta}(\mathbb C).
    \end{align}
\end{definition}

\begin{lemma}[Geometry of consecutive pair images]
    \label{lem:qca_consecutive_pair_geometry}
    \lean{SpinChain.disjoint_leftPair_rightPair,
        SpinChain.disjoint_leftPair_union_rightPair_rightPair_add_one,
        SpinChain.leftPair_union_rightPair_union_rightPair_add_one}
    \leanok
    \uses{def:qca_nearest_neighbor_pair_regions}
    Let $T_x=(L_x\cup P_x)\cup P_{x+1}$. The three regions
    $L_x,P_x,P_{x+1}$ are pairwise disjoint, and
    \begin{align}
        T_x & =\{2x-1,2x,2x+1,2x+2,2x+3,2x+4\}.
    \end{align}
\end{lemma}

\begin{proof}\leanok
    \uses{def:qca_nearest_neighbor_pair_regions}
    Substitute the three pair definitions and compare their six sites.
\end{proof}

\begin{theorem}[Canonical inclusions in tripartite coordinates]
    \label{thm:qca_consecutive_pair_tripartite_coordinates}
    \lean{SpinChain.tripartiteLocalAlgebraEquiv_localInclusion_left,
        SpinChain.tripartiteLocalAlgebraEquiv_localInclusion_right}
    \leanok
    \uses{def:qca_tripartite_local_coordinates,
        def:qca_local_inclusion}
    In the coordinates of Definition~\ref{def:qca_tripartite_local_coordinates},
    the canonical inclusion from $\Lambda\cup\Gamma$ is
    $X\mapsto X\otimes\mathbf 1$, and the canonical inclusion from
    $\Gamma\cup\Delta$ is $Y\mapsto\mathbf 1\otimes Y$, with the ordered
    factors $\Lambda,\Gamma,\Delta$.
\end{theorem}

\begin{proof}\leanok
    \uses{def:qca_tripartite_local_coordinates,
        def:qca_bipartite_local_algebra_coordinates,
        def:qca_local_inclusion}
    Restrict each reconstructed configuration to the retained region and its
    complement.  In the first inclusion the complementary coordinates are
    equal precisely on $\Delta$; in the second they are equal precisely on
    $\Lambda$.  The two matrix entries are therefore the left and right
    overlapping lifts, respectively.
\end{proof}

\begin{theorem}[Commutation across the common physical pair]
    \label{thm:qca_consecutive_support_algebras_commute}
    \lean{
        SpinChain.PropagatesWithin.embeddedOddSupport_commute_embeddedEvenSupport_add_one}
    \leanok
    \uses{def:qca_site_indexed_support_algebras}
    For every $x\in\mathbb Z$, the canonically embedded support algebras
    $\widehat{\mathcal R}_{2x+1}$ and
    $\widehat{\mathcal R}_{2x+2}$ commute elementwise.
\end{theorem}

\begin{proof}\leanok
    \uses{lem:qca_consecutive_pair_geometry,
        thm:qca_consecutive_pair_tripartite_coordinates,
        thm:qca_finite_region_restriction_embedding,
        thm:qca_disjoint_quasi_local_commute,
        thm:qca_overlapping_support_algebras_commute}
    The source regions $E_x$ and $E_{x+1}$ are disjoint, so their canonical
    quasi-local observables commute.  Multiplicativity of $\omega$ preserves
    this equality.  Embed the two finite images into $\mathcal A_{T_x}$ and
    use the tripartite coordinates above to obtain
    \begin{align}
        (X\otimes\mathbf 1)(\mathbf 1\otimes Y)
         & =(\mathbf 1\otimes Y)(X\otimes\mathbf 1).
    \end{align}
    The overlapping-support theorem, in the order
    $\operatorname{Spp}_{\mathrm R}(S_x)$ followed by
    $\operatorname{Spp}_{\mathrm L}(S_{x+1})$, gives the claim.
\end{proof}

\begin{definition}[Physical region of a site-indexed support algebra]
    \label{def:qca_support_algebra_region}
    \lean{SpinChain.supportRegion,
        SpinChain.supportRegion_even,
        SpinChain.supportRegion_odd}
    \leanok
    \uses{def:qca_nearest_neighbor_pair_regions}
    Define $Q_{2x}=L_x$ and $Q_{2x+1}=P_x$.
\end{definition}

\begin{theorem}[Unified physical localization]
    \label{thm:qca_unified_support_algebra_region}
    \lean{SpinChain.PropagatesWithin.embeddedSupportAlgebra_supportedIn}
    \leanok
    \uses{def:qca_support_algebra_region,
        def:qca_site_indexed_support_algebras}
    Every element of $\widehat{\mathcal R}_y$ is supported in $Q_y$.
\end{theorem}

\begin{proof}\leanok
    \uses{def:qca_support_algebra_region,
        thm:qca_site_indexed_support_regions}
    Split according to the parity of $y$ and apply the corresponding
    localization of $\widehat{\mathcal R}_{2x}$ or
    $\widehat{\mathcal R}_{2x+1}$.
\end{proof}

\begin{theorem}[Classification of overlapping support regions]
    \label{thm:qca_support_region_overlap_classification}
    \lean{SpinChain.PropagatesWithin.supportRegion_disjoint_or_consecutive}
    \leanok
    \uses{def:qca_support_algebra_region}
    If $y\ne z$, then $Q_y$ and $Q_z$ are disjoint unless
    \begin{align}
        (y,z) & =(2x+1,2x+2)
    \end{align}
    for some $x$, or unless this ordered pair is reversed.
\end{theorem}

\begin{proof}\leanok
    \uses{def:qca_support_algebra_region}
    Write each index in its even or odd form.  Two regions of equal parity
    are disjoint.  In the mixed-parity cases, the defining two-site sets
    intersect exactly for the displayed consecutive odd--even pair.
\end{proof}

\begin{theorem}[Pairwise commutation of distinct support algebras]
    \label{thm:qca_site_indexed_support_algebras_commute}
    \lean{SpinChain.PropagatesWithin.embeddedSupportAlgebra_pairwise_commute}
    \leanok
    \uses{def:qca_site_indexed_support_algebras}
    Distinct members of the canonically embedded family commute elementwise:
    \begin{align}
        [X,Y] & =0,
              & X\in\widehat{\mathcal R}_y,\quad
        Y\in\widehat{\mathcal R}_z,\quad y\ne z.
    \end{align}
    This is the homogeneous fixed-$d$ specialization of the
    consecutive/disjoint argument in
    \cite[lines~1270--1274]{Gross2012IndexTheory}, using Lemma
    \texttt{sppcomm} from lines~1221--1246; compare
    \cite[lines~1174--1194]{SchumacherWerner2004QCA}.  Cirac et al. provide the
    homogeneous four-site localization at
    \cite[lines~2313--2329]{Cirac2017MPU}, but do not state this commutation
    result.  The site-dependent GNVW statement is not claimed; see
    \cite{gap:gnvw12_site_dependent_adjacent_generation_scope}.
\end{theorem}

\begin{proof}\leanok
    \uses{thm:qca_support_region_overlap_classification,
        thm:qca_unified_support_algebra_region,
        thm:qca_consecutive_support_algebras_commute,
        thm:qca_disjoint_quasi_local_commute}
    In the exceptional odd--even order apply
    Theorem~\ref{thm:qca_consecutive_support_algebras_commute}; the reversed
    order follows by symmetry.  Every remaining ordered pair has disjoint
    physical supports, so disjoint quasi-locality applies.
\end{proof}

\begin{lemma}[Two-site translation of the pair regions]
    \label{lem:qca_pair_region_two_site_translation}
    \lean{SpinChain.translateRegion_evenPair,
        SpinChain.translateRegion_leftPair,
        SpinChain.translateRegion_rightPair}
    \leanok
    \uses{def:qca_nearest_neighbor_pair_regions}
    Translation by two sites advances each of the three two-site regions by
    one index:
    \begin{align}
        E_x+2 & =E_{x+1}, &
        L_x+2 & =L_{x+1}, &
        P_x+2 & =P_{x+1}.
        \label{eq:qca_pair_region_two_site_translation}
    \end{align}
\end{lemma}

\begin{proof}\leanok
    \uses{def:qca_nearest_neighbor_pair_regions}
    Add $2$ to each of the two points of (\ref{eq:qca_pair_regions}) and
    compare with the same set at the index $x+1$.
\end{proof}

\begin{lemma}[Matrix coefficients as compressions by matrix units]
    \label{lem:qca_support_coefficient_compression}
    \lean{Matrix.leftKroneckerEmbed_bipartiteSlice_eq_sum,
        Matrix.rightKroneckerEmbed_operatorBlock_eq_sum}
    \leanok
    \uses{def:qca_left_coefficient_set, def:qca_right_coefficient_set}
    Let $X\in M_{I\times J}(\mathbb C)$ and write $e_{rs}$ for the matrix
    units of either factor.  Then
    \begin{align}
        X^{\mathrm L}_{rs}\otimes\mathbf 1
         & =\sum_{t\in J}
        (\mathbf 1\otimes e_{tr})\,X\,(\mathbf 1\otimes e_{st}),
        \label{eq:qca_left_coefficient_compression} \\
        \mathbf 1\otimes X^{\mathrm R}_{ij}
         & =\sum_{k\in I}
        (e_{ki}\otimes\mathbf 1)\,X\,(e_{jk}\otimes\mathbf 1).
        \label{eq:qca_right_coefficient_compression}
    \end{align}
    Every coefficient of the expansion used at
    \cite[lines~1188--1191]{Gross2012IndexTheory} and
    \cite[lines~1137--1169]{SchumacherWerner2004QCA} is therefore a sum of
    products formed inside the ambient algebra.
\end{lemma}

\begin{proof}\leanok
    Compare entries.  On the left of
    (\ref{eq:qca_left_coefficient_compression}), the entry at
    $((i,t),(j,u))$ equals $X((i,r),(j,s))$ when $t=u$ and vanishes
    otherwise.  On the right, the summand indexed by $t$ contributes
    $X((i,r),(j,s))$ exactly when both second coordinates equal $t$, so the
    sum over $t$ reproduces the same entry.  The second identity follows by
    exchanging the roles of the two factors.
\end{proof}

\begin{lemma}[Quasi-local form of the coefficient compressions]
    \label{lem:qca_quasi_local_coefficient_compression}
    \lean{SpinChain.matrixToQuasiLocalObservable_bipartiteSlice_eq_sum,
        SpinChain.matrixToQuasiLocalObservable_operatorBlock_eq_sum}
    \leanok
    \uses{lem:qca_support_coefficient_compression,
        def:qca_bipartite_local_algebra_coordinates,
        def:qca_quasi_local_observable}
    Let $\Gamma$ and $\Delta$ be disjoint finite regions and let
    $A\in\mathcal A_{\Gamma\cup\Delta}$, written in the ordered
    $\Gamma\times\Delta$ coordinates.  For configurations
    $r,s\in\mathcal C_\Delta$ and $i,j\in\mathcal C_\Gamma$, the canonical
    quasi-local copies of the two coefficient families are
    \begin{align}
        \widehat{A^{\mathrm L}_{rs}}
         & =\sum_{t\in\mathcal C_\Delta}
        \widehat{e^\Delta_{tr}}\ \widehat A\ \widehat{e^\Delta_{st}},
        \label{eq:qca_quasi_local_left_compression} \\
        \widehat{A^{\mathrm R}_{ij}}
         & =\sum_{k\in\mathcal C_\Gamma}
        \widehat{e^\Gamma_{ki}}\ \widehat A\ \widehat{e^\Gamma_{jk}},
        \label{eq:qca_quasi_local_right_compression}
    \end{align}
    where a hat denotes the canonical image in the quasi-local algebra.  Both
    right-hand sides are formed inside the quasi-local algebra alone.
\end{lemma}

\begin{proof}\leanok
    \uses{lem:qca_support_coefficient_compression,
        thm:qca_bipartite_local_algebra_inclusions,
        lem:qca_quasi_local_observable}
    In the bipartite coordinates the inclusion of $\mathcal A_\Gamma$ is
    $X\mapsto X\otimes\mathbf 1$ and the inclusion of $\mathcal A_\Delta$ is
    $Y\mapsto\mathbf 1\otimes Y$.  Transport
    (\ref{eq:qca_left_coefficient_compression}) and
    (\ref{eq:qca_right_coefficient_compression}) through the inverse of the
    bipartite coordinate equivalence, then apply the canonical map of
    $\mathcal A_{\Gamma\cup\Delta}$ into the quasi-local algebra, under which
    a finite-region inclusion leaves the quasi-local observable unchanged.
\end{proof}

\begin{theorem}[Two-site translation covariance of the support algebras]
    \label{thm:qca_site_indexed_support_algebras_two_site_covariance}
    \lean{SpinChain.PropagatesWithin.embeddedEvenSupportAlgebra_map_quasiLocalTranslation_two,
        SpinChain.PropagatesWithin.embeddedOddSupportAlgebra_map_quasiLocalTranslation_two,
        SpinChain.PropagatesWithin.embeddedSupportAlgebra_map_quasiLocalTranslation_two,
        SpinChain.PropagatesWithin.embeddedSupportAlgebra_map_quasiLocalTranslation_two_even,
        SpinChain.PropagatesWithin.embeddedSupportAlgebra_map_quasiLocalTranslation_two_odd}
    \leanok
    \uses{def:qca_site_indexed_support_algebras, def:qca_translation_covariant,
        def:qca_quasi_local_translation}
    Let $\omega$ be a star-algebra automorphism of the homogeneous quasi-local
    algebra which propagates within $\{-1,0,1\}$ and commutes with every
    lattice translation.  Then, for every $y\in\mathbb Z$,
    \begin{align}
        \tau_2\bigl(\widehat{\mathcal R}_y\bigr)
         & =\widehat{\mathcal R}_{y+2},
        \label{eq:qca_two_site_covariance}
    \end{align}
    and in the two parities
    \begin{align}
        \tau_2\bigl(\widehat{\mathcal R}_{2x}\bigr)
         & =\widehat{\mathcal R}_{2x+2}, &
        \tau_2\bigl(\widehat{\mathcal R}_{2x+1}\bigr)
         & =\widehat{\mathcal R}_{2x+3}.
        \label{eq:qca_two_site_covariance_parities}
    \end{align}
    The displacement is two sites of the original chain, that is, one blocked
    cell.  The corresponding one-site formula
    $\tau_1(\widehat{\mathcal R}_y)=\widehat{\mathcal R}_{y+1}$ is false in
    general and is not asserted, because translation by one does not preserve
    the even pairing $E_x$.  For the shift automorphism at
    \cite[lines~1267--1268]{Gross2012IndexTheory} the even algebra is a
    two-site full matrix algebra while the odd algebra is the scalars, and no
    translation carries one onto the other.

    Covariance of the support algebras under translation is the identity
    displayed at
    \cite[lines~1218--1226]{SchumacherWerner2004QCA} for the algebras defined
    at \cite[lines~1199--1206]{SchumacherWerner2004QCA}, and the ordered even
    and odd factors are equation \texttt{RR2x} at
    \cite[lines~1251--1266]{Gross2012IndexTheory}.  The translation-covariant
    fixed-$d$ setting is that of
    \cite[lines~2292--2298]{Cirac2017MPU}.

    \textbf{Scope restriction (one-dimensional homogeneous automorphism):}
    Schumacher--Werner state the translated-support identity for a
    translation-commuting homomorphism on the cube lattice in arbitrary
    dimension with the cubic nearest-neighbour scheme,
    \cite[lines~309--318 and 920--941]{SchumacherWerner2004QCA}.  The statement
    above treats the one-dimensional invertible case, as recorded in
    \cite{gap:sw04_support_covariance_scope}.  GNVW allow
    site-dependent cell sizes and assume no translation covariance, so
    (\ref{eq:qca_two_site_covariance}) is not a theorem of that paper; the
    fixed-$d$ restriction of the family $\mathcal R_y$ is recorded in
    \cite{gap:gnvw12_site_dependent_adjacent_generation_scope}.
\end{theorem}

\begin{proof}\leanok
    \uses{lem:qca_quasi_local_coefficient_compression,
        lem:qca_pair_region_two_site_translation,
        def:qca_site_indexed_support_algebras,
        def:qca_quasi_local_translation,
        thm:qca_finite_region_restriction_embedding}
    Fix $x$ and put $\Lambda=E_x$, $\Gamma=L_x$, $\Delta=P_x$.  A
    finite-region observable belongs to the finite local image exactly when
    its quasi-local copy lies in the image of $\mathcal A_\Lambda$ under
    $\omega$.  Hence, by
    (\ref{eq:qca_quasi_local_left_compression}), the canonical quasi-local
    copy of $\mathcal R_{2x}$ is generated by the elements
    \begin{align}
        \sum_{t\in\mathcal C_\Delta}
        \widehat{e^\Delta_{tr}}\ \omega\bigl(\widehat B\bigr)\
        \widehat{e^\Delta_{st}},
        \qquad B\in\mathcal A_\Lambda,\quad r,s\in\mathcal C_\Delta,
        \label{eq:qca_left_support_generators}
    \end{align}
    and the copy of $\mathcal R_{2x+1}$ is generated by the analogous
    elements built from (\ref{eq:qca_quasi_local_right_compression}) with
    $\Gamma$ in place of $\Delta$.  Since a star-algebra automorphism carries
    the algebra generated by a set to the algebra generated by the image of
    that set, it suffices to translate the two generating families.

    Translation carries a matrix unit of a region to the matrix unit of the
    translated region at the translated configurations, and covariance gives
    $\tau_2(\omega(\widehat B))=\omega(\tau_2(\widehat B))$, whose right-hand
    side is $\omega$ applied to the quasi-local copy of the translated
    observable.  Reindexing the summation over $\mathcal C_\Delta$ along the
    configuration translation therefore turns
    (\ref{eq:qca_left_support_generators}) into the same family for the
    regions $\Lambda+2$ and $\Delta+2$, and the translation is a bijection
    between the two families.  By
    (\ref{eq:qca_pair_region_two_site_translation}) those regions are
    $E_{x+1}$ and $P_{x+1}$, so the translated family generates
    $\widehat{\mathcal R}_{2x+2}$.  The odd case is identical with $\Gamma$
    in place of $\Delta$ and $L_{x+1}$ in place of $P_{x+1}$.

    Finally, splitting $y$ by parity and using $2x+2=2(x+1)$ and
    $2x+3=2(x+1)+1$ gives (\ref{eq:qca_two_site_covariance}) and
    (\ref{eq:qca_two_site_covariance_parities}).
\end{proof}

Equation~(A1) in lines~2300--2306 of \cite{Cirac2017MPU} defines the
local action associated with an MPU by eventual finite-chain conjugation.  In
lines~2300--2306, eventual constancy is used to obtain a norm-preserving
star-algebra automorphism of $\mathcal A_{\mathrm{loc}}$, together with a
translation formula and the support bound
$\mathcal N=[-4D^4,4D^4]\cap\mathbb Z$, and then to extend this automorphism to
$\mathcal A$.  The results below isolate the direct-limit and completion
argument.  They begin with compatible forward and inverse finite-local actions
and assume a forward norm formula, a support bound, and a unit-translation
formula.  The inverse norm formula follows from forward isometry and mutual
invertibility.  The numerical radius above belongs to the MPU-specific
stabilization argument: the generic result below neither proves that radius nor
constructs the finite-local actions from an MPU.  It also does not provide the
blocked circuit or MPU standard form discussed at line~2308.

\begin{definition}[Compatible finite-local automorphism data]
    \label{def:qca_compatible_local_automorphism}
    \lean{SpinChain.CompatibleLocalAutomorphism}
    \leanok
    \uses{thm:qca_algebraic_local_algebra_universal}
    For each finite region $\Lambda$, let
    \begin{align}
        f_\Lambda,g_\Lambda\colon\mathcal A_\Lambda
         & \longrightarrow\mathcal A_{\mathrm{loc}}
    \end{align}
    be star-algebra homomorphisms such that, whenever
    $\Lambda\subseteq\Gamma$,
    \begin{align}
        f_\Gamma\circ\iota_{\Lambda,\Gamma} & =f_\Lambda, &
        g_\Gamma\circ\iota_{\Lambda,\Gamma} & =g_\Lambda.
    \end{align}
    Let $f,g\colon\mathcal A_{\mathrm{loc}}\to
        \mathcal A_{\mathrm{loc}}$ be the maps induced by these compatible
    families.  Compatible finite-local automorphism data require
    \begin{align}
        g(f_\Lambda(A)) & =\iota_\Lambda(A), &
        f(g_\Lambda(A)) & =\iota_\Lambda(A)
    \end{align}
    for every $A\in\mathcal A_\Lambda$.
\end{definition}

\begin{definition}[Induced algebraic-local homomorphisms]
    \label{def:qca_compatible_local_homs}
    \lean{SpinChain.CompatibleLocalAutomorphism.forwardHom,
        SpinChain.CompatibleLocalAutomorphism.inverseHom}
    \leanok
    \uses{def:qca_compatible_local_automorphism, thm:qca_algebraic_local_algebra_universal}
    The compatible families induce star-algebra homomorphisms
    \begin{align}
        f,g\colon\mathcal A_{\mathrm{loc}}
         & \longrightarrow\mathcal A_{\mathrm{loc}}.
    \end{align}
\end{definition}

\begin{proof}\leanok
    \uses{thm:qca_algebraic_local_algebra_universal}
    Apply the star-algebra universal property to each compatible family.
\end{proof}

\begin{lemma}[Evaluation of the induced homomorphisms]
    \label{lem:qca_compatible_local_homs_evaluation}
    \lean{SpinChain.CompatibleLocalAutomorphism.forwardHom_localObservable,
        SpinChain.CompatibleLocalAutomorphism.inverseHom_localObservable}
    \leanok
    \uses{def:qca_compatible_local_homs}
    For every finite region $\Lambda$ and $A\in\mathcal A_\Lambda$,
    \begin{align}
        f(\iota_\Lambda(A)) & =f_\Lambda(A), &
        g(\iota_\Lambda(A)) & =g_\Lambda(A).
    \end{align}
\end{lemma}

\begin{proof}\leanok
    \uses{thm:qca_algebraic_local_algebra_universal}
    These are the evaluation identities supplied by the universal property.
\end{proof}

\begin{definition}[Induced algebraic-local equivalence]
    \label{def:qca_compatible_algebraic_equiv}
    \lean{SpinChain.CompatibleLocalAutomorphism.algebraicEquiv}
    \leanok
    \uses{def:qca_compatible_local_automorphism, def:qca_compatible_local_homs}
    Compatible finite-local automorphism data determine a star-algebra
    equivalence
    \begin{align}
        F\colon\mathcal A_{\mathrm{loc}}
         & \simeq\mathcal A_{\mathrm{loc}}
    \end{align}
    whose forward and inverse maps are $f$ and $g$, respectively.
\end{definition}

\begin{proof}\leanok
    \uses{def:qca_compatible_local_automorphism, lem:qca_compatible_local_homs_evaluation}
    Induct on finite-region representatives in the algebraic direct limit.
    The two inverse assumptions give
    \begin{align}
        g(f(\iota_\Lambda(A))) & =g(f_\Lambda(A))=\iota_\Lambda(A), \\
        f(g(\iota_\Lambda(A))) & =f(g_\Lambda(A))=\iota_\Lambda(A).
    \end{align}
    Thus $f$ and $g$ are mutually inverse.
\end{proof}

\begin{lemma}[Evaluation of the algebraic equivalence]
    \label{lem:qca_compatible_algebraic_equiv_evaluation}
    \lean{SpinChain.CompatibleLocalAutomorphism.algebraicEquiv_apply,
        SpinChain.CompatibleLocalAutomorphism.algebraicEquiv_symm_apply}
    \leanok
    \uses{def:qca_compatible_algebraic_equiv, def:qca_compatible_local_homs}
    For every $A\in\mathcal A_{\mathrm{loc}}$,
    \begin{align}
        F(A) & =f(A), & F^{-1}(A) & =g(A).
    \end{align}
\end{lemma}

\begin{proof}\leanok
    \uses{def:qca_compatible_algebraic_equiv}
    Both identities follow from the defining forward and inverse maps of $F$.
\end{proof}

\begin{definition}[Forward local norm hypothesis]
    \label{def:qca_compatible_local_norm}
    \lean{SpinChain.CompatibleLocalAutomorphism.ForwardNormPreserving}
    \leanok
    \uses{def:qca_compatible_local_automorphism, def:qca_algebraic_local_operator_norm}
    Assume $d>0$.  The forward local norm hypothesis is
    \begin{align}
        \lVert f_\Lambda(A)\rVert & =\lVert A\rVert
    \end{align}
    for every finite region $\Lambda$ and every $A\in\mathcal A_\Lambda$.
\end{definition}

\begin{lemma}[Forward algebraic norm formula]
    \label{lem:qca_compatible_algebraic_norm}
    \lean{SpinChain.CompatibleLocalAutomorphism.norm_algebraicEquiv}
    \leanok
    \uses{def:qca_compatible_algebraic_equiv, def:qca_compatible_local_norm, def:qca_algebraic_local_operator_norm}
    Under the forward local norm hypothesis, every
    $A\in\mathcal A_{\mathrm{loc}}$ satisfies
    \begin{align}
        \lVert F(A)\rVert & =\lVert A\rVert.
    \end{align}
\end{lemma}

\begin{proof}\leanok
    \uses{lem:qca_compatible_algebraic_equiv_evaluation, lem:qca_compatible_local_homs_evaluation, def:qca_algebraic_local_operator_norm}
    Induct on finite-region representatives.  For
    $A=\iota_\Lambda(A_\Lambda)$,
    \begin{align}
        \lVert F(A)\rVert
         & =\lVert f_\Lambda(A_\Lambda)\rVert
        =\lVert A_\Lambda\rVert
        =\lVert A\rVert.
    \end{align}
\end{proof}

\begin{lemma}[Forward and inverse algebraic isometries]
    \label{lem:qca_compatible_algebraic_isometry}
    \lean{SpinChain.CompatibleLocalAutomorphism.algebraicEquiv_isometry,
        SpinChain.CompatibleLocalAutomorphism.algebraicEquiv_symm_isometry}
    \leanok
    \uses{lem:qca_compatible_algebraic_norm, def:qca_compatible_algebraic_equiv}
    Under the forward local norm hypothesis, for all
    $A,B\in\mathcal A_{\mathrm{loc}}$,
    \begin{align}
        \lVert F(A)-F(B)\rVert           & =\lVert A-B\rVert, &
        \lVert F^{-1}(A)-F^{-1}(B)\rVert & =\lVert A-B\rVert.
    \end{align}
\end{lemma}

\begin{proof}\leanok
    \uses{lem:qca_compatible_algebraic_norm, def:qca_compatible_algebraic_equiv}
    Apply the forward norm identity to $A-B$ to prove that $F$ is an isometry.
    Since $F\circ F^{-1}$ is the identity, the right-inverse law then makes
    $F^{-1}$ an isometry as well.
\end{proof}

\begin{lemma}[Inverse algebraic norm formula]
    \label{lem:qca_compatible_algebraic_inverse_norm}
    \lean{SpinChain.CompatibleLocalAutomorphism.norm_algebraicEquiv_symm}
    \leanok
    \uses{lem:qca_compatible_algebraic_isometry}
    Under the forward local norm hypothesis, every
    $A\in\mathcal A_{\mathrm{loc}}$ satisfies
    \begin{align}
        \lVert F^{-1}(A)\rVert & =\lVert A\rVert.
    \end{align}
\end{lemma}

\begin{proof}\leanok
    \uses{lem:qca_compatible_algebraic_isometry}
    Apply the inverse isometry to the distance from $A$ to $0$ and use
    $F^{-1}(0)=0$.
\end{proof}

\begin{definition}[Completed compatible local automorphism]
    \label{def:qca_compatible_quasi_local_equiv}
    \lean{SpinChain.CompatibleLocalAutomorphism.quasiLocalEquiv}
    \leanok
    \uses{def:qca_compatible_algebraic_equiv, lem:qca_compatible_algebraic_isometry, def:completion_star_alg_equiv}
    Under the forward local norm hypothesis, the algebraic equivalence extends
    to a star-algebra equivalence
    \begin{align}
        \overline F\colon\mathcal A & \simeq\mathcal A
    \end{align}
    of the quasi-local completion.
\end{definition}

\begin{proof}\leanok
    \uses{lem:qca_compatible_algebraic_isometry, def:completion_star_alg_equiv}
    The forward norm formula makes $F$ an isometry, and the inverse law makes
    $F^{-1}$ an isometry.  Apply the completion extension for star-algebra
    equivalences to these two continuous maps.
\end{proof}

\begin{lemma}[Dense evaluation of the completed automorphism]
    \label{lem:qca_compatible_quasi_local_dense_evaluation}
    \lean{SpinChain.CompatibleLocalAutomorphism.quasiLocalEquiv_algebraicToQuasiLocal,
        SpinChain.CompatibleLocalAutomorphism.quasiLocalEquiv_symm_algebraicToQuasiLocal}
    \leanok
    \uses{def:qca_compatible_quasi_local_equiv, def:completion_star_alg_equiv}
    Let $j\colon\mathcal A_{\mathrm{loc}}\to\mathcal A$ be the canonical
    completion map.  For every $A\in\mathcal A_{\mathrm{loc}}$,
    \begin{align}
        \overline F(j(A))      & =j(F(A)),      &
        \overline F^{-1}(j(A)) & =j(F^{-1}(A)).
    \end{align}
\end{lemma}

\begin{proof}\leanok
    \uses{def:qca_compatible_quasi_local_equiv, def:completion_star_alg_equiv}
    These are the canonical-image formulas for the continuous extension of a
    star-algebra equivalence and for its inverse.
\end{proof}

\begin{lemma}[Finite evaluation of the completed automorphism]
    \label{lem:qca_compatible_quasi_local_finite_evaluation}
    \lean{SpinChain.CompatibleLocalAutomorphism.quasiLocalEquiv_quasiLocalObservable,
        SpinChain.CompatibleLocalAutomorphism.quasiLocalEquiv_symm_quasiLocalObservable}
    \leanok
    \uses{lem:qca_compatible_quasi_local_dense_evaluation, lem:qca_compatible_algebraic_equiv_evaluation, lem:qca_compatible_local_homs_evaluation, def:qca_quasi_local_observable}
    For every finite region $\Lambda$ and $A\in\mathcal A_\Lambda$,
    \begin{align}
        \overline F(j_\Lambda(A))      & =j(f_\Lambda(A)), &
        \overline F^{-1}(j_\Lambda(A)) & =j(g_\Lambda(A)).
    \end{align}
\end{lemma}

\begin{proof}\leanok
    \uses{lem:qca_compatible_quasi_local_dense_evaluation, lem:qca_compatible_algebraic_equiv_evaluation, lem:qca_compatible_local_homs_evaluation}
    Substitute $j_\Lambda=j\circ\iota_\Lambda$ into the dense evaluation
    formulas and use
    $F(\iota_\Lambda(A))=f_\Lambda(A)$ and
    $F^{-1}(\iota_\Lambda(A))=g_\Lambda(A)$.
\end{proof}

\begin{definition}[Finite-local support bound]
    \label{def:qca_compatible_local_support_bound}
    \lean{SpinChain.CompatibleLocalAutomorphism.MapsSupportWithin}
    \leanok
    \uses{def:qca_compatible_local_automorphism, def:qca_supported_in, def:qca_region_sumset}
    For a finite neighborhood $\mathcal N\subset\mathbb Z$, the forward family
    has support bound $\mathcal N$ if
    \begin{align}
        f_\Lambda(A) & \in
        \iota_{\Lambda+\mathcal N}
        (\mathcal A_{\Lambda+\mathcal N})
    \end{align}
    for every finite region $\Lambda$ and every $A\in\mathcal A_\Lambda$.
\end{definition}

\begin{theorem}[Propagation from a finite-local support bound]
    \label{thm:qca_compatible_local_propagation}
    \lean{SpinChain.CompatibleLocalAutomorphism.propagatesWithin_quasiLocalEquiv}
    \leanok
    \uses{def:qca_compatible_local_support_bound, lem:qca_compatible_quasi_local_finite_evaluation, def:qca_propagates_within}
    If the forward family has support bound $\mathcal N$, then the completed
    automorphism $\overline F$ propagates within $\mathcal N$.
\end{theorem}

\begin{proof}\leanok
    \uses{def:qca_compatible_local_support_bound, lem:qca_compatible_quasi_local_finite_evaluation}
    Let $X=j_\Lambda(A)$ be supported in $\Lambda$.  Choose
    $B\in\mathcal A_{\Lambda+\mathcal N}$ such that
    $f_\Lambda(A)=\iota_{\Lambda+\mathcal N}(B)$.  Then
    \begin{align}
        \overline F(X)
         & =j(f_\Lambda(A))
        =j_{\Lambda+\mathcal N}(B),
    \end{align}
    so the image is supported in $\Lambda+\mathcal N$.
\end{proof}

\begin{definition}[Finite-local unit-translation formula]
    \label{def:qca_compatible_local_unit_translation}
    \lean{SpinChain.CompatibleLocalAutomorphism.CommutesWithUnitTranslation}
    \leanok
    \uses{def:qca_compatible_local_automorphism, def:qca_algebraic_local_translation}
    The forward family commutes with unit translation if, for every finite
    region $\Lambda$ and every $A\in\mathcal A_\Lambda$,
    \begin{align}
        f_{\Lambda+1}(\tau_{1,\Lambda}(A))
         & =\tau_1(f_\Lambda(A)).
    \end{align}
\end{definition}

\begin{lemma}[Algebraic commutation with unit translation]
    \label{lem:qca_compatible_algebraic_unit_translation}
    \lean{SpinChain.CompatibleLocalAutomorphism.algebraicEquiv_commutes_unitTranslation}
    \leanok
    \uses{def:qca_compatible_local_unit_translation, def:qca_compatible_algebraic_equiv, def:qca_algebraic_local_translation}
    If the finite-local unit-translation formula holds, then every
    $A\in\mathcal A_{\mathrm{loc}}$ satisfies
    \begin{align}
        F(\tau_1(A)) & =\tau_1(F(A)).
    \end{align}
\end{lemma}

\begin{proof}\leanok
    \uses{def:qca_compatible_local_unit_translation, lem:qca_compatible_algebraic_equiv_evaluation, lem:qca_compatible_local_homs_evaluation, def:qca_algebraic_local_translation}
    Induct on finite-region representatives and write
    $A=\iota_\Lambda(A_\Lambda)$.  Then
    \begin{align}
        F(\tau_1(A))
         & =f_{\Lambda+1}(\tau_{1,\Lambda}(A_\Lambda))
        =\tau_1(f_\Lambda(A_\Lambda))
        =\tau_1(F(A)).
    \end{align}
\end{proof}

\begin{theorem}[Translation covariance from unit translation]
    \label{thm:qca_compatible_translation_covariant}
    \lean{SpinChain.CompatibleLocalAutomorphism.translationCovariant_quasiLocalEquiv}
    \leanok
    \uses{lem:qca_compatible_algebraic_unit_translation, lem:qca_compatible_quasi_local_dense_evaluation, def:qca_quasi_local_translation, thm:qca_translation_covariant_iff_commute_one}
    If the finite-local unit-translation formula holds, then the completed
    automorphism $\overline F$ is translation covariant.
\end{theorem}

\begin{proof}\leanok
    \uses{lem:qca_compatible_algebraic_unit_translation, lem:qca_compatible_quasi_local_dense_evaluation, def:qca_quasi_local_translation, thm:qca_translation_covariant_iff_commute_one}
    By Theorem~\ref{thm:qca_translation_covariant_iff_commute_one}, it suffices
    to prove that $\overline F$ commutes with translation by one site.  Induct
    on the completion.  On every $A\in\mathcal A_{\mathrm{loc}}$,
    \begin{align}
        \overline F(\overline\tau_1(j(A)))
         & =j(F(\tau_1(A)))
        =j(\tau_1(F(A)))
        =\overline\tau_1(\overline F(j(A))).
    \end{align}
    The equality set is closed because both composites are continuous, so the
    identity extends from the canonical dense subalgebra to the completion.
\end{proof}

\begin{theorem}[QCA assembled from compatible finite-local actions]
    \label{thm:qca_compatible_local_automorphism}
    \lean{SpinChain.CompatibleLocalAutomorphism.isQCA_quasiLocalEquiv}
    \leanok
    \uses{def:qca_compatible_quasi_local_equiv, thm:qca_compatible_local_propagation, thm:qca_compatible_translation_covariant, def:qca}
    Let $d>0$.  Suppose the compatible forward family preserves the local
    operator norm, has a finite support bound, and satisfies the finite-local
    unit-translation formula.  Then the completed automorphism $\overline F$ is
    a one-dimensional QCA.
\end{theorem}

\begin{proof}\leanok
    \uses{thm:qca_compatible_local_propagation, thm:qca_compatible_translation_covariant}
    The support hypothesis supplies a finite forward propagation neighborhood,
    while the unit-translation formula supplies translation covariance.  These
    are exactly the two defining properties of a one-dimensional QCA.
\end{proof}
